% 本文件由 LaTeX 排版 Agent 根据有效 translation.json 自动生成。
% author=Socrates Scholasticus (c. 379-c. 450); work_id=2601; title=教会历史

\chapter{教会历史}

% structure: p0001 source=h1 target=omitted reason=page-title-already-rendered-by-parent

第一卷\newline{}第二卷\newline{}第三卷\newline{}第四卷\newline{}第五卷\newline{}第六卷\newline{}第七卷\par

\SourceRecord{Translated by A.C. Zenos.；Nicene and Post-Nicene Fathers, Second Series；Vol. 2.；Edited by Philip Schaff and Henry Wace.；Buffalo, NY: Christian Literature Publishing Co.,；1890.；<http://www.newadvent.org/fathers/2601.htm>.}

% structure: p0001 source=h1 target=section reason=page-title

\section{教会史（第一卷）}

% structure: p0002 source=h2 target=subsection reason=relative-heading-rank; base=h2

\subsection{第一章 作品序言}

欧瑟伯，别号庞斐禄，撰写了十卷\WorkTitle{教会史}，结束于君士坦丁皇帝时期，那时戴克里先对基督徒的迫害已告终止。同一作者在撰写\WorkTitle{君士坦丁生平}中，仅简略地处理了关于亚略的事，更专注于其修辞的华美和皇帝的颂赞，而非对事实的准确陈述。如今，我们打算记载自他那个时代直至我们今日在教会中所发生之事的细节，便从他省略的细目开始叙述；我们不会刻意炫耀辞藻，而是将我们从文献中能搜集到的，以及从熟悉这些事实的人所讲述中听到的，呈现在读者面前。鉴于此事与当前主题关系重大，宜乎简要叙述君士坦丁归信基督信仰的经过，以此事件为开端。\par

% structure: p0004 source=h2 target=subsection reason=relative-heading-rank; base=h2

\subsection{第二章 君士坦丁皇帝如何成为基督徒}

当戴克里先与马克西米安（别号赫丘利）经双方同意退位并隐退后，与其共同执政的马克西米安（别号伽勒里乌斯）进入意大利，任命了两名凯撒：马克西明治理帝国东部，塞维鲁治理意大利。然而，在不列颠，君士坦丁被拥立为皇帝，以代替其父君士坦提乌斯；君士坦提乌斯死于第二百七十一届奥林匹亚纪元第一年七月二十五日。在罗马，马克西米安·赫丘利的儿子马克森提乌斯被近卫军士兵立为僭主，而非皇帝。在此情形下，赫丘利被重新掌权的欲望所驱使，试图除掉其子马克森提乌斯；但士兵阻止了他，不久后他在奇里乞亚的塔尔苏斯去世。同时，塞维鲁凯撒被伽勒里乌斯·马克西米安派往罗马以擒拿马克森提乌斯，却因士兵背叛而被杀。最后，伽勒里乌斯·马克西米安，这位曾掌握最高权力的人也死了，此前他已任命其老友和战友李锡尼（达契亚人）为继承人。与此同时，马克森提乌斯残酷压迫罗马人民，与其说是国王，不如说是僭主；他无耻地侵犯贵族之妻，杀害许多无辜者，并犯下其他类似暴行。君士坦丁皇帝得知此事后，竭力使罗马人摆脱他的奴役，并立即思考如何推翻僭主。当他的心思专注于这一重大主题时，他考虑在战争中应祈求哪一位神明的帮助。他开始意识到戴克里先一派并未从他们所试图安抚的外教神祇中获得任何益处；但他自己的父亲君士坦提乌斯，因弃绝了希腊的各种宗教，却生活得远为昌盛。在此犹豫不决之际，当他率领军队行进时，一个超自然的异象，超越一切描述，向他显现。事实上，大约在太阳过午西斜的时刻，他看见天上有一道光柱，呈十字架形状，上面铭刻着这些字：\Quote{\Emph{以此克胜}}。这个记号的出现令皇帝惊愕，几乎不相信自己的眼睛，便问周围之人是否也看见同样的景象；当他们一致说看见时，皇帝的心因这神圣而奇异的显现而坚固。次日夜晚，他在睡梦中看见基督，基督指示他按照所见之像制作一面军旗，并用以对抗敌人，作为必胜的保证。遵从这神圣的启示，他令人制作了一面十字架形状的军旗，至今仍保存在皇宫中；他更热切地推进措施，在罗马城门前、穆尔维桥附近攻击敌人并战胜了他，马克森提乌斯本人溺死于河中。这场胜利是在征服者统治的第七年取得的。此后，当与他共同执政且是妹夫（因其妹君士坦提娅嫁给李锡尼）的李锡尼驻守东方时，君士坦丁皇帝鉴于所蒙受的巨大恩惠，向天主献上感恩作为报谢；这些感恩包括：解除基督徒的迫害，召回被流放者，释放被囚禁者，归还被剥夺者的财产；此外，他重建了教堂，并以极大热忱完成了这一切。大约此时，已退位的戴克里先死于达尔马提亚的萨洛纳。\par

% structure: p0006 source=h2 target=subsection reason=relative-heading-rank; base=h2

\subsection{第三章 君士坦丁优待基督徒，其同僚李锡尼却加以迫害}

现在，\Person{君士坦丁}皇帝既已拥抱基督信仰，便以他宣称的基督徒身份行事，重建圣堂，并以辉煌的奉献加以装饰；他还关闭或摧毁了外教人的庙宇，并使他们庙中的神像受到民众的蔑视。但他的同僚\Person{李锡尼乌斯}坚持其外教信条，仇视基督徒；虽然由于惧怕\Person{君士坦丁}皇帝，他避免激起公开的迫害，却设法暗中谋害他们，最后竟不加掩饰地骚扰他们。然而，这场迫害是地方性的，仅局限于\Person{李锡尼乌斯}本人所在的地区；但由于这些及其他公开的暴行未能长久瞒过\Person{君士坦丁}，\Person{李锡尼乌斯}发现后者对其行为感到愤怒，便求助于辩解。他这样平息了\Person{君士坦丁}后，便假装缔结友谊同盟，并以许多誓言承诺不再施行暴政。但他刚发誓便犯了伪善罪；因为他既没有改变其暴虐性情，也没有停止迫害基督徒。他甚至颁布法律，禁止主教们拜访未皈依的外教人，以免这成为使他们改信基督信仰的借口。于是这场迫害既众所周知又隐密。它名义上被认可，实际上却显而易见；因为那些遭受迫害的人在身心和财产上都遭受了最严厉的苦楚。\par

% structure: p0008 source=h2 target=subsection reason=relative-heading-rank; base=h2

\subsection{第4章. 因基督徒而起的\Person{君士坦丁}与\Person{李锡尼乌斯}之间的战争。}

此举招致了\Person{君士坦丁}皇帝最严重的不悦；他们成了敌人，他们之间虚假的友谊盟约已被破坏。不久之后，他们作为公开的敌人拿起武器互相攻击。经过数场海陆交战，\Person{李锡尼乌斯}最终在\Place{比提尼亚}的\Place{克勒索波利斯}（一个\Place{卡尔西顿}人的港口）附近被彻底击败，并向\Person{君士坦丁}投降。因此，\Person{君士坦丁}将他活捉后，以最大的人道对待他，绝不愿处死他，而是命令他定居在\Place{得撒洛尼}安静度日。然而，他安静了短时间后，又设法召集了一些蛮族雇佣兵，试图通过重新诉诸武力来弥补他最近的失败。皇帝获悉他的行动后，下令将他处死，命令得以执行。\Person{君士坦丁}因此获得了唯一的统治权，于是被宣布为至高独裁君主，并再次致力促进基督徒的福祉。他以多种方式这样做，因他的努力，基督信仰享有了持续不断的和平。但这种平静状态之后不久便发生了内部纷争，我现在将尽力描述其性质和起因。\par

% structure: p0010 source=h2 target=subsection reason=relative-heading-rank; base=h2

\subsection{第5章. \Person{亚略}与他的主教\Person{亚历山大}的争论。}

\Person{伯多禄}（亚历山大的主教）在\Person{戴克里先}统治下殉道之后，\Person{阿基拉斯}被任命为主教职务，\Person{亚历山大}在上文提到的和平时期继任了他。\Person{亚历山大}在一次无畏地履行其教导和管理教会职责时，有一天在长老团及其他神职人员面前，试图以可能过于哲学的细致来解释那个伟大的神学奥秘——\Emph{天主圣三的合一}。他辖下的一位名叫\Person{亚略}的长老，拥有不小的逻辑敏锐力，想象主教是在微妙地教导与\Person{利比亚人撒伯里乌}相同的观点，便由于好争论而采取了与\Person{撒伯里乌}相反的意见，并自以为有力地回应了主教所说的话。\Quote{如果，}他说，\Quote{圣父生了圣子，那么被生的就有了存在的开始；由此可见，曾有一个时期圣子不存在。因此必然得出，他的本体来自无有。}\par

% structure: p0012 source=h2 target=subsection reason=relative-heading-rank; base=h2

\subsection{第6章. 教会因这场争论开始分裂；\Person{亚历山大}主教将\Person{亚略}及其追随者开除教籍。}

从他那新颖的推理链中得出这一推论后，他激发了许多人思考这个问题；于是，从一点火星燃起了一场大火：因为始于亚历山大教会的邪恶，蔓延至整个埃及、利比亚和上底比斯，最终扩散到其余各省和城市。许多人也采纳了\Person{亚略}的意见，但尤其是\Person{欧瑟伯}热切地为之辩护：不是恺撒利亚的\Person{欧瑟伯}，而是那位以前曾担任\Place{贝鲁特}教会主教，当时不知何故占据着\Place{比提尼亚}的\Place{尼科美底亚}主教职务的\Person{欧瑟伯}。当\Person{亚历山大}意识到这些事（既通过他自己的观察，也通过传闻）时，他被激怒到极点，便召集了许多主教的会议，将\Person{亚略}及其异端的支持者开除教籍；同时他写信给各城市的主教，信函如下：——\par

\Emph{亚历山大主教的书信}\par

至各处天主教会我们亲爱的、最受尊敬的共事司铎们，\Person{亚历山大}在主内问安。\par

既然天主教会是一个身体，我们在圣经中被命令要维持\Quote{团结与和平的联系}\BibleRef{弗 4:3}，我们就有责任写信，相互告知我们各自的情况，以便\Quote{若一个肢体受苦或喜乐，我们可以彼此同情或一同喜乐}\BibleRef{格前 12:26}。因此，要知道在我们的教区最近出现了一些无法无天、反基督的人，他们教导背教，这样的人可被公正地视为并称为假基督的前驱。我本想把这种混乱掩盖起来，希望这邪恶只限于背教者本身，不至于蔓延到其他地区，污染某些单纯之人的耳朵。但是，既然现在在\Place{尼哥美狄亚}的\Person{欧瑟比}认为教会的事务在他控制之下，因为他确实擅自离开了他在\Place{贝里托}的职务，毫无顾忌地在\Place{尼哥美狄亚}教会掌握了权力，并且他自任这些背教者的首领，甚至胆敢为他们四处发送推荐信，企图诱使一些无知之人陷入这种最不虔敬、反基督的异端，我感到有义务不再沉默，因为知道法律上所写的，而要将这一切通知你们，好使你们了解背教者是谁，以及他们的异端是多么可鄙，并且不要理会\Person{欧瑟比}写给你们的任何东西。因为他现在想要重拾他以前的恶意（这恶意似乎已被时间遗忘），他假装为他们写信，而事实本身清楚地表明他这样做是为了推进他自己的目的。这些就是成为背教者的人：\Person{亚略}、\Person{阿基拉斯}、\Person{艾塔勒斯}、\Person{卡尔波内斯}、另一个\Person{亚略}、\Person{萨尔马特斯}、\Person{欧佐伊乌斯}、\Person{路济乌斯}、\Person{犹利安}、\Person{梅纳斯}、\Person{赫拉迪斯}和\Person{加犹斯}；此外还有\Person{塞根都斯}和\Person{特奥纳斯}，他们曾被称为主教。他们发明并断言、违背圣经的教义如下：天主并不一直是圣父，而是有一个时期他不是圣父；天主的圣言不是从永远存在的，而是从无中造成的；因为永远存在的天主（\Quote{我是}---永恒的那位）从无中造了先前不存在的那一位；因此，有一个时期他并不存在，因为圣子是受造物和作品。他既不像圣父那样具有本质，也不是本性上圣父的真圣言或真智慧，而只是他的作品和受造物之一，被错误地称为圣言和智慧，因为他自己是由天主的圣言和存在于天主内的智慧所造，天主藉此创造了万物和他自己。因此，他的本性是可变的、易受改变的，正如所有其他理性的受造物一样：故圣言与天主的本质是相异、不同的；圣父是圣子无法解释的，圣子也不能看见圣父，因为圣言既不完美、精确地认识圣父，也不能清晰地看见他。圣子不知道他自己本质的本性：因为他是为了我们而被造的，好使天主能藉着他如同一个工具而创造我们；若非天主愿意创造我们，他也绝不会存在。\par

于是有人问他们：\Quote{天主的圣言能否改变，正如魔鬼那样？}他们竟敢回答：\Quote{能，因为祂是受生的，所以能改变。}于是我们，与埃及和利比亚的主教们，近一百人一同集会，因亚略无耻地公开承认这些异端，将他连同所有赞成这些异端的人一起弃绝。然而欧瑟伯的党羽却接纳了他们，试图将虚假与真理、不敬与神圣混为一谈。但他们不会得逞，因为真理必要得胜；\BibleQuote{光明与黑暗不能相通，基督与贝里雅耳也不能和谐。}\BibleRef{格后 6:14}谁曾听过这样的亵渎？或有哪个虔诚人听到这些，不惊恐万分，掩耳不听，唯恐这些话的污秽玷污他的听觉？谁听到若望说：\BibleQuote{在起初已有圣言}，而不谴责那些说：\BibleQuote{曾有一段时间圣言不存在}的人呢？或者谁听到福音中说到\BibleQuote{独生子}，并说\BibleQuote{万物是藉着祂造成的}，而不憎恶那些宣称圣子是受造物之一的人呢？祂怎能是那由祂自己所造之物中的一件呢？如果祂被列入受造物中，又怎能是独生子呢？祂怎能从无中产生，因为圣父曾说：\Quote{我的心涌出美词}；\Quote{我在黎明前从怀中生了你}？或者，祂怎会与圣父的本体不同？祂是\BibleQuote{祂的完美肖像}\BibleRef{哥 1:15}，是\BibleQuote{祂光荣的反映}\BibleRef{希 1:3}，且说：\BibleQuote{看见了我，就是看见了父}？再者，如果圣子是圣言和天主的智慧，怎会有一段时间祂不存在呢？因为这就等于说天主曾一度没有圣言和智慧。祂怎能是可变的、能改变的？祂亲自说：\BibleQuote{我在父内，父也在我内}\BibleRef{若 14:10}；\BibleQuote{我与父原是一体}\BibleRef{若 10:30}；又通过先知说：\BibleQuote{请看，我是自有者，而且没有改变}\BibleRef{拉 3:6}。但即使有人将这说法应用在圣父身上，如今也更适合用在圣言身上，因为祂降生成人并未改变，反而如宗徒所说：\BibleQuote{耶稣基督，昨天、今天、直到永远，都是一样的}\BibleRef{希 13:8}。但保禄明明宣告：\BibleQuote{万物都是为祂而有，藉祂而存在}\BibleRef{希 2:10}，又有什么能说服他们说祂是因我们而受造的呢？确实，他们亵渎地声称圣子不完全认识圣父，这不足为奇；因为他们既已决意与基督作战，便连主亲口说的话也拒绝，祂说：\BibleQuote{正如父认识我，我也认识父}\BibleRef{若 10:15}。因此，如果父只是部分地认识子，显然子也只是部分地认识父。但如果断言这点是不妥的，而且承认父完全认识子，那么显然，正如父认识自己的圣言，圣言也认识自己的父，祂正是祂的圣言。我们陈述这些事，阐释神圣的圣经，经常驳倒了他们；但他们又如变色龙一般改变，努力将那经上所记的话用在自己身上：\Quote{恶人到了罪孽的深处，便变得轻慢}。在此之前已兴起许多异端，它们胆大妄为，超出了所有限度，陷入了完全的痴迷；但这些人，因试图在一切言论中推翻圣言的神性，更接近了假基督，反而相对减轻了先前异端的可憎。因此，他们被教会公开弃绝，并遭绝罚。我们确实为这些人的丧亡而悲伤，尤其因为他们先前曾受教会道理的教导，如今却背弃了这些道理。尽管如此，我们也不十分惊讶，因为依默纳约和非肋托也曾同样跌倒\BibleRef{弟后 2:17}；在他们之前，犹达斯原是救主的门徒，后来却离弃了祂，成了背叛者。关于这些人，我们也并非没有预先的警告：因为主亲自说：\Quote{你们要谨慎，免得有人欺骗你们：因为将有许多人冒我的名来，说：\InnerQuote{我是基督}，并且要欺骗许多人}\BibleRef{玛 24:4}；又说：\Quote{时期近了，你们不要跟从他们}\BibleRef{路 21:8}。保禄从救主那里得知这些事，便写道：\Quote{在末后的时期，有些人要背弃信德，听从欺诈的神灵和魔鬼的道理}，这些人歪曲真理。因此，既然我们的主和救主耶稣基督亲自命令了这事，又藉宗徒赐给我们关于这些人的警告，我们亲自听见他们的不敬，便因此弃绝了他们，正如我们先前所说，并宣布他们与天主教会和信德隔绝。此外，我们已将此告知你们的敬爱，亲爱的、最可敬的同工司铎，为使你们既不接待他们中的任何人（如果他们胆敢到你们那里），也不被诱使信任欧瑟伯或任何可能写信给你们的其他人。因为作为基督徒，我们有责任远离所有说话或心里反对基督的人，如同远离那些反抗天主和毁灭灵魂的人；我们甚至不应\Quote{问候这样的人}，正如圣若望所禁止的，\Quote{免得我们有时也分沾他们的罪过}。请问候与你们在一起的弟兄；与我在一起的弟兄也问候你们。\par

当\Person{亚历山大}向各城的主教们这样写信时，邪恶反而加剧，因为接受这消息的人被激发起争论。有些人完全赞同并签署了信中所表达的观点，而另一些人则相反。但尼苛美第亚的\Person{欧瑟伯}比其他人更被争论所激动，因为\Person{亚历山大}在信中对他进行了个人且苛责的暗示。此时，\Person{欧瑟伯}拥有很大影响力，因为皇帝居住在尼苛美第亚。事实上，\Person{戴克里先}不久前在那里建造了一座宫殿。因此，许多主教都向\Person{欧瑟伯}讨好。他多次写信给\Person{亚历山大}，要求他放弃已激起的讨论，重新接纳\Person{亚略}及其追随者进入共融；也写信给各城的主教，要求他们不要赞同\Person{亚历山大}的行动。这些导致到处混乱：不仅看到教会领袖在争论，百姓也分裂了，有的支持一方，有的支持另一方。这件事发展到了如此可耻的地步，以至于基督教甚至在剧院中成为大众嘲弄的对象。在亚历山大里亚的人就教义的最高点激烈争论，并派遣代表团到各教区的主教那里；而反对派也制造了类似的混乱。\par

亚略派与梅利提乌斯派混杂在一起，后者不久前已与教会分离。但现在必须说明这些[梅利提乌斯派]是谁。\par

亚历山大里亚的\Person{伯多禄}主教在\Person{戴克里先}统治时殉道，某个\Person{梅利提乌斯}，埃及某城的主教，因多项指控，尤其在迫害期间背信并献祭，被罢免。此人被剥夺了尊严，却仍有许多追随者，成为至今在埃及以他名字命名的梅利提乌斯派异端的领袖。由于他没有合理的借口离开教会，他便假装自己受了冤枉，对\Person{伯多禄}加以诽谤性的责备。\Person{伯多禄}在迫害中殉道，于是\Person{梅利提乌斯}转而辱骂继任\Person{伯多禄}主教的\Person{阿基拉斯}，此后又辱骂\Person{阿基拉斯}的继任者\Person{亚历山大}。在他们中间处于这种状况时，关于\Person{亚略}的讨论兴起了；\Person{梅利提乌斯}及其追随者与\Person{亚略}联合，与他共谋反对主教。但那些认为\Person{亚略}的观点站不住脚的人，赞同\Person{亚历山大}对他的决定，并认为那些支持他观点的人被公正地定罪。同时，尼苛美第亚的\Person{欧瑟伯}及其同党，以及赞同\Person{亚略}观点的人，写信要求撤销对\Person{亚略}的绝罚判决，并重新接纳被排除的人进入教会，因为他们没有持守错误的教义。这样，对立双方的书信被寄到亚历山大里亚的主教那里；\Person{亚略}收集了对自己有利的信件，而\Person{亚历山大}则收集了不利的信件。这因此为现今流行的亚略派、欧诺米派以及以\Person{玛策多纽}之名得名的派别提供了似是而非的辩护机会；因为这些派别分别利用这些书信来辩护自己的异端。\par

% structure: p0021 source=h2 target=subsection reason=relative-heading-rank; base=h2

\subsection{第七章 君士坦丁皇帝因教会的纷扰而忧伤，派遣西班牙人\Person{何西乌}前往亚历山大里亚，劝勉主教与\Person{亚略}和好与合一。}

当皇帝得知这些混乱时，他非常忧伤，将此事视为个人的不幸，立即努力扑灭已经点燃的火焰，并派一位名叫\Person{何西乌}（西班牙科尔杜瓦主教）的可靠之人送信给\Person{亚历山大}和\Person{亚略}。皇帝非常爱此人，并对他极为尊重。在此处引入这封信的一部分是合适的，整封信记载在\Person{欧瑟伯}所著的\WorkTitle{君士坦丁传}中。\par

\Emph{维克托·君士坦丁·马克西穆斯·奥古斯都致\Person{亚历山大}与\Person{亚略}}。\par

我得知你们目前的争论起因如下。当你，\Person{亚历山大}，询问你的司铎们对于圣经中某个不可解释的段落（更确切地说，是一个不适合讨论的问题）各自有何看法；而你，\Person{亚略}，则鲁莽地表达了一种观点，这种观点要么根本不应被构思，要么即使在你心中出现，你也应当将其埋于沉默。这个争论在你们中间被激起后，共融被拒绝；最神圣的百姓分裂为两派，脱离了共同身体的和谐。因此，你们每个人，彼此顾及，要倾听你们同仆的公正劝诫。他提供什么建议呢？起初提出这样一个问题既不谨慎，对所提出的问题给予回复也不谨慎：因为任何法律都不要求探究这类主题，而是闲暇时的无聊废话引起它们。即使它们是为了锻炼我们的自然能力而存在，我们也应将它们限制在我们自己的思考中，而不应轻率地在公共集会上提出，也不应轻率地让每个人都听到。实际上，有多少人能够充分阐述或准确理解如此宏大深奥之事的意义呢？\par

即使有人认为能够圆满达成此事，他又能说服多少民众呢？或者，谁能与这种研究的精微之处较量而不致陷入错误？因此，我们在这样的主题上必须谨言慎行，以免因我们本性的软弱而无法解释所提出的主题，或者听众的迟钝理解使他们无法清楚领会所教导的内容：而无论哪一种失败，民众都必然陷入亵渎或分裂。因此，你们当中任何一个人出于疏忽的问题和不加考虑的答案，都应当彼此同样宽恕。你们所引发的分歧并不涉及律法中所包含的任何重要诫命；你们也没有在敬拜天主方面引入任何新异端；相反，你们在这些问题上持守同一个判断，即信经。此外，当你们为了微小甚至几乎毫不重要的事情固执地彼此争论时，你们不宜照管如此众多的天主子民，因为你们意见分歧：这不仅不合适，而且被认为全然不合法律。\par

为了用一个较低层次的例子提醒你们的职责，我可以说：你们深知，即使哲学家们自己也团结在同一学派之下。然而他们在理论的某些部分往往彼此不同：但即使他们在最高科学分支上存在分歧，为了维持其整体的合一，他们仍同意联合起来。如今，若他们尚且如此，那么你们这些已被立为至高天主仆人的人，在如此宗教的宣认上彼此同心同德，岂不更加合理？但让我们以更详尽和更深入的关注来审视已经说过的话。因你们之间关于言辞的无谓而虚妄的争执，使弟兄与弟兄对峙；藉着我们为了如此不重要且绝非必要之事而彼此争斗，使尊贵的共融被不圣洁的分歧所扰乱，这合理吗？这些争吵是庸俗的，更符合孩童般的轻率，而不适合司铎和明智之士的智慧。我们应当自发地回避魔鬼的诱惑。我们众人的伟大天主和救主已将共同的光明扩展到所有人。在他的眷顾之下，请允许我，他的仆人，使我的此番努力成功结束；藉着我的劝勉、服务与恳切的告诫，我能引导你们，他的子民，回归共融的合一。因为，如我所说，你们中间只有一个信德，一个关于宗教的感受，并且律法的诫命在一切部分都将所有人的心灵合一于同一目的，所以请不要让这个已在你们中间引起纷争的意见分歧，在任何情况下造成冲突与分裂，因为它并不影响整个律法的效力。现在，我说这些话，并非强迫你们在这非常微不足道的争论主题上完全一致，无论它是什么；因为共融的尊贵可以不受影响地保持，与所有人的同一交谊可以保留，即使你们之间在无关紧要的事上存在一些感受上的差异。因为，当然，我们并非在每一方面都渴望同样的事物；我们之中也没有不变的本性或判断标准。因此，关于天主的眷顾，愿有一个信德、一个感受，以及一个天主的盟约：但你们之间如此精细地进行的那些细微研究，即使你们在它们上不达成一致看法，也应保留在你们自己的反思范围内，隐藏在内心的隐秘角落。那么，让一种不可言说且精选的普遍友谊的纽带，伴随着对真理的信德、对天主的敬畏，以及对他律法的虔诚遵守，在你们中间坚定不移。恢复相互的友谊与恩宠；向全体子民恢复他们惯常的亲切拥抱；而你们自己，在洗净了自己灵魂的力量上，再次彼此相认。因为友谊在敌意消除后往往变得更加甘甜。这样，请归还我平静的日子和无忧的夜晚；使我也能在纯净的光明中保有一些快乐，在我余生的岁月中拥有欢愉的宁静：否则，我必然叹息，全然被泪水浸透；我尘世生命的剩余时光也无法平静地维持。因为当天主的子民（我指的是我的同仆们）被如此不配且有害的争斗彼此分离时，我如何能保持我惯常的平静呢？但为了你们能对我因这不幸分歧而过度悲伤的程度有所了解，请听我即将陈述的话。近日我抵达\Place{尼科美底亚}城时，原本打算立即前往东方；但我正急速向你们而来，并在途中前进了一段相当距离后，此事的消息完全改变了我的意图，以免我被迫亲眼目睹那种我几乎无法忍受听闻的状况。因此，请通过你们从此处的和解，为我打开你们因彼此争斗而阻塞的通往东方之路：并请允许我尽快看见你们和全体其余的子民一同欢喜；并向天主表达我应得的感谢，因为各方的普遍和谐与自由，伴随着你们诚挚的赞美之声。\par

% structure: p0027 source=h2 target=subsection reason=relative-heading-rank; base=h2

\subsection{第八章 论在\Place{比提尼亚}的\Place{尼西亚}举行的会议，以及在那里所颁布的信经。}

皇帝的书信包含着如此明智而卓越的劝告。然而，邪恶的势力已变得过于强大，既非皇帝的劝勉所能战胜，也非携带书信者的权威所能遏制：因为无论是亚历山大还是阿略，都没有被这请求软化；此外，人民之间不断发生争执和骚乱。而且，那里先前已存在另一个地方性的不安定因素，困扰着教会——即关于逾越节的争议，这仅在东方的地区进行。这争议源于一些人希望更符合犹太人的习俗来庆祝该庆节；而另一些人则更喜欢普世基督徒普遍采用的庆祝方式。然而，这种差异并未妨碍他们的共融，尽管他们共同的喜乐必然受到了阻碍。因此，当皇帝看到教会因这两个原因而动荡不安时，他召集了一次大公会议，致信所有主教，请他们到比提尼亚的尼西亞与他相会。于是，主教们从各省各城聚集而来；关于他们，尤西比乌·潘菲卢斯在他的 \WorkTitle{君士坦丁传} 第三卷中逐字写道：\par

\Quote{因此，在遍布欧罗巴、阿非利加和亚细亚的所有教会中，最杰出的天主之仆们被召集起来。同一座神圣的建筑，仿佛是蒙天主扩展，在同一时刻容纳了叙利亚人、基利家人、腓尼基人、阿拉伯人和巴勒斯坦人，此外还有埃及人、底比斯人、利比亚人和来自美索不达米亚的人。在这次会议上，还有一位波斯主教出席，斯基泰人也没有缺席这场集会。蓬托斯、加拉提亚、旁非利亚、卡帕多细亚、亚细亚和弗里吉亚也提供了他们中最杰出的人。此外，还有色雷斯人、马其顿人、阿哈伊亚人和伊庇鲁斯人，甚至还有比这些更远地方的人，而最著名的西班牙人自己也坐在众人之中。帝国京城的首席主教因年迈缺席，但他的一些司铎在场，代他出席。这样的冠冕，作为和平的纽带，只有君士坦丁皇帝曾献给他的救主基督，作为胜敌后堪配天主的感恩祭，按照宗徒集会的样式为我们召集了这次会议。\BibleRef{宗 2:5} 因为据说在他们中间聚集了\BibleQuote{天下每一国家下的虔诚人，帕提亚人、玛待人、厄蓝人，以及住在美索不达米亚、犹太和卡帕多细亚、蓬托斯和亚细亚、弗里吉亚和旁非利亚、埃及和靠近基勒乃的利比亚部分的人，还有来自罗马的侨居者，包括犹太人和归依者，以及克里特人和阿拉伯人。}然而，那次集会在这一点上不如这次：所有在场者并非都是天主的仆人；而这次会议中，主教的人数超过三百位；而陪同他们的司铎、执事、辅礼者及其他人的数目几乎无法计数。这些天主的仆人中，有些因其智慧而卓越，有些因其生活的严谨和（遭受迫害时的）坚忍而卓越，还有一些人则将所有这些杰出品质集于一身；有些人因年迈而可敬，有些人因年轻和思想活力而引人注目，还有些人刚刚开始他们的牧职生涯。皇帝为他们所有人提供了丰富的日常食物供应。}\par

这就是尤西比乌关于这次与会者的记述。皇帝在完成了战胜李锡尼的庆节典礼后，也亲自来到了尼西亞。\par

主教中有两位特别著名：上底比斯的主教帕夫努提乌，以及塞浦路斯的主教斯皮里东。我为何特别提及这两位，将在以后说明。许多平信徒也出席了，他们精于推理艺术，每个人都急于为自己的派别辩护。如前所述，尼科米底亚的主教尤西比乌支持阿略的意见，与特奥格尼斯和玛里斯一起；前者是尼西亞的主教，玛里斯是比提尼亚的卡尔西顿的主教。他们受到亚历山大教会执事阿塔纳修的强力反对，阿塔纳修深受其主教亚历山大的尊重，并因此受到许多嫉妒，如以后将见到的。在主教们全体集会前不久，辩论者在群众面前进行了预备性的逻辑竞赛；当许多人因他们的话语兴趣而被吸引时，一位平信徒，\Emph{一位宣信者}，他具有单纯的理解力，斥责了这些推理者，告诉他们基督和他的宗徒们并没有教导我们\OriginalTerm{辩证法}{dialectics}、技巧或虚妄的\OriginalTerm{巧辩}{subtleties}，而是教导我们单纯的心志，这心志因信德和善行而得以保存。他说这话时，所有在场者都钦佩这位发言者，并赞同他言论的公正；辩论者们本人听了他对真理的简明陈述后，也表现出了更大的节制：于是，由这些逻辑辩论引起的骚动在当时被平息了。\par

第二天，所有主教聚集在同一地点；皇帝很快到达，一进门便站在他们中间，直到主教们鞠躬示意他入座，他才肯就座：这就是皇帝对这些人的尊敬和敬畏。当符合场合的静默得到遵守后，皇帝从他的座位上开始向他们发表劝勉和睦与团结的讲话，并恳求每个人放下个人的怨恨。因为前一天，其中几个人已互相控告，许多人甚至向皇帝递交了请愿书。但他将他们的注意力引向当前的问题以及他们聚集的目的，命令将这些请愿书烧掉，只指出：\Quote{基督吩咐那些渴望获得宽恕的人，要宽恕他的兄弟。}因此，当他极力强调保持和睦与和平之后，他再次批准了他们更深入探讨争议问题的意图。但最好听听尤西比乌在他的 \WorkTitle{君士坦丁传} 第三卷中关于此事的记述。他的话如下：\par

\Quote{各方提出了各种议题，从一开始就引起了诸多争论，皇帝耐心倾听了一切，慎重而公正地考虑所提出的一切。他部分支持双方所作的陈述，逐渐缓和了那些激烈对立者的尖刻态度，以他的温和与和蔼调解了每一方。由于他用希腊语向他们讲话（他对希腊语并非不熟悉），他既引人入胜又具说服力，使一些人心中信服，又用恳求说服了另一些人，他称赞那些发言良好的人。他激励所有人达到一致，最终成功地使他们在所有争议点上达成了相似的判断和一致的意见：不仅在信仰宣认上有了统一，而且在庆祝救恩节期的时间上也有了普遍的共识。此外，这些获得共同同意的教义由每个人签字确认。}\par

以上便是\Person{尤西比乌斯}以他本人的话在著述中留给我们的关于这些事的见证；我们并未不恰当地使用了它，而是将他所说的话视为权威，为了这部历史的忠实而将其引入此处。这样做也是出于以下目的：如果有人谴责在\Place{尼西亚}大公会议上所宣认的信仰为错误，我们可以不受其影响，也不信赖马其顿派的\Person{萨比努斯}，因为他称所有与会者为无知者和头脑简单者。这位\Person{萨比努斯}是\Place{色雷斯}地区\Place{赫拉克利亚}的马其顿派主教，他收集了各主教大公会议颁布的法令，尤其对\Place{尼西亚}会议的参与者加以轻蔑和嘲笑；他没有意识到，他这样做的同时也指责了\Person{尤西比乌斯}本人的无知，因为\Person{尤西比乌斯}经过最仔细的审查后作出了同样的宣认。事实上，他故意忽略了一些事情，歪曲了另一些事情，并且对所有事情都作出了有利于他自己观点的解释。然而，他称赞\Person{尤西比乌斯·庞非鲁斯}为可信赖的见证人，并赞扬皇帝在阐述基督教教义方面的能力；但他仍然指责\Place{尼西亚}所宣认的信仰是由无知者、对此事毫无见识的人所制定的。这样，他就自愿轻视了一个被他本人称为智慧而真实见证者的话：因为\Person{尤西比乌斯}声明，出席\Place{尼西亚}大公会议的天主仆人中，有些人以智慧之言著称，有些人以生活严谨著称；皇帝本人也在场，引导众人达到一致，在他们中间建立了统一的判断和一致的意见。至于\Person{萨比努斯}，我们将在适当的时候进一步提及。但\Place{尼西亚}大公会议中以高声欢呼所同意的信仰一致声明如下：\par

\Quote{我们信唯一的天主，全能的圣父，一切有形无形之物的创造者；——并信唯一的主耶稣基督，天主子，圣父的独生子，即出自圣父的本体；出自天主的天主，出自光的光；出自真天主的真天主；受生而非受造，与圣父同体；万物借他而造，无论在天上还是在地上；他为了我们人类，并为了我们的救恩，降下，成为血肉，并成为人；受难，第三日复活，升了天，并将再次来临审判生者死者。［我们也］信圣神。但圣而公的宗徒教会诅咒那些说\InnerQuote{曾有一个时期他不存在}、\InnerQuote{他未受生之前并不存在}和\InnerQuote{他是从不存在者中受造的}的人，以及那些断言他具有与圣父不同的本体或本质、或者他被造、或者他是可改变的人。}\par

这一信经得到了三百一十八位［主教］的承认和默许；并且，正如\Person{尤西比乌斯}所说，在表达和情感上一致，他们签署了它。只有五人不愿接受，他们反对\OriginalTerm{同体}{homoousios}或\OriginalTerm{同体}{consubstantial}这一术语：他们是\Place{尼科米底亚}主教\Person{尤西比乌斯}、\Place{尼西亚}的\Person{狄奥格尼斯}、\Place{迦克墩}的\Person{马里斯}、\Place{马尔马里卡}的\Person{塞奥纳斯}、\Place{托勒迈伊}的\Person{塞昆杜斯}。他们说：\InnerQuote{因为同体之物或是由于分有、衍生或萌发而来自另一者；萌发如同枝条从根生出；衍生如同子女从父母生出；分有如同两三个金器从一团金分出；而圣子从圣父而出皆非此类：因此他们宣布不能同意这一信经。}他们就这样嘲笑\OriginalTerm{同体}{consubstantial}一词，拒绝签署对\Person{阿略}的绝罚。于是，大公会议绝罚了\Person{阿略}以及所有坚持他意见的人，同时禁止他进入\Place{亚历山大城}。同时，皇帝的谕令将\Person{阿略}本人连同\Person{尤西比乌斯}、\Person{狄奥格尼斯}及其追随者流放；然而，\Person{尤西比乌斯}和\Person{狄奥格尼斯}在被放逐后不久，提交了一份书面声明，表示他们改变了观点，同意圣子与圣父同体的信仰，正如我们将在下文所述。\par

在此期间，当大公会议进行时，别号为庞非鲁斯的\Person{尤西比乌斯}，即\Place{巴勒斯坦}的\Place{凯撒勒雅}的主教，曾短暂地回避，经过深思熟虑是否应当接受这一信仰定义，最终同意了并与其他所有人一起签署了它：他还将信经的副本连同对\OriginalTerm{同体}{homoousios}一词的解释送给他辖下的子民，以免有人因他先前的犹豫而质疑他的动机。\Person{尤西比乌斯}所写的原文如下，他本人的话：\par

\Quote{亲爱的诸位，你们大概已听说为教会信仰而在尼西亚召开的大公会议所议决之事，因为传闻往往先于实情。但唯恐你们仅凭传闻而对此事产生不当的判断，我们认为有必要先向你们提交一份我们书面提出的信仰声明，随后再提交第二份，即经我们同意并作了若干措辞补充而颁布的声明。我们所提出的信仰声明，在我们最虔诚的皇帝面前宣读时，似乎获得普遍赞同，其内容如下：}\par

\Quote{我们既从前任主教们领受，无论是在真理知识的教育中，还是在我们受洗时，又正如我们亲自从圣经所学；并且我们无论在担任司铎职分时，还是在主教任内，都如此相信并教导。现今我们仍如此相信，并向你们明确宣认我们的信仰。内容如下：}\par

\Quote{我们信唯一的天主，全能的圣父，造成有形无形万物的主。并信唯一的主耶稣基督，天主的圣言、出自天主的天主、出自光明的光明、出自生命的生命、独生子，在一切受造之前所生，由天主圣父在万世之前所生，万物藉祂而造成；祂为了我们的救恩而成为肉身，居住人间；祂受难，第三日复活，升到父那里，并将在荣耀中再来审判生者死者。我们也信唯一圣神。我们相信每一位的存在和自立：圣父真是圣父，圣子真是圣子，圣神真是圣神；正如我们的主派遣门徒去传福音时所说：\BibleQuote{你们去教导万民，因父、及子、及圣神之名给他们授洗。}关于这些教义，我们坚定不移地持守其真理，并申明我们完全确信；我们迄今如此，而我们将至死不渝地持守这信仰，并在此信仰的坚定持守中，诅咒一切不虔敬的异端。在天主全能者和我们的主耶稣基督面前，我们证明：自从我们有了正确的自我认识以来，我们便从心底和灵魂相信并思想如此；现今我们所思所说的完全符合真理。我们更准备用无可辩驳的证据向你们证明并使你们信服：我们至今如此相信，如此宣讲。}\par

\Quote{当这些信条提出时，似乎没有反对的理由；不，我们最虔诚的皇帝亲自首先承认它们完全正确，并且他本人也持守这些条文中的思想；他劝勉所有在场的人赞同并签署这些信条，从而一致宣认它们，但需插入\Quote{\Emph{homoousios}}（\OriginalTerm{同质}{homoousios}）一词，皇帝亲自解释，该词并不表示肉体上的属性或性质；因此圣子并非由圣父分裂或切断而存有：他说，那非物质和非形体的本性不可能受任何肉体性质的影响；因此我们对这类事物的理解只能以神圣而奥秘的方式表达。这就是我们最有智慧、最虔诚的君王对此问题的哲学见解；主教们因\Emph{homoousios}一词而拟定这一信仰公式。}\par

\Emph{信经}\par

\Quote{我们信唯一的天主，全能的圣父，造成有形无形万物的主；并信唯一的主耶稣基督，天主的圣子，圣父所生的独生子，即出自圣父的本体；出自天主的天主，出自光明的光明，出自真天主的真天主；受生而非受造，与圣父同质；万物藉祂而造成，无论在天上或在地上；祂为了我们人类并为了我们的救恩，降下、成为肉身、成为人，受难并在第三日复活；祂升天，并将再来审判生者死者。〔我们〕也信圣神。但那些说\InnerQuote{曾有祂不存在的时期}，或\InnerQuote{祂在受生之前不存在}，或\InnerQuote{祂由无中造成}，或主张\InnerQuote{祂与圣父不同本体或本质}，或说天主子是受造的、可变的、或易改变的，天主教会和宗徒的教会都予以诅咒。}\par

当他们提出这项信仰宣认时，我们并未忽略探究这些表述的明确含义，即\Quote{出自圣父的本体，并与圣父同体}。于是，问题被提出并得到解答，这些术语的含义得到了清晰的定义；当时普遍承认，\OriginalTerm{本质（实体）}{ousias}（即本质或实体）仅仅意味着圣子确实出自圣父，但并非作为圣父的一部分而存在。对于这种解释神圣教义的说法——即圣子出自圣父，但并非圣父本体的一部分——我们认为是正确的，予以同意。因此我们自己也赞同这一阐释；我们也不挑剔\Quote{\OriginalTerm{同质}{homoousios}}这个词，出于对和平的考虑，并担心失去对该问题的正确理解。基于同样的理由，我们也接受了\Quote{受生而非受造}这个表述：\Quote{因为\Emph{受造},}他们说，\Quote{是一个普遍适用于所有由圣子所造的受造物的术语，而圣子与它们毫无相似之处。因此，他不是像那些由他所造的受造物一样的受造物，而是具有一种远超过任何受造物的本体；神圣训谕教导说，这一本体是由圣父所生，其生的方式无法被任何受造物解释甚或理解。}同样，关于\Quote{圣子与圣父同体}的声明在讨论后，大家一致同意，这不应被理解为物质意义上的，或以任何类似于会死之受造物的方式来理解；因为它既不是本体分割，也不是切割，也不是圣父的本体或能力有任何变化，因为圣父的非衍生本性与此一切不符。那么，圣子与圣父同体仅仅意味着，天主子与受造物毫无相似之处，而是在各方面只像生他的圣父；并且他不是出于任何其他本体或本质，而是出于圣父。对于以这种方式解释的教义，我们认为同意是正确的，尤其因为我们知道，在古代的一些杰出主教和博学作者已经在关于圣父和圣子本性的神学论述中使用了\Quote{\OriginalTerm{同质}{homoousios}}一词。这就是我要向你们陈述的，关于已经颁布的信德条文；我们都一致同意，并非未经适当考察，而是按照所指定的含义——这些含义是在我们至为敬爱的皇帝面前经过研究，并因上述理由而获得批准的。我们也认为他们在信仰宣认之后所宣布的绝罚并无冒犯之处；因为它禁止使用不合法的用语，而教会几乎所有的纷扰和骚动都源于此。因此，既然没有任何天主默感的圣经包含\Quote{从非存在的事物}和\Quote{曾有一个时期他是不存在的}这样的表述，以及其中所附的其他短语，我们认为说出并教导它们是不合理的；而且，这一决定获得了我们的认可，更是因为考虑到我们过去从未习惯使用这些术语。亲爱的，我们认为有责任让你们知道，我们在考察并同意这些事时所表现出的谨慎；以及我们基于正当理由，在最后时刻仍抵制引入某些有异议的表述，直到它们不被接受为止；而当我们在深思熟虑后审视这些词语的含义，发现它们与我们最初提出的正确信仰宣认相符时，我们便毫无争议地接受了它们。\par

这样的信是\Person{尤西比乌斯·潘菲鲁斯}写给在\Place{巴勒斯坦}的\Place{凯撒勒雅}的基督徒的。同时，大公会议本身也一致同意，写了以下的信，致\Place{亚历山太}教会及在\Place{埃及}、\Place{利比亚}和\Place{彭塔波利斯}的信徒。\par

% structure: p0046 source=h2 target=subsection reason=relative-heading-rank; base=h2

\subsection{第九章 大公会议的函件，论及其决定：谴责\Person{亚略}及赞同他的人。}

致蒙天主恩宠的圣而伟大的\Place{亚历山太}教会，及遍布\Place{埃及}、\Place{利比亚}和\Place{彭塔波利斯}的亲爱的弟兄们：在\Place{尼西亚}集会的主教们，组成伟大而神圣的大公会议，在主内问安。\par

既然，因天主的恩宠，一个伟大而神圣的大公会议已在\Place{尼西亚}召开，我们至为虔诚的君主\Person{君士坦丁}为此目的从各城各省召集了我们，我们觉得极有必要以神圣大公会议的名义写信给你们，好使你们知道哪些问题被提出并审查，以及最终决定并颁布了什么。\par

首先，我们最虔诚的皇帝\Person{君士坦丁}亲临审问了\Person{亚略}及其追随者的不敬与罪过；一致决定，应将他那不敬的主张以及他所发出的所有亵渎言辞——例如说\Quote{天主子出自无物}、\Quote{曾有一段时间他尚未存在}，以及说\Quote{天主子因拥有自由意志，既能行恶也能行善；并称他为受造物和作品}——统统处以绝罚。神圣主教会议对这些不敬的主张，或更确切地说，对这种疯狂和亵渎的言辞，几乎无法忍受听闻，便予以绝罚。至于我们针对他的诉讼结果，你们或许已经得知，或即将知晓；因为我们不愿践踏一个已因其罪行而受到惩罚的人。然而，他那有害的谬误竟如此具有传染性，以至于拖垮了\Place{马尔马里卡}的主教\Person{特奥纳斯}和\Place{托勒迈伊斯}的\Person{塞孔都斯}；他们遭受了与他相同的判决。当天主的恩宠将我们从那些可憎的教义及其全部不敬与亵渎中，以及从那些胆敢在原本和睦的人民中制造不和与分裂的人中解救出来后，仍有\Person{默利提乌斯}及其所祝圣之人的顽抗需要处理。现在，我们亲爱的弟兄们，我们向你们报告主教会议就此作出的决议。虽然严格来说\Person{默利提乌斯}完全不值得恩待，但主教会议出于极大的仁慈，决定他应留在自己的城内，但不得行使任何祝圣或提名祝圣的权力；不得以此借口出现在其他地区或城市，而仅保留名义上的尊位。那些由他任命的人，在经更合法的祝圣确认后，可按以下条件获准共融：他们应继续保有职衔和职务，但须在各方面自视为低于我们最尊敬的同事\Person{亚历山大}在各地区和各教会所祝圣与设立的人，以便未经\Person{亚历山大}下属的某位天主教会主教的同意，不得提名或推荐任何人，或做任何事。另一方面，那些因天主的恩宠和你们的祈祷而未曾陷入分裂，并一直忠信于天主教会的人，应有权提名和祝圣那些堪当圣职的人，并按教会法规与惯例行事。若任何任职于教会者去世，则这些新近获接纳的人可升任亡者的职位，条件是应显为堪当，且人民选举他们，并由\Person{亚历山大}主教批准其选择。这一特权确实授予所有其他人，但尤其不授予\Person{默利提乌斯}本人，因为他先前行为不端，且性格轻率鲁莽，以免授予他可能再次制造类似骚乱的权力或管治权。这些是特别涉及\Place{埃及}和最神圣的\Place{亚历山大}教会的事宜；若还有其他法规或条例确立，我们的主和最尊敬的同事及兄弟\Person{亚历山大}，因他参与并主导了所有事务，在返回你们那里后，会向你们提供更详细的说明。我们还有令人欣慰的消息告诉你们，关于最神圣的复活节的统一日期：因为这一问题也因你们的祈祷而得到圆满解决；以至于东方所有此前与犹太人一同庆祝此节日的弟兄，今后将遵行罗马人和我们，以及所有自始便按我们的时期庆祝复活节的人的做法。因此，为这些决议以及普遍的合一与和平，也为所有异端的铲除而欢欣，请以更大的尊敬和更丰盛的爱接待我们的同事及你们的主教\Person{亚历山大}，他的到来使我们大为喜悦，即使在年迈之际，他也为在你们中间重建和平而付出非凡的努力。请为我们众人祈祷，愿所决定的公正之事藉全能的天主、我们的主耶稣基督及圣神得以永保无虞；愿光荣归于祂，直到永远。阿们。\par

这份主教会议的书信表明，他们不仅绝罚了\Person{亚略}及其追随者，也绝罚了其主张的表述本身；并且他们在复活节的庆祝上达成了一致，重新接纳了异端首领\Person{默利提乌斯}进入共融，允许他保留主教的名位，但剥夺了他行使主教职权的所有权力。我想，正因如此，直至今日\Place{埃及}的\Person{默利提乌斯}派仍与教会分离，因为主教会议剥夺了\Person{默利提乌斯}的一切权力。还应注意，\Person{亚略}曾写下一篇阐述其主张的论文，题为\WorkTitle{塔利亚}；但其书性质松散放荡，风格与韵律类似\Person{索塔得斯}的歌曲。主教会议当时也谴责了这一作品。不仅主教会议不辞辛劳致函各教会宣告和平的恢复，而且皇帝\Person{君士坦丁}本人也亲自致函\Place{亚历山大}教会，发出以下训词。\par

\Emph{皇帝的信函。}\par

君士坦丁奥古斯都，致亚历山大城的天主教会。亲爱的弟兄们，你们好！我们已从天主眷顾中领受了不可估量的祝福，得以摆脱一切错谬，并团结于承认同一信德之中。魔鬼将不再能对我们施展任何权势，因为他奸诈图谋毁灭我们的一切，已从根本上被彻底推翻。天主的命令之下，真理的光辉驱散了那些纷争、分裂、骚乱，以及可谓不和的致命毒药。因此，我们全体敬拜唯一真天主，并相信他的存在。但为达成此事，我受天主的劝告，在尼西亚城召集了多数主教；我与他们一同，我自己也是你们中的一员，且非常欣喜成为你们的同仆，进行了对真理的探究。因此，所有因模棱两可而似乎为分歧提供借口之点，均经讨论并精确审查。愿天主至尊宽恕那些人对伟大救主、我们的生命与希望所无耻妄言的可怕亵渎之严重性；他们声称并认信与天主所启示的圣经相违之事。而超过三百位以节制和思维敏锐著称的主教，一致确认同一信德，这信德按照天主的真理和合法诠释，只能是\Emph{那}信德；唯独亚略被魔鬼的诡计所蒙蔽，被发现在你们中间首先传播这祸害，随后又以不圣洁的目的传播给其他人。因此，让我们拥抱全能者向我们提出的道理；让我们回到我们亲爱的弟兄们那里，一个不敬的魔鬼仆役曾使我们与他们分离；让我们赶紧回到共同的肢体和我们自身的骨肉。因为这与你们的明智、信德和圣德相宜；既然错谬已证明归咎于那真理之敌，你们应当重返天主之爱。因为那为三百位主教所判断认可的，只能是天主的道理；既然住在那许多尊贵者心中的圣神，已有效地光照他们明了天主圣意。因此，愿无人踌躇或徘徊，而愿所有人欣然回到无疑的职责之途；好使我到达你们中间时（这将尽快实现），能与你们一同向鉴察万物的天主献上应谢之恩，因为他揭示了纯正的信德，并将你们所祈求的爱德归还给你们。愿天主保护你们，亲爱的弟兄们。\par

皇帝就这样写信给亚历山大城的基督徒，向他们保证信德的阐释既非轻率也非随意，而是在大量研究、严格探究之后口授的；而且并非某些事被提及，另一些则被默然压下；而是凡能适当支持任何意见的一切都得以充分陈述。确实，没有任何事是仓促决定的，而是事先以极致的精确加以讨论；因此，凡似乎为含义模糊或意见分歧提供借口的每一点，都经彻底筛选，其疑难也得以消除。简而言之，他称所有与会者的思想为天主的思想，并且不怀疑如此众多杰出主教的一致是圣神所成就的。然而，马其顿异端首脑萨比诺却故意摒弃这些权威，称那些聚集者是无知无识的人；他甚至几乎指控凯撒勒雅的欧瑟比自己无知；他也不想想，即使组成那次会议的人都是平信徒，但他们既蒙天主光照并赐予圣神恩宠，就绝不可能背离真理。尽管如此，请继续听皇帝在另一道通谕中针对亚略及其支持者所颁布的法令；他将此通谕发送至各地的主教和人民。\par

\Emph{君士坦丁的另一封书信}\par

胜利者君士坦丁马克西穆斯奥古斯都，致主教及人民。——亚略既效法邪恶不敬之人，理当蒙受同样的耻辱。因此，正如那虔敬之敌波斐利，因撰写违背宗教的放肆论文而得到合适的报应，从而使他自此以后蒙受羞辱，为其不敬的著作遭受当得的谴责，他的著作也被销毁；同样，如今似乎适宜，使亚略及持其见解者被称为波斐利派，使他们因效法其行为而取其名号。此外，若发现亚略所著的任何论文，应将其付诸火焰，以使其邪恶的道理不仅被压制，而且他的任何纪念也丝毫不得存留。因此我如此规定：倘若有人被发现隐藏亚略所编之书，而不立即呈出焚毁，则此罪的刑罚为死刑；因为定罪之后，罪犯须立即受死。愿天主保佑你们！\par

\Emph{另一封书信}\par

君士坦丁奥古斯都，致各教会。\par

我从公共事务的兴旺状况中体验到了天主能力的恩宠何其伟大，因此判断首要关切之事莫过于在至可称颂的天主教会集会中，以真诚的爱和一致的虔诚维护对全能天主的同一信德。然而我认识到，除非全体或至少大部分主教能聚集在同一地点，并由他们以会议形式讨论我们最神圣宗教的每一点，否则这无法稳固持久地建立；因此，我尽最大可能召集了他们，我自己也如同你们中的一员出席了；因为我不会否认我特别引以为荣的事，即我是你们的同仆。随后所有细节得到详尽审查，直到一个为鉴察万物的那一位所悦纳的决定被颁布，以促进判断与实践的一致；使得今后在信德问题上再不留下任何分裂或争论。那里也审议了关于至圣的复活节日期的议题，大家一致同意各地在同一日庆祝它应该是合宜的。因为还有什么比这更合宜、更隆重呢？我们从此庆节中获得了永生的希望，理应一致遵循同一秩序来庆祝，其理由是显而易见的。首先，我们跟随犹太人的习俗来庆祝这至圣的庆节，似乎极不相称；这个民族双手沾满了最可憎的暴行，由此玷污了他们的灵魂，理应被蒙蔽。因此，我们撇开他们的做法，可以自由地确保今后此庆节的庆祝按照从受难的第一天直到现在所持守的更正确的秩序来进行。所以，不要与那最敌对的人民——犹太人——有任何共同之处。我们从救主那里领受了另一条路；因为在我们神圣的宗教中，为我们设定了一条既合法又准确的道路：让我们一致追求它，最可敬的弟兄们，使自己远离那可憎的团体。因为若没有他们的教导，我们就不能正确遵守这些事，他们如此夸耀实在是荒谬。因为他们在杀害了他们的主之后，已丧失了理智，不是被任何理性动机，而是被无法控制的冲动所引导，随其天生的狂怒而行，他们又能在什么事上做出正确判断呢？因此，即使在此事上他们也不认识真理，以至于他们不断地犯下极大的错误，不进行适当的改正，反而在同一年中第二次庆祝逾越节。那么，我们为何要效法那些被公认为染有严重错误的人呢？的确，我们绝不应容忍复活节在同一年中庆祝两次！但即使这些考量没有摆在你们面前，你们的明智也总当留意，通过勤勉和祈祷，确保你们灵魂的纯洁绝不与那些如此彻底败坏之人的习俗有任何相通，或显得相通。此外也当考虑，在如此重要且具有如此宗教意义的事上，最轻微的分歧也是最不敬的。因为我们的救主只留下了一天让我们纪念我们的得救，即他的至圣受难日；他也愿意他的公教会是一个整体；其成员无论散居各地，都被同一圣神所呵护，即天主的旨意。与你们的神圣身份相称的明智，请思量这是何等严重和不体面：同一些日子，有人守斋，有人庆节；复活节之后，有人纵情宴乐，有人则按规定守斋。因此，天主眷顾指示进行适当的纠正，并确立一致的作法，我想你们都知道。\par

既然当时希望这应如此修正，使我们与那弑亲者及杀害其主的民族毫无共同之处；并且这礼仪是合适的，为西方、南方、北方所有教会以及东方一些教会所遵守；基于这些考虑，目前众人认为妥当，我也承诺，以你们睿智的洞察力，会满意地接受在罗马城、整个意大利、非洲、全埃及、西班牙、法兰西、不列颠、利比亚、全希腊以及亚洲、本都、西里西亚各教区如此一致遵守的做法；也考虑到，不仅前述地方的教会数量更多，而且这一点尤为神圣的义务：凡理性所要求且与犹太人的背誓无关之事，众人都应共同向往。简而言之，共同决定应将至圣的逾越节庆贺在同一日隆重举行；因为在这等神圣的庆典中，任何差异都不合宜；采纳那不含异端错误、不偏离正道的意见更为可取。因此，既然这些事如此一致，请你们欣然接受这天上且真正神圣的命令：因为在主教们的神圣集会中所行的一切，皆归于天主的旨意。所以，当你们向所有亲爱的弟兄们指明这些规定后，理当公布上述文字，接受所提出的理由，并确立这至圣之日的遵守：这样，当我终于见到长久热切渴望的你们的秩序时，便能与你们在同一日庆贺这神圣的节庆；并因万事与你们一同喜乐，看见撒殚的残暴因我们的努力而被天主的能力挫败，而你们的信德、平安与和谐处处兴盛。愿天主保佑你们，亲爱的弟兄。\par

\Emph{致犹西比乌的另一封信}\par

\Person{胜利者君士坦丁·马克西穆斯·奥古斯都}，致\Person{犹西比乌}。\par

既然一种不虔的企图和暴政直到如今还迫害我们救主天主的仆人，我确实听说并深信，最亲爱的弟兄，我们所有的圣堂或因疏忽而破败，或因迫近的危险之恐惧而未得到应有的尊严装饰。但现在自由已恢复，那迫害的巨龙李锡尼已因至高天主的眷顾并藉我们的工具被清除出公共事务的管理，我想天主的能力已向众人彰显，同时那些或因恐惧或因不信而陷入任何罪恶的人，既已承认永生的天主，必将归向真实正确的生命之道。因此，吩咐你所管辖的教会，以及各地管辖的其他主教，连同你所知道的司铎和执事，要勤勉于圣堂事务，或修缮那些尚存的，或扩建，或在任何需要的地方新建。你自己也请求，并让其余人藉你向各省总督及禁卫军长官的官员请求必要的供应：因为已命令他们尽心执行你圣善的命令。愿天主保佑你，亲爱的弟兄。\par

这些关于建造圣堂的指示是由皇帝发给各省主教的；但他关于准备若干份圣经抄本写给巴勒斯坦的犹西比乌的内容，我们可以从信件本身得知：\par

\Person{胜利者君士坦丁·马克西穆斯·奥古斯都}，致\Person{凯撒勒雅的犹西比乌}。\par

在与我同名的城中，藉我们救主天主的助佑眷顾，极多的人已归附至圣教会，以致那里大得增长。因此，需要在该处增设更多圣堂：所以请你们极其诚挚地接纳我所怀的意愿。我认为当以此告示你的明智，应命令由精通其艺的合格抄写员在精心准备的羊皮纸上抄写五十份圣经，既清晰可读，又便于携带，其提供和使用你知道对教会的训导是必需的。朕的仁慈也已致函该教区的财务官，令其注意为抄写提供一切所需。为使这些抄本尽快准备好，这任务当由你的勤勉承担；并授权你凭此信使用两辆公共马车运送；这样，最令人满意的抄本可便于送来供朕检视，由你教会中的一位执事完成此任务；当他到达朕处时，将蒙受朕的恩赐。愿天主保佑你，亲爱的弟兄。\par

\Emph{致马卡里的另一封信}\par

君士坦丁·维克多·马克西穆斯·奥古斯都致耶路撒冷的马卡里乌斯。——我等的救主如此施予恩宠，以致任何言辞都不足以表达祂现今彰显的恩赐。因为祂至圣苦难的纪念物，长久埋藏于地下，竟被隐藏了如此多年，直到借着消灭万民公敌之际，才在祂的仆人重获自由后闪耀出来，这超越了一切赞叹。即使全地所有被视作智慧的人聚集一处，欲以言语相称此事，也必连最小的部分都无法企及；因为这奇事的领悟远超一切人性推理的本性，恰如天界之事强于人间。因此，这常是我特别的目标：既然真理的可信性每日以新的奇迹彰显自身，我们众灵魂就应更加勤勉于神圣律法，怀着谦逊与一致的热忱。但我愿你完全知晓我所想的大致众所周知之事：如今我的首要关切，乃是以宏伟的建筑妆点那块神圣之地——我奉天主之命，已将那最可耻的偶像附加物如同沉重负担般从那里除去；那地自起初本在祂的计划中分别为圣，但自从祂显明了救主苦难的证据，便更为明显地圣化了。因此，你的明智应当做出安排，并预备一切所需，使那里不仅建起一座本身超越任何别处的圣堂，而且其余部分也当如此，使各城所有最辉煌的建筑都被此所超越。至于墙壁的工艺与纯洁的建造，你知道我们将这些事托付给了我的朋友、最尊贵的近卫军长官代理德拉西利安，以及行省总督：朕的虔诚已下令，工匠和工人，以及你的明智能告知为结构所需的其他一切，都当通过他们的关照立即送来。关于石柱或大理石，凡你判断为更珍贵与合用者，你亲自视察计划后，请写信告诉我们；以便我们从你的信中了解需要多少以及何种物品时，可以从各地运送给你：因为世上最奇妙之地，理应与它的尊荣相配。但我想知道你认为大殿的穹顶应做成格子形，还是用其他方式建造：因为若是格子形，还可以用金装饰。最后，你的圣德应尽快告知前述官员，你需要多少工人和工匠，以及多少费用。同时请速速向我报告，不仅关于大理石和石柱，也关于格子穹顶，如果你确实认为这样更美。愿天主保佑你，亲爱的兄弟。\par

皇帝还写了其他几封更具修辞色彩的信，反对阿里乌及其追随者，并命人在各城广泛发布，嘲笑他，以讽刺挖苦他。此外，他在写给尼科米底亚人的信中反对优西比乌和特奥格尼斯，谴责优西比乌的不当行为，不仅因其阿里乌主义，还因他从前对统治者本怀好意，却背信弃义地阴谋反对他的事务。他随后劝他们改选另一位主教取代他。但我认为在此抄录这些关于这些事的信是多余的，因为篇幅过长：有意者可在别处找到并阅读。这些事就此说明足矣。\par

% structure: p0069 source=h2 target=subsection reason=relative-heading-rank; base=h2

\subsection{第十章 皇帝也召请诺瓦西安派主教阿基西乌出席主教会议。}

皇帝的勤勉促使我提及另一件表明其心意的事，足以显示他多么渴望和平。因为追求教会和谐，他也召请了诺瓦西安派的主教阿基西乌前来参加会议。当信德宣言已被书写并由主教会议签字后，皇帝问阿基西乌是否也同意这份信经以及关于复活节日期的决定。他回答说：\Quote{我的君王，主教会议并没有规定什么新东西：因为从起初，甚至从宗徒时代，我就传统地接受了信德的定义和复活节庆祝的时间。}因此，当皇帝进一步问他：\Quote{那你为何与教会其余的共融分离？}他便讲述了德西乌斯迫害期间所发生的事，并引用了那道严厉的教规，该教规宣称：在领洗后又犯了神圣圣经称为\Quote{至于死的罪}\BibleRef{若 5:16}的人，不应被认为配得参与圣事；他们确实应被劝告悔改，但不可期待从司铎那里得到赦免，而应从天主那里，因为祂有能力并有权柄赦罪。阿基西乌这样说后，皇帝对他说：\Quote{架起梯子，阿基西乌，独自爬上天去吧。}优西比乌·庞菲鲁斯或其他任何人都未曾提及这些事；但我从一个绝不轻信谎言、非常年长、并在叙述中单纯陈述了会议中所发生之事的人那里听到的。由此我猜想，那些对此事保持沉默的人，是受到了驱使许多其他史学家的同样动机的影响：因为他们常常出于对某些人的偏见或对另一些人的偏袒而压制重要事实。\par

% structure: p0071 source=h2 target=subsection reason=relative-heading-rank; base=h2

\subsection{第十一章 论主教帕夫努提乌斯。}

正如我们先前许诺要提及\Person{帕弗努提乌斯}和\Person{斯彼里东}，现在正是该论及他们的时候。\Person{帕弗努提乌斯}当时是\Place{上泰贝}地区一座城市的主教：他深受天主眷顾，以至于借他之手行了非凡的奇迹。在教难期间，他被剥夺了一只眼睛。皇帝极其敬重此人，常召他入宫，并亲吻那被挖掉眼睛的部位。\Person{君士坦丁}皇帝就是这样虔诚。关于\Person{帕弗努提乌斯}，只举这一件事就够了：我现在要说明另一件事，这件事因他的建议而发生，既有利于教会，也有益于圣职人员的荣誉。主教们认为宜在教会中引入一条新法规，即那些处于圣秩中的人——我指的是主教、司铎和执事——不得与他们尚为平信徒时所娶的妻子有夫妻关系。当此事即将讨论时，\Person{帕弗努提乌斯}在主教会议中站起来，恳切地请求他们不要将如此沉重的枷锁加在神职人员的身上；他声称：\BibleQuote{婚姻本身是尊贵的，床榻也是无玷的}\BibleRef{希 13:4}。他在天主面前力劝他们不应以过于严苛的限制损害教会。他说：\Quote{因为并非所有人都能承受严格的节欲实践；恐怕各人的妻子的贞洁也无法保持}：他将男人与合法妻子的关系称为贞洁。他认为，按照教会的古老传统，那些先前已进入神圣职务的人应放弃婚姻，这已足够了；但不应该将那些尚未被祝圣时所结合的妻子与之分离。他表达这些意见时，虽然他自己没有婚姻经验，而且坦率地说，从未认识过女人：因为他从小就在隐修院长大，并以贞洁著称于众人。全体圣职人员都赞同\Person{帕弗努提乌斯}的理由：因此他们就此点停止了进一步讨论，将是否愿意对妻子实行禁欲留给身为丈夫者自行斟酌。关于\Person{帕弗努提乌斯}，就说到这里。\par

% structure: p0073 source=h2 target=subsection reason=relative-heading-rank; base=h2

\subsection{第十二章 论\Person{斯彼里东}，\Place{塞浦路斯}主教}

关于\Person{斯彼里东}，他在做牧羊人时圣德如此崇高，以至于被认为配得上做人群的牧者；并被指派为\Place{塞浦路斯}一座名为\Place{特里米图斯}的城市的主教，但因他极其谦卑，在担任主教期间仍继续放牧他的羊群。关于他有许多非凡事迹；我不过记录一两个罢了，免得偏离主题。有一次约在半夜，盗贼偷偷进入他的羊圈，企图偷走一些羊。但保护牧者的天主也保护了他的羊；因为盗贼被一种无形的力量捆绑在羊圈里。天亮时，他来到羊群那里，发现那些人双手被反绑在身后，便明白了所发生的事。他祈祷后释放了盗贼，恳切地告诫和劝勉他们要靠诚实的劳动维持生计，不可妄取不义之财。然后给了他们一只公羊，打发他们走，幽默地补充说：\Quote{免得你们看起来整夜白守了。}这是与\Person{斯彼里东}有关的一个奇迹。另一个奇迹是这样：他有一个童贞女儿名叫\Person{依勒内}，与她父亲的虔诚有份。一位熟人托付她保管一件价值不菲的饰物；她为了更安全地保存，将托付之物藏在地下，不久后便去世了。随后，物主来索取；找不到那童贞女，便与父亲激动地争执起来，时而指责他企图欺诈，时而又恳求他归还寄存物。老人将此人的损失视为自己的不幸，来到女儿的坟墓前，呼求天主在适当的季节之前向他显示所应许的复活。他的希望没有落空：因为童贞女再次复活显现给父亲，向他指出藏饰物的地点，然后又离开了。如此的人物在\Person{君士坦丁}皇帝的时代装饰了教会。这些细节我从许多\Place{塞浦路斯}居民那里获得。我也找到一篇由司铎\Person{鲁菲努斯}用拉丁文撰写的论文，我从其中收集了这些以及一些其他将在后面引述的事。\par

% structure: p0075 source=h2 target=subsection reason=relative-heading-rank; base=h2

\subsection{第十三章 论隐修士\Person{欧提基安}}

我又听说关于\Person{欧提基安}的事，他是一位虔诚的人，大约在同一时期活跃；他也属于诺瓦替安教会，但因行类似的奇迹而受尊敬。我将明确说明这段叙述的权威来源，即使冒犯某些人，我也不试图隐瞒。这来自\Person{奥克萨农}，诺瓦替安教会的一位年事极高的司铎；他年轻时曾陪同\Person{阿凯修斯}前往尼西亚会议，并向我讲述了关于他的事。他的寿命从那时一直延续到\Person{小狄奥多西}统治时期；在我很年轻时，他向我详述了\Person{欧提基安}的事迹，大大夸赞了彰显在他身上的天主的恩宠：但他提到一件发生在君士坦丁统治时期的事，特别值得一提。那些被称为内卫（或近卫）军的侍卫中有一人，因涉嫌叛逆而逃命。愤怒的君主下令，无论他在何处被发现，都必被处死；此人在\Place{比提尼亚的奥林波山}被逮捕，被沉重痛苦的锁链捆绑，关押在\Place{奥林波山}附近，那里正是\Person{欧提基安}过着独居生活、并治愈许多人的身体和灵魂之处。年迈的\Person{奥克萨农}当时非常年轻，与他在一起，并受他训练隐修生活。许多人来到这位\Person{欧提基安}面前，恳求他代为向皇帝求情以释放囚犯。因为\Person{欧提基安}所行奇迹的名声已传到皇帝耳中。他欣然答应前往君主那里；但锁链造成了难以忍受的痛苦，那些为他求情的人宣称，锁链导致的死亡将先于皇帝的报复和为囚犯所做的任何求情。因此\Person{欧提基安}派人请求狱卒释放那人；但他们回答说释放罪犯会使自己陷入危险，于是\Person{欧提基安}亲自带着\Person{奥克萨农}前往监狱；当他们拒绝开门时，\Person{欧提基安}身上所赐的恩宠更加显扬：因为监狱的大门自动打开，而钥匙还在狱卒手中。\Person{欧提基安}与\Person{奥克萨农}一进入监狱，令在场所有人大为惊奇的是，锁链自动从囚犯的四肢脱落。然后他与\Person{奥克萨农}前往昔日称为\Place{拜占庭}、后来称为\Place{君士坦丁堡}的城，进入皇宫后，他救了那人免于死亡；因为皇帝对\Person{欧提基安}非常敬重，欣然答应了请求。这确实发生在稍晚一些的时候（相对于本段历史所涉及的时期）。\par

在尼西亚会议中集会的诸位主教，制定并记录了某些他们通常称为教规的教会规章后，又各自返回自己的城。因我认为这将受到学问爱好者的欣赏，我在此附上与会者的名单，就我所查明的，包括他们各自主持的省份和城，以及此次会议举行的日期。\Person{奥西西乌斯}，我以为是西班牙的\Place{科尔多瓦}的主教，如前所述。\Person{维托}和\Person{文森提乌斯}，\Place{罗马}的司铎，\Person{亚历山大}，\Place{埃及}的主教，\Person{尤斯塔提乌斯}，\Place{大安提阿}的主教，\Person{玛卡里乌斯}，\Place{耶路撒冷}的主教，以及\Person{哈波克拉提昂}，\Place{奇诺波利斯}的主教：其余人的名字详载于\Place{亚历山大里亚}主教\Person{阿塔纳修斯}所著的\WorkTitle{会议汇编}。此会议是在\Person{保利努斯}和\Person{尤利安}执政之年、5月20日、从\Place{马其顿}的\Person{亚历山大}统治起算的第636年召开的（我们是从会议记录前注的日期发现的）。因此会议的工作就完成了。应注意，会议之后，皇帝前往帝国的西部。\par

% structure: p0078 source=h2 target=subsection reason=relative-heading-rank; base=h2

\subsection{第十四章 \Place{尼科米底亚}主教\Person{尤西比乌斯}与\Place{尼西亚}主教\Person{狄奥格尼斯}因与\Person{亚略}见解一致而被放逐，在发表悔改声明并赞同信经后，恢复其主教座。}

\Emph{\Person{尤西比乌斯}和\Person{狄奥格尼斯}向主要主教们发送了悔改声明后，奉皇帝谕旨从流放地被召回，并恢复了自己的教会，取代了那些被任命在他们位置上的人——\Person{尤西比乌斯}取代\Person{安菲翁}，\Person{狄奥格尼斯}取代\Person{克雷斯图斯}。\newline{}以下是他们书面撤回声明的副本：}\par

\Quote{我们既曾于某时被陛下定罪，未经正式审判，理应对陛下神圣判决保持缄默。然而，若我们因缄默而助长对我们的诬告者，实为不当，为此我们声明：我们完全赞同你们的信仰；并且，在仔细考量了\OriginalTerm{同质}{consubstantial}一词的含义后，我们始终致力于和平，从未追随异端。为了教会的安全，我们提出了一切所思所想，并充分安抚了受我们影响的人，我们签署了信仰宣言；但我们没有签署绝罚令；并非反对信经，而是不相信被指控的一方如所述那般，这点我们从他给我们的信件以及亲自交谈中已得满足。但若你们的神圣会议确信，我们不反对而赞同你们的决定，则通过此声明我们给予完全同意和确认：我们这样做并非因厌倦流放，而是为了摆脱异端的嫌疑。因此，若你们如今认为宜恢复我们到你们面前，你们将在一切事上找到我们顺从并同意你们的法令；尤其是，既然陛下仁慈地对待甚至召回主要被指控者，我们若缄默而自证其罪，实为荒谬，因为看似责任人已被允许澄清对他的指控。请依照你们那爱基督的虔诚，提请我们最虔诚的皇帝，呈递我们的请求，并速速以合宜的方式为我们决定。}\par

以上是尤西比乌和特奥格尼斯的悔改之词；由此我推断他们签署了所颁布的信仰条款，但不愿参与对亚略的谴责。看来亚略在他们之前已被召回；然而，尽管这可能是真的，他却被禁止进入亚历山大里亚。这从以下事实可见：他后来通过伪造悔改，为自己找到了返回教会和亚历山大里亚的途径，我们将在适当的地方说明。\par

% structure: p0082 source=h2 target=subsection reason=relative-heading-rank; base=h2

\subsection{第十五章  大公会议后，亚历山大去世，亚大纳削被立为亚历山大里亚主教。}

此后不久，亚历山大里亚主教亚历山大去世，亚大纳削被委派管理该教会。鲁菲努斯记载，此人在幼年时曾与同龄人在圣游戏中玩耍：这游戏模仿了司祭职和献身者的等级。在这游戏中，亚大纳削被指定为主教席位，其他男孩分别扮演司铎或执事。孩子们在纪念殉道者兼主教伯多禄的节日里进行这项游戏。当时亚历山大里亚主教亚历山大恰巧经过，看到了他们的游戏，便召来孩子们，询问他们在游戏中各自扮演的角色，认为此事或许预兆着什么。然后他吩咐将孩子们带到教堂学习功课，尤其是亚大纳削；后来在他成年后祝圣他为执事，并在大公会议召开时带他到尼西亚协助辩论。鲁菲努斯在自己的著作中记载了亚大纳削的这段经历；此事并非不可能发生，因为这类事情常有。关于此事，以上所述已足够。\par

% structure: p0084 source=h2 target=subsection reason=relative-heading-rank; base=h2

\subsection{第十六章  君士坦丁皇帝扩建古拜占庭，称其为君士坦丁堡。}

大公会议后，皇帝花了一些时间休整，并在公开庆祝登基二十周年后，立即致力于修缮教堂。他在其他城市以及以他命名的城市都实施了此事；该城原名拜占庭，他加以扩建，围以坚固城墙，并用各种建筑装饰；使其与帝国罗马相提并论，命名为\Emph{君士坦丁堡}，并以法律确定其应被称为\Emph{新罗马}。此法律刻在石柱上，立于斯特拉特吉翁广场皇帝骑马雕像附近。他还在同一城市建造了两座教堂，一座名为\Emph{和平堂}，另一座名为\Emph{宗徒堂}。他不仅改善了基督徒的事务，如上所述，还摧毁了异教徒的迷信；因为他将他们的神像公开展示以装饰君士坦丁堡城，并将德尔斐的三脚鼎公开展示于竞技场。现在提及这些事似乎多余，因为它们尚未听闻已先目睹。但那时基督徒的事业得到了最大发展；因为天主眷顾在君士坦丁皇帝时代保存了许多其他事物。尤西比乌·庞非录以华美的言辞称颂了皇帝；我认为尽可能简洁地提及这些并非不合时宜。\par

% structure: p0086 source=h2 target=subsection reason=relative-heading-rank; base=h2

\subsection{第十七章  皇帝母亲赫肋纳来到耶路撒冷，寻找并找到了基督的十字架，建造了一座教堂。}

皇帝的母亲赫肋纳（君士坦丁皇帝曾将一度是村庄的德勒帕农扩建为城市，并以她的名字命名为赫肋诺波里）蒙天主指引，由梦境前往耶路撒冷。她发现那曾是的耶路撒冷，依先知所言，荒凉\InnerQuote{像秋季果实的保护区}，便仔细寻找基督的坟墓，即祂埋葬后复活之处。经过许多困难，在天主的助佑下，她找到了它。困难的缘由，我稍作解释。那些在祂受难后接受基督信仰的人，极为尊敬这座坟墓；但仇恨基督信仰的人，用土丘覆盖了这地方，并在其上建造了一座维纳斯庙，立上她的像，不顾及这地方的纪念。此事持续了很长时间；后来为皇帝的母亲所知。于是她命人推倒雕像，移除泥土，彻底清理地面，便在坟墓中发现了三个十字架：其中一个是那有福的十字架，基督曾悬于其上；另外两个是与他同钉的两个盗贼死去时所悬的。与这些一同被发现的还有比拉多的牌子，其上他用多种文字写着：那被钉的基督是犹太人的君王。然而，由于无法确定哪一个是他们所寻找的十字架，皇帝的母亲极为忧愁；但不久耶路撒冷主教玛加略解除了她的烦恼。他以信德解决了疑问，因为他向天主求了一个征兆，且获得了。征兆是这样的：邻村有一妇人，久病缠身，现已濒临死亡；于是主教安排将每个十字架依次带到那垂死的妇人面前，相信她一触到那宝贵的十字架便会痊愈。他的期望没有落空：因为那两个不是主的十字架接触后，妇人仍处于垂死状态；但当第三个，即真十字架，触到她时，她立即痊愈，恢复了先前的力气。就这样，真正的十字架被找到了。皇帝的母亲在坟墓之上建造了一座宏伟的圣堂，命名为\Place{新耶路撒冷}，使其面向那座古老而荒凉的城市。她在那里留下了一部分十字架，封在一银匣中，作为纪念，供愿意观看的人瞻仰；另一部分她送给了皇帝；皇帝深信，凡保存此圣物之处，城市必得完全安全，便私下将它封入自己的雕像中，该雕像矗立在君士坦丁堡称为君士坦丁广场的斑岩大柱上。我固然是根据传闻记述；但几乎所有的君士坦丁堡居民都断言此事属实。此外，基督双手被钉在十字架上的钉子（因为他的母亲也在坟墓中找到了这些钉子并送给了他），君士坦丁拿去制成了马衔铁和头盔，用于他的军事征战中。皇帝为建造圣堂提供了所有材料，并致信主教玛加略，催促加速这些建筑。皇帝的母亲完成\Place{新耶路撒冷}后，又在白冷——基督按肉身诞生之洞窟——上建造了另一座毫不逊色的圣堂：她并未止步，又在祂升天的山上建造了第三座。她对此事如此虔诚，以至于她与妇女们一同祈祷；并邀请登记在教会名册上的贞女们赴宴，亲自服侍她们，将菜肴端上桌。她也对圣堂和穷人非常慷慨；度过虔诚的一生后，约八十岁时去世。她的遗体被运往首都新罗马，安葬在皇家陵墓中。\par

% structure: p0088 source=h2 target=subsection reason=relative-heading-rank; base=h2

\subsection{第十八章 君士坦丁皇帝废除异教并在各地建立许多圣堂}

此后，皇帝对基督徒的利益日益关注，并摒弃了外教的迷信。他废止了角斗士的搏斗，并在庙宇中竖立了自己的雕像。由于外教人断言是\Person{塞拉皮斯}使尼罗河上涨以灌溉埃及，因为通常会将一腕尺带入他的庙中，他命\Person{亚历山大}将这腕尺转移到教堂。尽管他们预言因\Person{塞拉皮斯}之怒尼罗河将不再泛滥，然而次年及之后的河水泛滥均如期发生：事实证明了尼罗河的上涨并非出于他们的迷信，而是出于天主的眷顾。大约同一时期，那些蛮族萨尔马特人和哥特人侵入了罗马领土；然而皇帝对教会的热忱并未因此减弱，而是对这两件事均作了适当安排。他信靠基督旗帜，彻底战胜了敌人，甚至取消了前代皇帝惯于向蛮族进贡的金税；而蛮族自身，因出乎意料的失败而惊恐，于是首次接受了基督信仰，\Person{君士坦丁}正是藉此信仰得以护佑。他再次建造了其他教堂，其中一座建于\Place{曼默勒}橡树附近，圣经宣称\Person{亚巴郎}曾在那里款待天使。皇帝得知那棵橡树下曾设立祭台，并举行外教祭献，便致信\Place{凯撒勒雅}的\Person{欧瑟比}主教加以谴责，并下令拆除祭台，在橡树旁建一座祈祷所。他还指示在\Place{腓尼基}的\Place{赫略波利斯}建造另一座教堂，理由如下。我无法说明最初为\Place{赫略波利斯}居民立法的人是谁，但他的品格和道德可从该城的习俗来判断；因为该地的法律规定妇女为公共所有，因此出生的孩子血统不明，以致无法区分父亲和子女。他们的处女也被提供给到访的陌生人行淫。皇帝急于纠正这种长期存在于他们中间的恶习。他颁布了严肃的贞洁法律，清除了这一可耻的恶行，并规定家庭之间的承认。他在那里建立了教堂，并委任了主教和圣职人员。就这样，他改革了\Place{赫略波利斯}人的败坏风俗。他还摧毁了\Place{黎巴嫩山}上\Place{阿法卡}的维纳斯神庙，废除了那里举行的恶行。何需我描述他从\Place{基里基雅}驱逐皮同魔的事迹，他下令将那魔鬼潜藏的府邸夷为平地？皇帝对基督信仰的热忱如此之大，以至于当他即将对\Place{波斯}开战时，他准备了一座以绣花细麻布制成的幕帐，仿照教堂的样式，正如\Person{梅瑟}在旷野中所作的那样；\EditorialNote{出谷纪 35-40}且如此构造以便于搬迁，以便即使在最荒凉的地区也能拥有一座祈祷所。但当时由于皇帝之威而使敌人恐惧，战争并未进行。我想在这里描述皇帝重建城市、将许多村庄变为城市的勤勉并不恰当；例如他以其母亲的名字命名的\Place{德雷帕努姆}，以及以其姊妹命名的\Place{巴勒斯坦}的\Place{君士坦提亚}。因为我的任务不是列举皇帝的作为，而仅是那些与基督宗教相关的事，尤其是与教会相关的。因此，我将这些光荣成就留给更有能力的人去详述，因为它们属于另一个主题，需要单独的论述。如果教会没有受到分裂的干扰，我本人本会保持沉默：因为当主题不提供可叙述的内容时，就不需要叙述者。既然微妙而虚妄的辩论已经混乱并同时分散了基督宗教的宗徒信仰，我认为记录这些事情是可取的，以免教会的事迹湮没无闻。因为对这些问题的准确了解能为众人赢得声誉，同时使了解它们的人更加免于错误，并教导他不被任何可能听到的空洞诡辩所左右。\par

% structure: p0090 source=h2 target=subsection reason=relative-heading-rank; base=h2

\subsection{第十九章 —— 在\Person{君士坦丁}时代印度内陆各民族如何归化基督信仰。}

现在我们必须说明基督信仰在这位皇帝在位期间是如何传播的：因为正是在他的时代，内陆的印度人和伊比利亚人这两个民族最初接受了基督信仰。但我将简要解释为何使用了附加的表达\Emph{内陆}。当宗徒们通过抽签往各民族去时，\Person{多默}接受了帕提亚人的宗徒职分；\Person{玛窦}被分派到\Place{厄提约丕雅}；而\Person{巴尔托罗茂}则分到与\Place{厄提约丕雅}接壤的印度地区：但内陆印度——那里居住着许多使用不同语言的蛮族——在\Person{君士坦丁}时代之前尚未被基督信仰的教义所光照。现在我来讲述导致他们皈依基督信仰的原因。一位名叫\Person{梅罗皮乌斯}的哲学家，是\Place{提洛}人，受哲学家\Person{梅特罗多罗斯}榜样的激励——此人先前曾游历过印度地区——决定亲自了解印度人的国家。于是他带着两个与他有亲缘关系的年轻人——他们并非不谙希腊语——乘船抵达了该国；在考察了他所想考察的一切之后，他在一个拥有安全海港的地方停靠，以购置一些必需品。恰巧在那之前不久，罗马人和印度人之间的条约被违反了。因此，印度人抓了那位哲学家和与他同船的人，将他们全部杀死，只留下了他两个年轻的亲属；但出于对他们年幼的怜悯而饶恕了他们，并将他们作为礼物送给了印度人的国王。国王对这两个少年的外貌感到满意，便任命其中一个名叫\Person{埃德西乌斯}的为他的司酒官；另一个名叫\Person{福鲁孟提乌斯}的，则委托他管理王室档案。国王不久后去世，给他们留下了自由，政府则交由他的妻子和幼子管理。王后看到儿子尚在幼年，便恳请这两个年轻人承担照管他的责任，直到他成年。于是，这两个年轻人接受了任务，开始治理王国。因此，\Person{福鲁孟提乌斯}掌管一切，并将查问在那国经商的罗马商人中是否有基督徒作为一项任务：发现了一些人后，他向他们说明了自己的身份，并劝诫他们选择并占据一些适当的地方来举行基督信仰的敬礼。过了不久，他建造了一座祈祷所；在教导了一些印度人基督信仰的原理后，他们使他们适于参与敬礼。当年轻的国王达到成年时，\Person{福鲁孟提乌斯}和他的同伴将公共事务的管理权交还给他——他们曾光荣地履行了管理职责——并恳求允许返回自己的国家。国王和他的母亲都恳求他们留下；但他们渴望重返故土，未能被说服，于是离开了。\Person{埃德西乌斯}则赶回\Place{提洛}去见他的父母和亲属；但\Person{福鲁孟提乌斯}抵达\Place{亚历山大城}，将此事报告给主教\Person{阿塔纳修斯}——他刚被授予这一尊位不久——并向他详细讲述了自己漂泊的经历以及印度人接受基督信仰的希望。他还恳求他派遣一位主教和圣职人员到那里，切不可忽略那些如此可得救恩的人。\Person{阿塔纳修斯}经过考虑，认为此事如何能最有效地成就，便请求\Person{福鲁孟提乌斯}本人接受主教之职，宣称他找不到比他更合适的人。于是这件事就这样成了；\Person{福鲁孟提乌斯}被授予主教权柄，回到印度，在那里成为福音的宣讲者，建造了几座教堂，并藉天主恩宠的助佑，行了许多奇迹，在医治灵魂的同时也医治了许多人的身体疾病。\Person{鲁菲努}向我们保证，他是从\Person{埃德西乌斯}那里听到这些事实的，\Person{埃德西乌斯}后来在\Place{提洛}被祝圣为司铎。\par

% structure: p0092 source=h2 target=subsection reason=relative-heading-rank; base=h2

\subsection{第二十章 伊比利亚人如何皈依基督信仰。}

现在适合叙述伊比利人如何在几乎同一时期成为信德的皈依者。某位虔诚贞洁的女子，在天主的上智安排中，被伊比利人掳去。这些伊比利人居住在\Place{黑海}附近，是\Place{西班牙}伊比利人的殖民。因此，那女子在被掳中，在野蛮人中修习德行：她不仅保持最严格的节欲，而且花费大量时间禁食和祈祷。野蛮人见此，对她的奇特行为感到惊讶。那时，王后之子，当时还是婴孩，患病了；王后按当地习俗，将孩子送到其他妇女那里去医治，希望凭借她们的经验提供疗法。婴孩被乳母抱着四处求治，却没有从任何妇女那里得到缓解，最后被带到这被掳的女子面前。她不懂医术，也没有采用任何药物；她只是抱起孩子，放在她铺着马衣的床上，在别的妇女面前，只说了一句：\Quote{基督，曾治愈多人，也要治愈这孩子}；然后，在说了这句信德之言后又祈祷，呼求天主，男孩立刻痊愈，从此一直健康。这奇迹的消息在野蛮人的妇女中广为流传，很快就传到了王后那里，那被掳的女子因此名声大噪。不久之后，王后本人患病，派人去请那被掳的女子。但这位女子谦逊退隐，推辞不去，于是王后被人抬到她那里。那被掳的女子像先前为王子所做的那样对待王后，疾病立刻消除。王后向这位陌生人道谢，但她回答说：\Quote{这工作不是我的，而是基督的，祂是创造世界的圣子}；因此她劝王后呼求祂，并认识真天主。伊比利王惊异于妻子突然康复，想用礼物酬谢他所知的、促成这些治愈的女子；但她却说不需要财富，因为她拥有宗教的安慰作为财富；她认为他能献给她的最大礼物，就是承认她所敬拜和宣告的天主。说完这话，她把礼物退了回去。国王将这番话铭记在心。次日出去打猎时，发生了这样的事：浓雾和黑暗笼罩了他打猎的山顶和森林，以致狩猎受阻，路径难辨。在困境中，国王恳切地呼求他所敬拜的众神；但毫无作用，最后他决定求助于那被掳女子的天主；他刚一开始祈祷，由浓雾产生的黑暗就完全消散了。他惊奇于所发生的事，欢喜地回到王宫，告诉妻子所发生的事；他立即派人去请那被掳的陌生人，恳求她告诉他，她所敬拜的是哪位神。那女子到来后，使伊比利王成为基督的宣讲者：因为他藉着这位虔诚的女子信了基督，他召集了所有在他权下的伊比利人；不仅向他们宣告了他妻子和孩子的治愈，以及打猎时发生的事，还劝他们敬拜那被掳女子的天主。于是，国王和王后都成了基督的宣讲者，一个向男性臣民，另一个向女性臣民宣讲。此外，国王从被掳者那里获悉罗马人建造圣堂的样式，便下令建造一座圣堂，立即提供了所需的一切；建筑工程随即开始。但当他们竖起柱子时，天主眷顾介入，以坚固居民的信德；因为有一根柱子纹丝不动，没有任何方法能移动它；绳索断裂，机械破碎；最后工人们放弃努力，离开了。这时，那被掳女子的信德之真实以如下方式显明：她夜间无人知晓地来到那地方，整夜祈祷；因天主的德能，那柱子被举起，悬在空中，立于柱基之上，却不接触柱基。天亮时，明智的国王亲自来查看工程，看见柱子这样悬空而无支撑，他和随从都非常惊讶。不久，就在他们眼前，柱子降落在自己的柱基上，牢牢固定。百姓为此欢呼，见证国王信德的真确，歌颂那被掳女子的天主。他们从此相信，并热心地竖起其余柱子，整座建筑很快完工。之后，他们派遣使节到\Person{君士坦丁}皇帝那里，请求从此与罗马人结盟，并从他们那里接受一位主教和祝圣的圣职人员，因为他们真诚相信基督。\Person{鲁菲努斯}说，他从\Person{巴库留斯}那里得知这些事实；\Person{巴库留斯}原是伊比利人的小王之一，但后来投靠罗马人，被任命为\Place{巴勒斯坦}军队的将领；最后在对抗暴君\Person{马克西穆斯}的战争中受命担任最高统帅，协助\Person{德奥多西}皇帝。这样，在君士坦丁的日子里，伊比利人也归信了基督信仰。\par

% structure: p0094 source=h2 target=subsection reason=relative-heading-rank; base=h2

\subsection{第二十一章 论隐修士\Person{安东尼}}

隐修士\Person{安东尼}是怎样的一个人，他生活在同一时代的埃及旷野中，如何与魔鬼公开争战，识破他们的诡计和狡猾的作战方式，以及他行了许多奇迹，我们无需赘述；因\Person{亚大纳削}，\Place{亚历山大}主教，已先于我们，写了一整本书为他作传记。在\Person{君士坦丁}皇帝年间，曾有大量这样的善人。\par

% structure: p0096 source=h2 target=subsection reason=relative-heading-rank; base=h2

\subsection{第二十二章  \Person{摩尼}，摩尼教异端的创始者，及其起源。}

但在好麦子中，稗子常会生长；因为嫉妒喜爱阴险地谋害良善。因此，就在\Person{君士坦丁}时代前不久，一种异教化的基督教与真正的基督教一同出现；正如假先知在真先知中出现，假宗徒在真宗徒中出现一样。因为当时，藉着\Person{摩尼}，异教哲学家\Person{恩培多克勒}的教义披上了基督教教义的外衣。该撒勒亚的\Person{尤西比乌斯·旁斐利}确实在他的\WorkTitle{教会史}第七卷中提到了此人，但没有详细叙述有关他的细节。因此，我认为有义务补充一些他忽略了的细节：这样就会知道这个\Person{摩尼}是谁，他从何而来，以及他僭妄行事是怎样的性质。\par

一个名叫\Person{斯泰辛}的撒拉逊人娶了一个来自\Place{上底比斯}的女俘。因她之故，他居住在埃及，并精通了埃及人的学问，他巧妙地将\Person{恩培多克勒}和\Person{毕达哥拉斯}的理论引入基督信仰的教义中。他断言有两性，一善一恶，并像恩培多克勒那样，称后者为\OriginalTerm{不和}{Discord}，前者为\OriginalTerm{友爱}{Friendship}。这个\Person{斯泰辛}有一个门徒叫\Person{布达斯}，先前被称为\Person{特勒宾图斯}；他前往波斯人所居住的\Place{巴比伦}，说了许多关于自己的夸张言论，声称自己由童贞女所生，在山中长大。此人后来又写了四本书，一本名为\WorkTitle{奥秘}，另一本为\WorkTitle{福音}，第三本为\WorkTitle{宝藏}，第四本为\WorkTitle{纲要}；但当他假装举行某种神秘仪式时，被一种灵推下悬崖，便如此丧命。一位曾留宿他的妇人埋葬了他，并占有了他的财产，买了一个约七岁的男孩，名叫\Person{库布利库斯}；这妇人给了他自由，并给予他良好的教育，不久后去世，将\Person{特勒宾图斯}的一切所有都留给了他，包括他根据\Person{斯泰辛}所教导的原则写的那些书。这位获释的\Person{库布利库斯}带着这些东西，退入波斯地区，改名为\Person{摩尼}；并将\Person{布达斯}或\Person{特勒宾图斯}的那些书当作自己的著作，散布给受他迷惑的信众。这些论文的内容在表达上表面上与基督教相符，但在情感上是异教的：因为\Person{摩尼}是个无神论者，他煽动门徒承认多神，并教导他们崇拜太阳。他还引入了\OriginalTerm{命运}{Fate}的教义，否认人的\OriginalTerm{自由意志}{free-will}；并肯定\OriginalTerm{身体轮回}{transmutation of bodies}，明显追随\Person{恩培多克勒}、\Person{毕达哥拉斯}和埃及人的意见。他否认\Person{基督}曾存在于肉身中，声称祂只是一个\OriginalTerm{幻影}{apparition}；还拒绝了法律和先知，称自己为\Quote{护慰者}——所有这些教义都与教会的正统信仰完全相悖。在他的书信中，他甚至胆敢称自己为宗徒；但这种毫无根据的僭称使他以如下方式遭到了应得的报应。波斯君王的儿子患病，他的父亲为他的康复而焦虑，并尝试了一切方法；他听说\Person{摩尼}显奇迹，认为这些奇迹是真的，就派人请他作为宗徒前来，相信通过他儿子能得痊愈。于是他来到宫中，以他装模作样的方式着手治疗年轻的王子。但国王见孩子死在他手中，便将这骗子关进监狱，意图处死他。然而，他设法逃出，逃往\Place{美索不达米亚}；但波斯王发现他住在那，便派人强行将他带来，活活剥了他的皮，用糠秕填满他的皮囊，悬挂在城门前。我们陈述这些并非自己编造，而是从一本名为\WorkTitle{卡什卡拉主教阿尔凯劳斯辩论集}的书中收集而来（\Place{卡什卡拉}是美索不达米亚的一座城市）。因为\Person{阿尔凯劳斯}本人声称他曾与\Person{摩尼}面面对质，并提到了我们刚才提到的有关他生平的情况。因此，正如我们之前所说，嫉妒喜爱在繁荣局势中暗中作祟。但天主的良善为何允许此事发生，无论是祂愿意借此彰显教会原则的卓越，并彻底粉碎常与信德相连的自负心理，还是出于其他原因，这既是一个难题，也与当前的讨论无关。因为我们的目的既不是审视教义观点的正确性，也不是分析天主的眷顾和审判的神秘原因；而是尽可能忠实地详述发生在教会中的事件的历史。摩尼教的迷信在\Person{君士坦丁}时代前不久兴起的方式已经这样描述了；现在让我们回到这本历史书的主要对象——时代与事件。\par

% structure: p0099 source=h2 target=subsection reason=relative-heading-rank; base=h2

\subsection{第二十三章  \Place{尼科美底亚}主教\Person{尤西比乌斯}与\Place{尼西亚}主教\Person{狄奥格尼斯}重获信心，密谋反对\Person{亚大纳削}，企图颠覆\OriginalTerm{尼西亚信经}{Nicene Creed}。}

欧瑟比和德奥格尼的同党从流放地返回后，后两者得以恢复其教会，如我们所见，他们驱逐了那些代替他们被祝圣的人。此外，他们受到皇帝极大的重视，皇帝非常尊敬他们，视其为从谬误回归正统信仰的人。然而，他们滥用这赐予的自由，在世界中引起了比以前更大的骚动；这受两种原因驱使——一方面是他们先前所受的亚略异端感染，另一方面是对亚大纳削的强烈敌意，因为他在会议讨论信德条款时曾如此有力地抵制他们。首先，他们反对亚大纳削的祝圣，部分因为他不配主教职，部分因为他是由不合格的人选出的。但当亚大纳削证明自己高于此诽谤（因为他接管亚历山大教会后，热切捍卫尼西亚信经），欧瑟比便竭尽全力阴谋使亚大纳削被撤职，并让亚略返回亚历山大；因为他认为只有这样，他才能抹去同质论，并引入亚略主义。因此欧瑟比写信给亚大纳削，希望他重新接纳亚略及其追随者进入教会。他的信函语气虽是恳求，但公开威胁他。由于亚大纳削绝不答应，欧瑟比设法促使皇帝召见亚略，然后允许他返回亚历山大；至于他用什么手段达到目的，我将在适当的地方提及。与此同时，在这之前教会中又发生了另一场骚动。事实上，她自己的儿女再次扰乱了她的平安。欧瑟比·庞非利乌斯说，会议之后不久，埃及便因内部分裂而动荡：但他没有说明原因，因此他赢得了不诚实的名声，并避免说明这些纷争的原因，这出于他决心不赞同尼西亚会议的决定。然而，正如我们从会议后主教们彼此通信的信函中所发现的，\OriginalTerm{同质}{homoousios}一词困扰了他们中的一些人。因此，当他们过于细致地探究其含义时，便激起了彼此之间的争斗；这似乎不亚于一场黑暗中的竞赛；因为双方似乎都不清楚他们互相诽谤的依据。那些反对\OriginalTerm{同质}{homoousios}一词的人，认为赞成它的人赞同撒伯里乌和蒙丹的意见；因此他们称这些人为亵渎者，因为他们颠覆了天主圣子的存在。而支持这个词的人则指责对手为多神论者，抨击他们引入了异教迷信。安提约基雅主教尤斯塔修斯指责欧瑟比·庞非利乌斯歪曲尼西亚信经；欧瑟比则否认他违反那信德表述，并反责尤斯塔修斯是撒伯里乌意见的辩护者。由于这些误解，每个人都如同与对手争辩般写作：尽管双方都承认天主圣子有独立的位格和存在，并都承认有一个天主三位，但我无法揣摩是什么原因，他们彼此不能同意，因此完全无法保持和平。\par

% structure: p0101 source=h2 target=subsection reason=relative-heading-rank; base=h2

\subsection{第二十四章 论安提约基雅会议，该会议罢免了安提约基雅主教尤斯塔修斯，因他而爆发了一场几乎毁灭该城的骚乱。}

于是，他们在\Person{安提约基雅}召集了一次主教会议，撤职了\Person{尤斯塔修斯}，因为他是撒伯里乌异端的支持者，而非\Place{尼西亚}会议所订信理的支持者。据一些人断言，此举是出于其他且不充分的理由，尽管并未公开说明其他理由；这是常见之事：主教们惯于在所有情形中如此行事，指控并被撤职者说他们不虔敬，却不解释这样做有何依据。\Place{叙利亚}\Place{老底嘉}主教\Person{乔治}，是厌恶\OriginalTerm{同质}{homoousios}一词的人之一，他在其\WorkTitle{尤西比乌斯·埃米森纳斯颂词}中向我们保证，他们因\Person{尤斯塔修斯}偏向撒伯里乌主义而撤职了他，指控来自\Place{贝罗亚}主教\Person{西鲁斯}。关于\Person{尤西比乌斯·埃米森纳斯}，我们将在别处按顺序述说。\Person{乔治}写到\Person{尤斯塔修斯}，有些前后矛盾；因为先断言他被\Person{西鲁斯}指控持守撒伯里乌异端，随后又告诉我们\Person{西鲁斯}本人被定罪为同一错误并因此被贬。然而，\Person{西鲁斯}自己倾向撒伯里乌主义，怎能指控\Person{尤斯塔修斯}为撒伯里乌派？因此，\Person{尤斯塔修斯}很可能因其他理由而被定罪。然而，那时在\Place{安提约基雅}，因其撤职而爆发了危险的骚乱；因为当进行继任者选举时，激起了如此激烈的纷争，足以使整个城市面临毁灭。民众分为两派：一派强烈主张将\Person{尤西比乌斯·潘菲卢斯}从\Place{巴勒斯坦}的\Place{凯撒利亚}调任到\Place{安提约基雅}；另一派同样坚持恢复\Person{尤斯塔修斯}的职位。该城民众在这场基督徒间的纷争中感染了党争精神，双方都列阵军队，心怀敌意，若非\Person{天主}和皇帝之恐惧压制了群众的暴力，就会发生流血冲突。因为皇帝通过书信，而\Person{尤西比乌斯}通过拒绝接受主教职，平息了骚乱；为此，这位主教受到皇帝极大钦佩，皇帝写信赞扬他的明智决定，并祝贺他被视为不仅是一座城市，而是几乎整个世界的配得主教。因此，据说此后\Place{安提约基雅}教会的宗座空缺了连续八年；但最终，在那些旨在推翻尼西亚信经的人之努力下，\Person{欧弗洛尼乌斯}被正式立为主教。这就是我所了解的关于\Place{安提约基雅}会议因\Person{尤斯塔修斯}之事的情况。在此之后不久，早前离开\Place{贝里图斯}、当时主持\Place{尼科美底亚}教会的\Person{尤西比乌斯}，与他的同党竭力将\Person{亚略}带回\Place{亚历山大里亚}。但现在必须叙述他们如何设法做到这一点，以及皇帝是如何被说服接见\Person{亚略}和\Person{欧佐伊乌斯}二人的。\par

% structure: p0103 source=h2 target=subsection reason=relative-heading-rank; base=h2

\subsection{第二十五章 论为召回\Person{亚略}而奔走的司铎}

皇帝\Person{君士坦丁}有一位姊妹名\Person{君士坦提娅}，是\Person{李锡尼}的遗孀；\Person{李锡尼}曾与\Person{君士坦丁}共同执掌帝国大权，但后来篡夺专制权力而被处死。这位公主在其府邸中供养着一位心腹司铎，此人沾染了亚略主义的教义；\Person{尤西比乌斯}等人怂恿他，他在与\Person{君士坦提娅}的日常谈话中趁机暗示，会议对\Person{亚略}不公，且关于他的一般传闻并不确实。\Person{君士坦提娅}完全相信这位司铎的断言，但不敢向皇帝报告。恰巧她身患重病，她的兄弟每日探视她。当病情加重，她预期死亡时，便将这位司铎托付给皇帝，证明他的勤勉和虔诚，以及他对君主的忠贞。她不久便去世了，于是这位司铎成为皇帝身边最亲信的人之一；他逐渐放胆直言，将先前对他姊妹所说的话重复给皇帝，声称\Person{亚略}的见解与会议所宣称的意见别无二致；并且，如果他蒙准觐见，他将完全同意会议所颁布的法令；他还补充说，\Person{亚略}受到了不合理的诽谤。皇帝对司铎的话感到惊奇，说道：\Quote{若\Person{亚略}签名同意会议并持其观点，我不仅接见他，还要体面地送他回\Place{亚历山大里亚}。}说完这话，他立即写信给\Person{亚略}，信文如下：\par

\Quote{ \Emph{胜利者君士坦丁马克西穆斯奥古斯都，致\Person{亚略}。} }\par

\Quote{ 日前已通知阁下，你可前来朕的宫廷，以便获得与朕会面之机。朕颇感惊讶，你未立即照办。因此，请立刻乘公共驿车，速到朕的宫廷；这样，当你体验了朕的仁慈和关怀后，便可返回本国。愿天主保佑你，挚爱的。十一月二十五日。 }\par

这是皇帝给\Person{亚略}的信。我不能不赞叹这位君主对宗教的热忱；因为从这份文件中可以看出，他曾多次劝诫\Person{亚略}改变观点，因为他责备\Person{亚略}迟迟不肯回到真理，尽管他自己曾多次写信给\Person{亚略}。当\Person{亚略}接到此信后，他与\Person{欧佐伊乌斯}一同来到\Place{君士坦丁堡}；\Person{欧佐伊乌斯}是\Person{亚历山大}在绝罚\Person{亚略}及其同党时已免去其执事之职的。皇帝于是接见了他们，问他们是否同意信经。他们欣然同意后，皇帝命他们交出一份书面信仰声明。\par

% structure: p0108 source=h2 target=subsection reason=relative-heading-rank; base=h2

\subsection{第二十六章 \Person{亚略}被召回后向皇帝\Person{君士坦丁}呈递悔过书，并假装接受尼西亚信经}

他们拟定了如下内容的宣言，呈递给皇帝。\par

\Quote{ \Person{亚略}和\Person{欧佐伊乌斯}，致我们最虔诚、最敬主的皇帝\Person{君士坦丁}。 }\par

\Quote{遵照您虔诚敬主的命令，主上，我们宣认我们的信德，并在天主面前以书面声明：我们及我们的追随者相信如下：}\par

\Quote{我们相信唯一的天主，全能的圣父；并相信主耶稣基督，祂的圣子，祂在万世之前由圣父所生，天主圣言，万物藉祂而造成，包括天上的和地上的；祂降下，降生成人，受难，复活，升天，并将再次来临审判生者死者。我们也相信圣神，相信肉身的复活，来世的生命，天国，以及天主的唯一天主教会，从地极延展到地极。}\par

\Quote{此信德我们由圣福音领受，主在其中对祂的门徒说：\BibleRef{玛 28:9} \BibleQuote{你们要去使万民成为门徒，因圣父、圣子、圣神之名给他们授洗。}如果我们不相信并真实领受圣父、圣子、圣神，正如整个天主教会和圣经所教导的（我们在一切方面都相信它们），那么天主现在和将来的审判中都是我们的审判者。因此我们恳求您的虔诚，最虔诚的皇帝，我们这些奉献于圣职的人，持守教会和圣经的信德与情感，愿因您和平而虔诚的敬虔，得以重新与我们的母亲教会联合，避免一切多余的问题和争论：好使我们和整个教会都处于平安中，能共同为您的安宁统治和您的全家献上惯常的祈祷。}\par

% structure: p0114 source=h2 target=subsection reason=relative-heading-rank; base=h2

\subsection{第二十七章 \Person{亚流}得皇帝同意返回\Place{亚历山大里亚}，因未被\Person{亚大纳削}接纳，\Person{欧瑟伯}党羽在皇帝面前提出多项控诉反对\Person{亚大纳削}。}

\Person{亚流}既这样满足了皇帝，便返回了\Place{亚历山大里亚}。但他压制真理的诡计并未成功；因为他到了\Place{亚历山大里亚}时，\Person{亚大纳削}不肯接纳他，反而视他如瘟疫般避开，他便试图在该城散布其异端，重新煽起骚动。于是，\Person{欧瑟伯}本人写信，并劝说皇帝也写信，为使\Person{亚流}及其党羽能重新获准进入教会。然而\Person{亚大纳削}完全拒绝接纳他们，并回信告知皇帝：那些一旦背弃信德、并被绝罚的人，不可能在他们返回时重新获准共融。但皇帝被此答复激怒，以下列言辞威胁\Person{亚大纳削}：\par

\Quote{既然你已得知我的旨意，就应让所有渴望进入教会的人无障碍地进入教会。因为若有人向我报告，说你禁止任何声称要重新与教会联合的人，或阻碍他们进入，我将立即派人，按我的命令革除你的职务，并将你流放。}\par

皇帝如此写信，是出于增进公共福祉的愿望，并因他不愿看到教会分裂；因为他辛勤努力，要使他们全都和谐共处。那时，\Person{优西比乌}的党羽，对\Person{亚大纳削}心怀恶意，自以为找到了合时的机会，便欢迎皇帝的不悦，视作自己目的的助力：为此他们掀起大骚动，力图将他从主教职上逐出；因为他们抱有希望，认为只有撤去\Person{亚大纳削}，\OriginalTerm{亚略学说}{Arian doctrine}才会得胜。反对他的主要共谋者是\Place{尼科美底亚}主教\Person{优西比乌}、\Place{尼西亚}的\Person{特奥格尼斯}、\Place{加采东}的\Person{马里斯}、\Place{上默西亚} \Place{辛吉杜努姆}的\Person{乌尔萨基乌}，以及\Place{上潘诺尼亚} \Place{穆尔萨}的\Person{瓦伦斯}。这些人用贿赂收买了\OriginalTerm{梅里提安异端}{Melitian heresy}中的某些人，捏造各种指控反对\Person{亚大纳削}；首先他们通过梅里提安派的\Person{伊西翁}、\Person{欧德蒙}和\Person{卡利尼库}指控他，说他曾命令埃及人向\Place{亚历山大里亚}的教会缴纳一件麻布衣服作为税赋。但这诽谤立即被\Place{亚历山大里亚}教会的司铎\Person{阿利皮乌}和\Person{玛卡里乌}推翻，他们当时恰好在\Place{尼科美底亚}；他们使皇帝相信，这些有损\Person{亚大纳削}的陈述是虚假的。因此皇帝写信严厉谴责了他的控告者，但敦促\Person{亚大纳削}前来见他。但在他到来之前，\OriginalTerm{优西比乌派}{Eusebian faction}抢先行动，在以前的指控上又增加了一项更加严重的罪行指控；控告\Person{亚大纳削}谋反，并曾为叛国目的将一箱黄金送给一个名叫\Person{斐路梅努}的人。然而，当皇帝本人在\Place{普萨玛提亚}（\Place{尼科美底亚}的郊区）调查此事，并发现\Person{亚大纳削}无辜时，他体面地送他离开；并亲笔写信给\Place{亚历山大里亚}教会，向他们保证他们的主教是被诬告的。诚然，对\OriginalTerm{优西比乌派}{Eusebian faction}随后对\Person{亚大纳削}的攻击缄默不提，本是合适且可取的，以免因这些情况，基督的教会受到那些敌视教会利益者的不利判断。但既然这些事已被记载下来，变得尽人皆知，我因此认为有必要尽可能简略地提及这些事情，其细节则需要一篇专论。因此，我现在将简述诽谤指控的起源以及策划者的性质。\Place{马雷奥提斯}是\Place{亚历山大里亚}的一个地区；其中包含许多村庄，人口众多，拥有许多宏伟的教堂；这些教堂都受\Place{亚历山大里亚}主教的管辖，并隶属于他的城市，作为堂区。在这个地区有一个人名叫\Person{伊斯库拉斯}，他犯了一件死有余辜的行为；因为他从未获授神品，却胆敢自称司铎，并行使属于司铎职的神圣职能。但他在亵渎的行径中被察觉，便从那里逃脱，到\Place{尼科美底亚}寻求庇护，在那里他恳求\OriginalTerm{优西比乌派}{Eusebian faction}的保护；这些人出于对\Person{亚大纳削}的仇恨，不仅接纳他为司铎，甚至许诺授予他主教职的尊位，只要他肯捏造一个针对\Person{亚大纳削}的指控，并以此为借口听取\Person{伊斯库拉斯}编造的任何故事。因为他散布传言，说他因遭受袭击而痛苦不堪；并说\Person{玛卡里乌}狂暴地冲向祭台，推翻了桌子，打破了一个奥秘之杯：他还补充说，他焚烧了圣书。作为这一指控的奖赏，如前所述，\OriginalTerm{优西比乌派}{Eusebian faction}许诺给他一个主教职；他们预见针对\Person{玛卡里乌}的指控会连同被告一起牵连\Person{亚大纳削}，因为\Person{玛卡里乌}似乎是奉他的命令行事。但这一指控是后来提出的；在此之前，他们策划了另一个充满最恶毒敌意的指控，我现在就要提及。他们通过某种我不知道的方式，弄到一只人手；无论是他们自己杀了人砍下手，还是从尸体上切下，唯有天主和那事的主谋知道：但无论如何，他们公开展示它，作为梅里提安主教\Person{阿尔塞尼乌}的手，同时将声称的手的主人隐藏起来。他们断言，这只手曾被\Person{亚大纳削}用于行使某些巫术；因此它成为这些诽谤者共同策划的最严重的指控理由：但如常发生的那样，所有对\Person{亚大纳削}怀有怨恨的人，同时出面提出了各种其他指控。当皇帝得知这些行径时，他写信给侄子监察官\Person{达尔马提乌}，后者当时住在\Place{叙利亚}的\Place{安提约基雅}，指示他命令将被告带到面前，在适当调查后，对可能被定罪者施加惩罚。他还派了\Person{优西比乌}和\Person{特奥格尼斯}到那里，以便案件在他们面前审理。当\Person{亚大纳削}得知他将被传唤到监察官面前时，他派人去\Place{埃及}严密搜查\Person{阿尔塞尼乌}；他确实查明他被藏在那里，但无法抓获他，因为他屡次改变藏匿地点。与此同时，皇帝取消了原定在监察官面前进行的审判，原因如下。\par

% structure: p0118 source=h2 target=subsection reason=relative-heading-rank; base=h2

\subsection{第二十八章 因对\Person{亚大纳削}的指控，皇帝在\Place{提洛}召集主教会议}

皇帝曾下令在一次主教会议出席他于\Place{耶路撒冷}建造的圣堂的奉献典礼。因此，他指示，作为次要事务，他们应当在途中先聚集于\Place{提洛}，审查对\Person{阿塔纳修}的指控；以便在那里消除一切争执后，他们能更和平地举行奉献天主圣堂的祝圣仪式。这是\Person{君士坦丁}统治的第三十年，六十位主教应执政官\Person{狄奥尼西}之召从各地聚集于\Place{提洛}。至于司铎\Person{玛卡里乌斯}，他被戴上锁链，由军队护送从\Place{亚历山大}押来；而\Person{阿塔纳修}不愿前往，并非出于恐惧——因为他对自己被指控的罪名是无辜的——而是因为他担心对\Place{尼西亚}会议的决定作出任何更改；然而，在皇帝充满威胁的信件下，他被迫出席。因为给他写了信：如果他自愿不来，将被强制带来。\par

% structure: p0120 source=h2 target=subsection reason=relative-heading-rank; base=h2

\subsection{第二十九章　论\Person{阿尔塞尼乌斯}及其被指控被砍下的手}

天主特别的眷顾也将\Person{阿尔塞尼乌斯}引至\Place{提洛}；因为不顾他从贿赂他的控告者那里接受的指示，他乔装前往那里，想看会发生什么。不知何故，行省总督\Person{阿尔凯劳斯}的仆人们听到一些人在一家客栈断言，据报被谋杀的\Person{阿尔塞尼乌斯}被藏匿在一名市民的家中。他们听到这话并记住说这些话的人之后，将信息报告给他们的主人，主人立即下令严格搜查此人，发现并妥善扣留了他；之后他通知\Person{阿塔纳修}不必惊慌，因为\Person{阿尔塞尼乌斯}活着且在当场。\Person{阿尔塞尼乌斯}被捕后，起初否认他就是那个人；但\Place{提洛}主教\Person{保禄}以前认识他，确认了他的身份。天主眷顾如此安排后，不久\Person{阿塔纳修}被会议传唤；他一出现，诬告者就出示了那只手，并坚持指控。他非常审慎地处理此事，因为他询问在场的人以及他的控告者，谁认识\Person{阿尔塞尼乌斯}？好几个人回答说认识他，他让人将\Person{阿尔塞尼乌斯}带进来，双手用斗篷遮住。然后他再次问他们：\Quote{这就是那个失去一只手的人吗？}这一出乎意料的举动使所有人惊讶，除了那些知道那只手是从哪里砍下来的人；因为其余的人以为\Person{阿尔塞尼乌斯}真的少了一只手，并期望被告会以其他方式辩护。但\Person{阿塔纳修}将\Person{阿尔塞尼乌斯}的斗篷向一边掀开，露出他的一只手；再次，当有些人以为另一只手缺失时，让他们短暂犹豫片刻后，他又将斗篷向另一边掀开，露出另一只手。然后他对在场的人说：\Quote{如你们所见，\Person{阿尔塞尼乌斯}发现有两只手；让我的控告者出示第三只被砍下的手是从哪里来的。}\par

% structure: p0122 source=h2 target=subsection reason=relative-heading-rank; base=h2

\subsection{第三十章　\Person{阿塔纳修}被查明对其所控罪名无罪；他的控告者逃跑}

关于\Person{阿尔塞尼乌斯}的事发展到这一地步，这一骗局的策划者陷入困境；主要控告者之一\Person{阿哈布}（又称\Person{约翰}）在混乱中溜出法庭，得以逃脱。这样，\Person{阿塔纳修}未经任何辩护就澄清了这一指控；因为他确信仅看见\Person{阿尔塞尼乌斯}活着就会使诽谤者困惑。\par

% structure: p0124 source=h2 target=subsection reason=relative-heading-rank; base=h2

\subsection{第三十一章　当主教们不愿听\Person{阿塔纳修}对第二项指控的辩护时，他投奔皇帝}

但在反驳对\Person{玛卡里乌斯}的虚假指控时，他使用了法律形式；首先对\Person{尤西比乌斯}及其一派提出异议，视他们为仇敌，抗议任何人被对手审判的不公正。接着他坚持要求证明他的控告者\Person{伊斯库拉斯}确实获得了司铎的尊位；因为在起诉书中他就是这样被指称的。但由于法官们不允许这些反对意见中的任何一项，\Person{玛卡里乌斯}的案件被审理，而告发者缺乏证据，听证被推迟，直到某些人前往\Place{马雷奥蒂斯}，以便当场审查所有可疑之处。\Person{阿塔纳修}看到被派去的正是他所反对的那些人（被派去的是\Person{特奥格尼斯}、\Person{玛里斯}、\Person{特奥多鲁斯}、\Person{马其顿尼乌斯}、\Person{瓦伦斯}和\Person{乌尔萨基乌斯}），他喊道：\Quote{他们的程序既奸诈又欺诈；因为让司铎\Person{玛卡里乌斯}被锁链拘禁，而控告者连同作为他对手的法官被允许去进行\Emph{单方面}的证据收集是不公正的。}他在整个会议和行省总督\Person{狄奥尼西}面前作了这一抗议，发现无人理会他的上诉，便私下退出。因此，被派到\Place{马雷奥蒂斯}的人进行了\Emph{单方面}调查，认为控告者所说的是真实的。\par

% structure: p0126 source=h2 target=subsection reason=relative-heading-rank; base=h2

\subsection{第三十二章　\Person{阿塔纳修}离开后，组成会议的人投票决定罢免他}

于是\Person{阿塔纳修}离开了，赶往皇帝那里；会议首先在他缺席的情况下谴责了他；当在\Place{马雷奥蒂斯}进行的调查结果提交后，他们投票罢免了他；在他们的罢免判决中用辱骂性的言辞责难他，但对他控告者的谋杀指控的可耻失败完全保持沉默。此外，他们接纳了据报被谋杀的\Person{阿尔塞尼乌斯}进入共融；而这位先前曾是\OriginalTerm{梅利提安异端}{Melitian heresy}的主教，以\Place{伊普塞洛波利斯}城主教的身份签署了对\Person{阿塔纳修}的罢免。这样，通过一系列不寻常的情况，据称被\Person{阿塔纳修}暗杀的受害者被发现活着协助罢免他。\par

% structure: p0128 source=h2 target=subsection reason=relative-heading-rank; base=h2

\subsection{第三十三章　会议成员从提洛前往耶路撒冷，庆祝\Quote{新耶路撒冷}奉献礼，并接纳亚略及其追随者共融}

此时，皇帝的谕令送达，指示组成会议的人士火速前往\Quote{新耶路撒冷}。于是他们立刻离开提洛，全速赶赴耶路撒冷。在那里，他们为圣地的奉献举行了庆典，随后遵照皇帝，如他们所说，的意愿，重新接纳亚略及其追随者共融。皇帝在谕令中表示，他对亚略和欧佐依的信仰完全满意。此外，他们致函亚历山大里亚教会，声明所有嫉妒现已消除，教会事务已得平安；既然亚略藉其悔改承认了真理，他既已是教会成员，也理应从今以后被他们接纳，\EditorialNote{此处影射亚大纳削被放逐，他们在声明中称\Quote{所有嫉妒现已消除}。}同时，他们以几乎相同的措辞向皇帝禀报了所行之事。但正当主教们忙于这些事务时，皇帝另有一封预料之外的信函来到，暗示亚大纳削已逃往他处请求保护，因此他们有必要为他前往君士坦丁堡。皇帝这封突如其来的信函内容如下：\par

% structure: p0130 source=h2 target=subsection reason=relative-heading-rank; base=h2

\subsection{第三十四章　皇帝以书信召集会议至御前，以便当面对亚大纳削的指控进行仔细审查}

\Quote{胜利者君士坦丁·马克西穆斯·奥古斯都，致在提洛集会的诸位主教。}\par

朕确不知贵议会在如此动荡与风暴中所作的决定；但真理似乎已被某些纷乱无序的进程所歪曲，因为，也就是说，在你们彼此争辩的喜好中——你们似乎欲使其永存——你们忽视了那些中悦天主之事。然而，朕相信，天主眷顾之工一旦被察觉，将驱散由此嫉妒争斗所致的祸害，并向我们显明：你们这些与会者是否顾及了真理，你们对呈交判决之事所作的裁决是否摒弃了偏私或偏见。因此，你们所有人必须毫不延迟地前来朕之虔敬前，亲自严格汇报你们的行事。朕为何认为应当如此书写并召集你们至朕前，你们将由下文得知。当朕进入以朕之名命名的都城，即我们这最繁荣的家园君士坦丁堡时——恰巧朕正骑马前行——突然，亚大纳削主教带着他身边的一些神职人员，如此出乎意料地出现在朕的路上，以致引起惊慌。因为全知的天主可为朕作证：初见之时，朕并未认出他，直到朕的若干随从应朕之问询，极其自然地向朕告知他是何人，以及他遭受了何等不义。那时，朕既未与他交谈，也未与他有任何沟通。但他多次恳求晋见，朕不仅拒绝，而且几乎下令将他从朕前带走，他更大胆地说，他所求的不过是请你们被召至此，以便在朕面前，他因迫不得已而走上此路，能获得公平的机会诉说他的冤屈。因此，既然此请求看似合理，且符合朕治理的公正，朕欣然下令将这些话书写给你们。朕的命令是：所有在提洛集会的会议成员，应立即赶赴朕之仁慈的王庭，以便你们亲自在我——你们不得不承认是天主真实仆人者——面前，以事实本身显明你们裁决的纯洁与正直。正是因为朕对天主的虔敬服务之举，和平才处处统治，天主之名甚至在至今仍不知真理的蛮族中也得了真诚的敬畏。显然，不知真理者也不真正认识天主；然而，正如朕先前所言，连蛮族也因朕——身为天主真实的仆人——的缘故，承认并学习了敬拜那位他们已切实看到处处保护并照顾朕的神。因此，主要是由于对我们的畏惧，他们被引向认识并敬拜现在所敬拜的真神。然而，我们这些自称对其教会神圣奥迹怀有虔敬敬畏（朕不愿说守护）的人，我们，朕说，所做的一切无不是导致纷争与敌意，且直言不讳地说，是导致人类毁灭之事。但你们要尽快赶来，正如朕已说过的，全体到朕处来。你们要确信：朕将竭尽全力，使神律所包含的内容得以无玷保存，没有任何污点或责备能附着其上；当那些以神圣职业为掩护、引入众多多样亵渎之言的敌人被驱散、粉碎、彻底消灭时，此事必会成就。\par

% structure: p0133 source=h2 target=subsection reason=relative-heading-rank; base=h2

\subsection{第三十五章　会议未至皇帝处，欧瑟比党人指控亚大纳削威胁中断从亚历山大里亚供应君士坦丁堡的粮食；皇帝震怒，将亚大纳削放逐至高卢}

这封信使组成该主教会议的人非常恐惧，因此他们大多数人回到了各自的城市。但欧瑟伯、特奥格尼斯、马里斯、帕特罗菲鲁斯、乌尔撒基乌斯和瓦伦斯前往君士坦丁堡，不再允许对打碎的圣爵、掀翻的祭台和阿尔森纽斯的谋杀进行进一步调查；反而采用了另一种诽谤，告诉皇帝说，亚大纳削曾威胁要禁止通常从亚历山大运往君士坦丁堡的谷物运输。他们还断言，这些威胁是亚大纳削当着主教阿达曼提乌斯、阿努比昂、阿尔巴提翁和伯多禄的面说的，因为诽谤最盛行时，控告者似乎是一个可信的人。因此，皇帝被欺骗，因这一指控而对亚大纳削感到愤慨，立即判处他流放，命令他居住在高卢。现在有些人认为，皇帝做出这一决定是为了实现教会的合一，因为亚大纳削在拒绝与亚略及其追随者有任何共融方面毫不妥协。因此，他定居在特雷维斯，一座高卢的城市。\par

% structure: p0135 source=h2 target=subsection reason=relative-heading-rank; base=h2

\subsection{第36章。安基拉主教马塞鲁与诡辩家阿斯忒里。}

在君士坦丁堡集合的主教们也罢免了安基拉（小加拉达省的一个城市）的主教马塞鲁，原因如下。一个卡帕多细亚的修辞学家名叫阿斯忒里，他放弃了其技艺，声称自己皈依了基督信仰，撰写了一些至今仍存的论文，他在其中赞扬了亚略的教义；断言基督是天主的能力，如同梅瑟所说的蝗虫和蝻子是天主的能力一样\BibleRef{岳 2:25}，以及其他类似言论。阿斯忒里经常与主教们交往，尤其是那些不反对亚略教义的主教；他也参加他们的主教会议，希望暗中获得某个城市的主教职位；但由于他在迫害期间献过祭，未能获得圣职。因此，他走遍叙利亚各城，公开朗读他所写的书。马塞鲁得知此事，想抵消他的影响，在过于急切驳斥他时，陷入了完全相反的错误；因为他竟敢像撒摩沙特人那样说，基督只是一个凡人。当时在耶路撒冷聚会的主教们得知这些事后，没有理会阿斯忒里，因为他甚至没有被列入圣职人员的名册；但他们坚持认为马塞鲁作为圣职人员，应为他所写的书作出交代。发现他持有撒摩沙特保禄的观点后，他们要求他撤回自己的意见；他深感羞愧，答应烧掉他的书。但由于皇帝召他们到君士坦丁堡，主教会议仓促解散，欧瑟伯派的人到达该城后，再次审理了马塞鲁的案件；马塞鲁拒绝履行烧毁他那本不适时之书的承诺，在场的人便罢免了他，并派遣巴西尔代替他前往安基拉。此外，欧瑟伯写了一篇驳斥这部著作的三卷书，揭露了其中的错误教义。然而，马塞鲁后来在萨尔底加主教会议上恢复了主教职务，因为他声称自己的书被误解了，因此被视为支持撒摩沙特人的观点。但关于此事，我们将在适当的地方更详细地叙述。\par

% structure: p0137 source=h2 target=subsection reason=relative-heading-rank; base=h2

\subsection{第37章。亚大纳削被放逐后，皇帝召来亚略，亚略引起了对君士坦丁堡主教亚历山大的骚动。}

当这些事情发生时，君士坦丁统治的第三十年结束了。但是亚略和他的追随者回到亚历山大，再次搅乱了整个城市；因为亚历山大的民众对这位不可救药的异端者及其同伙的复职以及他们的主教亚大纳削被流放感到极其愤慨。当皇帝得知亚略的乖戾性情时，他再次命令他前往君士坦丁堡，说明他重新试图激起的骚动。那时，亚历山大（他此前已接替默特罗法内）主持君士坦丁堡的教会。这位主教是一位虔诚奉献的人，这一点从他与亚略的冲突中清楚地显现出来；因为亚略到达后，民众分裂为两派，整个城市陷入混乱：一些人坚持认为尼西亚信经决不可侵犯，而另一些人则争辩说亚略的意见合乎理性。在这种情况下，亚历山大陷入困境；尤其是尼科美底亚的欧瑟伯猛烈威胁说，除非他接纳亚略及其追随者进入共融，否则将立即罢免他。然而，亚历山大对自己的被罢免远不如此担忧，他更担心他们急于要实现的信仰原则的颠覆；他视自己为公认的教义和尼西亚议会所作决定的守护者，竭尽全力防止它们被违反或败坏。陷入如此绝境，他告别了所有逻辑手段，以天主为避难所，致力于持续禁食，从未停止祈祷。他没有向任何人透露他的意图，独自关在被称为\Emph{伊雷内}的教堂里；在那里他走上祭台，俯伏在圣体圣桌下的地上，流着泪倾吐他热切的祈祷；许多个连续的夜晚和白天，他不停地这样做。他如此恳切地向天主祈求的，他得到了；因为他的祈求是这样的：\Quote{如果亚略的意见是正确的，求主不要让他看到指定讨论它的那一天；但如果他自己持守真信仰，那么亚略作为这一切罪恶的制造者，应承受与他邪恶相应的惩罚。}\par

% structure: p0139 source=h2 target=subsection reason=relative-heading-rank; base=h2

\subsection{第39章。亚略的死亡。}

这就是亚历山大的恳求。同时，皇帝希望亲自审查阿里乌斯，便派人将他召至宫中，问他是否同意尼西亚会议的决定。他毫不犹豫地表示同意，并在皇帝面前签署了信德宣言，但心存虚伪。皇帝见他轻易顺从，便令他宣誓确认签名。他同样虚伪地做了。据我所闻，他躲避的方式是：他将自己的意见写在纸上，夹在腋下，这样他确实发誓他真正持有他所写下的意见。然而，这只是我根据传闻所写，但他确实在签署之后又加了宣誓，这是我从皇帝亲笔信函中查证到的。皇帝因此信服，便下令由君士坦丁堡主教亚历山大接纳他进入共融。那天是星期六，阿里乌斯期待在次日与教会聚集。但天主的惩罚临到了他的狂妄罪行。当他走出皇宫，由一群尤西比乌斯派党羽如护卫般簇拥时，他骄傲地穿过城市中央，吸引了所有人的注意。当他走近被称为君士坦丁广场的地方，那里立有斑岩柱时，良心的悔恨所产生的恐惧抓住了阿里乌斯，随着恐惧，突然发生剧烈的腹泻。于是他询问附近是否有方便之处，被指引到君士坦丁广场后方，他急忙前往。不久，一阵虚弱袭来，随着排泄物，他的肠子突出，接着是大量的出血，小肠也掉了下来；此外，他的脾和肝的碎片也在血中排出，以至于他几乎立即死亡。这场灾难的地点至今仍在君士坦丁堡展示，如我所说，在柱廊内的肉市后面。路人用手指着那个地方，这种异常死亡的记忆被永久保存。这样不幸的事件使尼科美底亚主教尤西比乌斯一党充满恐惧和惊慌；消息迅速传遍全城和整个世界。随着皇帝对基督宗教愈加热忱，并承认尼西亚的宣认是由天主所证实，他对所发生的事感到欣喜。他也为自己的三个儿子感到高兴，因为他已经宣布他们为凯撒，每一个都是在他在位每十年纪念日时依次被设立的。在第一个十年结束时，他将帝国西部的治理权分给长子，他称之为君士坦丁，以他自己的名字命名。第二个十年完成后，第二个儿子君士坦提乌斯，以他祖父的名字命名，被立为东部的凯撒。而最小的君士坦斯，在他统治的第三十年被授予同样的尊荣。\par

% structure: p0141 source=h2 target=subsection reason=relative-heading-rank; base=h2

\subsection{第三十九章　皇帝患病与逝世}

一年过去了，君士坦丁皇帝刚满六十五岁，便患了病。于是他离开君士坦丁堡，航海至海伦波利斯，想试试该城附近有益的温泉。但他发现病情加重，便推迟了使用温泉；从海伦波利斯迁往尼科美底亚，居住在城郊，并在那里领受了基督徒的圣洗。此后他变得开朗，立下遗嘱，指定他的三个儿子为帝国的继承人，照他生前所作安排，给他们每人分派一份。他还授予罗马和君士坦丁堡许多特权，并将遗嘱的保管委托给那位使阿里乌斯被召回的司铎（我们已提及此人），嘱咐他除了交给他的儿子君士坦提乌斯（他已将东方统治权授予他）外，不得交付任何人。立遗嘱后，他仅存活数日便去世了。他去世时儿子们无一在场。于是立即差遣信使前往东方，告知君士坦提乌斯其父的死讯。\par

% structure: p0143 source=h2 target=subsection reason=relative-heading-rank; base=h2

\subsection{第四十章　君士坦丁皇帝的葬礼}

皇帝的遗体由相关人士安放在金棺中，然后运往君士坦丁堡，在宫殿中陈列于高台之上，有卫兵环护，受到如同他生前般的尊敬，直到他一个儿子到来。当君士坦提乌斯从帝国东部来到时，遗体被赐予了皇帝的葬礼，并安放在称为\WorkTitle{宗徒}的教堂中。他正是为此目的而修建了这座教堂，使皇帝和主教们得以受到几乎不低于对宗徒圣髑所表示的敬礼。君士坦丁皇帝活了六十五岁，统治了三十一年。他死于费利西安和塔提安执政之年，五月二十二日，在第二百七十八届奥林匹亚的第二年。因此，本书涵盖了一段三十一年的时期。\par

\SourceRecord{Translated by A.C. Zenos.；Nicene and Post-Nicene Fathers, Second Series；Vol. 2.；Edited by Philip Schaff and Henry Wace.；Buffalo, NY: Christian Literature Publishing Co.,；1890.；<http://www.newadvent.org/fathers/26011.htm>.}

% structure: p0001 source=h1 target=section reason=page-title

\section{\WorkTitle{教会史（第二卷）}}

% structure: p0002 source=h2 target=subsection reason=relative-heading-rank; base=h2

\subsection{\Emph{第一章 绪论——作者修订第一、二卷的原因。}}

鲁菲努斯以拉丁文撰写了一部《教会史》，但在年代上犯了错误。因为他以为对阿塔纳修斯所行的一切发生在君士坦丁皇帝驾崩之后；他也不晓得他被流放至高卢以及其他许多情况。我们最初撰写这部史书的第一、二卷时，是跟随鲁菲努斯的；在撰写第三至七卷时，部分事实取自鲁菲努斯，部分取自其他作者，部分来自尚存者的叙述。然而此后，我们细读了阿塔纳修斯的著作，其中他描述了自己的苦难，以及他如何因欧瑟比派虚构的诬告而被放逐，并且断定，更应相信亲身受苦者以及那些目睹其事者，而非那些依赖猜测因而犯错的人。此外，我们获得了那个时期若干显赫人物的几封书信，也借助了这些书信来尽可能探究真相。为此，我们不得不修订这部史书的第一、二卷，不过，在鲁菲努斯显然不可能出错之处，我们仍使用他的证词。还应注意，在先前的版本中，既未收入对亚略的判决词，也未收入皇帝的书信，而只是叙述或事实，以免史书冗长，以繁琐细节使读者厌倦。但在本版中，为了你的缘故，\Person{特奥多禄}，\OriginalTerm{天主的人}{sanctus homo Dei}，我们作了这些修改和增补，使你不致不知晓诸王以自己的言词写了什么，以及主教们在各次主教会议中所作的决议，他们在其中不断更改信仰宣言。因此，凡我们认为必要者，均收入此后期版本。在第一卷中采取了这一方针，我们也将在史学接续部分——即第二卷——中努力同样行事。现在，让我们进入正题。\par

% structure: p0004 source=h2 target=subsection reason=relative-heading-rank; base=h2

\subsection{\Emph{第二章 尼苛米底亚主教欧瑟比及其一派再次试图引进亚略异端，在教会中引起骚动。}}

君士坦丁皇帝驾崩后，尼苛米底亚主教欧瑟比与尼凯亚主教特奥格尼斯以为有机可乘，便竭尽全力删除\OriginalTerm{同质}{homoousion}信理，而以亚略主义取而代之。然而，他们若想实现此目的，除非阿塔纳修斯不能返回亚历山大里亚。因此，为完成其计谋，他们寻求那位不久前帮助亚略自流放地被召回的长老的协助。此事如何成就，现宜叙述。该长老将已故君王的遗嘱与请求呈献给其子君士坦提乌斯；后者发现遗嘱中所安排的内容正是他最渴望的——东方的帝国已由父王遗命分派给他——于是厚待该长老，赐以恩宠，并准其自由进出宫殿与自己面前。这种自由很快使他得以与皇后及她的宦官亲密往来。当时有一位寝宫总宦官，名叫欧瑟比；长老说服他采纳了亚略的观点，此后其余的宦官也相继接受了同样的主张。不仅如此，皇后也在宦官和长老的影响下，对亚略的教义产生了好感；不久之后，此事也被引荐给皇帝本人。这样，它逐渐传遍朝廷及内外臣仆、侍卫之中，最终蔓延至全城居民。宫廷内侍与妇孺争论此教理；每个市民家中都有逻辑争辩。此外，祸害迅速扩展到其他省份和城市，这争论如同火星，起初微不足道，却在听众中激起争辩精神：凡询问骚乱缘由者，立即找到争论的题目，并在询问之际决心参与纠纷。由于这种普遍的争吵，一切秩序倾覆；然而骚动仅限于东方诸城，伊利里库姆和帝国西部地区则完全平静，因为它们不愿废除尼凯亚大公会议的决议。由于此事日益加剧，每况愈下，尼苛米底亚的欧瑟比及其一派视民众骚动为幸运之事。因为他们认为，只有这样，他们才能使自己教派的人担任亚历山大里亚主教。但此时阿塔纳修斯归来，挫败了他们的目的；因为他带着一位奥古斯都的书信回到那里，这封书信是小君士坦丁——他以父亲的名字命名——自高卢的特雷维斯城寄给亚历山大里亚居民的。该信副本附录于下。\par

% structure: p0006 source=h2 target=subsection reason=relative-heading-rank; base=h2

\subsection{\Emph{第三章 阿塔纳修斯受小君士坦丁书信鼓励，返回亚历山大里亚。}}

\Signature{君士坦丁凯撒}致亚历山大里亚天主教会成员。\par

我想，诸位虔诚的心灵不可能不知道，受人敬重的律法的阐释者亚大纳削，曾被暂时遣往高卢，以免因其嗜血敌人的邪恶而遭受不可弥补的伤害，他们的暴行不断危及他的圣洁生命。因此，为躲避这种邪恶，他从威胁他的人们口中被带到我的管辖之下的一个城市，在那里，在他被指定的居住期间，他获得了一切必需品的充足供应：尽管他的卓越美德因信赖天主的助佑，本可轻视更严酷命运的压迫。既然我们的君主，先君君士坦丁奥古斯都，因去世而未能实现将这位主教恢复其教座及你们至圣之虔诚的心愿，我认为继承他的任务并实现他的愿望是适宜的。他受到我们如何大的尊敬，你们在他到达时便会知道；没有人应对我给予他的荣誉感到惊讶，因为我既受对如此杰出人物应有之感的驱使，也深知你们对他的关怀之情。愿上智安排保护你们，亲爱的弟兄们。\par

凭着这封信，亚大纳削来到亚历山大里亚，并受到该城人民最欢乐的接待。然而，城中凡信奉亚略主义的人，联合起来，结党反对他，由此屡次激起骚动，为欧瑟伯派提供了借口，向皇帝控告他出于己意占领了亚历山大里亚教会，无视主教公会议的相反判决。他们确实极力指控，以致皇帝被激怒，将他从亚历山大里亚流放。此事究竟如何发生，我将在后面说明。\par

% structure: p0010 source=h2 target=subsection reason=relative-heading-rank; base=h2

\subsection{第四章 欧瑟伯·庞非勒之死，亚加爵继任凯撒勒雅主教}

此时，欧瑟伯，即巴勒斯坦凯撒勒雅的主教，别号庞非勒者，去世了。其门徒亚加爵继任主教。此人出版了几部著作，其中包括其师传记。\par

% structure: p0012 source=h2 target=subsection reason=relative-heading-rank; base=h2

\subsection{第五章 小君士坦丁之死}

此后不久，君士坦提乌斯皇帝的弟弟小君士坦丁，因与其弟君士坦斯皇帝管辖的帝国部分地区为由，与其弟的军队交战，被其军队所杀。此事发生在阿辛迪努斯和普罗克鲁斯执政官任期内。\par

% structure: p0014 source=h2 target=subsection reason=relative-heading-rank; base=h2

\subsection{第六章 君士坦丁堡主教亚历山大临终时，提议选举保罗或马其顿尼为其继任者}

大约同时，除我们已经记载的骚乱外，君士坦丁堡又发生了另一场骚乱，起因如下：在该城主持教会并坚决反对亚略的亚历山大，已离世，他担任主教二十三年，在世九十八年，未任命任何人继任。但他嘱咐适当的人选从他所提名的两人中选出一位：即，若他们想要一位擅长教导且卓越虔诚的人，就应选举保罗——他亲自祝圣为司铎，此人虽年轻，但智慧与审慎成熟；但若他们想要一位仅外表可敬、貌似圣洁的人，则可任命马其顿尼——他长期在他们中间任执事，且年事已高。因此，关于主教的选举爆发了巨大争论，严重困扰了教会；因为自人民分为两派以来，一派拥护亚略主义的教义，另一派坚持尼西亚会议所定义的信理，在亚历山大在世期间，主张同体论的人始终占优势，而亚略派内部意见分歧，争论不休。但在亚历山大大主教去世后，斗争的结果变得不确定：正统信仰的捍卫者坚持祝圣保罗，而全体亚略派则支持马其顿尼。因此，保罗在称为\Emph{依勒内}的教堂被祝圣为主教，该教堂位于伟大的索菲亚教堂附近；他的选举似乎更符合逝者的意愿。\par

% structure: p0016 source=h2 target=subsection reason=relative-heading-rank; base=h2

\subsection{第七章 君士坦提乌斯皇帝在保罗当选主教后将其驱逐，并派人从尼科米底亚召来欧瑟伯，授予其君士坦丁堡主教职}

不久之后，皇帝到达君士坦丁堡，对保罗的祝圣极为恼怒；他召集了一个亚略派主教会议，剥夺了保罗的尊位，并将欧瑟伯从尼科米底亚教区调离，立他为君士坦丁堡主教。做完此事后，皇帝便前往安提约基雅。\par

% structure: p0018 source=h2 target=subsection reason=relative-heading-rank; base=h2

\subsection{第八章 欧瑟伯在叙利亚安提约基雅召开另一次会议，并颁布新信经}

然而，欧瑟伯绝不能保持沉默，而是如俗语所说，不遗余力地实现他心中的目的。因此，他借口为奥古斯都之父所奠基、其子君士坦提乌斯在奠基后第十年完成的那座教堂举行奉献礼，在叙利亚的安提约基雅召集了一次主教会议，但其真实意图是推翻并废除\Emph{\OriginalTerm{同质}{homoousion}}的道理。来自各城的九十位主教出席了这次会议。然而，接替玛卡里乌斯的耶路撒冷主教玛克西穆没有出席，他记着自己曾受欺骗而签署了罢黜亚大纳削的决议。伟大的罗马主教儒略也没有出席，也没有派人代理，尽管教会法令规定，各教会不得违背罗马主教的意见作出任何规定。这次主教会议在奥古斯都之父君士坦丁去世后第五年、玛尔凯卢斯与普罗比努斯执政时期，在皇帝君士坦提乌斯面前于安提约基雅召开。当时，接替欧弗若尼乌斯的普拉基图斯（或称弗拉基卢斯）主持安提约基雅教会。欧瑟伯的同党先前已设计诽谤亚大纳削，首先指控他违反了当时制定的一条法令，即未经主教全体会议许可而擅自恢复主教职权，因为他从流放地回来后擅自占据了教会；又指控他进入时引发了骚乱，许多人在暴乱中丧生；此外，他还鞭打了一些人，并将另一些人送交法庭。他们还提出了在\Place{提洛}对亚大纳削作出的裁决。\par

% structure: p0020 source=h2 target=subsection reason=relative-heading-rank; base=h2

\subsection{第九章 论厄米萨的欧瑟伯}

基于这些指控，他们为亚历山大里亚教会提出了另一位主教，首先是绰号厄米森努斯的欧瑟伯。关于此人，当时在场的劳迪塞雅主教乔治告诉我们。因为他在自己所撰写的关于其生平的书卷中说，欧瑟伯出身于美索不达米亚厄德萨的贵族家庭，自幼学习圣经；后来在厄德萨的一位老师的教导下学习希腊文学；最后，帕特若斐鲁斯和欧瑟伯（后者主持凯撒勒雅教会，前者主持斯库托波利斯教会）为他解释了圣经。后来，当他住在安提约基雅时，欧斯塔丢因贝罗亚的齐禄指控他持有撒伯里乌斯的观点而被罢黜。之后，他又与欧斯塔丢的继任者欧弗若尼乌斯交往，并逃避主教职务，退隐到亚历山大里亚，专心研究哲学。返回安提约基雅后，他与欧弗若尼乌斯的继任者普拉基图斯（或称弗拉基卢斯）结为密友。最后，他被君士坦丁堡主教欧瑟伯祝圣为亚历山大里亚主教，但因该城人民对亚大纳削的依恋而未能前往，于是被派往厄米萨。由于厄米萨居民因他的任命而发起骚乱——因为他普遍被指控研究和实践占星术——他逃到劳迪塞雅，找到乔治，乔治提供了关于他的许多历史细节。乔治带他到安提约基雅，通过普拉基图斯和纳尔基苏斯使他再次被带回厄米萨；但后来他又被指控持有撒伯里乌斯观点。乔治更详细地描述了他祝圣的情况，并在结尾补充说，皇帝带他一起远征外邦人，以及他手中行了奇迹。我认为这里已充分转述了乔治关于厄米萨的欧瑟伯的记述。\par

% structure: p0022 source=h2 target=subsection reason=relative-heading-rank; base=h2

\subsection{第十章 安提约基雅的主教会议：厄米萨的欧瑟伯拒绝接受亚历山大里亚主教职后，祝圣了额我略，并更改了尼西亚信经的措辞。}

那时，欧瑟伯被推举后因害怕前往亚历山大里亚，安提约基雅的主教会议便指定额我略为那教会的主教。此事完成后，他们更改了信经；并非谴责尼西亚所订立的信经中的任何内容，而实际上决心通过频繁的会议和发表各种信仰阐明，逐步推翻并废除\OriginalTerm{同质}{consubstantiality}的道理，以建立亚略派的观点。这些事如何结局，我们将在叙述过程中表明；但当时颁布的关于信仰的书信如下：\par

\Quote{我们没有成为亚略的追随者——因为我们是主教，怎能受一个司铎引导呢？——我们也没有接受任何从起初订立的信仰以外的信仰。但我们作为他观点的审查者和判断者，承认其正确性，而不是从他那里采纳；你们将从我们即将陈述的内容中认识到这一点。我们从起初学会相信一位宇宙的天主，一切可思及和可感知事物的创造者和保存者；并相信一位天主的独生子，在万世之前就已存在，与生他的圣父共同存在，万物无论可见不可见都是藉着他而造的；他在末世按照圣父的美意，降来并从童贞圣母取了肉身，完全完成了圣父的旨意，即受苦、复活、升天，坐在圣父的右边，并且将要审判生者死者，永远为王并为天主。我们也相信圣神。如果有必要加上这一点，我们相信肉身的复活和永生。}\par

他们在第一封书信中如此写后，便派人送至各城的主教。但在安提约基雅停留一段时间后，仿佛要谴责前者，他们又发布了另一封信，内容如下：\par

\Quote{\Emph{另一篇信仰阐明。}}\par

依据福音和宗徒的传统，我们相信唯一的天主，全能的圣父，宇宙的创造者和形成者。并相信唯一的主耶稣基督，他的圣子，天主独生子，万物藉着他而受造：在万世之前由圣父所生，天主出自天主，全体出自全体，独一出自独一，完美出自完美，君王出自君王，主出自主；永生的圣言、智慧、生命、真光、真理之路、复活、牧者、门；不变不动摇；是圣父神性、实体、能力、旨意和荣耀的不变肖像；生在\Quote{一切受造之前}；他起初与天主同在，是天主圣言，正如\WorkTitle{福音}所记载：\BibleRef{若 1:1} \BibleQuote{圣言与天主同在，圣言就是天主}，万物藉着他而受造，万有在他内得以存立；他在末世自高天降下，按照圣经由童贞女所生，成为人，成为天主与人之间的中保，我们信德的宗徒，生命的元首，正如他所说：\BibleRef{若 6:38} \BibleQuote{我从天降下，不是为执行我的旨意，而是为执行派遣我者的旨意。}他为我们受苦，第三天为我们复活，升了天，坐在圣父的右边；并将带着荣耀和权能再来审判生者死者。我们也相信圣神，他赐给信徒，为安慰、圣化和成全他们；正如我们的主耶稣基督吩咐他的门徒说：\BibleRef{玛 28:19} \BibleQuote{你们去使万民成为门徒，因父及子及圣神之名给他们授洗}；即父是真实的父，子是真实的子，圣神是真实的圣神，这些称谓并非随意或无意义地使用，而是准确表达了所命名各位的独特存在、荣耀和次序：如此，位格有三，但和谐为一。因此，在天主和基督面前持守此信德，我们弃绝一切异端和虚假教义。若有人违反圣经纯正正确的信德而教导，声称在圣子存在之前有一个时期或时代，就让他受诅咒。若有人说圣子是受造物如同受造物之一，或是后裔如同后裔之一，并且不持守上述如神圣圣经所传授给我们的每一条教义；或若有人教导或宣讲任何其他与我们领受的相悖的教义，就让他受诅咒。因为我们真诚而无保留地相信并追随一切由先知和宗徒从神圣圣经传给我们的一切。\par

这就是当时在安提约基雅集会者所发表的信仰声明，\Person{额我略}作为亚历山大里亚主教也签署了，尽管他尚未进入该城。会议做完这些事并制定了一些其他法规后便解散了。此时，公共事务也发生了动荡。称为法兰克人的民族侵入高卢的罗马领土，同时东方发生了强烈地震，尤其在安提约基雅，那里整整一年持续受到震动。\par

% structure: p0029 source=h2 target=subsection reason=relative-heading-rank; base=h2

\subsection{第十一章  关于\Person{额我略}在军队护送下到达亚历山大里亚，\Person{亚大纳削}逃离}

此后，军事指挥官\Person{叙利安}和五千名重装士兵将\Person{额我略}护送到亚历山大里亚；城中那些有亚略派思想的人也加入他们。但在此处应叙述\Person{亚大纳削}在被迫离开教堂后如何逃脱抓捕者之手。当时是傍晚，人们在那里参加守夜礼，期待举行礼仪。指挥官到达，在教堂四周布置兵力列阵。\Person{亚大纳削}看到所发生的事，心中思量如何避免会众因他而遭受任何损害：于是他指示执事通知祈祷，之后命令诵读圣咏；当圣咏的优美歌声响起时，所有人都从教堂的一个门出去。在此期间，军队袖手旁观，\Person{亚大纳削}就这样在咏唱圣咏的人群中安然逃脱，并立即赶往罗马。于是\Person{额我略}在教堂得势；但亚历山大里亚人民对此义愤填膺，焚烧了称为\Person{狄约尼削}堂的教堂。此事就此为止。\Person{欧瑟伯}至此已实现其目的，便派遣代表团前往罗马主教\Person{儒略}处，请求他亲自审查对\Person{亚大纳削}的指控，并命令在他面前进行司法调查。\par

% structure: p0031 source=h2 target=subsection reason=relative-heading-rank; base=h2

\subsection{第十二章  君士坦丁堡人民在\Person{欧瑟伯}死后恢复\Person{保禄}的教座，而亚略派人选出\Person{玛克多尼}}

但\Person{欧瑟伯}未能活到得知\Person{儒略}关于\Person{亚大纳削}的决定，因为他在那次会议后不久就去世了。于是人民再次将\Person{保禄}引入君士坦丁堡教堂；但亚略派同时在那座献给\Person{保禄}的教堂内祝圣了\Person{玛克多尼}。这是那些曾与\Person{欧瑟伯}（那个扰乱公共安宁者）合作的人促成的，他们接管了他的一切权力。他们是尼开亚主教\Person{德奥尼}、加采东的\Person{玛里斯}、色雷斯赫拉克利亚的\Person{戴奥多}、上米西亚辛吉杜努姆的\Person{乌尔萨奇}和上潘诺尼亚穆尔萨的\Person{瓦伦斯}。\Person{乌尔萨奇}和\Person{瓦伦斯}后来确实改变了观点，向主教\Person{儒略}提交了书面撤回声明，因此在签署\OriginalTerm{同质}{consubstantiality}教义后重新获准共融；但那时他们热烈支持亚略派错误，是教会中最激烈冲突的煽动者，其中一次与\Person{玛克多尼}在君士坦丁堡有关。因基督徒间的内部战争，该城不断发生叛乱，许多人因此丧生。\par

% structure: p0033 source=h2 target=subsection reason=relative-heading-rank; base=h2

\subsection{第十三章  \Person{保禄}因\Person{君士坦修}的将军\Person{厄尔莫根}被杀，再次被\Person{君士坦修}逐出教堂}

\Emph{消息}传到了皇帝君士坦提乌斯耳中，当时他驻跸安提约基雅。于是，他命令已派往色雷斯的将军赫尔莫杰内斯在途中经过君士坦丁堡，将保禄逐出教会。赫尔莫杰内斯抵达君士坦丁堡后，试图驱逐主教们，使全城陷入混乱；因为民众立即起来暴动，急于保护主教。当赫尔莫杰内斯坚持用武力驱逐保禄时，民众如常激愤，猛烈攻击他，纵火烧了他的住宅，将他拖过全城，最后将他杀死。此事发生在两位奥古斯都担任执政官期间——即君士坦提乌斯第三次执政，君士坦斯第二次执政；当时君士坦斯已征服法兰克人，迫使他们与罗马人缔结和约。皇帝君士坦提乌斯得知赫尔莫杰内斯被杀的消息后，立即从安提约基雅骑马出发，抵达君士坦丁堡，立刻驱逐了保禄，然后通过对居民施加惩罚，取消了其父亲拨给免费分发给他\OriginalTerm{臣民}{citizens}的每日粮食配给中超过四万麦子量的份额；因为在此灾难之前，将近八万麦子量从亚历山大里亚运来的小麦已分发给市民。然而，他迟迟未能批准将马其顿尼任命为该城主教，不仅因为马其顿尼未经他同意而被祝圣，而且由于马其顿尼与保禄的竞争，他的将军赫尔莫杰内斯和许多其他人被杀。但皇帝允许马其顿尼在他被祝圣的教堂中供职，然后返回了安提约基雅。\par

% structure: p0035 source=h2 target=subsection reason=relative-heading-rank; base=h2

\subsection{亚流派将格列高利从亚历山大里亚主教座撤职，并任命乔治接替其职位。}

大约在同一时期，亚流派将格列高利从亚历山大里亚主教座撤职，理由是他不得人心，同时因为他放火烧了一座教堂，且在促进本派利益方面不够热心。因此，他们将乔治引入他的主教座；乔治是卡帕多细亚人，以其教义的有力辩护者而闻名。\par

% structure: p0037 source=h2 target=subsection reason=relative-heading-rank; base=h2

\subsection{亚他那修和保禄前往罗马，获得尤利乌斯主教的信函后，恢复了各自的教区。}

与此同时，亚他那修经过长途跋涉，终于抵达意大利。帝国西部当时由君士坦丁最小的儿子君士坦斯单独统治，其兄君士坦丁已被士兵杀害，如前所述。同时，君士坦丁堡主教保禄、迦萨的阿斯克勒帕斯、小加拉达安基拉的马尔塞鲁斯以及阿德里亚诺堡的路济乌斯，因各种指控被逐出各自的教会后，也抵达了帝国之都。在那里，每人将案件呈交罗马主教尤利乌斯。尤利乌斯凭借罗马教会的特殊特权，将他们送回东方，并附上推荐信加以强化；同时恢复每人自己的职位，并严厉斥责那些将他们免职的人。依靠尤利乌斯主教的签署，主教们离开罗马，重新接管自己的教会，并将信函转交收信人。这些人认为受到尤利乌斯斥责的侮辱，便在安提约基雅召开会议，口述了对尤利乌斯信函的回复，作为整个主教会议的一致意见。他们说，裁决任何被他们逐出教会的人不是他的职权；因为当诺瓦图斯被逐出教会时，他们并未反对他。东方教会的主教们将这些事告知了罗马主教尤利乌斯。但是，当亚他那修进入亚历山大里亚时，亚流派乔治的支持者引发骚乱，据称多人被杀；由于亚流派试图将此事的全部责任推给亚他那修，认为他是始作俑者，我们有必要对此稍加评论。万民的审判者天主知道这些混乱的真正原因；但有经验者不会不知道，此类致命事故大多是民众派系运动的伴随物。因此，亚他那修的诽谤者将责任归于他是徒劳的，尤其是马其顿派异端的主教撒彼努斯。因为如果后者反思亚他那修连同其他持\OriginalTerm{同质}{consubstantiality}论者所受亚流派的诸多不公，或因亚他那修之事而召开的多次会议对此事的诸多控诉，或简而言之，那位异端魁首马其顿尼本人在各教会所行之事，他要么完全沉默，要么被迫发言也会说更合理的话，而非这些责难。但他有意忽视这一切，故意歪曲事实。然而，他绝口不提异端魁首，极力隐瞒他所知此人犯下的滔天罪行。更令人惊异的是，他未曾对亚流派说过一句不利之言，尽管他远非赞同他们的观点。马其顿尼的祝圣典礼——他接受了其异端观点——他也缄默不语；因为如果他提及，必然也会记录其不敬之举，而这些在当日表现得极为明显。就此打住。\par

% structure: p0039 source=h2 target=subsection reason=relative-heading-rank; base=h2

\subsection{君士坦提乌斯皇帝通过命令近卫军长官斐理伯，确保保禄被流放，马其顿尼就任其主教座。}

当当时驻跸安提约基雅的君士坦提乌斯皇帝听说保禄再次获得了主教之位时，对他的僭越行为极为震怒。于是他发出一道书面命令给禁卫军长官菲利普——其权力超过其他行省总督，被称为仅次于皇帝的第二人——要他把保禄再次逐出教堂，并让马其顿尼取代他。这位长官菲利普害怕民众发生暴动，便用计谋诱捕主教：他保守皇帝的谕令秘密，前往名为齐克西普斯的公共浴场，假借处理一些公务，以极为尊敬的态度派人去请保禄，声称他的到场不可或缺。主教来了；当他应召前来时，长官立即向他出示皇帝的命令；主教耐心地接受了未经审讯的定罪。但菲利普害怕群众的暴力——因为许多人已经聚集在浴场周围，想看会发生什么事，他们的怀疑因流传的传闻而被激起——他命令打开一扇通往皇宫的浴场门，保禄被带出去，送上预先准备好的船只，立即被放逐。长官指示他去马其顿的首府得撒洛尼，那是他祖先的故乡；命令他居住在该城，但允许他访问伊利里库姆的其他城市，同时严格禁止他进入东罗马帝国的任何部分。这样，保禄出乎意料地立刻被逐出教堂和城市，再次被匆匆流放。皇帝的长官菲利普离开浴场，立刻前往教堂。与他同行的是马其顿尼，他坐在长官的同一辆战车上，与长官同坐，人人都看得见，周围有拔刀的卫兵。群众被这一景象完全震慑，阿里乌斯派和同质派都急忙赶往教堂，每个人都竭力想挤进去。当长官和马其顿尼靠近教堂时，一种非理性的恐慌抓住了群众甚至士兵自己；因为人群如此众多，没有空间让长官和马其顿尼通过，士兵试图用武力推开人群。但他们拥挤的空间使得无法后退，士兵们以为群众故意阻挡通路；于是他们开始拔出刀剑，砍倒挡路的人。据说这次事件中约有3150人被屠杀；其中大部分倒在士兵的武器下，其余的人在群众拼命躲避暴力的践踏中被踩死。取得了如此辉煌的成就之后，马其顿尼似乎没有造成任何灾难，对发生的一切毫无罪责，被长官安置在主教的座位上，而非依照教规。就这样，通过教堂里如此多的谋杀，马其顿尼和阿里乌斯派夺取了教会的统治权。大约在同一时期，皇帝建造了名为\Place{圣索菲亚}的大教堂，毗邻那座名为\Place{圣伊雷内}的教堂——后者原本规模较小，由皇帝的父亲扩大并装饰。如今两者都在同一围墙内，共有一个名称。\par

% structure: p0041 source=h2 target=subsection reason=relative-heading-rank; base=h2

\subsection{第十七章 亚大纳削因皇帝威胁而畏惧，再次返回罗马。}

此时，亚略派信徒又捏造了另一项针对\Person{阿塔纳修斯}的指控，他们编造了如下借口。两位奥古斯都的父亲早已赐予\Place{亚历山大}城教会谷物津贴，以救济贫困。他们断言，\Person{阿塔纳修斯}通常将这些谷物出售，并将收益据为己有。皇帝听信了这诽谤性报告，以死亡威胁\Person{阿塔纳修斯}作为惩罚；\Person{阿塔纳修斯}对这一威胁感到惊慌，便逃亡并躲藏起来。当\Place{罗马}主教\Person{尤利乌斯}得知亚略派信徒对\Person{阿塔纳修斯}的这些新阴谋，并收到了当时已故的\Person{尤西比乌斯}的信函后，他查明了\Person{阿塔纳修斯}藏身之处，邀请受迫害的\Person{阿塔纳修斯}前来。此前不久在\Place{安提约基雅}聚集的主教们的信函也刚刚到达他手中；同时还有来自\Place{埃及}主教们的信函，向他保证对\Person{阿塔纳修斯}的全部指控都是捏造的。收到这些相互矛盾的报告后，\Person{尤利乌斯}首先回复了从\Place{安提约基雅}写信给他的主教们，抱怨他们信中表现出的尖刻情绪，并指责他们违反教规，因为他们没有请求他出席主教会议，因为教会法律要求各教会不得通过违反\Place{罗马}主教意见的决议；然后他严厉谴责他们暗中试图歪曲信仰；此外，他们在\Place{提洛}的先前行为是欺诈性的，因为对\Place{马雷奥提斯}所发生之事的调查只听取了单方面陈述；不仅如此，关于\Person{阿尔塞尼乌斯}的指控明显被证明是诬告。\Person{尤利乌斯}在给聚集在\Place{安提约基雅}的主教的回信中写了诸如此类的意见；我们本应将它们连同写给\Person{尤利乌斯}的那些信函全文收录于此，若不是它们的冗长妨碍了我们的目的。然而，我们之前提到的\OriginalTerm{马其顿异端}{Macedonian heresy}的辩护者\Person{萨比努斯}并未将\Person{尤利乌斯}的信函收入其\WorkTitle{主教会议汇编}中；尽管他没有省略\Place{安提约基雅}主教们写给\Person{尤利乌斯}的信函。不过，这是他的一贯做法；他小心翼翼地收录那些不曾提及或完全排斥\OriginalTerm{同质}{homoousion}一词的信函，而有意对相反倾向的信函缄默不提。关于此事，说到这里就够了。此后不久，\Person{保禄}假称从\Place{得撒洛尼}前往\Place{格林多}，来到\Place{意大利}；于是两位主教向该地区的皇帝上诉，将各自的情况呈报给他。\par

% structure: p0043 source=h2 target=subsection reason=relative-heading-rank; base=h2

\subsection{第十八章。西方皇帝请求其兄弟派遣三人，以说明阿塔纳修斯与保禄被废黜的原因。所派之人发布了另一形式的信经。}

西方皇帝得知他们的情况后，同情他们的遭遇，并写信给他的兄弟\Person{君士坦提乌斯}，请求他派遣三位主教，以便向他说明废黜\Person{阿塔纳修斯}和\Person{保禄}的原因。应此请求，基利家人\Person{纳尔基苏斯}、色雷斯人\Person{特奥多雷}、\Place{加采东}的\Person{玛里斯}、叙利亚人\Person{马尔克}被委派执行此任务；他们抵达后，拒绝与\Person{阿塔纳修斯}或其友人进行任何交流，反而压制了在\Place{安提约基雅}颁布的信经，向皇帝\Person{君士坦斯}呈上了一份由他们自己拟定的另一份信仰宣言，内容如下：\par

\Emph{另一篇信仰阐述。}\par

我们信唯一的天主，全能的圣父，万物的创造者和制造者，诸天和地上的一切家族都藉祂得名，\BibleRef{弗 3:15}；并信祂的独生子，我们的主\Person{耶稣基督}，祂于万世之前由圣父所生；出自天主的天主；出自光的光；藉着祂，天上地下，有形无形的万物，皆被造成：祂是圣言、智慧、德能、生命、真光：祂在末期为了我们成为人，生于圣童贞；被钉十字架，受死；被埋葬，第三日从死者中复活，升天，坐在圣父的右边，将于时期的终结来临，审判生者死者，并按照各人的行为施行报应：祂的王国永存，延续至无穷的世代；因为祂坐在圣父的右边，不但在今世，也在来世。我们信圣神，即\OriginalTerm{护慰者}{Comforter}，主在升天后，按照祂的许诺，派遣给祂的宗徒们，为教导他们，并使他们想起一切：藉着祂，那些真心相信祂的人的灵魂得以圣化；凡主张圣子是由不存在之物造成，或出于另一实体而非出于天主，或曾说祂不存在的时候，\OriginalTerm{天主教会}{Catholic Church}视其为外人。\par

他们将这信经呈交皇帝，并向许多其他人展示后，便离开了，不再理会其他事情。但当西方与东方教会之间仍保持着不可分的共融时，在\Place{伊利里库姆}的\Place{西尔米乌姆}城又兴起另一个异端；因为\Person{福提努斯}——他主管该地区的教会，是\Place{小加拉太}人，且是被废黜的\Person{马尔凯鲁斯}的门徒——采纳了其师的说法，主张天主子只是一个凡人。不过，我们将在适当的地方更充分地讨论此事。\par

% structure: p0048 source=h2 target=subsection reason=relative-heading-rank; base=h2

\subsection{第十九章。东方主教们致意大利主教们的信经，称为长篇信经。}

在上述事件过去约三年后，东方主教们再次召开主教会议，并拟定了另一形式的信经，由当时\Place{盖尔马尼基亚}的主教\Person{尤多克修斯}、\Person{玛尔提里乌斯}和\Place{基利家}的\Place{莫普苏埃斯提亚}的主教\Person{马其顿尼乌斯}之手传送给意大利的主教们。这份信经的表述采用了更长的形式，包含了许多对之前信经的增加，其内容如下：\par

\Quote{我们信唯一的天主，全能的圣父，万物的创造者与造作者，天上和地上的一切家族都从他得名；也信他的独生子、我们的主耶稣基督，他在万世之前由圣父所生，是出自天主的天主，出自光明的光明；藉着他，天地间一切有形无形的事物都被造成。他是圣言、智慧、德能、生命、真光。他在末世为了我们而降生成人，由童贞玛利亚诞生，被钉十字架，受死，埋葬，第三日从死者中复活，升天，坐在圣父的右边，将在时代的终末来临，审判活人死人，按各人的行为予以报应：他的王国永久长存，延续无尽；因为他坐在圣父的右边，不仅在此世，也在来世。我们也信圣神，即护慰者，主在他升天之后按他的许诺派遣给宗徒们，教导他们、使他们记起一切；藉着圣神，那些真诚信主的人的灵魂得以圣化。但那些断言圣子是由不存在的事物、或由另一种\OriginalTerm{本体}{substance}、而非出自天主所造，或者断言曾有一个时间或时代他并不存在的人，圣而公教会视之为外人。圣而公教会也诅咒那些说有三个天主，或说基督在万世之前不是天主，或说他既不是基督也不是天主子，或说父、子、圣神是同一位置，或说圣子\OriginalTerm{非受生}{unbegotten}，或说父不是凭自己的意志或意愿生了圣子的人。断言圣子是从不存在的事物而有其存在，也是不安全的，因为这在天主默感的圣经中从未论及他。我们也未受教导说他从父以外的任何其他先在本体而有其存有，而是他真的单独由天主所生；因为神圣的圣言教导我们只有一个\OriginalTerm{非受生}{unbegotten}的\OriginalTerm{无始}{inoriginate}的元始，即基督的父。但那些未经圣经授权而鲁莽断言\InnerQuote{他曾经不存在}的人，不应预先设想任何时间间隔，而只应设想那无时间地生了他的天主；因为时间和时代都是藉着他造成的。然而不可认为圣子与圣父同\OriginalTerm{无始}{inoriginate}或同\OriginalTerm{非受生}{unbegotten}：因为对于同无始或同非受生者，严格来说并无父。但我们知道，唯有圣父是\OriginalTerm{无始}{inoriginate}、莫可名状的，他不可言说地、莫可名状地生了圣子，而圣子在诸时代之前受生，但不像父那样非受生，而是有一个开始，即生他的父，因为\Quote{基督的头是天主}\BibleRef{格前 11:3}。如今，虽然我们按圣经承认三个\OriginalTerm{位格}{person}或事物，即父、子、圣神，我们并不因此设立三个天主：因为我们知道只有一个天主，他本身是完美的、非受生的、无始的、不可见的，是独生子的天主和父，他独自从自身而有存在，并独自丰盛地赐予一切其他事物以存在。但我们既断言只有一个天主，即我们主耶稣基督的父、独生子，我们并不因此否认基督在万世之前是天主，像\Place{萨莫萨塔}的\Person{保禄}的追随者那样，他们声称他在降生成人后因被高举而神化，因为他本性上只是一个凡人。我们确实知道他服从于他的天主和父；然而他由天主所生，按本性是真实而完美的天主，并非后来由人变成天主；而是为了我们由天主变成人，且从未停止是天主。此外，我们憎恶并诅咒那些虚妄地称他仅为天主非实体的圣言，只在另一个人内存在，如同发出的话语或心中的意念；并声称在万世之前他既不是基督、天主子、中保，也不是天主的肖像；而是从他取了我们的肉体由童贞女出生之时，大约四百年前，他才成为基督和天主子。因为他们断言基督的王国从那时起开始，并在万物终结与审判之后将结束。这样的人是\Person{玛尔克路}和\Person{福提诺}的追随者，即安奇洛-迦拉达派，他们以确立其主权的借口，像犹太人一样废除了基督的永恒存在和神性，以及他王国的永久性。但我们知道他不仅仅是天主发出或意念中的圣言，而是自存的永生活圣言；是天主子、基督；他不仅在万世之前以临在的方式与父共存并交游，在创造一切有形无形万物时服务于父，更是父的\OriginalTerm{实体}{substance}的圣言，出自天主的天主：因为父曾对他说：\Quote{让我们按我们的肖像，按我们的模样造人。}他亲自显现于众族长，颁布法律，藉先知发言；最终降生成人，向众人彰显了他的父，并统治直到无穷世代。基督并未获得任何新的尊荣；但我们相信他从起初就是完美的，在一切事上肖似他的父；那些说父、子、圣神是同一位置，不敬地认为三个名称指向同一事物和位格的人，我们理当将他们逐出教会，因为藉着降生成人，他们使那莫可名状、不受苦难的父变得可理解、可受苦难。这些人中如罗马人所谓的\OriginalTerm{父受苦派}{Patropassians}，在我们则称为\OriginalTerm{撒伯流派}{Sabellians}。因为我们知道派遣的父留在他自己不变天主性的本性中；而被派遣的基督完成了降生成人的安排。同样地，那些不敬地断言基督不是凭他父的意愿和喜悦而受生，因而归于天主一种非出于选择的非自愿必然性，仿佛他强制地生了圣子，我们视之为最不敬神、与真理隔绝的人，因为他们竟敢论及他作出与我们对天主的普遍概念不一致、且确实违反天主默感圣经的含义的判断。因为我们知道天主自主、自为主宰，所以我们虔敬地主张他凭自己的意志和喜悦生了圣子。我们虔敬地相信关于他所说的话：\Quote{上主在造化之始创造了我，为他的化工开道。}然而我们并不认为他是像被造的受造物或作品一样被造。因为将创造者与他所造的受造物相比，或想象他具有与其本性完全不同的受造物同样的生出方式，是不敬且与教会信仰抵触的：因为圣经教导我们，那独一的独生子是真实而确实地受生的。当我们说圣子出于自身，像父一样生活并自存时，我们并不因此将他与父分离，仿佛我们设想他们被空间和距离的物质性间隔所分开。因为我们相信他们合一而无中介或间隔，彼此不可分离：整个父怀抱圣子，整个圣子依偎并永远安息于父的怀抱。因此，我们信从完全完美、至圣的天主圣三，断言父是天主，子也是天主，但我们不承认两个天主，而只承认一个，这是基于天主性的尊严以及王国的完美融合与合一：父统治万有，甚至统治圣子自己；圣子服从父，但除父之外，统治一切藉他而造并在他之后的事物；并由父的意志丰盛地将圣神的恩宠赐予圣徒们。因为神圣谕示告诉我们，基督所行使的主权特性就在于此。}\par

\Quote{我们自从出版先前的纲要以来，不得不对这信经作出更为详尽的阐述；这并非为了满足虚妄的野心，而是为了在那些不了解我们真实想法的人中，澄清关于我们信德的一切奇怪猜疑。并且，为了使西方居民既知晓异端党派的无耻歪曲，也认识东方主教们关于基督的教会观点——这观点已由神圣默感的圣经之不容置疑的见证所证实——在所有心智纯正的人中。}\par

% structure: p0052 source=h2 target=subsection reason=relative-heading-rank; base=h2

\subsection{第二十章 萨迪卡会议}

西方主教们因语言不同且不理解此阐述，拒绝接受；他们声称尼西亚信经已足够，不愿在此之外浪费时间。但皇帝再次写信要求恢复保禄和亚太纳削各自的教座，却因人民持续骚动而未果——这两位主教便要求召开另一次主教会议，以便他们的案件以及与信德相关的其他问题能由大公会议解决，因为他们显然表明，他们的被罢免无非是为了更容易歪曲信德。因此，在两位皇帝的共同授权下，另一届大公会议在萨迪卡——伊利里库姆的一座城市——召开；一位皇帝通过信件请求如此，另一位东方皇帝欣然同意。在两位奥古斯都之父去世后的第十一年，鲁菲努斯和欧塞比乌斯执政期间，萨迪卡主教会议召开。据亚太纳削所述，约有三百位来自帝国西部的主教出席；但萨比努斯说，只有七十位来自东部的主教出席，其中包括马里奥特的伊斯库拉斯，他是由那些罢免亚太纳削的人祝圣为该地主教。其余的人中，有的假装身体不适；有的抱怨通知时间太短，并将责任归咎于罗马主教尤利乌斯，尽管从会议召集时起已过去一年半；在此期间亚太纳削留在罗马等待主教会议召开。当他们最终在萨迪卡集会时，东方主教们拒绝与西方主教们会面或进行任何商议，除非他们先将亚太纳削和保禄排除在会议之外。但由于萨迪卡主教普洛托革尼斯和西班牙科尔多瓦主教奥西乌斯绝不允许他们缺席，东方主教们立即退出，返回色雷斯的腓立波波利斯召开了单独的会议，在那里他们公开诅咒\OriginalTerm{同质}{homoousios}一词；并将\OriginalTerm{不同质}{anomoion}的观点引入他们的信函中，向各处发送。另一方面，留在萨迪卡的人首先谴责了他们的离去，随后剥夺了亚太纳削控告者的尊严；然后确认了尼西亚信经，拒绝了\OriginalTerm{不同质}{anomoion}一词，更明确地承认了同质论，并将其写入致所有教会的信函中。双方都认为己方行为正确：东方主教们认为，西方主教们支持了他们所罢免的人；而西方主教们则不仅因为那些罢免他们的人在案件审查前就退出，还因为他们是尼西亚信德的捍卫者，而对方竟敢篡改它。因此，他们恢复了保禄和亚太纳削的教座，以及小加拉西亚安基拉的玛尔凯禄——他在很久以前就被罢免，正如我们在前书所述。当时，他确实竭尽全力争取撤销对他的判决，宣称他因误解其书中的某些表述而被怀疑持有撒摩萨塔的保禄的错误。但必须指出，欧塞比乌斯·庞非利乌斯写了三整卷书反对玛尔凯禄，其中他引用了该作者的原话，以证明他与利比亚的撒伯里乌斯和撒摩萨塔的保禄一样，断言主[耶稣]只是一个凡人。\par

% structure: p0054 source=h2 target=subsection reason=relative-heading-rank; base=h2

\subsection{第二十一章 为欧塞比乌斯·庞非利乌斯辩护}

但既然有人试图甚至污蔑欧塞比乌斯·庞非利乌斯本人在其著作中支持亚略的观点，在此略作评论或许并非无关紧要。首先，他曾出席尼西亚会议，该会议定义了\OriginalTerm{同质}{homoousion}的教义，并且他赞同了会议的决议。在\WorkTitle{君士坦丁传}第三卷中，他如此写道：\Quote{皇帝激励众人达成一致，直至使他们在先前有分歧的论点上观点一致；以至于他们在尼西亚会议上就信德问题完全达成了共识。}既然欧塞比乌斯在提及尼西亚会议时说所有分歧都已消除，所有人都达成了共识，那么有何理由假定他本人是亚略派呢？亚略派也当然被误导了，以为他是他们教义的支持者。但或许有人会说，在他的论述中，他似乎采纳了亚略的观点，因为他经常说\Quote{藉着基督}，对此我们应回答说，教会作者常用这种表达方式以及类似表示救主人性之经济（economy）的语句；而且，在所有这些人之前，宗徒就已使用了此类表达，却从未被视为谬误导师。此外，既然亚略竟敢说圣子是受造物，与其他的类似，请看欧塞比乌斯在其反对玛尔凯禄的第一卷书中对此如何论述：\par

\Quote{只有他，没有其他，被宣告为并确实是天主的独生子；因此，任何人都有理由谴责那些胆敢断言他是如同其他受造物一样从无中受造的受造物的人；因为那样他如何能是儿子？倘若他被赋予与其他受造物相同的本性……且他是众多受造物中的一员，既然他像它们一样，在这种情况下也从无中受造，那么他如何能是天主的独生子？但圣经并非这样教导我们。} 他稍后又补充道： \Quote{因此，凡是将圣子定义为由不存在之物所造，且是预先存在的从无中产生的受造物的人，就忘记了：当他承认圣子之名时，他实际上否认他是真正的儿子。因为从无中制造出来的，不能真正是天主子，正如其他被造之物一样；但真正的天主子，既是由圣父所生，就恰当地被称为圣父的独生子和爱子。为此缘故，他本身就是天主；因为天主的后裔，除了是他所生者的完美肖像之外，还能是什么？一位君主固然兴建一座城，却不生它；人们说他生一个儿子，而不是建造一个儿子。一位工匠也可被称为作品的制作者，却不是作品之父；而他绝不能被称为他所生者的制作者。同样，宇宙的天主是圣子的父亲；却可恰当地被称为世界的创造者和制造者。尽管在圣经中曾有一次说：\BibleRef{箴 8:22} \BibleQuote{主创造了我，作为他道路的开始，为他工作的缘故，} 然而我们应当考虑这句话的意义，我将在以后解释；而不应像马塞录那样，因一处经文就危及教会最重要的教义。}\par

欧瑟伯·庞菲鲁斯在\WorkTitle{反马塞录}第一卷中说了这些和许多其他类似的话；在他的第三卷中，他说明该如何理解\Emph{受造物}一词的含义，他说：\par

\Quote{因此，这些事情既然已经如此确立，那么就必须按照与前述相同的意义来理解这句话：\BibleQuote{主创造了我，作为他道路的开始，为他工作的缘故}，这必定是说过的话。因为尽管他说他被创造，但这并不等于说他从无进入存在，也不是说他像其他受造物一样从无中被造——这是某些人错误地假定的；而是说他是存有、活着、先存，并且在整个世界构成之前就已存在；并且被他的主和父任命来统治宇宙：\Quote{创造}一词在这里是代替\Quote{任命}或\Quote{设立}来使用的。确实，宗徒\BibleRef{伯前 2:13}明确地将人中的统治者和长官称为\Quote{受造物}，当他说：\BibleQuote{你们为主的缘故，要服从每个受造的人，无论是君王为最高的，或是执政权者作为他所派遣的。}先知也当他说：\BibleQuote{以色列，预备好呼求你的天主。因为看，那使雷声稳固的，创造圣神，向人宣告他的基督}时，……并没有使用\InnerQuote{创造者}一词来表示从无中制造的意思。因为天主在他向所有人宣告他的基督时，并没有创造圣神，既然\BibleRef{训 1:9} \BibleQuote{太阳之下无新事}；但圣神早已存在，早已有实体：他是在宗徒们聚集时被派来的，当时如雷声一般，\BibleQuote{从天上传来一阵响声，好像暴风刮来；他们就充满了圣神。}这样他们就向所有人宣告了天主的基督，符合那个预言所说的：\BibleRef{亚 4:13} \BibleQuote{看，那使雷声稳固的，创造圣神，向人宣告他的基督}：\InnerQuote{创造}一词是代替\InnerQuote{派遣}或\InnerQuote{任命}来使用的；而雷声在另一种图像中暗示着福音的宣讲。此外，那说\BibleQuote{天主，求你给我再造一颗纯洁的心}的人，并非说他无心，而是祈求他的心智得以净化。同样，也如此说：\BibleRef{弗 2:15} \BibleQuote{为把双方在自己身上造成一个新人}，意思是\Quote{结合}。也请思考这段经文是否同类：\BibleRef{弗 4:24} \BibleQuote{穿上那新人，是按照天主创造的}；以及这段：\BibleRef{格后 5:17} \BibleQuote{所以，若有人在基督里，他就是新的受造物}；以及任何仔细查考与神默感的圣经的人所能找到的其他类似表达。因此，一个人不应感到惊讶，如果在这段经文\BibleQuote{主创造了我，作为他道路的开始}中，\InnerQuote{创造}一词是比喻式地使用，代替\InnerQuote{任命}或\InnerQuote{设立}。}\par

欧瑟伯在他的\WorkTitle{反马塞录}中使用了这样的言辞；我们引用它们是因为那些诽谤性地试图诋毁和控告他的人。他们也无法证明欧瑟伯将存在的起始归于天主子，尽管他们可能发现他经常采用比喻的方式使用这些表达；尤其是因为他效仿并赞赏\Person{奥力金}的作品，在那些作品中，那些能理解奥力金著作深度的人，将会看到处处都声明圣子是由圣父所生。这些评论是顺便作出的，以反驳那些歪曲欧瑟伯的人。\par

% structure: p0060 source=h2 target=subsection reason=relative-heading-rank; base=h2

\subsection{第二十二章 萨底迦会议恢复保禄和亚大纳削的教座；因东方皇帝拒绝接纳他们，西方皇帝以战争威胁他}

凡在撒底迦聚集之人，以及在色雷斯斐理波波里组成另一会议之人，各自完成了他们认为必要之事后，便返回了各自的城邑。自那时起，西方教会与东方教会便分离了；二者共融的界限是称为苏西斯的那座山，它将伊利里亚人与色雷斯人分开。直到这座山，虽有信仰的差异，仍可无区别地共融；但过了此山，他们彼此便不再共融。这就是当时教会的动荡状况。这些事之后不久，西方帝国的皇帝将撒底迦发生之事告知其弟君士坦提乌斯，并恳求他恢复保禄和亚大纳削的教座。但君士坦提乌斯迟迟未予执行，西方皇帝再次写信给他，让他选择：要么恢复保禄和亚大纳削原先的尊位并将教会归还他们；要么，若他不这样做，就视他为敌人，并立即准备战争。他致其弟的信函如下：\par

\Quote{亚大纳削和保禄现在与我同在；经调查，我完全确信他们是为虔敬而受迫害。因此，你若承诺恢复他们的教座，并惩罚那些如此不义地伤害他们的人，我就将他们送给你；但你若拒绝这样做，你应确信，我将亲自前往那里，不顾你的反对，恢复他们自己的教座。}\par

% structure: p0063 source=h2 target=subsection reason=relative-heading-rank; base=h2

\subsection{第二十三章　君士坦提乌斯因惧怕其兄长的威胁，致信召回亚大纳削，并派他前往亚历山大里亚。}

东方皇帝收到此信后陷入困惑；他立即召来大部分东方主教，告知他们其兄弟向他提出的选择，并询问当如何行。他们回答说：宁可让出教会给亚大纳削，也不可发动内战。于是皇帝迫于无奈，召见亚大纳削及其朋友。与此同时，西方皇帝派保禄前往君士坦丁堡，带有两位主教和其他体面随从，并用他自己的信函及会议的函件加强了他的地位。但亚大纳削仍心存畏惧，犹豫不决，不敢前去见皇帝——他害怕诽谤者的阴谋——而东方皇帝不仅一次，而是第二次、第三次邀请他前来；这从其函件中可以清楚看出，这些函件从拉丁文翻译如下：\par

\Emph{君士坦提乌斯致亚大纳削的书信。}\par

君士坦提乌斯·胜利者·奥古斯都致亚大纳削主教。\par

我们富有同情的仁慈不能再容许你像被大海的汹涌波涛一样颠簸不安。我们不懈的虔敬并没有忘记你被逐出家乡、被剥夺财产、在无路可走的荒野中流浪。虽然我已太久没有写信告知你我的心意，期待你自愿前来寻求解决困境的方法；但既然可能恐惧阻碍了你的愿望实现，我们因此向你的尊驾送去充满宽仁的信函，使你能够无所畏惧地迅速来到我们面前，从而在体验到我们的善意之后，达成你的愿望，并恢复你的适当地位。为此，我请求我主和兄弟君士坦斯·胜利者·奥古斯都准许你前来，以便在双方的同意下，你能返回你的家乡，并获得我们眷顾的这一保证。\par

\Emph{另一封致亚大纳削的书信。}\par

君士坦提乌斯·胜利者·奥古斯都致亚大纳削主教。\par

虽然我们在前一封信中已充分告知，你可以放心前来我们的宫廷，因为我们极其渴望恢复你的适当位置，但我们又向你的尊驾发出了这封信。因此我们敦促你，不要有任何疑虑或恐惧，乘坐公车赶紧来到我们这里，以便你能得到你所期望的。\par

\Emph{另一封致亚大纳削的书信。}\par

君士坦提乌斯·胜利者·奥古斯都致亚大纳削主教。\par

当我们驻留在厄得撒时，你的司铎们也在场，我们乐意派他们中的一位到你那里去，以催促你尽快前来我们的宫廷，以便在引见给我们之后，你可以立即前往亚历山大里亚。但自从你收到我们的信以来，已经过了相当长的时间，而你还未到来，因此我们现提醒你迅速亲自前来见我们，这样你就能返回家乡，实现你的愿望。为进一步保证我们的意图，我们派了执事阿凯塔斯到你那里，从他那里你将领会我们对你心意，并且你将能确保你所期望的，即我们乐于促成你所追求的目标。\par

亚大纳削在阿奎莱亚收到这些信函时——因为他从撒底迦离开后便住在那里——他立即赶往罗马；并将这些信函给主教犹利乌斯看了，在罗马教会引起了极大的喜乐。因为东方皇帝似乎也承认了他们的信仰，既然他召回了亚大纳削。于是犹利乌斯代表亚大纳削写信给亚历山大里亚的圣职人员和教友如下：\par

\Emph{罗马主教犹利乌斯致亚历山大里亚人的书信。}\par

犹利乌斯主教致住在亚历山大里亚的司铎、执事和子民，亲爱的弟兄们，在主内问安。\par

我也与你们一同喜乐，亲爱的弟兄们，因为你们终于亲眼看见了你们信德的果实。事实上，任何人关于我的兄弟兼同为主教\Person{亚大纳削}，都能看出这一点：天主因他纯洁的生活并应你们的祈祷，已将他归还给你们。由此明显可见，你们向天主的祈求始终纯洁而充满爱德；因为你们牢记天主的应许以及与之相连的爱德——这是你们从我兄弟的教导中学到的——你们确实知道，并且按照你们内中纯正的信德清楚地预见到，你们的主教不会永远与你们分离，你们在虔诚的心中仿佛他常与你们同在。因此，我无需对你们多言，因为你们的信德已经先于我所能说的一切；你们众人的共同祈祷已按基督的恩宠得以成就。所以我与你们一同喜乐，并重申你们在信德中保全了你们的灵魂，无可战胜。我与我的兄弟\Person{亚大纳削}也一同喜乐，因为他在遭受诸多苦难时，从未忘记你们的爱德与渴望；虽然他表面上暂时离开你们，但在精神上却常与你们同在。此外，亲爱的弟兄们，我深信他所承受的一切试炼并非不光荣；因为你们的信德和他自己的信德都由此受到考验并向众人显明。如果他没有遭受如此多的磨难，谁会相信你们对这位卓越的主教怀有如此崇高的敬意与爱德，或者他拥有如此杰出的美德——并且也因着这些美德，他在天上的望德绝不会落空呢？因此，他获得了宣认的见证，在今世和来世都极其光荣。因为他经历了陆地和海上如此众多且多样化的试炼，他已践踏了亚略异端的种种阴谋；尽管屡次因嫉妒而面临危险，他却藐视死亡，蒙全能天主和我们的主耶稣基督保护，始终坚信他不仅能逃脱（敌人的）诡计，而且会为了安慰你们而得以恢复，并同时从你们自己的良心中带回更大的胜利。如此，他因纯洁的生活、坚定的目标、对天上教义的恒心而受到认可，这一切都有你们不变的敬爱与爱德作证，他因此得以声名远扬，直到地极，堪受光荣。因此，他回到你们身边，现在比离开时更为显赫。因为火能试炼贵金属（我说的是金银）以净化，那么对于这样一位卓越的人，我们该说什么才与他的价值相称呢？他在克服了如此多灾难与危险的火炼之后，如今被归还给你们，不仅由我们，而且由全体主教会议宣告无罪。所以，亲爱的弟兄们，请以虔敬的尊荣与喜乐，迎接你们的主教\Person{亚大纳削}，以及那些在磨难中与他同行的人。你们因达到祈祷的目标而喜乐，你们曾用支持的书信供给你们那饥渴慕义（几乎可以说）的牧者饮食，滋养他的精神福祉。事实上，当他在异乡寄居时，你们安慰了他；当他被阴谋陷害并遭受迫害时，你们以最忠信的情感养育了他。至于我，仅仅是凭想象描摹你们每个人在他归来时的喜悦、群众虔诚的问候、聚集迎接他的人的光荣庆典，以及那一天的全部景象——当我的兄弟被带回你们身边时，过去的苦难将终结，他珍贵而渴盼的归来将使所有心灵汇聚于最热烈的喜乐表达——这已令我幸福。这种感受要在最高程度上延及我们，我们视之为天主恩宠的标志，使我们有幸得以认识如此卓越的人物。因此，我们宜以祈祷结束这封信。愿全能天主和他的圣子我们的主救主耶稣基督不断赐予你们这项恩宠，以此奖赏你们对主教所表现出的令人钦佩的信德，并赐予你们显赫的见证：使那些最卓越的事物，即\BibleQuote{眼所未见，耳所未闻，人心所未想到的，就是天主为爱祂的人所预备的}\BibleRef{格前 2:9}，在来世等待你们和你们的亲人，以上主耶稣基督为凭借；愿光荣归于全能天主，至于无穷之世。阿们。亲爱的弟兄们，愿你们坚强。\par

亚大纳削倚赖这些信函，抵达了东方。当时君士坦提乌斯皇帝并非怀着敌意接待他；然而在亚略派信徒的唆使下，他仍设法图谋欺骗他，对他说：\Quote{你已按照主教会议的法令并经我们同意，恢复了你教座的职位。但因有些亚历山大城的人拒绝与你共融，望你允许他们在城内保留一所圣堂。}亚大纳削立即回应道：\Quote{陛下，您有权下令并执行您所喜悦的一切。因此，我也恳求您赐我一项恩惠。}皇帝欣然允诺后，亚大纳削立刻补充说，他希望皇帝也能赐予他自己所要求的事，即：在每个城市中，为那些拒绝与亚略派信徒共融的人保留一所圣堂。亚略派信徒看出亚大纳削的意图不利于他们，便说此事可改日再议；但他们任凭皇帝自行处置。于是皇帝将各自的教座归还给亚大纳削、保禄和玛尔塞路；同样也归还给迦萨的主教阿斯克勒帕斯，以及哈德良堡的路济安。因为这些人也已被萨迪卡主教会议所接纳：阿斯克勒帕斯出示了记录，证明欧瑟伯·庞非禄曾会同数人审理他的案件后，恢复了他先前的品秩；路济安则因控诉者逃逸而获接纳。于是皇帝的谕令被送至各城，命令居民欣然接纳他们。的确，在安其拉城，当巴西略被逐、玛尔塞路接替时，发生了相当大的骚乱，这成了敌人诽谤他的口实；但迦萨人民却欣然接纳了阿斯克勒帕斯。君士坦丁堡的马其顿纽暂时让位于保禄，自行在城内别一所圣堂内分别召集聚会。此外，皇帝还写信给主教、神职人员和教友，论及欣然接纳亚大纳削一事；同时另以信件下令，凡在法庭上对他所作的任何判决均应废除。有关这两件事的文书如下：\par

\WorkTitle{君士坦提乌斯为亚大纳削辩护之信函}\par

凯旋者君士坦提乌斯·马克西穆斯·奥古斯都，致天主教会的主教与司铎。\par

至可敬的亚大纳削主教并未被天主的恩宠遗弃。虽然他在人前暂时受到了考验，却从全知的天主眷顾中获得了应有的豁免；他已被天主的旨意和我们的决定，恢复到他先前受天主许可所治理的家乡和教会。因此，我们的仁慈理应对与此相符之事予以适当关注：凡先前对与他共融者所作的一切裁决，如今均应废除；今后对他的一切猜疑都应止息；与他同在的圣职人员的原有豁免权，理当予以确认。此外，我们以为还应在我们对他的恩惠中增加一点，即整个教会团体应明白，所有依附于他的人——无论主教或其他圣职人员——都受到保护；与他共融即足以证明各人意图纯正。故此，我们已下令，按照先前眷顾的相似方式，凡有智慧以更健全的判断加入正确一方并选择与他共融者，均可享有我等如今依天主旨意所赐之宽免。\par

\WorkTitle{另一封致亚历山大城人的信函}\par

凯旋者君士坦提乌斯·马克西穆斯·奥古斯都，致亚历山大城天主教会的人民。\par

我们以你们各方面的良好秩序为目标，并知悉你们久已失去主教照管，遂认为理应将你们的亚大纳削主教——此人以其生活与举止之正直圣洁而为众人所周知——再次送回给你们。你们当以常有的合宜礼遇接待他，并立他为你们向天主祈祷的助手，务须按照教会法典时时保持和睦与平安，这既能使你们自己蒙受荣耀，也令我们欣慰。因为若在你们中间激起任何分裂或党派之争，与当今繁荣时代相悖，实不合理；我们相信这样的不幸必将从你们中间完全消除。因此，我们劝勉你们，如我们先前所说，藉他的帮助，恒心持守你们惯常的热心祈祷；如此，当你们的这个决定广为人知，甚至那些仍奴役于无知偶像崇拜的外邦人，因进入众人的祈祷中，也会急于寻求我们神圣宗教的知识，最亲爱的亚历山大城人啊。我们再次劝勉你们留意这些事：衷心欢迎你们的主教，视他为按天主旨意和我们谕令所任命者；当以你们心灵的一切热情拥抱他，视他堪受此爱戴。因为此事符合你们身份，也与我们的仁慈一致。然而为抑制好乱之徒的煽动和骚乱倾向，已对我们的法官下令，凡发现煽动者，应施以法律严惩。因此，你们当留意我们的决定和天主的决定，以及我们为促使你们和谐所怀的焦虑，也纪念对无序者将施的惩罚，务使你们的行为符合神圣宗教的制裁，以一切尊敬荣耀你们的主教；好使你们能与他一同向万有的天主及圣父呈献祈祷，为你们自己，也为全人类的治理井然有序。\par

\WorkTitle{论废除针对亚大纳削的判决之信函}\par

凯旋者君士坦提乌斯·奥古斯都，致涅斯托里乌斯,并同样致奥古斯塔姆尼卡、特拜斯和利比亚的诸总督。\par

如果发现以往任何时候曾通过任何不利于或贬损与主教\Person{亚他那修}共融者的法令，则我们之意愿是，该令应如今完全废除；其圣职人员应重新享有从前授予他们的同样豁免。我们命令严格遵守此令，以期既然主教\Person{亚他那修}已恢复其教会，所有与他共融者皆可拥有如前之特权，并一如其他神职人员现今所享有者；如此，他们的事务得到妥善安排，亦可分享普遍的繁荣。\par

% structure: p0088 source=h2 target=subsection reason=relative-heading-rank; base=h2

\subsection{第二十四章. \Person{亚他那修}经\Place{耶路撒冷}返回\Place{亚历山大}，被\Person{玛克西穆}接纳共融；在该城召集的主教会议确认\WorkTitle{尼西亚信经}。}

主教\Person{亚他那修}有了这等信函作保障，便穿过\Place{叙利亚}，来到\Place{巴勒斯坦}。到达\Place{耶路撒冷}后，他将\WorkTitle{撒尔底加会议}所行之事以及\Person{君士坦提乌}皇帝已确认其决议之事告知主教\Person{玛克西穆}；随后他提议应在那里召开一次主教会议。于是\Person{玛克西穆}立即派人去请\Place{叙利亚}和\Place{巴勒斯坦}的几位主教，召集了一次会议，恢复了\Person{亚他那修}的共融及其先前的地位。此后，该会议致函\Place{亚历山大}城及\Place{埃及}与\Place{利比亚}的所有主教，报告了关于\Person{亚他那修}的决定。因此，\Person{亚他那修}的对手们大大嘲笑\Person{玛克西穆}，因为他先前协助罢黜他，却忽然改变主意，仿佛先前什么事都未发生，投票恢复他的共融和品级。\Person{乌撒基乌}和\Person{瓦伦斯}——他们曾是亚略主义的激烈支持者——得知这些事后，谴责了他们从前的热忱，前往\Place{罗马}，向主教\Person{尤利乌}递交了他们的撤回书，并赞同\OriginalTerm{同质}{consubstantiality}的道理；他们还写信给\Person{亚他那修}，表示今后愿意与他共融。于是\Person{乌撒基乌}和\Person{瓦伦斯}当时因\Person{亚他那修}的好运而屈服，被诱导承认正统信仰。\Person{亚他那修}在往\Place{亚历山大}的途中经过\Place{佩鲁西翁}，劝告各城的居民警惕亚略派，只接纳那些宣认\OriginalTerm{同质}{homoousios}信仰的人。他在一些教会中还施行了圣秩祝圣，这给了他另一个指控的理由，因为他竟然在别人的教区中施行祝圣。这就是那时关于\Person{亚他那修}的事务的进展。\par

% structure: p0090 source=h2 target=subsection reason=relative-heading-rank; base=h2

\subsection{第二十五章. 论篡位者\Person{玛格嫩提乌}和\Person{维特拉尼奥}。}

大约此时，一场非同寻常的骚动震撼了整个国家，我们将其主要概要略述，认为有必要不完全忽略它们。我们在第一卷中已提到，在\Place{君士坦丁堡}建城者去世之后，他的三个儿子继承了帝国；现在也必须说明，他们的一位亲属\Person{达尔马提乌}（从其父得名）与他们共享了皇帝权力。此人与他们共同执掌主权仅仅很短时间，士兵便杀死了他，\Person{君士坦提乌}既未命令处死他，也未加以禁止。\Person{小君士坦丁}如何也因侵入属于其弟的帝国部分而被士兵杀死，此前已不只一次记载。他死后，对罗马人发动了波斯战争，\Person{君士坦提乌}在其中一事无成；因为双方边界上的一次夜战中，\Person{波斯人}稍占优势。而此时基督徒的事务也变得同样不稳定，因\Person{亚他那修}和\OriginalTerm{同质}{homoousion}一词，众教会大乱。事情到了这般地步，在西部出现了一个叫做\Person{玛格嫩提乌}的篡位者，他通过背叛杀死了西部帝国皇帝\Person{君士坦斯}（当时驻在\Place{高卢}）。此事发生后，一场激烈的内战爆发，\Person{玛格嫩提乌}控制了整个\Place{意大利}，将\Place{非洲}和\Place{利比亚}置于其权力之下，甚至占有了\Place{高卢}。但在\Place{伊里利库姆}的\Place{西尔米乌姆}城，军队扶立了另一个名叫\Person{维特拉尼奥}的篡位者；同时一场新的麻烦使\Place{罗马}本身陷入骚动。原来\Person{君士坦丁}有一个侄子名叫\Person{内波提安}，他依靠一群角斗士的支持，在那里自立为君。但他被\Person{玛格嫩提乌}的一些军官杀死，\Person{玛格嫩提乌}本人则入侵西部省份，处处制造荒凉。\par

% structure: p0092 source=h2 target=subsection reason=relative-heading-rank; base=h2

\subsection{第二十六章. 西方皇帝\Person{君士坦斯}死后，\Person{保罗}与\Person{亚他那修}再次被逐出主教座：前者在流放途中被杀，后者则逃遁得脱。}

这些灾难性事件在短时间内集中发生；它们发生在萨迪卡会议后第四年，在塞尔吉乌斯和尼格里尼安执政期间。当这些情况被公布后，整个帝国的最高统治权似乎全部落在君士坦提乌斯一人身上，他在东方被宣布为唯一的独裁者，并以最积极的准备对付篡位者。这时，亚他那修的反对者们认为有利的时机已经到来，在亚他那修抵达亚历山大里亚之前，又对他提出了最恶毒的指控，向君士坦提乌斯皇帝保证说，他正在颠覆整个埃及和利比亚。而亚他那修竟在自己教区范围之外施行圣职，这也在一定程度上使对他的指控更加可信。与此同时，在此关键时刻，亚他那修进入了亚历山大里亚；他召集了埃及的主教会议，他们一致投票确认了萨迪卡会议和马克西穆斯在耶路撒冷召集的会议所作的决定。但是，早已浸染亚略学说的皇帝推翻了他不久前才决定的全部宽容措施。他首先下令将君士坦丁堡主教保禄流放；押送者在卡帕多西亚的库库苏斯将他勒死。玛策禄也被驱逐，巴西略再次被立为安西拉教会的主宰。阿德里安堡的路济乌斯被戴上镣铐，死于狱中。关于亚他那修的报告如此影响了皇帝的心，以致他在无法控制的愤怒中下令，无论何处找到亚他那修，都要将他处死；他还将色雷斯教会的主持者德奥杜禄和奥林比乌斯列入同一放逐名单。然而，亚他那修对皇帝的意图并非不知；他得知后再次选择逃亡，从而逃脱了皇帝的威胁。亚略派人士谴责这次退隐是犯罪，特别是基里基雅尼禄尼亚斯主教纳齐苏斯、劳迪基雅的乔治，以及当时治理安提约基雅教会的莱昂提乌斯。这最后一人在担任司铎时曾被剥夺品级，因为为了消除他与一名叫欧斯托利乌姆的女人有不正当交往的一切嫌疑——他大部分时间与她在一起——他阉割了自己，此后便与她更无拘束地生活，理由是再也没有恶意猜疑的余地了。后来，在君士坦提乌斯皇帝的迫切要求下，他被立为安提约基雅教会的主教，接替普拉基图斯的继任者斯德望。关于此事，就谈到这里。\par

% structure: p0094 source=h2 target=subsection reason=relative-heading-rank; base=h2

\subsection{第二十七章 马其顿尼占据君士坦丁堡主教座后，对与其分歧者大加迫害。}

那时，保禄按上述方式被除去后，马其顿尼成为君士坦丁堡教会的主宰；他获得了对皇帝的极大影响力，在基督徒中挑起了一场战争，其残酷程度不亚于篡位者自身正在进行的战争。因为他促使他的君主与他合作摧毁教会，他设法使他决定采取的任何有害措施都得到法律批准。为此，在各城市颁布了一道谕令，并指派军队执行皇帝的法令。因此，承认\OriginalTerm{本体同一}{consubstantiality}教义的人不仅被赶出教会，而且被赶出城市。起初，他们只满足于驱逐；但随着邪恶的增长，他们采取了更极端的措施，强迫他人与他们共融，几乎不顾及这种对教会的亵渎。他们的暴力几乎不亚于从前迫使基督徒敬拜偶像的人；因为他们施加了各种鞭打、各种折磨，并没收财产。许多人被流放惩罚；有的在酷刑下死去；有的在被押往流放途中被杀。这些暴行在所有东方城市施行，尤其是在君士坦丁堡；马其顿尼一旦获得主教职位，便猛烈加剧了先前尚属轻微的内部纷争。然而，希腊和伊利里库姆的城市，以及西方地区的城市，仍然享有安宁；因为它们彼此保持和谐，并继续持守尼西亞会议颁布的信德准则。\par

% structure: p0096 source=h2 target=subsection reason=relative-heading-rank; base=h2

\subsection{第二十八章 亚他那修记载亚略派者乔治在亚历山大里亚施行的暴行。}

乔治同时期在亚历山大里亚施行的残暴行径，可从亚他那修的叙述中得知，他亲身经历并目睹了这些事件。他在\WorkTitle{为逃亡辩护}中论及这些事时，如此写道：\par

\Quote{此外，他们又来到\Place{亚历山大里亚}，企图再次加害于我；这次他们的行径比以往更加恶劣。因为士兵突然包围了教堂，于是响起了战争的喧嚣，而非祈祷的声音。后来，在四旬期期间，从\Place{卡帕多西亚}派来的\Person{乔治}到达后，更助长了他奉命要做的恶事。复活节周过后，贞女们被投入监狱，主教们被军人用锁链押送，甚至连孤儿寡妇的住所也遭强行闯入，粮食被洗劫一空。基督徒在夜间被暗杀；房屋被查封；圣职人员的亲属因他们而陷入危险。这些暴行已是骇人听闻，但随后的行为更为可怕。因为在圣神降临节后的那一周，百姓禁食后前往墓园祈祷，因为所有人都厌恶与\Person{乔治}共融；这个最邪恶的人得知后，便煽动一个身为摩尼教徒的军官\Person{塞巴斯蒂安}来攻击他们。于是，\Person{塞巴斯蒂安}率领一队持刀、弓和标枪的士兵，在主日向百姓进攻；发现只有少数人在祈祷——因时间已晚，大部分人已离去——他便做出了他们所能预期的暴行。他点起火，把贞女们带到火旁，强迫她们宣称自己属于亚流派信仰。但看到她们坚守立场，不畏火焰，他便剥去她们的衣服，猛击她们的脸，以至于之后很长时间都难以认出她们。他又抓住大约四十个男人，以异常的方式鞭打他们：他用刚从棕榈树上砍下的、还带着刺的枝条抽打他们的背，以致有些人不得不反复求医，将刺从肉中取出，另一些人因无法忍受痛苦而死于鞭刑。所有幸存者连同一位贞女被流放到\Place{大绿洲}。至于死者的尸体，他们甚至不肯交给亲属，而是拒绝给予他们埋葬的礼仪，随意藏匿，以免他们残忍的罪证显露。他们这样做如同疯子一般。因为死者的亲友为其宣认而欢喜，却因遗体未能安葬而哀伤，这些行为的邪恶不人道反而更加昭彰。因为不久之后，他们从\Place{埃及}和两个\Place{利比亚}流放了以下主教：\Person{亚摩尼乌斯}、\Person{特姆伊斯}、\Person{卡尤斯}、\Person{斐禄}、\Person{赫尔墨斯}、\Person{普林尼}、\Person{普塞诺西里斯}、\Person{尼拉蒙}、\Person{阿加托}、\Person{阿纳甘普斯}、\Person{马尔谷}、\Person{亚摩尼乌斯}（另一人）、\Person{马尔谷}（另一人）、\Person{德拉孔提乌斯}、\Person{亚德尔斐乌斯}、\Person{雅典诺多鲁斯}；以及司铎\Person{希耶拉克斯}和\Person{狄奥斯哥鲁}。他们在押送中对待他们如此苛刻，以至于有人在途中死去，有些则在流放地去世。他们就这样除掉了三十多位主教，因为亚流派的急切愿望，如同\Person{阿哈布}一样，是要尽可能消灭真理。}\par

以上是\Person{亚大纳削}关于\Person{乔治}在\Place{亚历山大里亚}所行暴行的记述。与此同时，\Person{君士坦提乌斯}皇帝率领军队进入\Place{伊利里库姆}。因为那里有公共事务急需他到场，特别是军队宣布\Person{维特拉尼奥}为皇帝。到达\Place{西尔米乌姆}后，他在休战期间与\Person{维特拉尼奥}举行了一次会谈；他巧妙地使先前支持\Person{维特拉尼奥}的士兵倒戈，只向\Person{君士坦提乌斯}一人致敬为奥古斯都和至高君主。因此，在欢呼声中，\Person{维特拉尼奥}被完全忽略。\Person{维特拉尼奥}发现自己被抛弃，立即跪倒在皇帝脚下；\Person{君士坦提乌斯}取走他的皇冠和紫袍，以极大的仁慈对待他，并建议他以平民身份安静地度过余生，指出在他这个高龄，宁静的生活远比伴有焦虑和忧虑的尊位更合适。\Person{维特拉尼奥}的事就此了结；皇帝下令从国库中给予他丰厚的供养。后来\Person{维特拉尼奥}住在\Place{比提尼亚的普鲁萨}时，多次写信给皇帝，向他保证，使他摆脱帝权不可避免的烦扰，是赐予他最大的祝福。并补充说，他自己没有明智地享用他在隐居中所赐予他的幸福。此事就此为止。此后，\Person{君士坦提乌斯}皇帝立他的亲戚\Person{加卢斯}为凯撒，并赐予他自己的名字，派他前往\Place{叙利亚的安提约基雅}，以保卫东部地区。当\Person{加卢斯}进入该城时，救主的记号在东方显现：一颗十字架形的柱子出现在天上，使观看者大为惊奇。皇帝派遣他其余的将领率领相当兵力去对付\Person{马格嫩提乌斯}，自己则留在\Place{西尔米乌姆}等待事态发展。\par

% structure: p0100 source=h2 target=subsection reason=relative-heading-rank; base=h2

\subsection{第29章 论异端创始人\Person{福提诺}。}

在此期间，当时主持该城教会的\Person{福提诺}更为公开地宣扬他所创立的信经；因此引起骚动，皇帝下令在\Place{西尔米乌姆}召开主教会议。于是，东方主教中有\Place{阿雷苏萨}的\Person{马尔谷}、\Place{亚历山大里亚}的\Person{乔治}（如我先前所说，亚流派在\Person{格雷戈里}被撤职后派他管理该教区）、\Place{安基拉}教会的\Person{巴西略}（在\Person{马塞鲁斯}被驱逐后主持该教会）、\Place{佩鲁西翁}的\Person{潘克拉提乌斯}和\Place{赫拉克勒亚}的\Person{希帕提安}出席会议。西方主教中有\Place{穆尔萨}的\Person{瓦伦斯}和当时著名的\Place{西班牙科尔多瓦}的\Person{奥西乌}，后者十分不情愿地出席。他们在\Person{塞尔吉乌斯}和\Person{尼格里尼安}任执政官之后的年份在\Place{西尔米乌姆}会面；那一年因战争动荡，没有执政官举行惯常的开幕典礼。他们会面后，发现\Person{福提诺}持有利比亚的\Person{撒伯里乌}和\Place{萨莫萨塔}的\Person{保禄}的异端，便立即将他撤职。这一决定在当时和后来都普遍被视为光荣和公正；但那些留在那里的人随后的行为却远非如此得到普遍赞同。\par

% structure: p0102 source=h2 target=subsection reason=relative-heading-rank; base=h2

\subsection{第30章 在君士坦提乌斯皇帝面前于\Place{西尔米乌姆}公布的信经}

仿佛他们要撤销先前关于信仰的决定，他们又公布了新的信经阐述：即一份由\Person{阿雷图萨的马尔谷}用希腊文写成的；另有数份用拉丁文写成的，这些既在措辞上互不一致，也与阿雷图萨的主教所口述的那一份不相符。我在此将附上其中一份用拉丁文起草的，与马尔谷用希腊文写成的相对照；另一份后来在西尔米乌姆宣读的，将在描述\Place{阿里米努姆}事件时给出。但必须理解的是，这两份拉丁文本都被译成了希腊文。\Person{马尔谷}所提出的信仰宣言如下：\par

我们相信唯一的天主，全能的圣父，万有的创造者和造主，天上地下一切家族都由他得名，\BibleRef{弗 3:15}并相信他的独生子，我们的主耶稣基督，他在万世之前由圣父所生，天主出自天主，光出自光，藉着他，天上地下一切有形无形的万物都受造；他就是圣言、智慧、真光与生命；他在末世为了我们而成为人，由童贞玛利亚诞生，被钉十字架而死，被埋葬，第三日从死者中复活，被接升天，坐在圣父的右边，在时代完成之时来临，审判生者死者，按各人的行为报应各人；他的国度永远长存，延续无穷岁月；因为他不仅在此世，也要在来世，坐在圣父的右边。我们也相信圣神，即护慰者，主在升天之后，曾应许将他赐给宗徒们，教导他们，使他们想起一切事；藉着他，那些真诚相信主的人的灵魂得以圣化。但那些断言圣子来自虚无，或来自另一本体，而非来自天主，并且说他曾有一个时期或时代不存在的人，圣而公教会视之为外人。因此我们再说，若有人声称父与子是两位天主，当受绝罚。若有人承认基督是天主，且在万世之前是天主子，却不承认他在创造万物时辅佐圣父，当受绝罚。若有人胆敢断言无始者或其一部分由玛利亚所生，当受绝罚。若有人说圣子是依预知由玛利亚所生，而不是在万世之前由圣父所生、与天主同在，且万物不是藉着他造的，当受绝罚。若有人断言天主的本质是可扩展或可收缩的，当受绝罚。若有人说天主扩展的本质产生圣子，或称圣子为天主本质的扩展，当受绝罚。若有人称天主子为内在的言语或发出的言语，当受绝罚。若有人宣称由玛利亚所生的圣子仅仅是个人，当受绝罚。若有人承认由玛利亚所生的是天主而人，却暗示他就是无始的天主自己，当受绝罚。若有人理解经句\BibleQuote{我是首先的，我是末后的，除我以外没有天主，}\BibleRef{依 44:6}（这话是为消灭偶像和假神而说的）按犹太人的理解，以为是为推翻万世之前天主的独生子而说的，当受绝罚。若有人听见\BibleQuote{圣言成了血肉，}\BibleRef{若 1:14}竟以为圣言变成了血肉，或以为他在取血肉时经历了任何变化，当受绝罚。若有人听见天主的独生子被钉十字架，竟说他的神性经历了任何朽坏、苦难、变化、减损或毁灭，当受绝罚。若有人断言圣父没有对圣子说\BibleQuote{让我们造人，}\BibleRef{创 1:26}而是天主对自己说话，当受绝罚。若有人说亚巴郎所见到的不是圣子，而是无始的天主或其一部分，当受绝罚。若有人说以人形与雅各伯角力的不是圣子，而是无始的天主或其一部分，当受绝罚。若有人理解\Quote{主从天主降下雨来}不是指圣父与圣子，而说是他由自己降下雨来，当受绝罚：因为主圣子是从主圣父降下雨来。若有人听见\Quote{主圣父和主圣子}，却称圣父和圣子都是主，并说\Quote{从主来的主}，从而断言有两位天主，当受绝罚。因为我们并不将圣子与圣父并列，而是认为他隶属于圣父。因为他并非没有圣父的旨意就降世为身体；也不是由自己降下雨来，而是从行使至高权柄的\Emph{主}（即圣父）降雨；也不是自己坐在圣父的右边，而是服从圣父的命令：\Quote{你坐在我的右边}【当受绝罚】。若有人声称圣父、圣子、圣神是同一主体，当受绝罚。若有人称护慰者圣神为无始的天主，当受绝罚。若有人，如他所教导我们的，不承认护慰者不同于圣子，而主自己曾说：\Quote{圣父，我要求他，他就会派遣另一位护慰者给你们，}当受绝罚。若有人肯定圣神是圣父和圣子的一部分，当受绝罚。若有人说圣父、圣子、圣神是三位神，当受绝罚。若有人说天主子是由天主的旨意造为受造物之一，当受绝罚。若有人说圣子是在圣父不愿意的情况下受生的，当受绝罚：因为圣父并非由于某种自然必然性的强迫，在不情愿的时候生了圣子；而是他一愿意，就表明他在无时间、无情感中生了他。若有人说圣子是无始的、没有开始的，暗示有两个无始者，从而造出两位天主，当受绝罚：因为圣子是万物的头和开始；但\BibleQuote{基督的头是天主。}\BibleRef{格前 11:3}因此我们虔诚地藉着圣子将一切追溯至那无始的万有之源。此外，为了给出基督教义的准确概念，我们再说，若有人不宣认耶稣基督在万世之前就是天主子，并在创造万物时辅佐圣父；反而声称他仅从玛利亚出生时才被称为子和基督，并且那时才获得神性的开端，当受绝罚，如同撒摩撒塔人一样。\par

\Emph{另一个在\Place{西尔米乌姆}以拉丁文提出、后来被译为希腊文的信仰阐明。}\par

既然认为应当对信仰进行一些慎思，所有要点都在\Place{西尔米乌姆}，当着\Person{瓦伦斯}、\Person{乌尔撒西乌斯}、\Person{革尔米尼乌斯}等人的面，经过仔细考察和讨论。\par

显然，只有一位天主，全能的天主圣父，正如普世所宣认的；以及祂的独生子耶稣基督，我们的主、天主和救主，由祂在万世之前所生。但我们不应说有两个天主，因为主亲自说过：\BibleQuote{我往我父和你们父那里去，往我天主和你们天主那里去。}\BibleRef{若 20:17}因此，祂也是万有的天主，正如宗徒所教导的：\BibleQuote{难道天主只是犹太人的天主吗？不也是外邦人的天主吗？是的，也是外邦人的天主；因为只有一位天主，祂要使受割损的因信德而成义。}\BibleRef{罗 3:29}在其他所有事上也是一致的，没有任何模棱两可。但关于拉丁文称为\OriginalTerm{实体}{substantia}、希腊文称为\OriginalTerm{本质}{ousia}的东西，也就是说，为了更准确地表达意义，\OriginalTerm{同质}{homoousion}或\OriginalTerm{似质}{homoiousion}这两个词使许多人困惑，所以完全不宜提及这些词；也不应在教会中宣讲它们，因为圣经中对此毫无记载，并且这些事超越人类知识和能力，无人能解释圣子的受生，正如经上所载：\BibleQuote{谁能述说祂的世代？}\BibleRef{依 53:5}显然，只有圣父知道祂如何生了圣子；同样，圣子知道祂如何由圣父所生。但无人怀疑圣父在尊荣、威严和神性上，以及在天主的名称上更为伟大；圣子亲自作证：\BibleQuote{派遣我来的父比我大。}\BibleRef{若 14:28}并且无人不知这也是公教教义：圣父和圣子具有两个位格，圣父更为伟大；但圣子与圣父所交给祂的一切事物一同受于圣父。圣父无始、无形、不死、不苦；圣子则由圣父所生，从天主的天主，从光的光；无人能理解祂的受生，如前所述，只有圣父知道。圣子本人，我们的主和天主，取了血肉或身体，即人性，正如天使所报喜讯；并如全部圣经，尤其是外邦人之大宗徒所教导，基督由童贞玛利亚取人性，并借此受苦。全部信仰的总结和确认是：应始终持守\OriginalTerm{天主圣三}{Trinity}的教义，正如我们在福音中所读：\BibleQuote{你们去使万民成为门徒，因父、及子、及圣神之名给他们授洗。}\BibleRef{玛 28:19}因此，天主圣三的数目是圆满而完美的。现在，护慰者圣神，由圣子所派遣，按祂的应许而来，为使宗徒和所有信徒获得圣化和教导。\par

他们甚至在其被免职后，还试图劝诱\Person{福提努斯}同意并签署这些条款，承诺如果他通过撤回声明谴责他自己发明的教义，并采纳他们的意见，就恢复他的主教职位。但他不接受他们的提议，反而向他们挑战辩论：皇帝安排了一个日子，在场的诸位主教以及皇帝指示旁听听讨论的不少元老都聚集了。在他们面前，当时主持\Place{安基拉}教会的\Person{巴西略}被指定与\Person{福提努斯}辩论，速记员记录了他们各自的发言。双方的论证交锋非常激烈；但\Person{福提努斯}败北，被定罪，余生都在流亡中度过，期间他用两种语言——他在拉丁文上并非不精通——撰写了反对一切异端并支持他自己观点的论著。关于\Person{福提努斯}，说这些就够了。\par

当时在\Place{西尔米乌姆}聚集的主教们，后来对他们以拉丁文颁布的那份信经格式感到不满；因为在其公布之后，他们觉得其中包含许多矛盾。因此他们试图将其从抄写者手中收回；但由于许多人将其藏匿，皇帝通过谕令下令搜寻该版本，威胁若有人被发现藏匿将受到惩罚。然而，这些威胁无法压制已经落入许多人手中的东西。关于这些事，说这些就够了。\par

% structure: p0110 source=h2 target=subsection reason=relative-heading-rank; base=h2

\subsection{第三十一章 \Place{科尔多瓦}主教\Person{奥西乌斯}}

既然我们注意到西班牙人\Person{奥西乌斯}是违背自己意愿出席[\Place{西尔米乌姆}会议]的，有必要对他作一些简要介绍。不久前，他因亚略派人士的阴谋而被流放；但在\Place{西尔米乌姆}与会者的恳求下，皇帝召他前往那里，希望他能通过劝服或强迫赞同他们的做法；因为如果此事能够实现，他们认为这将大大增加他们意见的权威。基于这个原因，因此，如我所说，他非常不情愿地被迫出席：当他拒绝同意他们时，那位老人受到鞭挞和酷刑。因此，他被迫顺从并签署了他们的信仰阐明。这就是当时在\Place{西尔米乌姆}发生之事的结局。但在此之后，皇帝\Person{君士坦提乌斯}仍留在该地，在那里等待反对\Person{马格嫩提乌斯}的战争结果。\par

% structure: p0112 source=h2 target=subsection reason=relative-heading-rank; base=h2

\subsection{第三十二章 僭越者\Person{马格嫩提乌斯}的覆灭}

在此期间，马格嫩提乌斯控制了帝国首都罗马城，处死了许多元老院成员以及大量平民。但君士坦提乌斯的将领们一集结起罗马军队并向其进军，他便离开罗马，退入高卢地区。在那里发生了多次战斗，有时一方占优，有时另一方占优：但最终马格嫩提乌斯在穆尔萨——高卢的一座堡垒——附近被击败，并被严密围困。据说在此地发生了以下引人注目的事件。马格嫩提乌斯想提振因最近战败而气馁的士兵的勇气，为此登上一座高台。士兵们本想像往常一样发出向皇帝致敬的欢呼，却违背本意地同时高呼君士坦提乌斯·奥古斯都之名，而非马格嫩提乌斯。马格嫩提乌斯认为这是对自身不利的兆头，立即撤离堡垒，退往高卢最偏远的地区。君士坦提乌斯的将领们急速追击。在塞琉古山附近再次交战后，马格嫩提乌斯全军溃败，独自逃往高卢的里昂城，该城距离穆尔萨堡垒有三天的路程。马格嫩提乌斯到达该城后，先杀死了自己的母亲；然后又杀死了被他立为凯撒的兄弟，最后伏剑自杀。此事发生在君士坦提乌斯第六次任执政官、君士坦提乌斯·加卢斯第二次任执政官期间，八月十五日。不久之后，马格嫩提乌斯的另一个兄弟德森提乌斯上吊自尽。马格嫩提乌斯的野心就此终结。帝国的事务并未完全平息；因为此后不久，又出现了一个名为西尔瓦努斯的篡位者：但君士坦提乌斯的将领们在他于高卢挑起骚乱时迅速将其铲除。\par

% structure: p0114 source=h2 target=subsection reason=relative-heading-rank; base=h2

\subsection{第三十三章 论居住在巴勒斯坦凯撒利亚腓立比的犹太人}

大约同时，东方又发生了另一场内乱：因为居住在巴勒斯坦凯撒利亚腓立比的犹太人拿起武器反抗罗马人，并开始劫掠邻近地区。但也被称为君士坦提乌斯的加卢斯（皇帝立他为凯撒后派往东方）派遣军队攻打他们，并将其完全击败；之后他下令将凯撒利亚腓立比城夷为平地。\par

% structure: p0116 source=h2 target=subsection reason=relative-heading-rank; base=h2

\subsection{第三十四章 论凯撒加卢斯}

加卢斯完成这些事后，未能以节制的心态承受成功；而是立即企图反对立他为凯撒者的权威，自己觊觎君权。然而，他的意图很快被君士坦提乌斯察觉：因为他竟敢擅自处死当时的东方近卫总长多米提安和财务官马格努斯，而未向皇帝表明其意图。君士坦提乌斯对此极为愤怒，召加卢斯前来；加卢斯极为恐惧，极不情愿地前往；当他到达西部地区，并抵达弗拉诺纳岛时，君士坦提乌斯下令将其处死。但不久之后，他立加卢斯的兄弟尤利安为凯撒，并派他前往高卢对抗蛮族。君士坦提乌斯第七次任执政官、加卢斯（别号君士坦提乌斯）第三次任执政官时，加卢斯被处死；次年，阿尔贝提翁和洛利安任执政官时，尤利安于十一月六日被立为凯撒；关于他，我们将在下一卷中进一步提及。君士坦提乌斯解除了困扰他的忧虑后，再次将注意力转向教会纷争。因此，他从锡尔米乌姆前往帝国首都罗马，再次召集主教会议，召请部分东方主教赶往意大利，并安排西方主教在那里与他们会面。在东方为此做准备期间，罗马主教尤利乌斯去世，他治理该地教会十五年，由利贝里乌斯继任主教之职。\par

% structure: p0118 source=h2 target=subsection reason=relative-heading-rank; base=h2

\subsection{第三十五章 论欧诺米乌斯的导师、叙利亚人埃提乌斯}

在叙利亚的\Place{安提约基雅}，又兴起了一位异端首领，即被称为\OriginalTerm{无神论者}{Atheus}的\Person{阿埃提乌斯}。他在教义上与\Person{亚略}一致，并持守相同的意见；但他脱离了亚略派，因为他们接纳了\Person{亚略}共融。因为\Person{亚略}，正如我先前所述，心中怀有一种看法，口里却声称另一种；他假装同意并签署了尼西亚会议所制定的信经，以欺骗当时的皇帝。因此，\Person{阿埃提乌斯}便脱离了亚略派。然而，他早先已是一个异端者，并且是亚略观点的热心提倡者。在\Place{亚历山大里亚}接受了一些非常浅薄的教育后，他离开那里，来到\Place{叙利亚}的\Place{安提约基雅}，那是他的故乡，由当时该城的主教\Person{良提乌}祝圣为执事。自此，他开始因言论的奇特而令与他交谈的人惊讶。他这样做，是依赖\Person{亚里士多德}的\WorkTitle{范畴篇}的原则；那书以此命名，其要旨他既未亲自领会，也未通过与有学识者的交流而得到启发：以致他几乎不知道自己在构建荒谬的论证来困扰和欺骗自己。因为\Person{亚里士多德}撰写这部作品，是为了锻炼年轻弟子的才智，并用精妙的论证来困惑那些假装嘲笑哲学的诡辩家。因此，阐释\Person{柏拉图}和\Person{普罗提诺}著作的\OriginalTerm{怀疑派学者}{Ephectic academicians}，指责\Person{亚里士多德}在那本书中表现出的虚妄的精细；但从未有过学院导师指导的\Person{阿埃提乌斯}，却坚守\WorkTitle{范畴篇}的诡辩。为此缘故，他无法理解如何能有不受开始而受生的事，以及那受生者如何能与生者同永恒。事实上，\Person{阿埃提乌斯}是如此肤浅，对圣经知之甚少，又极其好争辩——这是任何粗野之人都能做到的——以至于他从未仔细研读过那些阐释基督神谕的古代作者；他完全拒绝\Person{克莱孟}、\Person{阿非利加努斯}和\Person{奥利金}，这些人以其在文学和科学各领域的博学而著称。但他在写给皇帝\Person{君士坦提乌}和其他一些人的信中，穿插了冗长的争论以展示他的诡辩。因此他被冠以\OriginalTerm{无神论者}{Atheus}之名。然而，尽管他的教义陈述与亚略派相似，但由于其三段论的晦涩难懂——亚略派同人无法理解——他们便宣称他为异端者。为此缘故，他被逐出他们的教会，便假装脱离与他们的共融。即使在今日，仍有一些人从前因他而被称为\OriginalTerm{阿埃提乌斯派}{Aëtians}，但现在称为\OriginalTerm{欧诺米乌斯派}{Eunomians}。因为过了一些时候，曾是他的文书\Person{欧诺米乌斯}，在从师父学习了这种异端的推理方式后，后来成了该派的领袖。至于\Person{欧诺米乌斯}，我们将在适当的地方更详细地论述。\par

% structure: p0120 source=h2 target=subsection reason=relative-heading-rank; base=h2

\subsection{第三十六章 米兰会议}

其时，主教们在\Place{意大利}会集，来自东方的甚少，多数因年迈或路途遥远未能前来；但来自西方的超过三百位。皇帝命令他们在\Place{米兰}聚集。会面时，东方的主教们以要求与会者一致通过谴责\Person{亚他那修}的判决来开启会议；其目的是，使他从此完全被逐出\Place{亚历山大里亚}。但\Place{高卢}的\Place{特雷维}的主教\Person{保利努}、\Place{阿尔巴}的主教\Person{狄奥尼西}（\Place{阿尔巴}为意大利首府），以及\Place{利古里亚}的\Place{韦尔切利}的主教\Person{欧瑟比}，察觉到东方主教们通过要求批准对\Person{亚他那修}的判决，意在颠覆信仰，便起立大声呼喊说：\Quote{这个提议表明了一个针对基督真理原则的隐蔽阴谋。因为他们坚持认为，对\Person{亚他那修}的指控是毫无根据的，只是控告者为了败坏信仰而捏造的。}他们以这种激烈的态度提出抗议后，主教会议便解散了。\par

% structure: p0122 source=h2 target=subsection reason=relative-heading-rank; base=h2

\subsection{第三十七章 阿里米农会议及其公布的信经}

皇帝得知所发生之事后，将这三位主教放逐；并决定召开一次大公会议，打算将东方所有主教都召集到西方，以便尽可能使他们达成一致。但经过考虑，旅程遥远似乎构成严重障碍，他便指示议会分为两部分：允许当时在\Place{米兰}的人前往意大利的\Place{阿里米努姆}集会；又通过信函指示东方主教们在\Place{比提尼亚}的\Place{尼科美底亚}集会。皇帝这样安排，旨在达成意见上的普遍一致；但结果却出乎他的意料。因为两个会议本身都不和谐，各自分成了对立派别：在\Place{阿里米努姆}召集的人彼此无法达成一致；而东方主教们在\Place{伊索利亚}的\Place{塞琉西亚}集会，又造成了另一分裂。两处发生的细节，我们将在本史书中叙述，但首先应对\Person{欧多克修斯}略作评述。大约此时，曾祝圣异端者\Person{艾提乌斯}为执事的\Person{莱昂提乌斯}已去世，\Place{日尔曼尼西亚}（此城在\Place{叙利亚}）主教\Person{欧多克修斯}当时正在\Place{罗马}，认为时间紧迫，便虚伪地向皇帝表示，他管辖的城市需要他的咨询和照管，请求允许立即返回那里。皇帝轻易就答应了，没有怀疑其隐秘意图：\Person{欧多克修斯}以皇帝寝宫中的几位首席官员为助手，放弃了自己原有的教区，并欺骗性地占据了\Place{安提约基雅}教区。他的首要愿望是恢复\Person{艾提乌斯}的地位；于是他召集了一个主教议会，旨在重新授予\Person{艾提乌斯}执事职的尊位。但这无论如何也无法实现，因为人们对\Person{艾提乌斯}的憎恶远甚于\Person{欧多克修斯}为他所作的努力。当主教们在\Place{阿里米努姆}聚集时，东方的主教们宣布他们愿意对\Person{亚大纳削}的案子默不作声：这一决议得到了\Person{乌尔萨西乌斯}与\Person{瓦伦斯}的热烈支持，这两人先前曾坚持\Person{亚略}的教理；但正如我已指出的，他们后来向罗马主教呈交了撤销他们意见的声明，并公开承认赞同同体教义。因为这些人总是倾向于与占主导地位的一方为伍。\Person{格米尼乌斯}、\Person{奥克森提乌斯}、\Person{德莫菲鲁斯}和\Person{加尤斯}也发表了关于\Person{亚大纳削}的相同声明。因此，当一些人在主教会议中力图提出某种提议，另一些人提出另一种时，\Person{乌尔萨西乌斯}与\Person{瓦伦斯}说，所有先前的信经草案都应视为置之不理，唯有最后一份，即他们最近在\Place{西尔米乌姆}大会上所拟定的，应被视为权威文件。随后他们让人宣读拿在手中的一份文件，其中包含另一种信经格式：此格式确实是在\Place{西尔米乌姆}拟定的，但正如我们先前所说，一直隐藏着，直到他们在\Place{阿里米努姆}当前会议上才公布。此信经由拉丁文译为希腊文，内容如下：\par

\Quote{大公信仰于\Place{西尔米乌姆}，在我们主\Person{君士坦提乌斯}面前，在卓越的\Person{弗拉维乌斯·欧塞比乌斯}与\Person{希帕提乌斯}担任执政官时，五月二十三日，予以阐释。}\par

\Quote{我们相信唯一真实的天主，全能圣父，万物的创造者与塑造者；又相信天主的独生子，在万世之前，在一切起始之前，在一切可想象的时间之前，在一切可理解的思想之前，无痛苦地受生；万世由他而受造，万物由他而受造；他作为圣父的独生子而受生，唯一者从唯一者，天主从天主，相似那生他的圣父，按照圣经：他的受生无人知晓，唯独生他的圣父知道。我们知道，他的这位独生子因圣父的同意从天降下，为除去罪恶，由\Person{童贞玛利亚}诞生，与他的门徒往来，并按圣父的旨意完成了每一项职分：被钉十字架而死，降入地下深处，并安排了那里的一切；阴间的（守门者）一见他就战栗；第三天复活后，他又与门徒往来，满四十天后升天，坐在圣父的右边；在末日，他将带着圣父的光荣来临，按照各人的行为施行报应。[我们相信]也相信圣神，天主子耶稣基督曾亲自应许将他作为护慰者赐给人类，正如经上记载：\InnerQuote{\BibleQuote{I go away to my Father, and will ask him, and he will send you another Comforter, the Spirit of truth. He shall receive of mine, and shall teach you, and bring all things to your remembrance.}}至于\InnerQuote{\InnerQuote{本体}}一词，我们的先辈为更简明起见曾使用过，但因百姓不理解，且圣经未包含此词，已引起反感，因此认为该词应予完全废除，今后不得再论及关乎天主的本体，因为神圣圣经从未谈论过圣父与圣子的本体。但我们要说，圣子在一切事上\Emph{相似}圣父，正如圣经所肯定和教导的。}\par

这些陈述宣读之后，那些对它们感到不满的人站起来说：\Quote{我们来此，并非因为我们缺少信经；我们完整保存了从起初所领受的，未曾更改。但我们现在聚集在此，是为了阻止可能对它进行的任何革新。因此，如果刚才宣读的内容没有引入新事物，那么现在就公开谴责亚流异端，如同教会古老法则谴责所有亵渎的异端一样。因为全世界都清楚，亚流的不虔信条引起了教会的纷争，以及直到如今仍然存在的困扰。}这个提议未被\Person{乌尔萨基乌斯}、\Person{瓦伦斯}、\Person{革尔弥尼乌斯}、\Person{奥克森提乌斯}、\Person{德摩斐卢斯}和\Person{加约斯}采纳，彻底分裂了教会：因为这些主教坚持当时在\Place{里米尼}教务会议上所宣读的内容；而另一些人则再次确认了尼西亚信经。他们还嘲笑所宣读的那份信经的标题；尤其是\Person{亚大纳削}，在他写给朋友的一封信中，如此表达自己：\par

\Quote{公教会的敬虔还缺少什么教义要点，以至于他们现在要就信仰进行探究，并将当代的执政官名放在他们发表的信仰阐述上？因为\Person{乌尔萨基乌斯}、\Person{瓦伦斯}和\Person{革尔弥尼乌斯}做了一件在基督徒中从未有过、甚至闻所未闻的事：他们编撰了一份信经，正如他们自己所愿意相信的，并在前面注明当前的执政官、月份和日期，以此向所有明辨的人证明：他们的信经并非古老的信仰，而是在当今皇帝\Person{君士坦提乌斯}统治下才产生的。此外，他们所写的一切都为了自己的异端：不仅如此，他们还假装为论述主，却称另一位\InnerQuote{主}为他们的主，即\Person{君士坦提乌斯}，他纵容了他们的不虔，以至于那些否认圣子永恒的人，竟称他为永远皇帝。这样，他们因自己的亵渎而被证明是基督的敌人。但或许圣先知记录时间为他们注意执政官提供了先例！即使他们敢以此为借口，也会极其明显地暴露出自己的无知。这些圣人的预言确实标记了时间。\Person{依撒意亚}和\Person{欧瑟亚}生活在\Person{乌齐雅}、\Person{约堂}、\Person{阿哈次}、\Person{希则克雅}的时代；\Person{耶肋米亚}在\Person{约史雅}时代；\BibleRef{耶 1:2} \Person{厄则克耳}和\Person{达内尔}在\Person{居鲁士}和\Person{达理阿}王朝；其他人在别的时期发出预言。然而他们那时并没有奠定宗教的基础。宗教在他们之前就已存在，且一直存在，甚至在创世之前，天主就在基督内为我们预备了它。他们也没有指明自己信仰的开端；因为他们自己先前已有信仰；但他们指明了通过他们所赐的应许的时间。这些应许主要指向我们救主的来临；所有关于以色列和外邦人未来事件进程的预言都是附带和次要的。因此，所提到的时期并非如我之前所说，表明他们信仰的开端，而是这些先知生活和预言这些事情的时间。但我们当代的这些贤哲，既不编纂历史，也不预示未来事件，在写了\InnerQuote{公教信仰得以公布}之后，立即附上执政官、月份和日期：正如圣先知记录了他们的著作和他们自己职务的日期一样，这些人暗示他们自己信仰的时代。但愿他们只写了关于\Emph{他们自己的}信仰的事（既然他们现在开始相信），而不是着手写关于公教信仰的事。因为他们没有写\InnerQuote{我们如此相信}，而是写了\InnerQuote{公教信仰得以公布}。此处所表现的轻率目的证明了他们的无知；而他们编造的文件中所发现的新颖表达表明它与亚流异端相同。以这种方式写，他们宣布了自己何时开始相信，以及他们希望从何时起他们的信仰被传扬。正如福音作者\Person{路加}说：\BibleRef{路 2:1} \InnerQuote{登记令得以公布}，他指的是一个先前不存在、而在那时生效、并由颁布者所写的法令；同样，这些人通过写\InnerQuote{信仰现已公布}，宣布了他们异端的教义是现代发明，以前并不存在。但既然他们将\InnerQuote{公教}一词应用其上，他们似乎无意中陷入了卡塔弗里吉亚人的极端主张，如同他们所说\InnerQuote{基督信仰首先向我们启示，并从我们开始}。正如那些被称为\Person{马克西米拉}和\Person{蒙塔努斯}的人一样，这些人称\Person{君士坦提乌斯}为他们的主，而非基督。但若按他们的说法，信仰是从当前的执政官时期开始的，那么教父们和真福殉道者们该怎么办？此外，他们自己如何对待那些曾受他们教导宗教原则、并在这一执政官时期之前去世的人？他们将用什么方法使这些人复生，以便抹去他们似乎曾教导过的东西，并用他们所公布的那些新发现取而代之？他们如此愚蠢，只能编造借口，且这些借口既不相称又不合理，并带有自身的驳斥。}\par

\Person{亚大纳削}这样写信给他的朋友们；有兴趣的人若通读他的整封信，将会看出他多么有力地处理了这个问题；但为简略起见，我们在此只摘录了其中一部分。教务会议罢免了\Person{瓦伦斯}、\Person{乌尔萨基乌斯}、\Person{奥克森提乌斯}、\Person{革尔弥尼乌斯}、\Person{加约斯}和\Person{德摩斐卢斯}，因为他们拒绝谴责亚流主义；这些人对自己被罢免极为愤怒，立即赶往皇帝那里，带去了在教务会议上所宣读的信仰陈述。教务会议也以公函将他们的决定告知皇帝，该公函由拉丁文译成希腊文，内容如下：\par

\Emph{\Place{里米尼}教务会议致\Person{君士坦提乌斯}皇帝书}\par

我们相信，先前颁布的法令之所以得以执行，既是出于天主的安排，也是出于陛下虔敬之命。因此，我们西方的主教们从各地来到\Place{阿里米诺}，为使公教会之信德得以彰显，并使持相反观点者得以显明。因我们审慎审议所有要点后，决定持守那由先知、福音和宗徒藉我们的主耶稣基督——陛下帝国的守护者和陛下本人的保护者——所启示的古老信德，这信德我们也一直持守。我们认为，删减任何由那些与陛下之父、荣光纪念的\Person{君士坦丁}一同出席\Place{尼西亚}大公会议的人所公正而正确地批准的事，都是不当且不敬的。他们的教义和观点已注入人心，并宣讲于民众耳中，且被证明强有力地反对甚至击败了\OriginalTerm{亚略异端}{Arian heresy}。不仅此异端，所有其他异端也为其所制服。因此，若在那时所确立的信条上有所增添或删减，将是危险的；因为若发生其中任何一事，敌人便将胆大妄为，为所欲为。\par

因此，\Person{乌尔撒基乌斯}和\Person{瓦伦斯}先前因涉嫌持有亚略派观点而被暂停共融；但为恢复共融，他们作了道歉，并声称已悔改其过失，正如其书面撤回声明所证明的；他们因而获得了宽恕和完全的赦免。\par

这些事发生之时，正值会议在\Place{米兰}举行，当时罗马教会的司铎也在场。\par

同时，我们得知\Person{君士坦丁}——即使在其死后也值得称誉——曾以恰当的精确性阐明了信德，但他作为人所生，领受了圣洗，并去世归入他应得的平安，我们从而认为，在他之后有所创新，而忽视如此众多圣洁的宣信者和殉道者——他们也是这宣认的作者，并持守他们对公教会古老体系的信德——是不适当的。他们的信德，天主藉我们的主耶稣基督延续到陛下统治的年代，藉其恩宠，你得以巩固你的统治，以至治理世界之一部分。\par

然而，这些迷妄而可怜之人，秉性不幸，竟再次胆敢自称是谬误教义的传播者，甚至试图颠覆教会的体制。因为当陛下的虔敬书信命令我们聚集审查信德时，他们便暴露了其意图，剥去了欺骗的外衣。因他们以某种诡计和混乱企图提出革新，并在此事上以\Person{格尔米尼乌斯}、\Person{奥克森提乌斯}和\Person{加约斯}为盟友，这些人不断引起争斗与不和，且他们单一的教训已超越了整个亵渎之体。但当他们察觉我们在其虚假观点上并无与他们相同的意向或心志时，他们便在我们的会议期间改变了主意，声称应提出另一信德表述。确实，很短的时间便使他们信服了自己观点的虚假。\par

因此，为使教会事务不再持续陷入同样境况，并为使麻烦和骚动不再不断发生并混乱一切，保持先前确定的观点坚定而不变，并将上述人员从我们共融中分离，似乎是稳妥的；为此，我们已派遣代表至陛下仁德，他们将向您传达会议的意见。我们特别嘱咐我们的代表，首要的是，他们应凭据古代和正确的决定，确认真理。他们将告知陛下之圣德，和平不会如\Person{乌尔撒基乌斯}和\Person{瓦伦斯}所言，在正确的一点被推翻时建立。因为那些破坏和平的人如何能和平？反之，这些事在\Place{罗马}教会中，也如同在其他城市中一样，将引起争斗和骚动。因此，现在我们恳求陛下仁德，以平静的目光看待我们的代表，并予以友善倾听，不允许任何事物被改变，从而侮辱已故之人，而应准许我们继续持守由我们祖先所界定和立法之事；我们应说，他们是凭着精明、智慧及圣神行事的。因为他们现今所引进的革新，使信徒充满不信任，使不信者充满残忍。我们进一步恳求您，吩咐住在远方的、受年老体弱和贫困之苦的主教们，得以顺利而迅速地返回自己的家乡，以免各教会失去其主教。此外，我们还恳求您这一点：从陛下虔敬之父的时代直至如今所存留的信德条文，不要删减，也不要增添；而应使它们保持不受损害。求您不许我们继续劳苦和忍受，也不许我们与我们的教区分离，而使我们能与自己的子民一同平安地有时日献上祈祷和感谢，为您的安康和统治的延续祈求——愿神性将这一切永远赐予您。我们的代表携带着主教们的署名和问候。这些代表将亲自从神圣圣经本身教导陛下之虔敬。\par

于是，该主教会议如此写信，并通过所选的主教们将信件送交皇帝。但\Person{乌尔撒基乌斯}和\Person{瓦伦斯}的同党先于他们到达，竭力诽谤会议，并展示他们所带来的信德说明。皇帝因事先已偏向\OriginalTerm{亚略主义}{Arianism}，对会议极为恼怒，却给予\Person{瓦伦斯}、\Person{乌尔撒基乌斯}及其朋友极大尊荣。因此，会议所委派的代表被扣留了相当长时间，未能得到答复；最终，皇帝通过那些到他那里的人作了如下答复：\par

\Quote{\Person{君士坦提乌斯} \OriginalTerm{胜利者与凯旋者}{Victor and Triumphator} \OriginalTerm{奥古斯都}{Augustus} 致在 \Place{里米尼} 集会的诸位主教。}\par

\Quote{连您们的圣善也并非不知，我们始终特别关怀神圣可敬的法律。然而至今我们仍未能接见您们派来的二十位主教使团，因为对蛮族的远征已不可避免。既然，如您们所承认，有关神圣法律的事务当以毫无挂虑的心情着手处理，故我已命这些主教等候我们返回 \Place{阿德里安堡}，待一切公务妥善处理后，我们方可听取并考虑他们的提议。在此期间，请勿因等候他们返回而觉得有烦扰您的尊驾；因为他们将我们的决定告知您们后，您们便可准备实施最有利于天主教会福祉的措施。}\par

主教们收到此信后，复函如下：\par

\Quote{我们收到了陛下——至蒙天主宠爱的君主——的来信，您在信中告知我们，国务的紧迫使您至今未能接见我们的代表；您命我们等候他们返回，好使您的虔敬能得知我们依照祖先传统所作的决定。但我们再次通过此信声明，我们绝不能背离初衷；我们也已委托代表陈述此事。因此我们恳求您，既以平和的面容命人诵读此封卑微的书信，也请欣然垂听代表们的陈述。您的温和无疑与我们一样明白，因在您这最蒙福的时代，如此众多教会失去主教，忧伤与悲痛何其深重。因此我们再次恳求陛下——至蒙天主宠爱的君主——若合您虔敬之意，请在寒冬来临前命我们返回各自的教会，使我们能与众信友一起，照常向全能的天主和我们的主、救主耶稣基督——祂的独生子——献上为陛下国运昌隆的祈祷，正如我们素常所做，如今也在祈祷中如此行。}\par

此信发出后，主教们一同等候了一段时间，但皇帝未予答复，于是他们各自返回本城。其实皇帝早就想在各地教会传播\OriginalTerm{亚略派}{Arian}学说，并急于使之占据优势；因此他假借主教们擅自离去为侮慢之举，声称他们藐视他，违反他的意愿解散了会议。于是他给予\Person{乌萨西乌斯}的同党在教会事务中为所欲为的无限自由，并下令将在\Place{里米尼}宣读的修订版信经送往意大利各地的教会，规定凡不签署者均应被逐出教座，另选他人接替。首先，罗马主教\Person{利贝里乌斯}因拒绝赞同该信经而被流放；\Person{乌萨西乌斯}的同党指派\Person{菲力克斯}继任，他原是该教会的一位执事，但接受亚略派异端后升任主教。然而有些人断言他并不赞同此见解，乃是迫于压力接受了主教祝圣。此后，西方各地充满动荡与骚乱，一些人被逐出和流放，另一些人则被扶植取代。这些事都是凭借皇帝谕令的权威以暴力实行的，这些谕令也发往东方各地。不久之后，\Person{利贝里乌斯}果然被召回，恢复其教座；因罗马民众发动骚乱，将\Person{菲力克斯}逐出教堂，皇帝即使不情愿也同意了。\Person{乌萨西乌斯}的同党离开意大利，经过东方各地，抵达色雷斯的\Place{尼斯}城，在那里暂住并召开了另一次主教会议，将曾在\Place{里米尼}宣读的信经译成希腊文，重新予以确认并公布，如前面所引的形式，并称之为大公会议，企图以此利用名称的相似性欺骗较为单纯的人，将他们在色雷斯\Place{尼斯}所炮制的信经冒充为在比提尼亚\Place{尼西亚}所颁布的信经。但这诡计对他们并无多大好处；因为很快就被识破，他们成为笑柄。关于西方所发生的事，所说已足够；现在我们必须叙述同时期在东方所发生的事。\par

% structure: p0142 source=h2 target=subsection reason=relative-heading-rank; base=h2

\subsection{第三十八章 \Person{马其顿尼乌斯}的残酷行径及其引发的骚乱}

亚流派的主教们因皇帝谕令而更加自信。他们以何种方式召集了一次主教会议，我们将在稍后说明。现在让我们先简要提及他们先前的一些行为。阿卡西乌和帕特罗斐禄将耶路撒冷主教玛克西穆逐出后，立济利禄接替其位。马其顿尼颠覆了君士坦丁堡邻近城省的正常秩序，将与他共同阴谋对抗教会的助手提升至教会荣誉职位。他任命埃勒乌西乌为齐济库斯主教，任命马拉托尼乌为尼科美底亚主教；后者先前曾在马其顿尼本人手下任执事，并在建立男女隐修院方面颇为活跃。但我们现在必须说明马其顿尼如何蹂躏君士坦丁堡周边城省的教会。此人，如我已说过的，夺取主教职位后，给那些不愿接受他观点的人带来无数灾难。他的迫害不仅限于被承认爲大公教会成员的人，也扩展到诺洼天派，因为他们知道他持守\OriginalTerm{同质}{homoousion}的道理；因此他们与其他人一同遭受了最难以忍受的苦难，但他们的主教安杰利乌却逃脱了。许多虔诚杰出的人被逮捕并受酷刑，因为他们拒绝与他共融；酷刑之后，他们强迫男人领受神圣奥迹，用木片撬开他们的口，将祝圣过的圣体塞进去。受到如此对待的人认为这是比其他一切更严厉的惩罚。此外，他们抓住妇女和儿童，强迫他们领受圣洗；如果有人抵抗或反对，立刻就是鞭打，鞭打之后是捆绑、监禁和其他暴力手段。我在此讲述一两件实例，使读者可以大致了解马其顿尼和当时掌权者所施行的严厉和残忍程度。他们先将那些不愿与他们共融的妇女的乳房压入盒子中，然后锯掉。其他妇女的同一身体部位，他们一部分用铁烫，一部分用火中加热至极热的鸡蛋。这种酷刑甚至异教徒都不知道，是由那些自称基督徒的人发明的。这些事是年长的诺洼天教会长老奥克萨农告诉我的，我在第一卷书中提到过他。他还说，在他晋升长老之前，他曾遭受亚流派不少严酷对待；他被投入监狱，与他的隐修伴侣帕夫拉戈尼亚的亚历山大一同多次受鞭打。他补充说，他自己能忍受这些酷刑，但亚历山大却因受刑而死于狱中。他现在葬在驶入称为凯拉斯的君士坦丁堡海湾的人的右侧，靠近河流，那里有一座以亚历山大命名的诺洼天派教堂。此外，亚流派在马其顿尼的煽动下，拆毁了各城市的许多教堂，其中包括靠近佩拉古斯的君士坦丁堡诺洼天派教堂。我为何特别提及这座教堂，可从与它相关的非凡情况看出，正如那位年长的奥克萨农所证实的。皇帝的谕令和马其顿尼的暴力已注定要摧毁那些持守\OriginalTerm{同质}{consubstantiality}道理之人的教会；这道谕令和暴力也波及了这座教堂，那些负责执行命令的人也已到场。我不禁赞叹诺洼天派在此场合所表现的热忱，以及他们从那些当时被亚流派驱逐、但现在平安拥有自己教会的人那里所得到的同情。因为当敌人的使者急于完成毁灭时，一大群诺洼天派信徒，在众多持同样情感之人的协助下，聚集在这座献身的教堂周围，将其拆毁，并把材料运到另一个地方：这地方在城对面，称为斯凯，是君士坦丁堡城的第十三区。这次迁移在极短的时间内完成，因为参与的人数众多，热情非凡：有人搬瓦片，有人搬石头，有人搬木料；各人肩负不同的东西。甚至妇女和儿童也协助工作，视之为他们最佳愿望的实现，认为被视作献给天主之物的忠实守护者是最大的光荣。就这样，当时诺洼天派的教堂被迁到了斯凯。许久之后，君士坦提乌去世，皇帝尤利安下令恢复其原址，并允许他们在那里重建。因此，民众如前一样将材料搬回，在原址重建了教堂；由于这一情况，以及其建筑和装饰的巨大改进，他们恰当地称之为\OriginalTerm{亚纳斯塔西亚}{Anastasia}。如我们先前所说，这教堂后来在尤利安统治时期得到重建。但那时，大公教会信徒和诺洼天派信徒都同样遭受迫害：因为前者厌恶在亚流派聚集的教堂中敬拜，却常去另外三座教堂——这是诺洼天派在城中所拥有的教堂数目——并与他们一同参与敬拜。事实上，若不是诺洼天派出于对古老规条的尊重而拒绝，他们早已完全合一。但在其他方面，他们彼此保持如此程度的真诚与友爱，以至于愿意为对方舍命。因此，双方不加区别地受到迫害，不仅在君士坦丁堡，也在其他省份和城市。在齐济库斯，该地的主教埃勒乌西乌对那里的基督徒施以同样的暴行，正如马其顿尼在其他地方所做的那样，骚扰并驱散他们，并且（其中一件事）他彻底拆毁了齐济库斯的诺洼天派教堂。但马其顿尼以下列方式完成了他的邪恶。他听说帕夫拉戈尼亚省有大量诺洼天派信徒，尤其是在曼提尼翁，并意识到这样众多的群体无法单凭教会人士赶离家园，便经皇帝许可，派遣四队士兵进入帕夫拉戈尼亚，以便他们因畏惧军队而接受亚流派的意见。但曼提尼翁的居民因宗教热忱而绝望，装备了长柄镰刀、斧头和手边任何可用的武器，出去迎击军队；结果发生冲突，确实有许多帕夫拉戈尼亚人被殺，但几乎全部士兵都被消灭了。我从一位帕夫拉戈尼亚农夫那里得知这些事，他说他当时在场；该省的许多其他人也证实了这一说法。这就是马其顿尼为基督教所建的功勋，包括谋杀、战斗、监禁和内战：这些行径不仅使他受到迫害对象的憎恶，甚至也为他自己的派别所憎恶。他也因此为皇帝所厌恶，尤其由于我即将叙述的情况。停放皇帝君士坦丁遗骸的棺材的教堂有倒塌的危险。因此，进入教堂的人和习惯在那里祈祷的人都十分恐惧。因此，马其顿尼希望移动皇帝的遗骸，以免棺材被废墟损坏。民众得知此事，竭力阻止，坚持说：\Quote{不应搅动皇帝的遗骨，因为这样的挖掘相当于将其挖出}；但许多人断言，移动遗骸不会损害尸体，于是这一问题形成了两派；持守\OriginalTerm{同质}{consubstantiality}道理的人与那些因其不敬而反对的人联合起来。马其顿尼完全不顾这些顾虑，令人将皇帝的遗骸运往安放殉道者阿卡西乌遗体的教堂。于是，一大群人分成两个敌对阵营冲向那座建筑，猛烈地互相攻击，造成大量伤亡，以致教堂庭院被血泊覆盖，教堂内的水井也溢出鲜血，流到邻近的门廊，甚至流到大街上。当皇帝得知这一不幸事件后，他对马其顿尼大为恼怒，既因为他造成了这场屠杀，也因为他竟敢未经请示就移动他父亲的遗体。因此，他留下西泽尤利安管理西方地区，自己动身前往东方。至于马其顿尼不久后被罢黜，从而为其臭名昭著的罪行遭受了极不相称的惩罚，我将在后面叙述。\par

% structure: p0144 source=h2 target=subsection reason=relative-heading-rank; base=h2

\subsection{Chapter 39. 论在伊索里亚的塞琉西亚所召开的会议。}

但我现在必须叙述另一场会议，即皇帝敕令在东方召集的、作为阿利米努姆会议对手的会议。起初决定主教们应在比提尼亚的尼科米底亚集合；但一场大地震几乎摧毁了该城，使得他们无法在那里集会。这事发生在塔提安和塞雷阿利斯任执政官期间，八月二十八日。因此他们计划将会议迁往附近的尼西亚城；但这个计划又被改变了，因为在西里西亚的塔尔苏斯集会似乎更为方便。他们对这个安排也不满意，最后终于聚集在伊索里亚的塞琉西亚（别号\Quote{阿斯佩拉}）。这事发生在同一年（阿利米努姆会议召开之年），即尤西比乌斯和伊帕提乌斯任执政官期间，与会人数约一百六十人。出席此次会议的有莱奥纳斯，一位宫廷显贵，皇帝敕令命在其面前进行关于信仰的讨论。伊索里亚的军队总司令劳里基乌斯也被命令到场，为主教们提供他们所需的服务。因此，在场这些显贵面前，主教们于九月二十七日聚集，并立即开始根据公开记录进行讨论，有速记员在场记录每个人的发言。凡欲了解各次发言详情者，可在萨比努斯的汇编中找到大量细节；我们只提及更重要的要点。在他们聚集的第一天，莱奥纳斯命令每人提出他认为合适的内容：但与会者说，在尚未到达的那些主教缺席的情况下，不应提出任何问题；因为君士坦丁堡主教马其顿尼、安基拉的巴西尔以及其他一些因担心行为被弹劾而尚未露面的人尚未到场。马其顿尼以身体不适为由未能出席；帕特罗菲鲁斯说他眼睛有毛病，因此需要留在塞琉西亚郊区；其余人则提出各种借口解释缺席原因。然而，当莱奥纳斯宣布，尽管这些人缺席，他们聚集起来要讨论的主题必须着手进行时，主教们回答说，在调查受指控者的生活和行为之前，他们无法着手讨论任何问题：因为耶路撒冷的济利禄、亚美尼亚塞巴斯特的尤斯塔提乌斯以及一些其他人很久以前就因各种理由被指控行为不当。由于这一异议，一场激烈的争论爆发了：一些人坚称应当首先审理所有这些控告，另一些人则否认任何事应当优先于信仰问题。皇帝的命令也加剧了这一争论，因为有人出示了他的信件，时而督促这个，时而督促那个必须优先考虑。由此产生争论后，分裂便发生了，塞琉西亚会议分为两派：一派由巴勒斯坦凯撒勒雅的阿卡西乌、亚历山大的乔治、提洛的乌拉尼乌和安提约基雅的尤多克西乌为首，他们仅获得约三十二名其他主教的支持。另一派人数多得多，其主要人物是叙利亚劳迪塞雅的乔治、帕弗拉戈尼亚庞培波利斯的索弗洛尼乌和基齐库斯的厄肋乌西乌。由于多数派决定先审查教义问题，阿卡西乌一派公开反对\WorkTitle{尼西亚信经}，并希望引入另一份信经取而代之。另一派人数相当多，赞同尼西亚会议的所有决定，但批评其采用了\OriginalTerm{同性同体}{homoousion}一词。因此他们就此点辩论，双方各抒己见，直到傍晚时分，塔尔苏斯教会的监督西尔瓦努斯以非常激烈的态度坚持说：\Quote{不需要一份新的信仰声明；他们的责任反而是确认在安提约基雅为那座教堂祝圣时所公布的那份。}听到这话，阿卡西乌及其同党私下退出了会议；其他人则取出在安提约基雅撰写的信经，念了一遍，然后那日散会。次日，他们在塞琉西亚的教堂集会，关上门后，再次念了同一份信经，并予以签署批准。当时，在场的诵经员和执事代表某些缺席但已表示同意其内容的主教签署了信经。\par

% structure: p0146 source=h2 target=subsection reason=relative-heading-rank; base=h2

\subsection{Chapter 40. 凯撒勒雅主教阿卡西乌在塞琉西亚会议上提出一份新的信经格式。}

\Person{阿卡奇乌斯}及其追随者批评所做的事：即他们关闭了教堂的大门，并以此方式附上他们的签署；声称\Quote{所有这些秘密的交易都理应受到怀疑，毫无效力可言。}他提出这些反对意见，是因为他急于提出另一份由他自己草拟的信仰陈述，这份陈述他已提交给总督\Person{勒翁纳斯}和\Person{劳瑞基乌斯}，现在他试图让这份陈述单独得到确认和确立，以取代已经签署的那份。第二天就这样被占用，除了他为此目标所做的努力外，别无他事。第三天，\Person{勒翁纳斯}试图促成双方的和好会议；\Place{君士坦丁堡}的\Person{马其顿尼乌斯}和\Place{安基拉}的\Person{巴西尔}在会议期间到达。但当阿卡奇乌斯派发现双方已达成相同立场时，他们拒绝会面，声称\Quote{不仅那些先前被革职的人，而且任何目前受到指控的人，都应被排除在会议之外。}经过双方多次争论，这一意见占了上风：那些受到任何指控的人退出了会议，阿卡奇乌斯派进入了他们的位置。\Person{勒翁纳斯}随后说，\Person{阿卡奇乌斯}交给他一份文件，他希望提请他们注意这份文件；但他没有说明这是一份信经草案，其中某些部分隐晦地，另一些部分明确地与前一份相矛盾。当在场的人都沉默时，以为文件内容除了信经陈述外还有别的东西，于是宣读了以下由\Person{阿卡奇乌斯}撰写的信经及其序言。\par

\Quote{我们于昨日奉皇帝之命，在\Place{伊苏里亚}的城市\Place{塞琉西亚}集会，时值九月二十七日，尽最大努力以一切温和方式维护教会的和平，并根据先知和福音的权威决定教义问题，以便在教会信仰宣认中不认可任何与圣经相违背的内容，正如我们最蒙天主喜爱的皇帝\Person{君士坦提乌斯}所命令的。但是，鉴于会议中的某些人对我们中的数人行为不公，阻止一些人有表达意见的自由，并违背其意愿将其他人排除在会议之外；同时，他们引入了已被革职的人以及违背教会法规被祝圣的人，以致会议呈现出一片混乱无序的景象，最显赫的伯爵\Person{勒翁纳斯}和最杰出的本省总督\Person{劳瑞基乌斯}都亲眼目睹。因此，我们有必要作出以下声明：我们不拒绝在\Place{安提约基雅}教堂祝圣时批准的信德，因为我们明确优先选择它，因为它得到了我们那些为讨论某些争议问题而聚集在\Place{安提约基雅}的神长们的一致同意。然而，既然\OriginalTerm{同质}{homoousion}和\OriginalTerm{类似本质}{homoiousion}这两个术语在过去曾使许多人困惑，并且仍持续困扰他们；此外，最近有人新创了一个术语，主张圣子与圣父\OriginalTerm{不同本质}{anomoion}：我们拒绝前两个术语，因为它们在圣经中找不到；但我们彻底诅咒最后一个，并将那些支持使用它的人视为与教会分离。我们明确认信圣子与圣父\OriginalTerm{相似}{homoion}，这与宗徒关于他所宣告的相符合：\BibleQuote{他是不可见的天主的肖像。}\BibleRef{哥 1:15}}\par

\Quote{我们信从并相信一位天主，全能的圣父，天地的创造者，以及一切有形无形之物的创造者。我们也相信他的圣子，我们的主耶稣基督，他在万世之前由圣父不受苦难而生，天主圣言，天主的独生子，光、生命、真理、智慧：藉着他，天上地上一切有形无形之物都得以造成。我们相信他在末世由圣童贞玛利亚取得肉身，为要消灭罪恶；他成为人，为我们的罪受苦，复活了，被提升天，坐在圣父的右边，将来必在荣耀中再来审判生者死者。我们也相信圣神，我们的主救主称他为安慰者，并在升天后按照他的应许派遣给门徒：藉着圣神，他在教会中圣化所有信徒，这些信徒因父及子及圣神之名受洗。凡传讲与此信经相悖的人，我们视之为与公教会分离的人。}\par

这是\Person{阿卡西乌斯}所提出的信仰宣言，由他本人及所有赞同他观点的人签署，其人数我们已提及。当这宣言被宣读后，\Person{索弗罗尼乌斯}（\Place{帕夫拉戈尼亚}的\Place{庞培奥波利斯}主教）发言说：\Quote{如果日日表达不同意见能被视为信仰的阐述，我们便永远无法达到对真理的准确理解。}这是\Person{索弗罗尼乌斯}的话。我坚信，如果这些主教的前任和继任者，对这\WorkTitle{尼西亚信经}抱有类似的情感，那么一切论战本可避免；教会也不致被如此猛烈而不合理的骚动所搅扰。然而，让那些能理解这些事情的人来判断吧。当时，在围绕这一教义声明以及被控各方发表了多方面的言论后，会议解散了。第四天，他们又在同一地点聚会，以从前那种争执的精神继续议事。这次，\Person{阿卡西乌斯}发言说：\Quote{既然\WorkTitle{尼西亚信经}不止一次，而是多次被修改，我们现在再发表另一信经并无阻碍。}\Person{库基库斯}主教\Person{厄勒乌西乌斯}回答说：\Quote{本次会议现在召开，不是为了学习先前不知道的事，也不是为了接受先前未赞同的信经，而是为了确认父老们的信仰，无论在生前还是死后，决不可背离此信仰。}\Person{厄勒乌西乌斯}这样反对\Person{阿卡西乌斯}，他所说的\Quote{父老们的信仰}指的是在\Place{安提约基雅}颁布的那份信经。但他也完全可以被如此合理地回应：\Quote{\Person{厄勒乌西乌斯}啊，你怎么称那些在\Place{安提约基雅}集会的人为\InnerQuote{父老们}，而你不承认那些是他们父老的人呢？制定\WorkTitle{尼西亚信经}、承认\Emph{同质（homoousios）}信仰的人，更有资格享有\InnerQuote{父老们}的称号；既因为他们时间上在先，也因为那些在\Place{安提约基雅}集会的人是因他们而被授予司祭职的。如果那些在\Place{安提约基雅}的人否认了自己的父老，那么跟随他们的人便不知不觉地追随了弒父者。此外，他们怎能从那些他们宣称信仰不正、不虔的人那里获得合法的圣职呢？然而，如果组成\Place{尼西亚}会议的人没有通过按手而赐予的\OriginalTerm{圣神}{Holy Spirit}，那么\Place{安提约基雅}的人就没有妥善接受司祭职：因为他们怎能从没有授予能力的人那里接受呢？}这些考虑本可以提交给\Person{厄勒乌西乌斯}以回应他的反对。但他们随后转入另一个问题，与\Person{阿卡西乌斯}在其信仰阐述中所作的断言相关，即\Quote{圣子相似圣父}；他们彼此询问这种相似在于什么。\Person{阿卡西乌斯}派肯定圣子只在意志上相似圣父，而非他的\Quote{实体}或\Quote{本质}；但其余的人坚持相似既包括本质也包括意志。在这一点的争论上，整整一天耗尽了；\Person{阿卡西乌斯}被他自己出版的著作所驳斥，他在其中曾断言\Quote{圣子在一切方面相似圣父}，他的对手于是问他：\Quote{你现在如何否认圣子与圣父在\InnerQuote{本质}上的相似呢？}\Person{阿卡西乌斯}回答说：\Quote{没有哪位作者，无论是古代还是现代的，曾因自己的著作而被定罪。}他们继续就此事进行冗长乏味的讨论，充满尖锐的情绪和微妙的论辩，但毫无达成一致判断的迹象，于是\Person{莱奥纳斯}起身解散了会议：这就是\Place{塞琉西亚}会议的结局。因为在次日，\Person{莱奥纳斯}虽被敦促，却不愿再与他们聚会。他说：\Quote{我受皇帝委派，参加一个期望达成一致的会议；但既然你们完全无法相互理解，我不能再出席了：如果你们愿意，就到教堂去，在那里沉溺于无聊的空谈吧。}\Person{阿卡西乌斯}派认为这一决定对他们有利，也拒绝与他人会面。反对派独自在教堂聚集，并要求\Person{阿卡西乌斯}的追随者出席，以便审理\Place{耶路撒冷}主教\Person{西里尔}的案件：因为这位主教早在很久以前就受到控告，但究竟基于什么理由，我无法说明。他甚至已被撤职，因为出于恐惧，他在两年整的时间里没有露面，尽管他曾多次被传唤以便调查对他的指控。然而，当他被撤职时，他给那些定他罪的人发送了一份书面通知，声称他将上诉到更高一级的司法机构；皇帝\Person{君士坦提乌斯}批准了这一上诉。\Person{西里尔}是第一位，也是唯一一位敢于打破教会惯例，成为上诉人（如同世俗法庭通常的做法）的神职人员；他现在在\Place{塞琉西亚}，准备接受审判；因此其他主教邀请\Person{阿卡西乌斯}派在集会中就座，以便在一次全体会议上对那些被控告者的案件作出明确的判决：因为他们同时也传唤了其他被控各种不法行为的人到庭，这些人为了自保，已逃到\Person{阿卡西乌斯}的追随者之中。因此，当那一派在多次被召之后仍坚持拒绝会面时，主教们撤职了\Person{阿卡西乌斯}本人，以及\Place{亚历山大里亚}的\Person{乔治}、\Place{提洛}的\Person{乌拉尼乌斯}、\Place{弗吕吉亚}的\Place{凯雷塔皮}的\Person{德奥杜卢斯}、\Place{吕底亚}的\Place{菲拉德尔斐亚}的\Person{德奥多西乌斯}、\Place{米提利尼岛}的\Person{厄瓦格里乌斯}、\Place{吕底亚}的\Place{特里波利斯}的\Person{良提乌斯}，以及\Person{尤多克西乌斯}，他原为\Place{日耳曼尼卡}的主教，但后来钻营到\Place{叙利亚}的\Place{安提约基雅}主教之位。他们也撤职了\Person{帕特罗斐卢斯}，因其藐视法庭，没有出面回答一位名叫\Person{多罗特}的司铎对他的指控。这些人他们撤职了；他们还绝罚了\Person{阿斯特里乌斯}、\Person{欧塞比乌斯}、\Person{阿布加鲁斯}、\Person{巴西利库斯}、\Person{斐玻斯}、\Person{斐德利斯}、\Person{欧提基乌斯}、\Person{玛格努斯}和\Person{欧斯塔提乌斯}；决定他们不得恢复共融，直到他们作出辩护，洗清他们所背负的指控为止。做完这些后，他们向每个主教被撤职的教会寄发了说明信。随后\Person{阿尼亚努斯}被任命为\Place{安提约基雅}主教，取代\Person{尤多克西乌斯}；但\Person{阿卡西乌斯}派很快就逮捕了他，将他交到\Person{莱奥纳斯}和\Person{劳里基乌斯}手中，由他们将他流放。祝圣他的主教们对此感到愤怒，向\Person{莱奥纳斯}和\Person{劳里基乌斯}提出抗议，公开指控他们违反了会议的决定。发现通过这种方式无法获得纠正，他们便前往\Place{君士坦丁堡}，将整个事情呈报给皇帝。\newline{}\par

% structure: p0151 source=h2 target=subsection reason=relative-heading-rank; base=h2

\subsection{第四十一章 论皇帝从西方返回，\OriginalTerm{阿卡奇派}{Acacians}在\Place{君士坦丁堡}聚集，并确认了\Place{阿里米努姆}信经，且对其作了一些增补。}

现在，皇帝从西方返回，并任命了一名叫\Person{奥诺拉图斯}的城长官管理\Place{君士坦丁堡}，同时废除了总督的职位。但\OriginalTerm{阿卡奇派}{Acacians}抢先于主教们，向皇帝控告他们，说服皇帝不要接受他们提出的信经。这使皇帝非常恼怒，决定驱散他们；因此他颁布了一项法令，命令那些应担任某些公职的人不再免于履行其职责。因为他们中有几个人可能会被召去担任各种公职，既与城市行政长官有关，也隶属于各省的省长和总督。当这些人受到骚扰时，\Person{阿卡奇}的追随者在\Place{君士坦丁堡}停留了相当长的时间，并召开了另一次主教会议。他们派人去请\Place{比提尼亚}的主教们，这次大约有五十位主教聚集，其中包括\Place{加采东}的主教\Person{玛里斯}：他们确认了在\Place{阿里米努姆}宣读的信经，该信经前面已附有执政官的名字。如果不是对它作了一些增补，在这里重复它是没有必要的；但既然已经做了，也许取录它的新形式是合适的。\par

\Quote{我们信惟一的天主，全能的圣父，万物由他而来。并信天主的独生子，在万世之前、在一切起始之先由天主所生，借着祂一切有形无形的受造物得以造成。祂是圣父所生的独生子，惟一者出自惟一者，天主出自天主，与生祂的圣父相似，符合圣经；祂的受生除生祂的圣父外无人知晓。我们知道天主的这个独生子，如经上所载，由圣父派遣，从天降下，为消灭罪恶与死亡；祂按肉体由圣神和童贞玛利亚所生，如经上所载，并与祂的门徒来往；当按圣父的旨意完成了每一项安排之后，祂被钉十字架且死去，被埋葬并下降到地下的深处，阴府本身在祂面前战栗。祂在第三日从死者中复活，再次与门徒交谈，满了四十天后被提升天，坐在圣父的右边，在末日，即复活之日，将带着圣父的荣耀降临，按每个人的行为予以报应。［我们信］也信圣神，就是天主的独生子、我们的主和天主基督，应许要差遣给人类作护慰者的那一位，正如经上所载，\BibleRef{若 15:26} \InnerQuote{真理之神}；祂在升天之后将圣神赐给了他们。但因教父们所使用的\OriginalTerm{ousia}{本质}一词，原意简单明了，但因百姓不理解而成为绊脚石，我们认为应弃绝此词，因为它甚至不包含在圣经中；今后不应再提及，因为圣经从未提过圣父和圣子的本质。也不应再提圣父、圣子和圣神的\Quote{自立}。但我们肯定圣子相似圣父，如同圣经所宣告和教导的那样。愿所有反对这一信仰宣示的异端，无论是已被谴责的还是近来兴起的，都受弃绝。}\par

这些事当时在\Place{君士坦丁堡}得到了认可。现在，我们终于穿过了各种信仰形式的迷宫，让我们数一数它们的数目。在\Place{尼西亚}颁布的那份之后，在\Place{安提约基雅}教堂奉献礼上又提出了两份。第三份由\Person{纳尔奇苏}及其同行者在\Place{高卢}呈给皇帝。第四份由\Person{欧多克西乌斯}送到\Place{意大利}。在\Place{西尔米乌姆}公布了三份信经，其中一份附有执政官名字的在\Place{阿里米努姆}宣读。\OriginalTerm{阿卡奇派}{Acacian party}在\Place{塞琉西亚}提出了第八份。最后一份是\Place{君士坦丁堡}的，包含了禁止提及关于天主的\Quote{本质}或\Quote{自立}的条款。哥特人的主教\Person{乌尔菲拉斯}同意了这份信经，尽管他先前赞同\Place{尼西亚}的；因为他是哥特人的主教\Person{特奥斐洛}的门徒，后者曾出席尼西亚大公会议，并签署了那里的决定。关于这些事，到此为止。\par

% structure: p0155 source=h2 target=subsection reason=relative-heading-rank; base=h2

\subsection{第四十二章 论\Person{马其顿尼}的被废，\Person{欧多克西乌斯}获得\Place{君士坦丁堡}主教席位。}

\Person{阿卡西乌斯}、\Person{欧多克修斯}以及那些在\Place{君士坦丁堡}与他们联合的人，变得极其焦虑，希望他们这一方也能罢免对方的一些人。现在应当注意到，两派在罢免他人时都没有受到宗教考虑的支配，而是出于其他动机：因为虽然他们在信仰上不一致，但他们相互罢免的理由并非教义错误。因此，阿卡西亚派利用皇帝对他人、尤其是对\Person{马其顿尼乌斯}的愤怒——皇帝正在积蓄这种愤怒并急于发泄——首先罢免了\Person{马其顿尼乌斯}，既因为他造成了如此多的杀戮，也因为他接纳了一个犯有奸淫罪的执事共融。然后他们罢免了\Place{基齐库斯}主教\Person{埃琉修斯}，因为他为\Place{提尔}的一个名叫\Person{赫拉克利乌斯}的赫拉克勒斯祭司施洗并随后授予他执事职，此人以行巫术闻名。对\Place{安基拉}的\Person{马塞尔鲁斯}的继任者\Person{巴西尔}（或称\Person{巴西拉斯}）也宣判了类似的判决：谴责的理由是，他不公正地监禁了某人，给他戴上镣铐并施以酷刑；他诽谤了一些人；并且他通过书信扰乱了非洲的教会。\Person{德拉孔提乌斯}也被罢免，因为他离开了\Place{加拉太}教会去了\Place{佩尔加莫斯}。此外，他们以各种借口罢免了\Place{塞琉西亚}（召开主教会议的城市）主教\Person{内奥纳斯}、\Place{帕夫拉戈尼亚}的\Place{庞培城}主教\Person{索弗洛尼乌斯}、\Place{马其顿}的\Place{萨塔拉}主教\Person{埃尔皮迪乌斯}、\Place{耶路撒冷}的\Person{西里尔}等人，各有不同的理由。\par

% structure: p0157 source=h2 target=subsection reason=relative-heading-rank; base=h2

\subsection{论\Place{塞巴斯特}主教\Person{欧斯塔修斯}}

但\Place{亚美尼亚}的\Place{塞巴斯特}主教\Person{欧斯塔修斯}甚至不被允许为自己辩护；因为他很久以前就被自己的父亲、\Place{卡帕多细亚}的\Place{凯撒利亚}主教\Person{欧拉留斯}以衣着不合司祭职份为由罢免了。应当注意，\Person{梅利提乌斯}被任命为他的继任者，关于此人我们将在后面叙述。\Person{欧斯塔修斯}后来确实在\Place{帕夫拉戈尼亚}的\Place{甘格拉}为他召开的主教会议上被定罪；他在\Place{凯撒利亚}会议罢免他之后，做了许多违反教会法规的事。因为他\Quote{禁止婚姻}，并主张应戒除肉类；他甚至使许多人离开他们的妻子，并说服那些不喜欢在教会中聚会的人在家中共融。在虔诚的伪装下，他也引诱仆人离开主人。他自己穿着哲学家的服饰，并引导他的追随者采用一种新的、奇特的服装，命令妇女剪短头发。他允许忽略规定的斋戒，但建议在星期日斋戒。简言之，他禁止在有结婚的人家中祈祷；并宣称，一个在还是平信徒时合法结婚、并继续与妻子同居的司铎，其祝福和共融都应被视为可憎而加以回避。因为他做了并教导这些事以及许多类似的事，正如我们所说，在\Place{帕夫拉戈尼亚}的\Place{甘格拉}召开的主教会议罢免了他，并谴责他的意见。但这是后来发生的事。但在\Person{马其顿尼乌斯}被逐出\Place{君士坦丁堡}教区后，\Person{欧多克修斯}——他现在认为\Place{安提约基雅}教区已属次要——被提升到空缺的主教职位；由阿卡西亚派祝圣，他们在此事上并不考虑这与他们先前的行动不一致。因为他们曾因\Person{德拉孔提乌斯}从\Place{加拉太}转到\Place{佩尔加莫斯}而罢免他，却在祝圣\Person{欧多克修斯}时明显违反了自己的原则和决定，因为后者当时做了第二次变动。此后，他们将经过修正和补充的信仰宣言送往\Place{阿利米农}，命令所有拒绝签署的人根据皇帝的法令流放。他们还告知东方其他与他们意见一致的主教他们所做之事；尤其是\Place{西托波利斯}主教\Person{帕特罗菲鲁斯}，他离开\Place{塞琉西亚}后直接回到了自己的城市。\Person{欧多克修斯}被立为帝都的主教后，名为\WorkTitle{索菲亚}的大教堂在当时被祝圣，时在\Person{君士坦提乌斯}第十次执政、\Person{尤利安·凯撒}第三次执政之年，2月15日。正是在\Person{欧多克修斯}占据这个教区时，他首先说出了那句至今仍在流传的话：\Quote{圣父是不恭敬的，圣子是恭敬的。}当人们对此表述似乎感到震惊，并开始骚动时，他说：\Quote{不要为刚才所说的话烦恼：因为圣父是不恭敬的，因为他不崇敬任何人；但圣子是恭敬的，因为他崇敬圣父。} \Person{欧多克修斯}说完这话后，骚动平息了，教堂里爆发出大笑：他的这句话直到今天仍是一个笑柄。异端首领们确实经常发明这样微妙的短语，并以此撕裂教会。\Place{君士坦丁堡}的主教会议就这样结束了。\par

% structure: p0159 source=h2 target=subsection reason=relative-heading-rank; base=h2

\subsection{论\Place{安提约基雅}主教\Person{梅利提乌斯}}

现在该谈论默肋提乌斯（\Person{Meletius}）了。我们不久前曾提到，在欧斯塔提乌斯（\Person{Eustathius}）被罢免后，他被立为亚美尼亚的\Place{Sebastia}主教；从\Place{Sebastia}他又被调往叙利亚的\Place{Berœa}。他出席了塞琉西亚（\Place{Seleucia}）的会议，签署了阿卡基乌斯（\Person{Acacius}）在那里提出的信经，随即返回\Place{Berœa}。当君士坦丁堡（\Place{Constantinople}）的会议召开时，安提约基雅（\Place{Antioch}）的民众发现尤多克西乌斯（\Person{Eudoxius}）因被君士坦丁堡主教座的辉煌所吸引，而轻视了他们的教会，于是他们请来默肋提乌斯，授予他安提约基雅教会的主教职。起初他回避一切教义问题，只讲论道德主题；但后来他向听众阐释了\WorkTitle{尼西亚信经}，并主张\OriginalTerm{同一实体}{\GreekText{ὁμοούσιον}}的教义。皇帝得知此事后，下令将他流放，并让先前与亚略（\Person{Arius}）一同被罢免的尤佐伊乌斯（\Person{Euzoïus}）接替他担任安提约基雅主教。然而，那些依恋默肋提乌斯的人脱离了亚略派集会，单独举行聚会；尽管如此，最初接受\OriginalTerm{主张同一实体者}{\GreekText{ὁμοούσιον}}意见的人不肯与他们共融，因为默肋提乌斯是由亚略派按立的，他的追随者也由亚略派施洗。于是安提约基雅教会内部产生了分裂，即使在信仰观点完全一致的人之间也是如此。与此同时，皇帝得知波斯人正准备对罗马人发动另一场战争，便匆忙赶往\Place{Antioch}。\par

% structure: p0161 source=h2 target=subsection reason=relative-heading-rank; base=h2

\subsection{第四十五章 马其顿纽斯的异端}

马其顿纽斯（\Person{Macedonius}）被逐出君士坦丁堡后，不堪忍受对他的判决，变得不安分；因此他加入了在塞琉西亚罢免阿卡基乌斯及其同党的另一个派别，并派遣代表团前往索弗若尼乌斯（\Person{Sophronius}）和厄肋乌西乌斯（\Person{Eleusius}）处，鼓励他们坚持最初在安提约基雅公布、后在塞琉西亚确认的信经，并提议给它一个冒名：\OriginalTerm{同质相似}{\GreekText{ὁμοιούσιον}}信经。借此他吸引了大量追随者，这些人从他得名，至今仍被称为\Quote{马其顿纽斯派}。尽管在塞琉西亚会议中与阿卡基乌派意见不同的人此前并未使用\OriginalTerm{同质相似}{\GreekText{ὁμοιούσιος}}一词，但从那时起他们明确主张此说。然而有传言说，这个词并非源于马其顿纽斯，而是出自玛拉托纽斯（\Person{Marathonius}）的创造，此人不久前被派往尼科美底亚（\Place{Nicomedia}）教会；因此，持守这一教义的人也被称为\Quote{玛拉托纽斯派}。欧斯塔提乌斯（\Person{Eustathius}）也加入了这一派，他因前述原因被逐出塞巴斯蒂亚教会。但当马其顿纽斯开始否认天主圣三中圣神的神性时，欧斯塔提乌斯说：\Quote{我既不能承认圣神是天主，也不敢断言祂是受造物。}因此，那些主张圣子\OriginalTerm{同一实体}{\GreekText{ὁμοούσιον}}的人称这些异端者为\OriginalTerm{圣神仇敌}{\GreekText{Πνευματομάχοι}}。这些马其顿纽斯派在赫勒斯滂（\Place{Hellespont}）何以变得如此众多，我将在适当之处叙述。与此同时，阿卡基乌派变得极为急切，希望在安提约基雅召开另一次会议，因为他们改变了先前关于圣子与圣父\Quote{在一切方面}相似的看法。于是他们中的少数人在随后的执政年——即托鲁斯（\Person{Taurus}）和弗洛伦提乌斯（\Person{Florentius}）担任执政官的那一年——聚集在叙利亚的\Place{Antioch}，当时皇帝正驻跸于此，尤佐伊乌斯（\Person{Euzoïus}）为主教。他们重新讨论了先前决定的某些点，在讨论中他们宣布应从阿里米农（\Place{Ariminum}）和君士坦丁堡（\Place{Constantinople}）公布的信德格式中删除\OriginalTerm{相似}{\GreekText{ὅμοιος}}一词；他们不再掩饰，公开宣称圣子与圣父全然不相似，不仅在本质方面，甚至在本意方面也是如此；还像亚略（\Person{Arius}）早已做的那样，大胆断言圣子是从无中受造的。该城中支持艾提乌斯（\Person{Aëtius}）异端的人赞同这一意见；因此，除了一般称为亚略派之外，在安提约基雅那些主张\OriginalTerm{同一实体}{\GreekText{ὁμοούσιον}}的人（他们当时因默肋提乌斯而内部分裂，如前所述）也称他们为\OriginalTerm{非相似派}{\GreekText{Ἀνόμοιοι}}和\OriginalTerm{从无而生者派}{\GreekText{Ἐξ} \GreekText{οὐκ} \GreekText{ὄντων}}。因此，当主张同一实体者问他们，既然在过去信经中承认圣子是\Quote{出于天主的天主}，何以还敢断言圣子与圣父不相似且从无中获取存在时，他们用这样似是而非的遁词来回避异议：\Quote{\Quote{\InnerQuote{天主出于天主}}这个说法，他们说，应与宗徒的话——\BibleRef{格前 11:12}——\Quote{\InnerQuote{但万有出于天主}}作同样理解。因此，圣子是\Quote{\InnerQuote{出于天主}}的，作为其中之一；正因如此，在信经草案中加了\Quote{\InnerQuote{照经上所说}}这句话。}这一诡辩的作者是老底嘉（\Place{Laodicea}）的乔治（\Person{George}）主教，他不谙此类措辞，不知道奥利振（\Person{Origen}）当初如何彻底探究了此事，并对宗徒的这些独特说法作了解释。但尽管有这些逃避的巧辩，他们仍无法承受由此招来的指责和羞辱，于是退回先前在君士坦丁堡提出的信经；这样，每个人各自回到自己的管区。乔治返回亚历山大里亚（\Place{Alexandria}），重新掌握那里教会的主权，亚大纳削（\Person{Athanasius}）仍未露面。他迫害城中反对他意见的人，行事严厉残酷，使自己极受民众憎恶。在耶路撒冷（\Place{Jerusalem}），阿雷尼乌斯（\Person{Arrenius}）代替济利禄（\Person{Cyril}）被立为教会领袖；我们也可以说，在他之后，赫拉克利乌斯（\Person{Heraclius}）被祝圣为主教，之后是希拉略（\Person{Hilary}）。最后，济利禄回到耶路撒冷，再度受任管理那里的教会。大约同一时期，又兴起另一种异端，源于以下情况。\par

% structure: p0163 source=h2 target=subsection reason=relative-heading-rank; base=h2

\subsection{第四十六章 亚波里拿留派及其异端}

在叙利亚的\Place{劳狄刻雅}，有同名两人，一为父，一为子；其名均为\Person{阿波利纳里斯}。前者为该教会的\OriginalTerm{长老}{presbyter}，后者为\OriginalTerm{读经员}{reader}。二人均教授希腊文学，父亲教授文法，儿子教授修辞。父亲原籍\Place{亚历山大里亚}，起初在\Place{贝鲁特}任教，后迁至\Place{劳狄刻雅}，在那里成婚，生下小\Person{阿波利纳里斯}。他们是修辞学家\Person{埃皮法尼乌斯}的同时代人，因真诚相待，彼此成为密友；但\Place{劳狄刻雅}的主教\Person{塞奥多图斯}担心这种交往会败坏他们的信仰，引他们陷入异教，遂禁止他们与他来往；然而他们对此禁令几乎不加理会，与\Person{埃皮法尼乌斯}的亲密关系依然如故。\Person{塞奥多图斯}的继任者\Person{乔治}也试图阻止他们与\Person{埃皮法尼乌斯}交谈；但无法以任何方式说服他们放弃此事，便对他们施行了\OriginalTerm{绝罚}{excommunicate}。小\Person{阿波利纳里斯}认为这一严厉处置不公，凭借其修辞诡辩的才能，创立了一个以发明者命名的\OriginalTerm{异端}{heresy}，至今仍有许多支持者。尽管如此，有些人断言，他们与\Person{乔治}分离并非出于上述原因，而是因为他们看到其信仰宣认的动摇与前后矛盾；因为他有时主张圣子与圣父相似，符合\Place{塞琉西亚}会议的决定，有时又附和\OriginalTerm{亚流主义}{Arian}观点。因此他们把这一点作为与他分离的借口：但因无人效法他们，他们便引入了一种新的教义形式，起初他们主张，在\OriginalTerm{降生成人}{incarnation}的\OriginalTerm{降生奥迹}{economy of the incarnation}中，\OriginalTerm{天主圣言}{God the Word}取了一个没有灵魂的人性身体。后来，仿佛改变了主意，他们收回此说，承认他确实取了灵魂，但那是一个\OriginalTerm{无理性灵魂}{irrational soul}，而天主圣言自身代替了理智。那些追随他们、至今以他们之名称呼的人断言，这是他们[与公教会]的唯一区别；因为他们承认\OriginalTerm{天主圣三}{Trinity}中\OriginalTerm{位格}{person}的\OriginalTerm{同体}{consubstantiality}。我们将在适当之处进一步提及这两位\Person{阿波利纳里斯}。\par

% structure: p0165 source=h2 target=subsection reason=relative-heading-rank; base=h2

\subsection{第四十七章 尤利安的成功；君士坦提乌斯皇帝的驾崩。}

当皇帝\Person{君士坦提乌斯}继续驻跸\Place{安提约基雅}时，\Person{尤利安}凯撒在\Place{高卢}与一支庞大的\OriginalTerm{蛮族}{barbarians}军队交战，击败了他们，因此在军旅中极受欢迎，并被他们拥立为皇帝。此事传开后，皇帝\Person{君士坦提乌斯}深感痛苦；于是他被\Person{尤佐伊乌斯}\OriginalTerm{受洗}{baptized}，并立即准备出征讨伐\Person{尤利安}。当他抵达\Place{卡帕多奇雅}与\Place{奇里乞亚}边境时，因过度心神不宁而引发中风，在\Place{摩普苏克莱尼}结束了他的生命，时当\OriginalTerm{执政官}{consulate}\Person{陶鲁斯}与\Person{弗洛伦提乌斯}之任期，十一月三日。这是第285\OriginalTerm{奥林匹亚德}{Olympiad}的第一年。\Person{君士坦提乌斯}享年四十五岁，在位三十八年；其中十三年他与父亲共同执政，父亲去世后又独自治理二十五年，后一时期的历史已包含在此卷中。\par

\SourceRecord{Translated by A.C. Zenos.；Nicene and Post-Nicene Fathers, Second Series；Vol. 2.；Edited by Philip Schaff and Henry Wace.；Buffalo, NY: Christian Literature Publishing Co.,；1890.；<http://www.newadvent.org/fathers/26012.htm>.}

% structure: p0001 source=h1 target=section reason=page-title

\section{\WorkTitle{教会史（第三卷）}}

% structure: p0002 source=h2 target=subsection reason=relative-heading-rank; base=h2

\subsection{第一章  论尤利安：其出身、教育、登基及背教归入异教}

君士坦提乌斯皇帝于十一月三日，在陶鲁斯与弗洛伦提乌斯执政期间，死于西里西亚边境；尤利安于同年十二月十一日左右离开帝国西部，在相同执政官治下前往君士坦丁堡，在那里被拥立为皇帝。既然我必须论及这位以学识著称的君王的品格，请他的仰慕者勿期望我采用华丽浮夸的修辞风格，仿佛描述必须与主题的尊严相称；因为我旨在编纂一部基督宗教史，为使易于理解并符合我起初的目的，保持谦卑朴素的笔调既是合宜的，也是一贯的。然而，描述其外貌、出身、教育及获得主权的方式是合适的；为此，有必要先叙述一些背景。为君士坦丁堡命名的君士坦丁，有两位同父异母的兄弟，即达尔马提乌斯和君士坦提乌斯。前者有一个与他同名的儿子；后者有两个儿子，加卢斯和尤利安。在君士坦丁堡建立者君士坦丁去世时，士兵杀死了其弟小达尔马提乌斯，他的两个孤儿也性命堪忧；但加卢斯因一场危及生命的疾病而免遭其父谋杀者的暴力；而尤利安因年幼——当时他只有八岁——而得以保全。皇帝对他们的猜忌逐渐平息后，加卢斯到伊奥尼亚的以弗所就学，那里留给他们相当可观的祖产。尤利安长大后，在君士坦丁堡继续学业，他常去宫殿（当时学校设在那里），衣着朴素，由太监马尔多尼乌斯监督。文法方面，拉刻代蒙人尼科克勒斯是他的教师；修辞则师从当时身为基督徒的智者埃克博利乌斯；因为皇帝曾规定他不得有异教教师，以免他被引诱至异教迷信。起初尤利安是一名基督徒。他的学业进步迅速，很快便流传说他能治理罗马帝国；这一民间传闻广泛散布，使皇帝内心大为不安，于是将他从大城调至尼科米底亚，同时禁止他前往叙利亚人利巴尼乌斯（一位智者）的学校。因为利巴尼乌斯当时被君士坦丁堡的教师们联合驱逐，退居尼科米底亚，并在那里开设学校。他在一篇针对这些教师所写的文章中表达了对他们的愤慨。然而，尤利安被禁止成为他的听众，因为利巴尼乌斯在宗教上是异教徒；但尤利安私下获取了他的演说稿，不仅极为赞赏，而且经常专心研读。他在修辞艺术上日渐精湛时，哲学家马克西穆斯来到尼科米底亚（不是欧几里得的父亲拜占庭人马克西穆斯，而是以弗所人马克西穆斯，皇帝瓦伦提尼安后来以施行巫术的罪名处死了他）。那是后来的事；当时吸引他到尼科米底亚的唯有尤利安的名声。尤利安从他那里，除了哲学原理之外，还接受了他自己的宗教情感，以及获取帝国的渴望。当这些事传到皇帝耳中时，尤利安介于希望与恐惧之间，急于平息已被激起的怀疑，于是开始表现出他从前实际所是的外在伪装。他剃光头发，假装过隐修生活；私下研习哲学，公开阅读基督徒的圣经，并且被任命为尼科米底亚教堂的读经员。通过这些花言巧语，他成功避免了皇帝的不悦。他做这一切是出于恐惧，但并未放弃希望；他告诉朋友们，更幸福的时刻不远了，他将拥有帝国统治权。在此情形下，他的兄弟加卢斯被立为凯撒，在前往东方途中到尼科米底亚看望他。但不久后加卢斯被杀，尤利安受到皇帝怀疑；因此皇帝下令监视他；但他很快逃脱，四处躲避，设法安全无恙。最后，皇后欧西比亚发现了他的藏身之处，说服皇帝不加害于他，并允许他去雅典继续研习哲学。简言之，皇帝从那里召回了他，立他为凯撒；此外还将自己的妹妹海伦许配给他，派他去对付蛮族。因为君士坦提乌斯皇帝曾雇佣来对抗僭主玛格嫩提乌斯的蛮族援军，在对篡位者毫无用处后，开始劫掠罗马城市。由于他年轻，皇帝命他在未咨询其他军事将领的情况下不得擅自行动。\par

这些将领获得如此权力后，便松懈了职责，蛮族因此强盛起来。\Person{朱利安}察觉此事，允许将领们纵情享乐、狂欢作乐，但自己则努力激发士兵的勇气，规定杀死一名蛮族者可得约定赏金。这一措施有效地削弱了敌人，同时赢得了军队的爱戴。据传，当他进入一座城镇时，悬于两柱之间的公民冠冕落在他头上，恰好合其头围；在场者无不欢呼赞叹，视之为他日后成为皇帝的预兆。有人断言，\Person{君士坦提乌斯}派他攻打蛮族，是希望他在交战中阵亡。我不知道说这话的人是否属实；但他既已与他缔结如此亲密的联盟，却又寻求他的毁灭，损害自身利益，这确实不太可能。各人可自行判断。\Person{朱利安}向皇帝抱怨其军官的懈怠，结果获得了一位更符合自己热忱的副手；两人合力对蛮族发起猛攻，以致蛮族派来使节，保证他们是奉皇帝的信函（信函被出示）之命入侵罗马领土。但\Person{朱利安}将使者投入监狱，并猛烈攻击敌军，大获全胜；俘虏了他们的国王，将其活生生送给\Person{君士坦提乌斯}。这一辉煌胜利之后，他立即被士兵拥立为皇帝；因当时没有皇冠，一名侍卫取下自己颈上的链子，系在\Person{朱利安}头上。\Person{朱利安}就这样成了皇帝；至于他后来是否举止如哲学家，让读者判断吧。因为他既未派遣使节与\Person{君士坦提乌斯}沟通，也未对他以往的恩惠表示丝毫敬意；反而任命了其他总督管理行省，一切随心所欲。此外，他还试图使\Person{君士坦提乌斯}蒙羞，在每座城市公开宣读他写给蛮族的信件；从而使得居民不满，轻易从\Person{君士坦提乌斯}那里倒向他自己。此后，他不再伪装基督徒，而是到处开放异教神庙，向偶像献祭；并自称为\OriginalTerm{大司祭}{Pontifex Maximus}，允许愿意者庆祝其迷信节日。以此方式，他挑起了对抗\Person{君士坦提乌斯}的内战；就他而言，这几乎使帝国陷入战争的一切灾难后果。因为这位哲学家的目标若不流血就无法实现；然而天主，以他自身主权之意愿，藉着除去其中之一，在无损国家的情况下制止了这两位对手的愤怒。当\Person{朱利安}抵达色雷斯时，传来\Person{君士坦提乌斯}已死的消息；如此，罗马帝国当时得以幸免于迫在眉睫的内部纷争。\Person{朱利安}随即进入\Place{君士坦丁堡}，思忖如何最好地安抚民众、赢得民心。于是他采取了以下措施：他知道\Person{君士坦提乌斯}因将拥护\OriginalTerm{同质}{homoousian}信仰者逐出教堂、放逐其主教而招致仇恨。他也知道异教徒因禁止献祭、渴望开放神庙、自由举行偶像崇拜礼仪而极度不满。事实上，他感觉到这两派都对其前任心怀怨恨，而一般民众对太监的暴行，尤其是对御前总管\Person{欧瑟伯}的贪婪极为愤慨。在此情况下，他诡诈地对待各方：对一些人他装作不知，对另一些人则施以恩惠来笼络（因为他喜欢表现仁慈），但总体上他显明自己对异教崇拜的偏爱。首先，为了玷污\Person{君士坦提乌斯}的记忆，使其显得对臣民残暴，他召回了被放逐的主教，归还了他们被没收的财产。接着，他命令适当官员立刻开放异教神庙。然后，他指示那些曾被太监勒索的人取回被掠夺的财产。御前总管\Person{欧瑟伯}被他处以死刑，不仅因为他加害于人，更因他确信正是此人暗中策划害死了他的兄弟\Person{加卢斯}。他对\Person{君士坦提乌斯}的遗体施以帝王葬礼，却将太监、理发师和厨师逐出宫廷。他辞退太监，是因为妻子已故，他决定不再婚娶，太监已无必要；辞退厨师，是因为他饮食极为简朴；辞退理发师，是因为他说一人足以为许多人理发。他以此理由辞退了他们；他还将大部分秘书恢复原职，并为其留下者指定了与职务相称的薪俸。他还改革了公务旅行和物资运输方式，废除了使用骡、牛、驴，只允许用马匹。这些削减措施虽为少数人所称赞，但其他所有人都强烈谴责，认为它们剥夺了那种对平民思想有强大影响的排场和奢华点缀，从而有损皇威。不仅如此，他夜间常熬夜撰写演说辞，随后在元老院发表；事实上，自\Person{尤利乌斯·凯撒}以来，他是第一位也是唯一一位在该集会演讲的皇帝。他极力庇护文学造诣卓越者，尤其那些职业哲学家；因此，大批此类假学者从各地涌入宫廷，披着希腊式外套，衣着比学问更引人注目。这些骗子总是附和君主的宗教情感，对基督徒的福祉满怀敌意；而\Person{朱利安}本人，因极度虚荣，在一本题为\WorkTitle{凯撒}的书中嘲讽所有前任，同样因这种傲慢习性而写文章反对基督徒。辞退厨师和理发师，在某种意义上固然合乎哲学家风度，却不合皇帝身份；而嘲笑和丑化他人，既非哲学家所为，也非皇帝所为：因为此类人物应超越嫉妒和诋毁的影响。一个皇帝可以在节制和自制方面成为哲学家；但若哲学家试图模仿皇帝应有的行为，则常会背离其自身原则。我们如此简略地讲述了\Person{朱利安}皇帝的事迹，追溯其出身、教育、性情以及他获得皇权的方式。\par

% structure: p0005 source=h2 target=subsection reason=relative-heading-rank; base=h2

\subsection{第二章 亚历山大城之动乱及乔治被杀。}

现在宜提及在同一皇帝统治下教会中所发生的事。亚历山大城因以下情形发生了一场大骚动。该城中有一处地方，久已荒废污秽，异教徒曾在此举行他们的神秘仪式，并向\Person{密特拉}献人祭。这地方空置无用，\Person{君士坦提乌斯}便将其授予亚历山大教会；\Person{乔治}想在原址上建造一座教堂，遂下令清理该地。在清理过程中，发现了一个极深的\OriginalTerm{内殿}{adytum}，揭示了他们异教仪式的本质：因为在那里发现了许多不同年龄之人的头骨，据称在异教徒举行这些以及类似的魔法艺术，借此蛊惑人心时，他们为了通过查看内脏来占卜而把人献祭。基督徒发现\OriginalTerm{密特拉神殿}{Mithreum}内殿中的这些可憎之物后，急切地出去将它们展示给众人看并招致咒骂；于是他们带着头骨在全城游行，如同凯旋队伍，供人们观看。当亚历山大城的异教徒看到这一切时，无法忍受这种侮辱性的行为，变得极为愤怒，他们用随手捡起的任何武器攻击基督徒，在狂怒中用以多种方式杀害了许多人：有的用剑杀死，有的用棍棒和石头打死；有的用绳索勒死，有的被钉在十字架上，故意施加这种死亡方式以示对基督十字架的轻蔑；他们刺伤了大多数人；在这种情形下，通常连朋友和亲属也不能幸免，朋友、兄弟、父母、子女都亲手沾上彼此的鲜血。因此基督徒停止了清理密特拉神殿；与此同时，异教徒将\Person{乔治}从教堂拖出，绑在骆驼上，将他撕成碎片后，连同骆驼一起烧掉。\par

% structure: p0007 source=h2 target=subsection reason=relative-heading-rank; base=h2

\subsection{第三章 皇帝对乔治被杀感到愤慨，以书信斥责亚历山大市民。}

皇帝对\Person{乔治}被刺杀极为愤怒，写信给亚历山大市民，用最强烈的言辞斥责他们的暴行。有传闻说，那些因\Person{亚大纳削}而憎恨\Person{乔治}的人，对\Person{乔治}施加了这种暴行；但我认为，无可置疑的是，那些对特定个人怀有敌意的人常常出现在民众骚乱中；然而皇帝的信显然将责任归于民众，而非基督徒中的任何人。不过，\Person{乔治}当时——且在那之前已有一段时间——极为众矢之的，这足以解释民众对他的强烈愤怒。皇帝将罪行归咎于人民，可从他的信中看出，这封信内容如下。\par

\Emph{皇帝凯撒\Person{犹利安}\OriginalTerm{马克西穆斯·奥古斯都}{Maximus Augustus}致\Place{亚历山大}市民。}\par

即使你们不尊敬你们城市的奠基人亚历山大，更不尊敬那伟大而最神圣的神塞拉皮斯，但你们怎能不仅不顾人类与社会秩序的共同要求，而且不顾对我们应尽的义务？——所有神，特别是伟大的塞拉皮斯，已将世界帝国托付给我们，因此你们本应将一切公共不法之事保留给我们裁决。然而，也许是愤怒和愤慨的冲动占据了心灵，常常刺激人做出最残忍的行为，使你们误入歧途。不过，看来当你们的狂怒稍有缓和时，你们又因在冲动时刻已犯的罪行之上再加上一项最严重的过错而加重了罪责：虽然你们不过是平民，却毫不羞愧地做出那些你们因此正当憎恶的行为。我以塞拉皮斯之名恳求你们告诉我，你们对乔治如此愤怒，是因什么不义之事？你们或许会回答，是因为他激怒了有福记忆的君士坦提乌斯反对你们；是因为他带领军队进入圣城；是因为埃及总督因此掠夺了那位神的最神圣圣殿中的神像、还愿奉献物以及其中其他祝圣的器具；当你们无法忍受看到如此可憎的亵渎，试图保护那位神免遭渎神之手，或者更确切地说，阻止抢劫那些为侍奉神而祝圣的物品时，他却违背一切正义、法律和虔诚，竟敢派遣武装队伍攻击你们。他这样做，或许是因为他畏惧乔治胜过君士坦提乌斯；但如果他没有犯下这种暴虐行为，而是坚持先前对你们的温和态度，他本会更有利于自己的安全。由于所有这些原因，你们对乔治作为众神的敌对者感到愤怒，再次玷污了你们的圣城；而你们本应将他告上法官。因为如果你们这样做，既不会发生谋杀，也不会发生任何其他不法行为；但若公正得到公平施行，将会使你们免于这些可耻的过犯，同时使他受到其不敬神罪行应得的惩罚。总之，这样也会遏制那些藐视众神、既不尊重如此伟大的城市也不尊重如此繁荣的人口、却将其对众神施行的野蛮行为作为其权力行使之序幕的人的狂妄。因此，请将我现在的这封信与我前不久写给你们的信相比较。当时我是何等高度赞扬你们！但现在，以不朽之神的名义，我虽有同样赞扬你们的意愿，却因你们严重的恶行而无法做到。民众竟然胆大包天，像狗一样将一个人撕碎；事后他们对此残忍行径毫无羞愧，也不愿洗净手上沾染的污秽，以便在众神面前伸出未沾血污的手。你们无疑会准备说乔治受到此惩罚是公正的；我们也可能倾向于承认他应受更严厉的折磨。如果你们进一步断言，为了你们的缘故，他应受这些苦难，这点或许也可同意。但如果你们接着说，应由你们来施加他罪有应得的惩罚，那我绝不能同意；因为你们有法律，你们每个人都必须服从法律，并在公开和私下都表现出对法律的尊重。如果有人违反那些最初为公共利益而制定的明智而有益的规定，这是否免除其他人遵守它们的义务？你们亚历山大居民很幸运，这样的暴行发生在我们的统治时期，因我们对众神的敬畏，以及因我们的祖父和叔父——我们继承其名号、曾统治埃及和你们城市——的缘故，我们仍对你们怀有兄弟般的感情。确实，那不容许被轻视的权力，以及拥有强健健康体魄的政府，不能纵容其臣民如此无法无天的放纵，而不通过施用足够强效的疗法毫不留情地清除这一危险病症。然而，对你们，出于已陈述的原因，我们将仅限于使用劝诫和告诫这种较为温和的治疗方法；我们相信你们会更乐于接受这种处理方式，因为据我们所知，你们原是希腊血统，并且在记忆和品格中仍保留着祖先光荣的痕迹。这封信要向我们的亚历山大市民公布。\par

这就是皇帝的信函。\par

% structure: p0012 source=h2 target=subsection reason=relative-heading-rank; base=h2

\subsection{第四章　论乔治之死，阿塔纳修斯返回亚历山大，重获其主教座}

此后不久，阿塔纳修斯从流放地返回，受到亚历山大居民的极大欢迎。他们当时将亚略派人逐出教堂，并将教堂归还阿塔纳修斯占有。同时，亚略派人在低矮隐蔽的建筑物中集会，祝圣了路求斯以填补乔治的位置。当时亚历山大城的状况便是如此。\par

% structure: p0014 source=h2 target=subsection reason=relative-heading-rank; base=h2

\subsection{第五章　论路济弗尔与尤西比乌斯}

大约同一时期，路济弗尔和尤西比乌斯奉皇帝诏令，从上特拜斯被召回，结束流放；前者是撒丁岛卡拉拉城的主教，后者是意大利利古里亚地区韦尔切利城的主教，如我先前所述。因此，这两位主教共同商议最有效的方法，以防止被忽略的教会法规和纪律今后再被违反和藐视。\par

% structure: p0016 source=h2 target=subsection reason=relative-heading-rank; base=h2

\subsection{第六章　路济弗尔前往安提约基雅，祝圣了保利诺}

因此决定，路济弗尔应前往叙利亞的安提约基雅，尤西比乌则前往亚历山大，以便与亚大纳削一同召集教务会议，确认教会的教义。路济弗尔派了一位执事作为代表，藉此承诺遵从会议的一切决议；但他本人却去了安提约基雅，发现那里的教会极其混乱，民众彼此不和。因为不仅由尤佐伊乌斯引入的亚略异端分裂了教会，而且，如前所述，梅利提乌斯的追随者也因对老师的依恋，从那些与他们意见一致的人中分离出来。因此，路济弗尔任命了保利诺作为他们的主教后，便再次离开了。\par

% structure: p0018 source=h2 target=subsection reason=relative-heading-rank; base=h2

\subsection{第七章　尤西比乌与亚大纳削合作，在亚历山大召开教务会议，宣布天主圣三为\OriginalTerm{同质}{Consubstantial}。}

尤西比乌一抵达亚历山大，便立即与亚大纳削共同召集了一次教务会议。这次从各城聚集而来的主教们，审议了许多至关重要的问题。他们确立了圣神的天主性，并将祂纳入同质的天主圣三中；他们还宣告，圣言在降生成人时，不仅取了肉身，也取了灵魂，这与早期教会人士的观点一致。因为他们并没有将任何自己臆造的新教义引入教会，而是满足于记录赞同那些从起初就为教会传统所坚持、并为明智的基督徒所明确教导的要点。这些情感是古代教父在所有论战著作中一贯持守的。依勒内、克肋孟、希拉波利的阿波利纳里以及治理安提约基雅教会的塞拉彼雍，在他们各自的著作中向我们保证，普遍接受的观点是基督在降生成人时具有灵魂。此外，为处理阿拉比亚的费拉德尔菲亚主教贝里禄而召开的教务会议，在其致该位主教的信函中也承认了同一教义。奥力振也在其现存的所有著作中各处承认，降生成人的天主取了人的灵魂。但他在\WorkTitle{创世纪注释}第九卷中更详细地解释了这一奥迹，其中表明亚当和厄娃是基督与教会的预象。那位圣人庞菲利，以及以他命名的尤西比乌，在这一主题上是可靠的见证人：这两位见证人在他们合著的\WorkTitle{奥力振传}及为回应那些对奥力振怀有偏见之人而作的精彩辩护中，证明他并非首先作出此宣告的人，而只是教会奥秘传统的阐释者。参与亚历山大会议的人士还非常细致地审视了关于\Quote{本质}或\Quote{实体}，\Quote{存在}、\Quote{自立}或\Quote{位格}的问题。因为西班牙科尔多瓦的主教霍修斯——他先前曾被提及受君士坦丁皇帝派遣平息亚略引起的骚动——在热切推翻利比亚的撒伯里乌的教条时，引发了关于这些术语的争论。然而，随后不久召开的尼西亚大公会议并未激化这一争议；但此后围绕着它产生了争论，因此在亚历山大问题被自由讨论。会上决定，诸如\OriginalTerm{本体}{ousia}和\OriginalTerm{位格}{hypostasis}等表达不应用于天主；因为他们论证说\Quote{本体}一词在圣经中任何地方都未被使用；而宗徒之所以\Quote{误用}\OriginalTerm{位格}{hypostasis}一词\BibleRef{希 1:3}，是出于教义本质所带来的不可避免的必要性。尽管如此，他们决定，在驳斥撒伯里乌谬误时，由于缺乏更恰当的用语，可以采纳这些术语，以免人们以为一个事物由三重名称所标示；而实际上我们更应相信，天主圣三中每一位都是在自己的适当位格中为天主。这些便是该会议的决议。若我们表达自己对实体与位格的看法，我们认为希腊哲学家们对\Quote{本体}给出了各种定义，却丝毫未注意\Quote{位格}。语法学家依勒内在他的字母顺序\WorkTitle{阿提卡风格}中甚至宣称它是一个蛮族术语；因为它不见于任何古人著作中，除非偶尔在与其今日含义完全不同的意义上使用。例如索福克勒斯在其悲剧\WorkTitle{菲尼克斯}中，用它表示\Quote{背叛}；在米南德那里，它意为\Quote{调味汁}；仿佛有人将酒桶底的\Quote{沉淀物}称为\Quote{位格}。尽管古代哲学作家几乎不注意这个词，近代作者却常以它代替\Quote{本体}。正如我们此前所述，此术语有不同定义；但那能被定义所限定的事物，怎么可能适用于不可理解的天主呢？厄瓦格留在他的\WorkTitle{隐修士论}中告诫我们，在谈论天主时要避免轻率不慎的言辞，禁止一切界定天主的企图，因为天主在其本性上完全是单纯的：\Quote{因为，}他说，\Quote{定义只属于复合的事物。}该作者进一步补充说：\Quote{每一命题都有所谓\InnerQuote{种}、\InnerQuote{属}、\InnerQuote{种差}、\InnerQuote{固有属性}、\InnerQuote{偶性}，或这些的组合；但这些中没有一项能被设想存在于神圣的天主圣三中。因此，不可言说的当在沉默中受钦崇。}这便是厄瓦格留的推理，我们以后还会提到他。这里我们确实做了题外话，但这有助于阐明所讨论的主题。\par

% structure: p0020 source=h2 target=subsection reason=relative-heading-rank; base=h2

\subsection{第八章　引自亚大纳削的\WorkTitle{逃亡辩护}}

此时，\Person{阿塔纳修}向在场者朗读了他此前为辩护自己逃跑而写的\WorkTitle{辩护}；从其中摘录几段也许有助于此处介绍，全文过长，不便转录，留待好学之士自行查阅研读。请看这些不敬之人胆大包天的恶行！他们的行径便是如此：然而他们非但不因先前对我们拙劣的阴谋感到羞愧，反而现在辱骂我们得以逃脱他们那谋杀之手；甚至，他们因未能将我们完全除掉而深感懊恼。简而言之，他们忽略了这一事实：当他们假装以\Quote{懦弱}责备我们时，实际上是在指控自己：因为逃跑若是可耻，追赶就更可耻，前者只是设法避免被谋杀，后者却在寻求行凶。但圣经本身也吩咐我们逃跑：\BibleRef{玛 10:23} 而那些迫害人至死的人，企图违反法律，迫使我们不得不逃亡。因此，他们更应对自己的迫害感到羞愧，而非指责我们寻求逃避；让他们停止骚扰，逃亡者也将停止。然而他们恶意毫无限度，用尽一切手段陷害我们，深知被迫害者的逃亡是对迫害者最强烈的谴责：因为没有人会逃避一个温和善良的人，而是逃避性情野蛮残酷的人。因此，\Quote{凡受窘迫的、欠债的}都从撒乌耳逃往达味。因此，我们的仇敌也同样想要杀死那些隐藏自己的人，好使没有任何证据能控告他们的邪恶。但在此事上，这些迷途之人也大错特错地自欺：因为逃避他们的努力越明显，他们有预谋的屠杀和流放就越昭然若揭。如果他们扮演刺客的角色，流出的血的声音将更大声地呼喊反对他们；如果他们判处流放，他们将在各处树立起自己不义和压迫的活的纪念碑。诚然，除非他们心智不健全，否则他们会看出自己的计谋使他们陷入两难。但由于他们失去了健全的判断，当他们消失时，他们的愚昧便暴露出来，而当他们试图停留时，却又看不见自己的邪恶。但如果他们指责那些成功隐藏自己以躲避嗜血仇敌恶意的人，并且辱骂那些从迫害者面前逃跑的人，那么对于雅各伯躲避他兄弟厄撒乌的怒气，\EditorialNote{创世纪第二十八章}以及梅瑟因害怕法老而退往米德杨地，他们又能说什么呢？\BibleRef{出 2:15} 这些多嘴的人又有什么借口为达味从撒乌耳面前逃跑辩护？\BibleRef{撒上 19:12} 当时撒乌耳从自己家中派使者去杀他；以及达味假装癫狂，从阿彼默肋客的诡计中脱身后，躲藏在山洞里？当有人提醒他们关于大先知厄里亚的事时——他藉呼求天主曾使死人复活，却因害怕阿哈布而隐藏自己，并因依则贝尔的威胁而逃跑——这些轻率地断言一切以合其目的的人又能回答什么呢？\BibleRef{列上 19:3} 那时，先知弟子们也被搜寻欲杀，便退避，由敖巴狄雅藏在山洞里；\BibleRef{列上 18:4} 还是因为年代久远，他们不知道这些事例？他们也忘记了福音所记载的，门徒们因害怕犹太人而退避躲藏？\BibleRef{玛 26:56} 保禄，当大马士革的官长搜捕他时，\BibleQuote{被人从城墙上用筐子缒下去，就逃脱了那搜寻他的手。}\BibleRef{格后 11:32} 既然圣经记载了圣人这些事迹，他们还能为自己的鲁莽编造什么借口呢？如果他们指责我们\Quote{懦弱}，那完全是麻木不仁，看不出这指责恰恰是对他们自己的谴责。如果他们诋毁这些圣人，声称他们的行为违背了天主的旨意，那就显明他们对圣经的无知。因为法律中曾命令设立\Quote{逃城}，\BibleRef{户 35:11} 以此来使那些被迫杀的人有保全自己的方法。再者，在时代的终结，圣父的圣言——先前曾藉梅瑟发言——亲自来到世上，他明确下令：\BibleQuote{几时人们在这城迫害你们，你们就逃往另一城。}\BibleRef{玛 10:23} 不久之后又说：\BibleQuote{所以，当你们看见先知达尼尔所说的\InnerQuote{招致荒凉的可憎之物}立在圣地（诵读的应当明白），那时，在犹太的，要逃往山上；在屋顶上的，不要下来取屋内的东西；在田里的，也不要回去取自己的衣服。}\BibleRef{玛 24:15} 圣人既知道这些诫命，他们行事就有了这样的训练：因为主那时所命令的，在他降生成人之前，已藉他的仆人宣示过了。这就是为人普遍而走向成全的规则：\Quote{实践天主所吩咐的一切。}为此，圣言自己为了我们而降生成人，当人搜寻他时，他屈尊隐藏了自己；\BibleRef{若 8:59} 当他又受迫害时，他俯就退避，以避开对他的阴谋。因为他正应当藉着饥饿、干渴和其他苦难，来证明他确实成了人。因为一开始，他刚出生，就藉天使吩咐若瑟：\Quote{起来，带着婴孩和他的母亲，逃往埃及，因为黑落德要寻找婴孩，杀害他。}黑落德死后，他似乎因害怕他的儿子阿赫劳斯而退往纳匝肋。随后，他藉治愈枯手给出毋庸置疑的神性证据时，\BibleQuote{法利塞人出去商订怎样陷害他，耶稣知道他们的恶意，就离开了那里。}\BibleRef{玛 12:14} 此外，当他复活了拉匝禄，他们更加决意要杀害他时，〔我们被告知〕\BibleQuote{耶稣不再公开地在犹太人中间往来，却退往旷野附近的地方。}\BibleRef{若 11:53} 再者，当救主说：\BibleQuote{在亚巴郎出现以前，我就存在。}\BibleRef{若 8:58} 犹太人便拿起石头要砸他；耶稣隐藏了自己，从他们中间穿行，走出圣殿，就从那里离去，这样逃脱了。既然他们看见这些事，或者更确切地说，理解这些事（因为他们不愿看见），那么，他们行事说话如此明显地违背主的一切所做和所教导的，按照经上的话，岂不应受火烧吗？最后，当若翰殉道，他的门徒埋葬了他的身体，耶稣听了这事，就独自乘船退往一处荒野。\BibleRef{玛 14:12} 主做了这些事，也这样教导。但愿我所说的这些人有节制，仅对人鲁莽，而不敢犯这样的疯狂，竟指责救主本人\Quote{懦弱}，尤其是在他们已经说了亵渎他的话之后。但即使他们癫狂，也不会被容忍，他们对福音的无知将被众人识破。在这种情况下退避和逃跑的缘由是合理而有效的，福音作者们为我们提供了我们救主本人行为中的先例；由此可以推断，圣人一向受到同一原则的正确影响，因为凡记载他作为人的事，都适用于全人类。因为他取了我们的性体，在自己身上显示了我们软弱的情绪，若翰这样指出：\BibleQuote{他们想要捉拿他，但没有人下手，因为他的时辰还没有到。}\BibleRef{若 7:30} 此外，在那时辰来到之前，他亲自对他的母亲说：\Quote{我的时辰还没有到。}并对那些被称为他弟兄的人说：\Quote{我的时候还没有到。}及至时候到了，他对门徒说：\BibleQuote{现在你们继续睡觉吧，安歇吧！看，时辰到了，人子要被交在罪人手里。}\BibleRef{玛 26:45} ……因此，他既不容许时辰未到就被捕，也不在时辰到时隐藏自己，而是甘愿把自己交在那些图谋害他的人手中。……有福的殉道者们也是如此在迫害时期保全自己：他们被迫害时就逃跑，隐藏自己；但一旦被发现，就殉道致命。\par

此乃亚大纳削在其为自己逃亡所作的辩护中的推理。\par

% structure: p0023 source=h2 target=subsection reason=relative-heading-rank; base=h2

\subsection{第九章　亚历山大会议后，欧瑟伯前往安提约基雅，发现天主教徒因保理诺的祝圣而意见分歧；他竭力调解无效，遂离去；路济弗尔之愤慨及以其命名的派别之起源。}

亚历山大会议一解散，韦尔切利主教欧瑟伯便从亚历山大前往安提约基雅；他在那里发现保理诺已由路济弗尔祝圣，而民众之间意见不合——因为默肋提乌斯的追随者分开举行集会——他对这次选举中的不和谐感到非常痛心，并在内心不赞成所发生的事。然而，他对路济弗尔的尊重使他对此保持缄默，离开时他承诺一切将通过一次主教会议得到纠正。随后，他极其热心地努力联合持异议者，但未成功。与此同时，默肋提乌斯从流放中归来；发现他的追随者离开他人单独举行集会，他便自任首领。但亚流异端的首领欧佐依乌斯占据着各教堂：保理诺只保留城内一座小教堂，欧佐依乌斯出于对其个人的尊重并未将他驱逐。而默肋提乌斯在城门外召集他的拥护者。欧瑟伯就在这种情况下当时离开了安提约基雅。当路济弗尔得知他对保禄的祝圣未得到欧瑟伯赞同时，视此为侮辱，勃然大怒；不仅与欧瑟伯断绝共融，而且开始以好斗的精神谴责会议的决定。这些事发生在严重混乱的时期，使许多人远离教会；因为许多人依附路济弗尔，从而形成了一个以\Quote{路济弗尔派}为名的独立派别。然而，路济弗尔未能充分表达他的愤怒，因为他已通过其执事承诺同意会议所决定的一切。因此他坚持教会的教义，回到撒丁尼亚自己的教区；但最初认同他争执的那些人，至今仍与教会分离。另一方面，欧瑟伯如同一位好医生，巡游东方各省份，完全恢复了那些在信仰上软弱的人，教导并确立他们的教会原则。此后，他前往伊里利库姆，再由此去意大利，在那里他进行了类似的工作。\par

% structure: p0025 source=h2 target=subsection reason=relative-heading-rank; base=h2

\subsection{第十章　关于普瓦捷主教依拉略。}

然而，普瓦捷（属第二阿基坦省的一座城）主教依拉略已先采取行动，他先前已确认了意大利和高卢的主教们持守正统信仰的教义；因为他首先从流放地返回这些地区。因此两人都高尚地联合他们的力量维护信仰：依拉略是一位极有口才的人，在他用拉丁文写的著作中强有力地维护了\OriginalTerm{同质}{homoousion}的教义。在这些著作中，他给[该教义]提供了充分的支持，并无可辩驳地驳斥了亚流的主张。这些事发生在被放逐者被召回后不久。但须注意，同时马其顿尼、厄勒乌息、欧斯塔提和索弗罗尼，连同他们所有的支持者，只有一个共同称呼\Quote{马其顿派}，在各地频繁召开会议。他们召集了色娄基雅中赞同他们观点的人，咒骂另一派的主教，即阿卡奇派；并摒弃了阿里米农信经，确认了在色娄基雅读过的信经。如前书所述，这信经与先前在安提约基雅颁布的相同。当有人问他们：\Quote{你们这些迄今被称为马其顿派的人，既然持有不同见解，为何仍与阿卡奇派保持共融，好像你们意见一致？}他们便通过帕夫拉戈尼亚的庞培城主教索弗罗尼如此回答：\Quote{在西方的人，}他说，\Quote{如同被疾病感染，沾染了同质派谬误；厄提乌则在东方引入声明\InnerQuote{实体不相像}的学说，玷污了信仰的纯洁。如今这两种教义都是非法的；因为前者鲁莽地将父子不同的位格混为一谈，用那邪恶的绳索——\OriginalTerm{同质}{homoousion}一词——将它们绑在一起；而厄提乌则用\OriginalTerm{不同质}{anomoion}一词，即关于实体或本质不相像，完全割裂了圣子与圣父之间的本质关联。既然这两种见解都陷入了完全相反的极端，我们看来，二者之间的中间道路更符合真理和虔诚：因此我们主张圣子\InnerQuote{在实体上与圣父相像}。}\par

这便是马其顿派经由索弗罗尼对那个问题所作的答复，正如撒彼努斯在其\WorkTitle{会议文件汇编}中所确认的那样。但他们将厄提乌而非阿卡奇贬斥为\OriginalTerm{不同质}{anomoion}学说的始作俑者，便公然掩盖了事实，以便显得一方面与亚流派同样疏远，另一方面与同质派同样疏远：因为他们自己的话就证明他们与两者分离，仅仅出于喜爱革新。我们就此结束对这些人的记述。\par

% structure: p0028 source=h2 target=subsection reason=relative-heading-rank; base=h2

\subsection{第十一章　尤利安皇帝从基督徒手中勒索钱财。}

虽然皇帝尤利安在其统治初期对所有人都表现得温和，但随着时间推移，他未能继续维持同样的平和。他确实非常乐意满足基督徒的请求，只要这些请求有助于对君士坦提乌斯的记忆施加谴责；但若缺乏这种动机，他便毫不掩饰对基督徒普遍怀有的敌意。因此，他很快下令重建尤佐伊乌斯彻底拆毁的基齐库斯诺瓦替安派教堂，并对该城主教厄勒乌西乌斯处以重罚，若他未能在两个月内自费完成该建筑。此外，他以全部权威支持外教迷信；外教神庙被开放，正如我们之前所述；但他本人也公开在设有福尔图娜像的主教座堂中向君士坦丁堡的福尔图娜女神献祭。\par

% structure: p0030 source=h2 target=subsection reason=relative-heading-rank; base=h2

\subsection{第十二章．论迦克墩主教马里斯；尤利安禁止基督徒从事文教。}

大约此时，比提尼亚迦克墩主教马里斯被人搀扶着来到皇帝面前——因他年事极高，患了眼疾，名为\Quote{白内障}——他严厉斥责了皇帝的亵渎、背教和无神论。尤利安以连串辱骂之词回应他的责备；并用言词称他为\Quote{瞎子}反击。\Quote{你这瞎眼的老傻瓜，}他说，\Quote{你们这位加里肋亚人的天主永远不会医治你。}因为他惯常称基督为\Quote{加里肋亚人}，称基督徒为加里肋亚人。马里斯更为大胆地回答说：\Quote{我感谢天主使我失明，因为这样我就不必目睹一个陷入如此可怕亵渎之人的面孔。}皇帝当时未再追究此事，但后来进行了报复。他观察到戴克里先统治下殉道者深受基督徒尊敬，并知道许多人热切渴望成为殉道者，于是决定以其他方式发泄愤怒。因此，他虽未采用戴克里先时期的极端酷刑，但也并非完全停止迫害（因为任何骚扰和困扰的措施，我都视为迫害）。他制定的计划如下：他颁布法律，禁止基督徒从事文学修养；\Quote{免得，}他说，\Quote{当他们磨利了舌头，就能更容易地反驳外教人的论证。}\par

% structure: p0032 source=h2 target=subsection reason=relative-heading-rank; base=h2

\subsection{第十三章．论外教人对基督徒所犯的暴行。}

此外，他禁止那些不放弃基督信仰并献祭给偶像者在宫廷任职；也不允许基督徒担任行省总督；\Quote{因为，}他说，\Quote{他们的法律禁止他们对犯有死罪的罪犯动用刀剑。}他还以谄媚和礼物诱使许多人献祭。立刻，仿佛在炉中受炼，所有人都看清了谁是真正的基督徒，谁只是挂名的。心地纯全的基督徒非常乐意辞去职务，宁愿忍受一切也不否认基督。其中就有约维安、瓦伦提尼安和瓦伦斯，他们后来都成为皇帝。但其他心怀不轨之人，偏爱今世的财富和荣誉胜过真福，毫不犹豫地献祭。其中之一是君士坦丁堡的诡辩家厄刻玻利乌斯，他迎合皇帝的性情，在君士坦提乌斯统治时期假装是热心的基督徒；而在尤利安时代，他又表现得同样狂热的外教人；尤利安死后，他再次声称信奉基督教。他俯伏在教堂门前，喊道：\Quote{践踏我吧，因为我成了失了味的盐。}此人在其整个历史时期都如此反复无常。大约此时，皇帝欲对波斯人进行报复，因他们在君士坦提乌斯统治时期屡次侵犯罗马领土，便以极快的速度穿越亚细亚进军东方。但他深知战争带来的一系列灾祸，以及成功进行战争所需的巨额资源，且没有这些就无法进行，于是他诡计多端地设计了一种从基督徒那里勒索钱财的计划。他对所有拒绝献祭的人课以重税，并严厉地向真正的基督徒征收，每人按其财产比例缴纳。通过这些不义的手段，皇帝很快积累了巨额财富；因为这项法律在尤利安亲临之处及他不在之处均被执行。同时，外教人攻击基督徒；大批自称\Quote{哲学家}的人聚集。他们继而举行某些可憎的秘仪；宰杀纯洁的男女儿童，检查其内脏，甚至品尝其肉。这些丑恶的仪式在其他城市也有进行，但尤其在雅典和亚历山大里亚；在后者，有人对主教亚大纳削提出诽谤指控，皇帝被告知他意图摧毁该城乃至整个埃及，只有将他驱逐出境才能拯救该地。因此，亚历山大里亚总督奉皇帝谕令逮捕他。\par

% structure: p0034 source=h2 target=subsection reason=relative-heading-rank; base=h2

\subsection{第十四章．亚大纳削的逃亡。}

但\Person{亚大纳削}再次逃走，对他亲近的人说：\Quote{朋友们，我们退避一会儿吧；这不过是一片小云，很快就会过去。}他随即上船，渡过\Place{尼罗河}，全速赶往\Place{埃及}，追捕他的人紧追不舍。当他得知追兵已近，随行的人催促他再次退入旷野，但他用计逃脱。他说服同行的人回转迎敌，他们立刻照办；接近追兵时，对方只问\Quote{他们在哪里见过\Person{亚大纳削}}，他们回答说\Quote{他离得不远}，并且\Quote{如果他们赶快，很快就能追上他}。追兵受骗，加速追赶，却徒劳无功；\Person{亚大纳削}安全退去，秘密返回\Place{亚历山大里亚}，在那里隐藏，直到迫害结束。这就是\Place{亚历山大里亚}主教所接连遭遇的危难，这些来自外教人的危难，发生在他先前所受来自基督徒的危难之后。此外，各省总督利用皇帝的迷信以满足自己的贪欲，对基督徒施加比皇帝授权更严重的暴行：有时勒索超出应得的款项，有时对他们施以肉刑。皇帝得知这些暴行，却纵容不管；当受苦者向他上诉控告压迫者时，他嘲讽地说：\Quote{你们理应耐心忍受这些苦难，因为这是你们天主的命令。}\par

% structure: p0036 source=h2 target=subsection reason=relative-heading-rank; base=h2

\subsection{第十五章　\Person{尤利安}治下\Place{夫黎基雅}省\Place{梅隆}城的殉道者。}

\Place{夫黎基雅}省总督\Person{阿玛基乌斯}下令打开该省\Place{梅隆}城的庙宇，清除年久积存的污秽，并重新打磨其中供奉的雕像。此事付诸实施，令基督徒极为痛心。有一位名叫\Person{玛色多尼乌斯}、\Person{德奥杜禄}和\Person{塔提安}的人，不堪忍受他们的宗教所受的侮辱，激于对德行的热忱，夜间冲进庙宇，打碎了神像。总督大怒，本欲处死城中许多完全无辜的人，但肇事者自愿投案，宁可为真理而死，也不愿看着别人替他们受死。总督逮捕了他们，命令他们献祭以赎所犯之罪；他们拒绝后，法官以酷刑威胁，但他们蔑视其恐吓，勇气非凡，声明愿受任何苦难，也不愿因献祭而玷污自己。总督对他们施以各种酷刑，最后将他们放在火上的铁格架上，这样杀害了他们。但即使在最后时刻，他们也表现出了最英勇的坚毅，对无情的总督说：\Quote{阿玛基乌斯，你若想吃烤肉，就把我们翻过来，免得我们烤得半生不熟，不合你的口味。}这些殉道者就这样结束了生命。\par

% structure: p0038 source=h2 target=subsection reason=relative-heading-rank; base=h2

\subsection{第十六章　论两位\Person{阿波利纳里斯}父子的文学劳作，以及皇帝禁止基督徒受教于希腊文学。}

禁止基督徒学习希腊文学的帝国法令，使我们上面提到的两位\Person{阿波利拿里}比以前更加著名。因为两人都精通文雅之学，父亲是文法学家，儿子是修辞学家，他们在此危急关头为基督徒做出了贡献。前者作为文法学家，撰写了一部符合基督信仰的文法书；他还将\WorkTitle{梅瑟五书}译成英雄诗体，并用韵文意译了旧约所有历史书，部分采用长短格，部分改编成戏剧悲剧。他特意使用了各种诗体，以使希腊语言特有的表达形式在基督徒中不被陌生或未闻。年轻的\Person{阿波利拿里}在修辞学上训练有素，以对话方式阐释福音和宗徒的教导，正如希腊人中的\Person{柏拉图}所做的那样。因此，他们对基督教事业显示出助益，通过自己的辛劳战胜了皇帝的诡计。但天主的眷顾比他们的辛劳或皇帝的计谋更为有力：因为不久之后，以我们将要解释的方式，那法令完全失效；这些人的著作如今已无关紧要，仿佛从未被写过一般。但或许有人会激烈反驳说：\Quote{你凭什么肯定这两件事都是天主眷顾的结果？皇帝的突然死亡对基督教非常有利，这确实是显而易见的：但两位阿波利拿里的基督教作品被弃置，基督徒重新开始以异教徒的哲学熏陶心智，这对基督教毫无益处，因为异教哲学教导多神论，且有害于真宗教的推广。}对于这一异议，我将根据目前想到的理由予以回应。希腊文学确实从未被基督或其宗徒承认为神圣默感之作，但另一方面也未被完全拒斥为有害。我认为他们这样做并非轻率。因为希腊人中许多哲学家离认识天主并不远；事实上，他们受逻辑学训练，坚决反对否认天主眷顾的伊壁鸠鲁学派和其他好斗的诡辩派，驳斥他们的无知。因此，他们对所有热爱真虔诚的人都有用：然而他们自己却不熟悉真宗教的元首，因为对基督的奥迹一无所知，这奥迹\BibleQuote{从世世代代以来是隐藏的}\BibleRef{哥 1:26}。事实如此，宗徒在致罗马人书中这样宣告：\BibleRef{罗 1:18}\BibleQuote{天主的忿怒，从天上发显在一切不虔敬和不义的人身上，就是那些以不义压制真理的人。因为认识天主为他们是很明显的事，原来天主已将自己显示给他们了。其实，自从天主创世以来，他那看不见的美善，即他永远的大能和他为神的本性，都可凭他所造的万物，辨认洞察出来，以致人无可推诿。他们虽然认识了天主，却没有当作天主光荣他。}从这些话看来，他们拥有真理的知识，这是天主启示给他们的；但他们的罪在于，当认识了天主时，却没有当作天主光荣他。因此，没有禁止学习希腊人的学问，是任由愿意的人自行决定。这是我们为所持立场辩护的第一个论据：另一个论据可以这样提出：默感的圣经无疑教导了本身可敬且具有天上性质的道理；它们也卓越地促成那些受其规诫引导的人产生虔诚和正直的生活，指出一条为天主高度赞许的信德之路。但它们并不教导我们推理的艺术，使我们能够成功地抵抗那些反对真理的人。此外，当我们能用对手自己的武器攻击他们时，最容易挫败他们。但这种能力并非由两位阿波利拿里的作品提供给了基督徒。\Person{尤利安}心中明白这一点，因此他通过法律禁止基督徒接受希腊文学教育，因为他很清楚其中包含的寓言会暴露整个异教体系——他已成为该体系的捍卫者——使之受嘲笑和轻蔑。就连他们最著名的哲学家\Person{苏格拉底}也鄙视这些荒谬之事，并因此被定罪，仿佛他试图侵犯神明的尊严。此外，基督和其宗徒吩咐我们要\Quote{成为识别的兑换银钱者}，以便\BibleQuote{考验一切，保持好的}\BibleRef{得前 5:21}，还指示我们要\BibleQuote{谨慎，免得有人用哲学和虚妄的诡计把我们掳去}\BibleRef{哥 2:8}。但若非我们掌握了对手的武器，我们无法做到这一点：要小心在获取这武器时，不要采纳他们的观点，而是考验它们，弃绝邪恶，保留一切善与真：因为善无论在哪里，都是真理的财产。若有人认为，我们这样说是曲解圣经的正当含义，请记住：宗徒不仅没有禁止我们受希腊学问的教育，而且他自己似乎也绝非忽视它，因为他知道许多希腊人的说法。他从哪里得到这句：\BibleQuote{克里特人常说谎，是恶兽，是贪吃的懒汉}\BibleRef{铎 1:12}？难道不是从阅读克里特啓蒙者\Person{埃庇米尼得斯}的\WorkTitle{神谕}得来的吗？或者他怎么会知道\BibleQuote{我们也是他的子孙}\BibleRef{宗 17:28}，若非他熟悉天文学家\Person{阿拉托斯}的\WorkTitle{现象}？此外，这句\BibleQuote{坏交往败坏好品行}\BibleRef{格前 15:33}足以证明他熟悉\Person{欧里庇得斯}的悲剧。但何须详述这一点？众所周知，在古代，教会的圣师们在无阻碍的惯例中，习惯操练希腊学问，直到年事已高；他们这样做是为了提高自己的口才，加强和磨砺心智，同时能够驳斥异教徒的错误。关于两位阿波利拿里所引发的这一主题，以上论述已足够。\par

% structure: p0040 source=h2 target=subsection reason=relative-heading-rank; base=h2

\subsection{第十七章。皇帝准备远征波斯人，抵达\Place{安提约基雅}，因居民嘲笑他，他以讽刺著作题为\WorkTitle{Misopogon}（或称《恨须者》）回击。}

皇帝从基督徒身上榨取了巨额金钱后，加速远征波斯人，抵达叙利亚的\Place{安提约基雅}。在那里，他想要向市民炫耀自己追求荣耀，便不当地压低了商品价格；既不考虑当时的形势，也不思考军队驻留会给各省居民带来不便，并必然减少对城市的粮食供应。因此商人和零售商都停止交易，无法承受帝国诏令带来的损失；于是生活必需品短缺。安提约基雅人无法忍受这种侮辱——因为他们天生无法忍受侮辱——立刻对尤里安发出指责；还讽刺他的胡子，他的胡子非常长，说应该剪掉做成绳索。他们还说，铸在他钱币上的公牛象征着他蹂躏了世界。因为皇帝极其迷信，不停地在他的偶像祭台上献祭公牛，并下令在他的钱币上铸造公牛和祭台的图案。被这些嘲弄激怒后，他威胁要惩罚安提约基雅城，然后返回\Place{西里西亚}的\Place{塔尔索}，命令尽快准备离开。于是诡辩家\Person{利巴尼奥斯}趁机写了两篇演说，一篇向皇帝为安提约基雅人求情，另一篇向安提约基雅居民讲述皇帝的愤怒。然而据称这些作品只是写成，从未公开诵读。尤里安放弃了最初以伤害行为报复讽刺者的意图，转而以互相辱骂来发泄怒气；他写了一个小册子反对他们，题为\WorkTitle{Antiochicus, or Misopogon}，从而给该城及其居民留下了不可磨灭的污点。但现在我们必须讲述他给安提约基雅的基督徒带来的祸害。\par

% structure: p0042 source=h2 target=subsection reason=relative-heading-rank; base=h2

\subsection{第十八章。皇帝咨询神谕，魔鬼因敬畏殉道者\Person{巴比亚拉}的临近而不回应。}

他命令打开安提约基雅的外教神庙后，非常渴望从达佛涅的\Person{阿波罗}那里获得神谕。但占据那庙的魔鬼因畏惧他的邻居殉道者\Person{巴比亚拉}而保持沉默；因为那位圣人的棺材就在附近。当皇帝得知此事，他命令立即移走棺材；于是安提约基雅的基督徒，包括妇女和儿童，以庄严的欢庆和圣咏歌唱，将棺材从\Place{达佛涅}运到城中。所唱的圣咏都是斥责外教神祇以及那些信赖他们和他们偶像的人。\par

% structure: p0044 source=h2 target=subsection reason=relative-heading-rank; base=h2

\subsection{第十九章。皇帝的愤怒和宣信者\Person{塞奥多洛}的坚定。}

那时皇帝的真实脾气和性情——他此前一直尽量不让人察觉——完全显露出来：因为他如此夸耀自己的哲学，却再也无法自制；反而被这些羞辱性的赞美诗激得几乎疯狂，准备对基督徒施加\Person{戴克里先}的爪牙从前曾施加的同样残酷行为。然而，由于他对波斯远征的担忧使他无暇亲自执行其意愿，他命令禁卫军长官\Person{萨路斯特}逮捕那些在唱圣咏方面最为热心的人，以儆效尤。这位长官虽然是外教人，却一点也不喜欢这个差事；但由于不敢违背，他下令逮捕了一些基督徒，并将其中一些人关入监牢。有一个名叫\Person{塞奥多洛}的年轻人，外教人将他带到长官面前，长官对他施以各种酷刑，使他的身体被撕裂，直到他认为此人不可能再活下去时才停止进一步的惩罚；然而天主保全了这位受苦者，使他长期存活于那次宣信之后。拉丁文\WorkTitle{教会史}的作者\Person{鲁菲努斯}说，他在相当长的时间后亲自与这位塞奥多洛交谈过；并问他，在鞭打和拷问的过程中，他是否感到最剧烈的疼痛；他回答说，他感受到的酷刑之痛非常短暂；有一个年轻人站在他旁边，既拭去因所经受的剧烈考验而产生的汗水，同时又坚固他的心灵，使他将这段考验时刻视为狂喜而非受苦的时期。关于最奇妙的塞奥多洛，就说这么多。大约此时，波斯使者来到皇帝那里，请求他在某些特定条件下结束战争。但尤里安突然打发他们离去，说：\Quote{你们不久就会亲眼见到我，因此不必派使团了。}\par

% structure: p0046 source=h2 target=subsection reason=relative-heading-rank; base=h2

\subsection{第二十章。犹太人受皇帝煽动试图重建圣殿，但因奇迹干涉而受挫。}

皇帝又试图骚扰基督徒，暴露了他的迷信。他喜爱献祭，不仅自己以牺牲之血为乐，而且认为他人若不如此便是对他的侮辱。他发觉这样的人寥寥无几，便召来犹太人，询问他们为何不献祭，既然梅瑟法律命令他们这样做。他们回答说，除了在耶路撒冷，他们不可在任何其他地方行此事。他立即下令他们重建撒罗满的圣殿。与此同时，他亲自出发远征波斯。犹太人久已渴望获得重建圣殿的良机，以便在其中献祭，因此十分积极地投入工作。此外，他们对基督徒极其无礼，威胁要加给他们与自己所受的罗马人迫害同等的伤害。皇帝下令从国库拨付这项工程的费用，所有材料很快备齐，如木材、石料、烧砖、粘土、石灰以及一切其他必需的建材。此时，耶路撒冷的主教济利禄想起达尼尔的预言——基督在神圣福音中也曾确认——并在众人面前预言，那时刻确实已到，\Quote{在这圣殿中，将没有一块石头留在另一块上}，但救主的预言宣告将完全应验。主教的话如此。次日夜间，一场大地震掀起圣殿旧地基的石头，将它们连同邻近的建筑一并震散。犹太人因此对这事心生恐惧；消息传到远方，许多人聚集到现场。当大批民众聚集时，又发生了一件异事。有火从天降下，烧毁了所有建筑工具：火焰吞噬了木槌、磨石和抛光石头的铁器、锯子、斧头、扁斧，简言之，工人们为工程所需而准备的各种工具；火在这些东西上燃烧了一整天。犹太人确实惊恐万分，不情愿地承认基督，称他为天主；但他们并未遵行他的旨意，反而因根深蒂固的偏见仍固守犹太教。即使后来发生的第三个奇迹也未能引导他们相信真理。因为第二天夜里，他们的衣服上出现了发光的十字架印记，天亮后他们想擦洗或洗掉，却徒劳无功。因此，\Quote{他们眼目昏花}，正如宗徒所说的，并且抛弃了手中的好处：这样，圣殿非但没有重建，反而在那时被彻底摧毁。\par

% structure: p0048 source=h2 target=subsection reason=relative-heading-rank; base=h2

\subsection{第二十一章 皇帝入侵波斯及驾崩。}

其时，皇帝在初春之前便入侵了波斯人的国土，因为他得知波斯各民族在冬季极为衰弱、全无斗志。由于他们不耐寒冷，故在那季节免除兵役，有一句谚语说：\Quote{米底人此时绝不会将手从斗篷下伸出。}他深知罗马人惯于抵御一切气候的严酷，便将他们释放到该国。在蹂躏了一大片地区，包括无数村庄和堡垒之后，他们随即围攻城市；他包围了大城\Place{泰西封}，将波斯国王逼入绝境，后者多次派遣使节到皇帝处，提议割让部分领土，条件是他撤离该国并结束战争。但\Person{尤利安}不为这些屈服所动，毫不怜悯求饶的敌人；他也不曾记得这句格言：\Quote{赢得胜利是光荣的，但超过胜利者则招致嫉妒。}他信任与他不断来往的哲学家\Person{马克西穆斯}的占卜，受其欺骗，相信自己的事迹不仅将等同，而且会超越\Person{马其顿的亚历山大}；因此他傲慢地拒绝了波斯国王的恳求。他甚至按照\Person{毕达哥拉斯}和\Person{柏拉图}关于灵魂转世的教导，认为自己拥有\Person{亚历山大}的灵魂，或者更确切地说，他本人就是\Person{亚历山大}在另一个身体中的化身。这一荒谬的幻想迷惑了他，导致他拒绝了波斯国王提出的和平谈判。因此，后者确信谈判无用，被迫准备交战；于是，在拒绝其使节的次日，他将他所有的军队列成战阵。罗马人确实批评他们的君主，本可有利地避免交战却不这样做；尽管如此，他们还是攻击了那些抵抗他们的人，再次将敌人击溃。皇帝亲临战马，在战斗中鼓励士兵；但他仅凭胜利的希望，未穿铠甲。在这种无防护的状态下，一支不知谁投出的标枪刺穿了他的手臂，进入他的肋部，造成了伤口。他因这伤而死。有人说是一个波斯人投掷了标枪，然后逃走了；另一些人则断言是他自己手下的人所为，这其实是最可靠、最流行的说法。但他的侍卫之一\Person{卡利斯托斯}，曾用英雄诗体歌颂这位皇帝的功绩，在叙述这场战争的细节时说，他因此致命的那一伤是由一个恶魔造成的。这或许仅仅是诗人的虚构，或者确有其事；因为复仇的愤怒无疑毁灭了许多人。但无论情况如何，可以肯定的是，他天性中的热忱使他粗心大意，他的学识使他虚荣，他的宽仁矫饰使人轻视。于是，正如我们所说，\Person{尤利安}在波斯结束了生命，时值他第四次执政（与同僚\Person{萨卢斯特}共同执政）。这事发生在六月二十六日，他统治的第三年，也是\Person{君士坦提乌斯}立他为凯撒之后的第七年，当时他三十一岁。\par

% structure: p0050 source=h2 target=subsection reason=relative-heading-rank; base=h2

\subsection{第二十二章　\Person{约维安}被拥立为皇帝。}

这突发事件使士兵们陷入极大的困惑，他们次日便毫不犹豫地拥立\Person{约维安}为皇帝，此人以其英勇和出身而同样出众。当时\Person{尤利安}发布敕令，给军官们一个选择：要么祭献，要么辞去军职；\Person{约维安}是一名军事保民官，他宁愿辞去职务，也不愿服从一位不虔敬君主的命令。然而，\Person{尤利安}迫于眼前的战争紧迫，仍将他留作将领之一。当士兵向他呼喊致敬时，他坚决拒绝接受君主权力；士兵们强行把他推出来时，他声明：\Quote{我是一个基督徒，我不愿统治一群选择以异教为宗教的人民。}他们全体齐声回答说他们也是基督徒；于是他接受了皇帝尊严。他发觉自己突然处于极其困难的境况，在波斯领土中间，军队因缺乏必需品而面临灭亡的危险，于是同意结束战争，即使条件对罗马的荣誉并不光彩，但危急形势使之成为必要。因此，他忍受了丧失叙利亚统治权的损失，并放弃了美索不达米亚的\Place{尼西比斯}城，撤离了波斯领土。这些消息的宣布给基督徒带来了新的希望；而异教徒则痛惜\Person{尤利安}之死。尽管如此，全军谴责他的暴躁脾气，并将丧失所失领土的耻辱归咎于他轻信一个波斯逃兵的狡猾报告；因为他被这个逃亡者的陈述所欺骗，烧毁了为他们提供水运补给的船只，从而使他们遭受饥饿的一切恐怖。随后，\Person{利巴尼奥斯}为他撰写了一篇葬礼演说，题为\WorkTitle{尤利安，或墓志铭}，其中他以极高的赞词歌颂了他几乎所有的行为；但在提到\Person{尤利安}所写的反对基督徒的书籍时，他说他已在其中清楚证明了他们神圣经卷的可笑和琐碎特性。倘若这位诡辩家仅限于赞扬皇帝的其他行为，我会平静地继续我的历史叙述；但既然这位著名的修辞学家认为有必要借机攻击基督信仰的圣经，我们也打算稍作停顿，在简短的回顾中考量他的话。\par

% structure: p0052 source=h2 target=subsection reason=relative-heading-rank; base=h2

\subsection{第二十三章　驳斥诡辩家\Person{利巴尼奥斯}有关\Person{尤利安}的言论。}

他说：\Quote{冬天加长了夜晚的时候，皇帝对那些使\Place{巴勒斯坦}人成为天主和天主圣子的书籍进行了\InnerQuote{攻击}；经过一连串的论证，他证明了这些被基督徒如此尊崇的著作是荒谬且毫无根据的，从而表明他自己比\Place{提尔}老人更聪明、更熟练。但愿这位\Place{提尔}的智者对我仁慈，温和地容忍已被断言的事情，因为他已被儿子超越了！}这就是诡辩家\Person{利巴尼乌斯}的话。但我承认，他确实是一位杰出的修辞学家，但我深信，如果他与皇帝在宗教情感上不一致，他不仅会表达基督徒对他所说的一切，而且会像一个修辞学家自然会做的那样，夸大每一个谴责的理由。因为当\Person{君士坦提乌斯}活着的时候，他写诗赞美他；但在他死后，他却对他提出了最具侮辱性和辱骂性的指控。所以，如果\Person{波斐利}是皇帝，\Person{利巴尼乌斯}肯定会更喜欢他的书而不是\Person{尤利安}的书；如果\Person{尤利安}只是一个诡辩家，他会称他为非常平庸的诡辩家，就像他在\WorkTitle{尤利安挽歌}中对待\Person{埃凯博利乌斯}那样。既然他是以异教徒、诡辩家和所称赞之人的朋友的精神说话，我们将尽力回应他所提出的观点，尽我们所能。首先，他说皇帝在漫长的冬夜里\Quote{攻击}这些书籍。现在，\Quote{攻击}意味着将反驳这些书的写作作为一项任务，就像诡辩家们在教授他们艺术的基本原理时通常所做的那样；因为他很久以前就已经阅读过这些书，但此时才攻击它们。但在整个长期的论战中，他并没有像\Person{利巴尼乌斯}所声称的那样试图用合理的论证来反驳任何事情，而是在缺乏真理的情况下求助于他非常喜欢的嘲笑和轻蔑的玩笑；因此，他试图嘲讽那些坚不可摧、无法推翻的东西。因为每一个与他人争论的人，有时试图歪曲真理，有时又试图隐藏真理，用尽一切可能的手段篡改对立者的立场。敌人不仅对与他有分歧的人做恶意行为，而且还会对他说坏话，并将他自己意识到的缺点归咎于他所不喜欢的人。\Person{利巴尼乌斯}称之为\Quote{\Place{提尔}老人}的\Person{尤利安}和\Person{波斐利}都非常喜欢嘲笑，这是显而易见的，从他们自己的著作中可以看出。因为\Person{波斐利}在他的\WorkTitle{哲学家历史}中嘲笑了所有哲学家中最杰出的\Person{苏格拉底}的生活，对他发表了连他的控告者\Person{梅勒图斯}或\Person{安尼图斯}都不敢说的话；我说的是\Person{苏格拉底}，他被所有希腊人赞赏他的谦虚、正义和其他美德；\Person{柏拉图}（他们中最令人钦佩的）、\Person{色诺芬}和其他哲学群体不仅尊他为天主所爱的人，而且习惯于认为他被赋予了超人的智慧。而\Person{尤利安}，模仿他的\Quote{父亲}，在他的书\WorkTitle{凯撒们}中表现出类似的精神病态，在其中他诽谤他所有的皇帝前任，甚至不放过\Person{马可}哲学家。因此，他们自己的著作表明他们都喜欢嘲讽和辱骂；我无需用大量巧妙的表达来证明这一点；但关于他们在这方面的心境，刚才所说的已经足够。现在写这些，我以每人的演说作为他们性情的见证，但\Person{尤利安}尤其如此，\Person{纳齐盎的额我略}在他的\WorkTitle{第二篇反异教徒演说}中所说的是以下内容：\par

\Quote{这些事情经经验向其他人显明了，在他获得了帝国权力并可以随心所欲之后：但自从我在雅典认识他时，我就预见到了这一切。他是在他兄弟的命运改变后不久，得到皇帝许可来到那里的。他这次访问有两个动机：一个理由对他来说是光荣的，即参观希腊并进入那里的学校；另一个则更隐秘，几乎无人知晓，因为他的不敬尚未敢于公开表露，即有机会就自己的命运咨询献祭者和其他骗子。我清楚地记得，即使那时我对这个人也不是一个糟糕的占卜者，尽管我绝不声称自己是精通占卜术的人之一：但他性格的反复无常和心智的难以置信的放纵使我具有了预言能力；如果真有\InnerQuote{猜中事件的最好的先知}的话。因为在我看来，一个很少稳定的脖子、频繁耸肩、恶狠狠且不停转动的眼睛，伴随疯狂的表情；不平稳且摇摆的步态，只散发轻蔑和侮辱的鼻子，带着同样意味的荒谬扭曲面容；过分且非常大声的笑声，点头如同同意，以及摇头如同否定，而没有任何明显原因；说话犹豫且被呼吸打断；混乱且无意义的问题，回答也好不到哪里去，一切都混杂在一起，毫无连贯或条理。我何必细说？我预见他会如此，正如我后来从经验中发现的那样。如果当时在场并听到我说话的人中有人现在在这里，他们会乐意作证，当我观察到这些预兆时，我喊道：\InnerQuote{啊！罗马帝国给自己养成了多大的祸害啊！}当我说完这些话时，我祈求天主让我成为一个假先知。因为远为更好的是[我应被证明做出了错误的判断]，而不是世界充满如此多的灾难，并且出现了前所未见的这样一个怪物：尽管记载了许多洪水和大火，许多地震和裂缝，并描述了众多凶残和不人道的人，以及由不同种族混合而成的野兽怪物，大自然产生了不寻常的形态。他的结局确实与他疯狂的生涯相称。}\par

这便是额我略留给我们的尤利安画像。此外，许多人已证明，他们在各自编著中竭力对真理施暴，有时篡改圣经经文，有时增添原话并按自己目的曲解其意，从而驳斥其诡辩并揭露其谬误。尤其是奥利振，他远在尤利安时代之前，便亲自对圣经中看似困扰某些读者的章节提出异议并充分解答，从而堵塞了轻率者的恶意叫嚣。倘若尤利安和波斐利曾认真诚实地研读他的著作，他们定会谈论其他话题，而不是转向编造亵渎的诡辩。显而易见，这位皇帝在演说中意在迷惑无知者，并未向那些拥有圣经所呈真理之\Quote{形状}的人讲话。因他将圣经中按经济安排、更依人方式谈论天主的种种表达汇集一处，如此评论：\InnerQuote{这些说法无一不充满对天主的亵渎，除非其中含有某种隐秘神秘的含义——这倒是我所能设想的。}这是他在第三卷\WorkTitle{驳基督徒}中的原话。然而在其著作\WorkTitle{论犬儒哲学}中，他说明宗教题材可虚构寓言到何种程度时，声称此类真理必须隐藏：\InnerQuote{因为，}引其原话，\InnerQuote{自然喜爱隐藏；诸神隐藏的本质不能以赤裸裸的言辞投入受污染的耳朵。}由此显然，皇帝对圣经抱有这样的观念，即它们是一些神秘言谈，内含深奥意义。他对众人不持相同见解也极为恼怒，并抨击那些按更字面意义理解圣经的基督徒。但他如此激烈地辱骂平民的单纯，并为此对圣经傲慢无礼，实在不当；他也无权因他人正确领会了事物而按自己所不愿的方式理解，就对其避而远之。然而，看来一种类似的厌恶原因也作用于他，正如同影响波斐利的那种；波斐利在巴勒斯坦的凯撒利亚被一些基督徒殴打，因无法忍受（这种待遇），被无法抑制的愤怒驱使而背弃了基督教；并因仇恨那些殴打他的人而开始撰写反对基督徒的亵渎作品——正如尤西比乌·庞菲鲁斯所证明，他同时驳斥了其著作。同样，皇帝在一群无知民众面前对基督徒口出轻蔑之辞，因同样的病态心态而陷入波斐利的亵渎。既然他们二人都故意堕入不敬，便因罪疚之意识而受惩罚。但诡辩家利巴尼乌斯嘲弄地说基督徒把\Quote{一个巴勒斯坦人既当作天主又当作天主子}，他似乎忘记了自己在演说结尾已神化尤利安：\InnerQuote{他们几乎杀死了，}他说，\InnerQuote{第一个报告他死讯的人，仿佛此人曾对一位神祇撒谎。}稍后他又补充：\InnerQuote{啊，众神所眷爱者！众神的门徒！众神的伙伴！}如今，尽管利巴尼乌斯或许另有他意，但既然他未能避免一个有时带贬义的双关语，他似乎说出了基督徒曾以讽刺方式说过的话。若他本意是赞美，就应避免歧义用语；正如他另一次受到批评时，避开某个词并将之从其作品中删除一样。再者，那人在基督内与神性结合，以致他虽然外表是人，却是不可见的天主，这两点最为真实，基督徒的神圣典籍清楚教导。但外邦人在相信之前无法理解，因为神圣神谕宣告：\BibleQuote{除非你们相信，你们必定不能理解。}因此，他们不羞于将许多人列于他们的众神之中；但愿他们至少是将善良、正义、节制者列于其中，而非不洁、不义、耽于酗酒者，例如赫拉克勒斯、狄俄尼索斯和埃斯库拉庇乌斯——利巴尼乌斯在其演说中常以这些神祇起誓而毫不脸红。若我试图一一列举这些神祇的反常淫行和可耻通奸，离题将过于冗长；但渴望了解此事的读者，\WorkTitle{亚里士多德的披风}、\WorkTitle{狄俄尼索斯的花冠}、\WorkTitle{瑞吉努斯的波利姆尼蒙}以及所有诗人便足以证明外邦神学是一堆荒诞不经之物。我们确实可以多种实例表明，将人神化的做法在外邦人中绝非罕见，甚至他们毫不犹豫地这样做：仅举几例即可。罗得岛人因某公共灾祸求问神谕，答复指示他们应朝拜阿提斯——一位在弗里吉亚创制狂热仪式的异教祭司。神谕如此说：\par

\Quote{平息阿提斯，伟大的神，贞洁的阿多尼斯，有福的、金发的、恩赐丰厚的狄俄尼索斯。}\par

此处阿提斯，因情欲癫狂而自阉者，被神谕称为阿多尼斯和巴克斯。\par

再者，当马其顿王亚历山大进入亚洲时，近邻同盟会议讨好他，女祭司皮提亚发出此神谕：\par

\Quote{向至尊神宙斯和雅典娜·特里托革尼亚致敬，并向那隐藏于凡人形态中的神圣君王致敬，宙斯生他以荣耀，为凡人中的保护者与正义分配者——亚历山大王。}\par

这是德尔斐的魔鬼所说的话，他为了谄媚权贵，毫不犹豫地将他们列在众神之中。其动机或许是为了以奉承来和解；但对于拳击手\Person{克莱奥梅德斯}的情况，他们在这个神谕中将他列在众神之中，又能说什么呢？\par

\Quote{最后的英雄是阿斯提帕莱亚的\Person{克莱奥梅德斯}。以祭献尊崇他，因为他已不再是凡人。}\par

因为这个神谕，犬儒\Person{狄奥革涅斯}和哲学家\Person{俄诺玛俄斯}严厉谴责\Person{阿波罗}。\Place{库齐库斯}的居民宣称\Person{哈德良}为第十三位神；而哈德良本人也将他的娈童\Person{安提诺乌斯}奉为神。\Person{利巴尼乌斯}并未称这些为\Quote{可笑可鄙的荒谬之言}，尽管他熟悉这些神谕，也熟悉\Person{阿德里亚斯}关于\Person{亚历山大}（\Place{帕夫拉戈尼亚}的假先知）生平的著作；而且他本人也不犹豫以类似的方式尊崇\Person{波菲利}，在将\Person{尤利安}的著作置于波菲利之上后，他说：\Quote{愿叙利亚人对我仁慈。}这一题外话足以反驳这位诡辩家的嘲弄，无需进一步追随他提出的观点；因为进行彻底的反驳需要一部专门的著作。因此，我们将继续我们的历史。\par

% structure: p0063 source=h2 target=subsection reason=relative-heading-rank; base=h2

\subsection{第24章 主教们聚集在\Person{约维安}周围，各自试图将他拉向自己的信经。}

\Person{约维安}从\Place{波斯}返回后，教会的骚动再次掀起：因为那些治理教会的人争先恐后，希望皇帝依从他们自己的主张。然而他从一开始就持守\OriginalTerm{同质}{homoousios}信仰，并公开宣称他偏爱此信仰胜过其他一切。此外，他写信鼓励\Place{亚历山大里亚}的主教\Person{亚太纳削}，后者在\Person{尤利安}死后立即恢复了亚历山大里亚教会，那时因着这些书信而获得信心，摆脱了一切恐惧。皇帝还将\Person{君士坦提乌斯}放逐、且未被\Person{尤利安}重新安置的所有主教召回。异教庙宇再次关闭，他们尽可能地隐藏起来。哲学家也脱下披风，穿上普通服装。那种在\Person{尤利安}统治期间大量甚至令人厌恶地献上牺牲之血的公众污染，现在也被消除了。\par

% structure: p0065 source=h2 target=subsection reason=relative-heading-rank; base=h2

\subsection{第25章 马其顿派与阿卡齐乌派在\Place{安提阿}会合，宣布赞同\WorkTitle{尼西亚信经}。}

与此同时，教会的情况绝不安宁；因为各教派的首领殷勤地向皇帝献媚，认为保护自己就意味着获得对抗公认对手的权力。首先，马其顿派向他呈递了一份请愿书，恳求将所有主张圣子与圣父不相似的人逐出教会，并允许他们取代其位。这份请愿书由\Place{安基拉}的\Person{巴西略}、\Place{塔尔苏斯}的\Person{西尔瓦努斯}、\Place{庞培波利斯}的\Person{索弗洛尼乌斯}、\Place{泽莱}的\Person{帕西尼库斯}、\Place{科马纳}的\Person{莱昂提乌斯}、\Place{克劳迪奥波利斯}的\Person{卡利克拉特斯}和\Place{卡斯塔巴拉}的\Person{特奥菲卢斯}呈递。皇帝阅览后，没有给他们任何其他答复，只说：\Quote{我憎恶纷争；但我爱并尊敬那些努力促进和睦的人。}当这番话广为人知后，抑制了那些渴望争执之人的激烈情绪，于是皇帝的意图得以实现。此时，阿卡齐乌派的真实精神及其随时调整意见以适应最高掌权者的态度，比以往更加明显。因为他们在\Place{叙利亚}的\Place{安提阿}聚集，与\Person{美利提乌斯}（他不久前已脱离他们，接受了\OriginalTerm{同质}{homoousios}观点）举行会议。他们这样做是因为看见美利提乌斯深受当时驻在安提阿的皇帝器重；因此，他们达成一致，起草了一份声明，承认\Emph{同质}并批准\WorkTitle{尼西亚信经}，呈递给皇帝。其内容如下。\par

\Quote{从各省份在\Place{安提阿}召集的主教会议，致最虔诚、蒙天主眷爱的我主\Person{约维安}·维克多·奥古斯都：}\par

\Quote{陛下之虔诚，首要目标在建立教会的和平与和谐，对此，最深虔的皇帝陛下，我们完全知晓。我们也深知，您明智地判断，承认正统信仰乃是这合一的精髓。因此，为免我们被列入那些篡改真理教义者之列，我们谨向陛下声明：我们接纳并坚定持守昔日尼西亚圣公会议所宣认的信仰。尤其因为\OriginalTerm{同性同体}{homoousios}一词，虽在某些人看来似乎新奇不当，但教父们已明智地解释，该词表明圣子是由圣父的实体所生，且祂在实体上与圣父相似。这并不是说在那不可言喻的受生中应理解为任何情感。而且教父们使用\OriginalTerm{实体}{ousia}（即\Quote{substance}）一词时，并未取其希腊人常用的任何意义；而是为了推翻亚略对基督所妄加的不虔断言，即：祂是由\InnerQuote{不存在之物}造成的。近来兴起的\OriginalTerm{异质派}{Anomoeans}更加胆大妄为地坚持这一异端，致使教会合一全然破坏。因此，我们在此声明后附上尼西亚主教会议所订的信仰文本，我们也完全满意于此。其文如下：\InnerQuote{我们信唯一的天主，全能的圣父}，以及信经的其余全文。我们这些签署人在呈交此声明时，最衷心地赞同其内容。梅里丢斯、安提约基雅主教，欧塞比、萨莫萨塔主教，埃瓦格留、西西里主教，乌拉尼、阿帕美雅主教，左伊禄、拉里撒主教，阿加丘、凯撒勒雅主教，安提帕特、罗苏斯主教，阿布拉米、乌里米主教，阿里斯托尼科、贝鲁斯河畔塞琉西亚主教，巴尔拉梅诺、帕加玛主教，乌拉尼、梅里提纳主教，玛格努、卡尔西顿主教，欧提基、厄琉特罗城主教，伊萨科奇、大亚美尼亚主教，提托、波斯特拉主教，伯多禄、西皮主教，佩拉纠、劳迪刻雅主教，阿拉比安、安特罗斯主教，皮索、阿达纳主教（由拉米德里翁司铎代表），萨比尼安、泽格马主教，阿塔纳修、安基拉主教（由奥菲图和艾提乌司铎代表），依勒尼翁、迦萨主教，皮索、奥古斯塔主教，帕特里丘、帕尔图斯主教（由拉米里翁司铎代表），阿纳托利、贝洛雅主教，德奥提摩、阿拉伯人主教，以及路济安、阿尔卡主教。}\par

我们发现此声明载于撒彼诺那部题为\WorkTitle{主教会议法案汇编}的著作中。皇帝曾决意尽可能以温和的态度和说服的语言平息双方争执者的好斗精神；他声明\Quote{不会因宗教信仰骚扰任何人，且应当爱护和高度敬重那些热心促进教会合一的人}。哲学家特米斯提在为他\Quote{执政官任期}所写的演说词中证实了这样的行为。因为他赞扬皇帝允许每个人凭良心自由敬拜天主，从而战胜了谄媚者的诡计。并暗讽那些阿谀奉承者所受的打击，他诙谐地指出，经验表明这些人\Quote{崇拜的是紫袍而非天主；他们好似变幻无常的尤里普斯海峡，时而波浪漫向一方，时而流向完全相反的方向}。\par

% structure: p0070 source=h2 target=subsection reason=relative-heading-rank; base=h2

\subsection{第二十六章　若维安皇帝驾崩}

皇帝就这样在当时抑制了那些好争辩者的急躁；随后立即离开安提约基雅，前往基里基雅的塔尔索，在那里妥善地为尤里安举行了葬礼，之后他被宣布为执政官。由此直接前往君士坦丁堡，到达了一个名为达达斯塔纳的地方，位于加拉提雅与比提尼亚交界处。在那里，哲学家特米斯提与其他元老院成员迎接他，并在他面前宣读了执政官演说，之后他又在君士坦丁堡向民众重述。确实，罗马帝国得此卓越君主，本应极度繁荣，很可能民事与教会事务都会得到良好治理，若非他的突然去世使国家失去如此杰出的人物。因为某种阻塞引起的疾病在上述地方冬季发作，他于二月十七日，在其本人与其子瓦洛尼安担任执政官之年，享年三十三岁，在位七个月后驾崩。\par

本书记载了两年零五个月期间发生的事件。\par

\SourceRecord{Translated by A.C. Zenos.；Nicene and Post-Nicene Fathers, Second Series；Vol. 2.；Edited by Philip Schaff and Henry Wace.；Buffalo, NY: Christian Literature Publishing Co.,；1890.；<http://www.newadvent.org/fathers/26013.htm>.}

% structure: p0001 source=h1 target=section reason=page-title

\section{教会史（第四卷）}

% structure: p0002 source=h2 target=subsection reason=relative-heading-rank; base=h2

\subsection{第一章 若维安死后，瓦伦提尼安被拥立为帝，并以其弟瓦伦斯为同朝皇帝；瓦伦提尼安持守正统信仰，而瓦伦斯则为亚流派信徒。}

皇帝\Person{若维安}于其本人与儿子\Person{瓦洛尼安}的执政官任期内2月17日在\Place{达达斯塔纳}去世，如我们所述，军队于七日内从\Place{加拉提亚}抵达\Place{俾提尼亚}的\Place{尼西亚}，并于同一年任期内2月25日一致拥立\Person{瓦伦提尼安}为帝。他出身潘诺尼亚人，原籍\Place{基巴利斯}城，担任军事指挥，在战术上展现卓越才能。此外，他禀性宽宏，以至于无论获得何种荣誉，他似乎总是超然其上。他被拥立为帝后，立即启程前往\Place{君士坦丁堡}；在其即位后三十天，他使其弟\Person{瓦伦斯}成为帝国同僚。二人均自称信奉基督教，但所持信条不同：\Person{瓦伦提尼安}尊重尼西亚信经，而\Person{瓦伦斯}则偏爱亚流观点。这种偏见源于他被\Place{君士坦丁堡}主教\Person{欧多克西乌斯}施洗。二人均热心于各自党派的观点；但一旦掌握最高权力，便表现出非常不同的性情。因为在\Person{尤利安}统治时期，\Person{瓦伦提尼安}担任军事护民官，\Person{瓦伦斯}在皇帝卫队中指挥，二人都证明了自己对信仰的热忱；因被迫献祭，他们宁愿放弃军阶也不愿献祭而背弃基督教。然而，\Person{尤利安}深知此二人对国家不可或缺，便保留了他们的职位，继承帝位的\Person{若维安}亦然。后来，当他们获得帝权后，在公共事务上意见一致，但如前所述，在基督教方面，他们的行为大相径庭：\Person{瓦伦提尼安}虽偏爱与自己观点一致者，却未对亚流派施以暴力；而\Person{瓦伦斯}则热心推广亚流派事业，严重骚扰那些与他们观点相异者，正如我们历史进程所示。当时，\Person{利贝里乌斯}统领罗马教会；在\Place{亚历山大里亚}，\Person{亚大纳削}是\OriginalTerm{同质派}{Homoousians}的主教，而\Person{路济乌斯}被亚流派立为\Person{乔治}的继任者。在\Place{安提约基雅}，\Person{尤佐伊乌斯}是亚流派的领袖；但\OriginalTerm{同质派}{Homoousians}分为两派，一派以\Person{保利努斯}为首，另一派以\Person{梅利提乌斯}为首。\Person{济利禄}再次被立为\Place{耶路撒冷}教会的主教。\Place{君士坦丁堡}的教会由\Person{欧多克西乌斯}治理，他公开教导亚流主义的教义，但\OriginalTerm{同质派}{Homoousians}在城中仅有一所小教堂可供集会。属于马其顿派异端者，即那些在\Place{塞琉西亚}与阿卡西派分裂者，当时仍在各城保有他们的教堂。当时教会事务的状况就是如此。\par

% structure: p0004 source=h2 target=subsection reason=relative-heading-rank; base=h2

\subsection{第二章 \Person{瓦伦提尼安}前往西方；\Person{瓦伦斯}留在\Place{君士坦丁堡}，应允马其顿派举行会议的请求，却迫害拥护\OriginalTerm{同质}{Homoousion}者。}

两位皇帝中，\Person{瓦伦提尼安}迅速前往帝国西部；因为当时局势需要他出现在那里。与此同时，\Person{瓦伦斯}留在\Place{君士坦丁堡}，多数马其顿派异端主教向他请求召开另一次会议以修正信经。皇帝以为他们与\Person{欧多克西乌斯}和\Person{阿卡西乌斯}意见一致，便允准了他们；于是他们准备在\Place{兰普萨库斯}城集会。但\Person{瓦伦斯}急速赶往叙利亚的\Place{安提约基雅}，担心波斯人破坏在\Person{若维安}统治时期签订的为期三十年的条约，入侵罗马领土。然而波斯人保持安静；\Person{瓦伦斯}利用这段外部平静的时期，向所有承认\OriginalTerm{同质}{homoousion}者发动了一场消灭战争。他们的主教\Person{保利努斯}因卓越的虔诚，他未予骚扰；\Person{梅利提乌斯}则处以流放；其余所有拒绝与\Person{尤佐伊乌斯}共融的人，他都赶出安提约基雅的教会，并使他们遭受各种损失和惩罚。甚至据说他使许多人在流经该城的\Place{奥龙特斯河}中溺死。\par

% structure: p0006 source=h2 target=subsection reason=relative-heading-rank; base=h2

\subsection{第三章 当\Person{瓦伦斯}在东方迫害正统基督徒时，一名叫\Person{普罗科皮乌斯}的篡位者在\Place{君士坦丁堡}兴起；同时发生地震和洪水，毁坏数城。}

当\Person{瓦伦斯}在\Place{叙利亚}忙于此事时，一名叫\Person{普罗科皮乌斯}的篡位者在\Place{君士坦丁堡}兴起；他很快聚集了大批军队，策划远征对抗皇帝。这一消息使皇帝心中极度忧虑，并暂时制止了他对一切敢于在意见上与他不同者发起的迫害。正当人们痛苦地预料内战的骚动时，一场地震发生，对许多城市造成严重破坏。海水也改变了惯常的边界，在某些地方泛滥到如此程度，以至于船只可以在以前是道路的地方航行；而在其他地方则退去甚远，使地面干涸。这些事件发生在两位皇帝的第一次执政官任期内。\par

% structure: p0008 source=h2 target=subsection reason=relative-heading-rank; base=h2

\subsection{第四章 马其顿派在\Place{兰普萨库斯}举行会议，正值世俗与教会动荡时期；在确认\Place{安提约基雅}信经并诅咒\Place{阿里米努姆}所颁布的信经后，他们再次批准废除\Person{阿卡西乌斯}和\Person{欧多克西乌斯}。}

當這些事件發生時，無論在教會或國家都不得安寧。此時，在剛剛提及的執政官任期內，受皇帝委派召開大公會議的人聚集於\Place{朗普萨科斯}；這是在\Place{塞琉西亚}大公會議之後七年。在那裡，他們確認了在塞琉西亚所簽署的\WorkTitle{安提约基雅信經}後，就譴責了他們從前的同道在\Place{阿里米努姆}所頒布的信經。他們還再次譴責了\Person{阿卡西乌斯}與\Person{欧多克西乌斯}一派，並宣佈他們的撤職是正當的。當時即將爆發的內戰使\Place{君士坦丁堡}主教\Person{欧多克西乌斯}無法反駁或報復這些決定。因此，\Place{基齊库斯}主教\Person{厄勒乌西乌斯}及其追隨者一度佔了上風；因為他們支持\Person{马其顿尼}的觀點，這些觀點雖然之前鮮為人知，卻因朗普萨科斯主教會議而廣為人知。我認為，這次主教會議是赫勒斯滂地區馬其頓派增長的原因；因為朗普萨科斯位於赫勒斯滂的狹窄海灣之一。這次會議的結果就是如此。\par

% structure: p0010 source=h2 target=subsection reason=relative-heading-rank; base=h2

\subsection{第五章　\Person{瓦伦斯}與僭主\Person{普罗科皮乌斯}在\Place{弗里吉亚}的\Place{纳科利亚}附近交戰；其後，僭主被其將領出賣，與他們一同被處死。}

次年，在\Person{格拉提安}與\Person{达伽莱弗斯}擔任執政官期間，戰爭開始了。僭主\Person{普罗科皮乌斯}一離開\Place{君士坦丁堡}，率領軍隊向皇帝進發，\Person{瓦伦斯}便從\Place{安提约基雅}趕來，在\Place{弗里吉亚}一個名為\Place{纳科利亚}的城市附近與他交戰。初次交鋒時，瓦倫斯戰敗；但不久後，他藉著兩位將軍\Person{阿吉洛}和\Person{戈马利乌斯}的背叛活捉了普罗科皮乌斯，並對他們處以極其嚴厲的刑罰。他下令將這兩名叛徒鋸成兩半，不顧他曾向他們起誓。兩棵相近的樹被強力彎下，將僭主的一條腿分別綁在每棵樹上，然後突然鬆開樹木，使其恢復直立，樹木的彈力將暴君撕成兩半；僭主就這樣被撕裂而死。\par

% structure: p0012 source=h2 target=subsection reason=relative-heading-rank; base=h2

\subsection{第六章　\Person{普罗科皮乌斯}死後，\Person{瓦伦斯}強迫組成主教會議的人及所有基督徒信奉亞略異端。}

皇帝如此成功地結束了衝突，立即開始對付基督徒，企圖使每一個異端派別都皈依亞略異端。但他尤其惱怒那些組成\Place{朗普萨科斯}主教會議的人，不僅因為他們撤換了亞略派主教，還因為他們譴責了在\Place{阿里米努姆}公佈的信經。因此，當他到達\Place{比提尼亚}的\Place{尼科美底亞}時，他派人找來\Place{基齊库斯}主教\Person{厄勒乌西乌斯}，此人如前所述，緊隨\Person{马其顿尼}的觀點。於是，皇帝召集了一次亞略派主教會議，命令厄勒乌西乌斯同意他們的信仰。起初他拒絕了，但被流放和沒收財產的威脅所恐嚇，他因恐懼而同意了亞略派的信仰。然而，他隨即後悔了；回到基齊庫斯後，他在全體百姓面前痛苦地抱怨，聲稱他的順從是出於強迫，而非自願。他然後勸他們為自己另找一位主教，因為他被迫放棄了自己的意見。但基齊庫斯的居民太愛戴和尊敬他，不願失去他；因此他們拒絕服從任何其他主教，也不許他離開自己的教會：於是在他的監督下，他們繼續堅守自己的異端。\par

% structure: p0014 source=h2 target=subsection reason=relative-heading-rank; base=h2

\subsection{第七章　\Person{欧诺米乌斯}取代馬其頓派\Person{厄勒乌西乌斯}，擔任\Place{基齊库斯}主教；他的出身及模仿\Person{埃提乌斯}，曾是他的抄寫員。}

君士坦丁堡主教得知这些情况后，便立欧诺弥乌斯为基齐库斯主教，因他是一位有能力以其口才赢得众人归附其思想的人。欧诺弥乌斯抵达基齐库斯时，颁布了一道帝国敕令，命令驱逐厄肋乌息乌，由欧诺弥乌斯接替其位。此事执行后，追随厄肋乌息乌的人在城外建造了一座圣堂，与他一同在那里集会。关于厄肋乌息乌的话已足矣；现在让我们略述欧诺弥乌斯。他曾是阿厄提乌斯（绰号\Quote{无神者}，前文已提及）的秘书，并从他那里学会了模仿其诡辩式的推理方式；却未察觉，自己在练习制造谬误论证和使用某些微不足道的术语时，其实是在欺骗自己。然而这种习惯使他充满骄傲，陷入亵渎的异端，成为亚略教义的拥护者，并以多种方式成为真理教义的敌人。由于他对圣经字句只有非常浅薄的认识，完全无法领会其精神。但他言辞丰富，习惯于用不同的术语重复同样的思想，却从未能清楚解释自己所提出的问题。他的七卷著作\WorkTitle{论宗徒致罗马人书}便是明证，他在其上投入了大量徒劳的功夫：尽管他用了无数言辞试图解释，却远未理解这封书信的宗旨和目的。他现存的其他作品也与此类似，任何愿意花功夫检查的人都会在其中发现，在丰富的辞藻中缺乏意义。这位欧诺弥乌斯由厄乌多克修斯擢升为基齐库斯主教；他到达后，以他那\Quote{辩证术}的非凡展示令听众震惊，于是在基齐库斯引起了巨大轰动。最终，人们再也无法忍受他空洞而自负的言辞排场，将他逐出城。于是他退往君士坦丁堡，与厄乌多克修斯同住，被视为一位名义主教。但为了避免我们似乎出于毁谤而说这些话，让我们听听欧诺弥乌斯本人在其诡辩言谈中关于天主自身的胆大妄语，因为他用了以下言辞：\Quote{天主对自己的本质所知并不比我们更多；这本质对他来说并不比我们更为可知：凡我们关于神圣本质所知道的，天主恰恰也知道；反之，凡他所知道的，你也会发现在我们身上毫无差别。}欧诺弥乌斯还习惯于在对其自身愚钝的完全麻木中编造许多其他类似冗长而荒谬的谬误。他后来因何原故与亚略派分离，我们将在适当的地方叙述。\par

% structure: p0016 source=h2 target=subsection reason=relative-heading-rank; base=h2

\subsection{第八章　关于瓦伦斯皇帝下令拆毁卡尔西顿城墙时在一块石头上发现的谶语}

皇帝下令拆毁与拜占庭相对的城市卡尔西顿的城墙：因为他曾发誓，在征服僭越者之后必行此事，因为卡尔西顿人曾支持僭越者，对瓦伦斯说了侮辱的话，当他经过该城时关闭城门不让他进入。因此，根据帝国敕令，城墙被夷平，石头被运往君士坦丁堡，用于建造被称为\Quote{君士坦提亚纳}的公共浴场。在其中一块石头上发现了一条刻写的谶语，隐藏已久，预言当城市获得充足水源时，城墙将用于浴场；无数蛮族部落蹂躏罗马帝国各省，造成巨大破坏后，终将自身灭亡。我们在此录下这条谶语，以飨好学之士：\par

\Quote{当水神以湿足踏过骄傲的拜占庭的大街，\newline{}跳起神秘的舞蹈；\newline{}当愤怒的城墙被推倒，\newline{}其石块将去围成浴场：\newline{}那时蛮荒之地将派出无数蜂群，\newline{}饰以金发和闪亮的武器，\newline{}他们越过伊斯忒河的银色水流，\newline{}将蹂躏斯基泰田野和默西亚的草地。\newline{}但当他们征服了色雷斯，满载着胜利的喜悦，\newline{}命运将在那里赐予他们安葬之地。}\par

这便是那预言。事实上后来发生的事是，当瓦伦斯修建引水道使君士坦丁堡拥有充足水源时，蛮族部落进行了多次入侵，正如我们后文将要看到的。但也有人另有解释。因为当那条引水道完工时，城市长官克勒阿尔库斯建造了一座宏伟的浴场，命名为\Quote{丰水浴场}，位于现今称为德奥多西广场的地方；因此人们欢庆节日，他们认为这应验了谶语中那些话：\par

\Quote{他们以湿足跳起神秘的舞蹈，\newline{}踏过骄傲的拜占庭的大街。}\par

但预言的应验发生在后来。当拆毁进行时，君士坦丁堡居民恳求皇帝暂停毁墙；尼科美底亚和尼西亚的居民也从未提尼亚派人到君士坦丁堡提出同样请求。但皇帝对卡尔西顿人极为恼怒，勉强听取了这些请愿；但为了履行誓言，他命令拆除城墙，同时用其他小石头填补缺口。因此，今天在城墙的某些部分可以看到，极为巨大的石块上覆盖着极其劣质的材料，形成当时修补的难看补丁。关于卡尔西顿城墙，这就足够了。\par

% structure: p0022 source=h2 target=subsection reason=relative-heading-rank; base=h2

\subsection{第九章　瓦伦斯因诺瓦替派接受正统信仰而迫害他们}

然而那位皇帝并未停止迫害那些信奉\OriginalTerm{同质}{homoousion}道理的人，反而将他们逐出君士坦丁堡；诺瓦提安派既承认同一信仰，也遭受了类似的对待。他下令封闭他们的教堂，并将他们的主教放逐。该主教名叫阿格里乌（\Person{Agelius}），自君士坦丁时代起就一直管理他们的教会，度着宗徒般的生活：因为他总是赤足行走，只穿一件外衣，遵守福音的训示\BibleRef{玛 10:10}。但皇帝对这一教派的恼怒，因一位虔诚而雄辩之人名叫玛尔西安（\Person{Marcian}）的努力而得以缓和；此人曾在皇宫服兵役，但当时是诺瓦提安教会的一位司铎，并教授瓦伦斯皇帝的女儿阿纳斯塔西娅（\Person{Anastasia}）和卡罗萨（\Person{Carosa}）文法；前者之名被用于瓦伦斯在君士坦丁堡所建的、至今尚存的公共浴场。因此，出于对这位人物的尊重，那些关闭了一段时间的诺瓦提安教堂再次开放。然而，亚略派人士不愿让这些人安宁，因为他们不喜欢诺瓦提安派对\OriginalTerm{同质者}{Homoousians}所表现的同情与爱，诺瓦提安派在见解上与他们一致。这就是当时的情况。我们在此可以提及，反对篡位者普罗科皮乌（\Person{Procopius}）的战争大约在五月末，格拉提安（\Person{Gratian}）与达加拉伊夫斯（\Person{Dagalaïfus}）执政期间结束了。\par

% structure: p0024 source=h2 target=subsection reason=relative-heading-rank; base=h2

\subsection{第十章 小瓦伦提尼安的诞生}

这场战争结束后不久，同一执政年间，西方皇帝瓦伦提尼安得了一个儿子，以父亲的名字命名。因为格拉提安在他登基之前就已经出生了。\par

% structure: p0026 source=h2 target=subsection reason=relative-heading-rank; base=h2

\subsection{第十一章 异常大小的冰雹；比提尼亚与赫勒斯滂的地震}

次年6月2日，卢皮奇努（\Person{Lupicin}）与约维安（\Person{Jovian}）执政期间，君士坦丁堡降下冰雹，其大如人手。许多人断言这场冰雹是出于天主的愤怒，因为皇帝放逐了数位从事神圣职务的人，即那些拒绝与欧多克西乌（\Person{Eudoxius}）共融的人。同一年执政期间，8月24日，皇帝瓦伦提尼安立其子格拉提安为奥古斯都。次年，瓦伦提尼安与瓦伦斯再次担任执政官，10月11日，比提尼亚发生地震，摧毁了尼西亚城。这大约是在尼科美底亚遭受类似灾难之后十二年。不久之后，赫勒斯滂的革尔玛（\Place{Germa}）大部分地区因另一场地震而沦为废墟。然而，这些事件并未触动亚略派主教欧多克西乌或皇帝瓦伦斯的心；因为他们并未停止对那些在信仰上与他们意见相左之人的残酷迫害。与此同时，这些大地的震动被视为教会动荡的预兆：正如我们所说，许多神职人员被放逐；唯独巴西略（\Person{Basil}）与额我略（\Person{Gregory}）因天主上智的特殊安排，因其卓越的虔诚而免于此罚。前者是卡帕多西亚的凯撒勒雅主教；后者则管理凯撒勒雅附近的小城纳齐盎。但我们今后在叙述中还要提及巴西略和额我略。\par

% structure: p0028 source=h2 target=subsection reason=relative-heading-rank; base=h2

\subsection{第十二章 马其顿派因受皇帝压迫，遣使往罗马主教利贝里乌处，并签署尼西亚信经}

当坚持\Quote{同质}道理的人遭受如此严厉的对待并被赶散后，迫害者又开始骚扰马其顿派；他们出于恐惧而非暴力，从一个城市到另一个城市彼此派遣代表团，宣布有必要向皇帝的兄弟以及罗马主教利贝里乌（\Person{Liberius}）申诉：并且对他们而言，接受他们的信仰远胜于与欧多克西乌一派共融。为此，他们派遣了塞巴斯特的主教欧斯塔修（\Person{Eustathius}）——此人曾多次被免职——、奇里乞亚塔尔苏斯的西尔瓦诺斯（\Person{Silvanus}），以及同省的卡斯塔巴拉的特奥斐禄（\Person{Theophilus}）；委托他们在信仰上不得与利贝里乌有任何分歧，而要进入与罗马教会的共融，并确认\Quote{同质}的道理。于是这些人前往旧罗马，随身携带了那些在塞琉西亚与阿卡西乌（\Person{Acacius}）分离之人的信件。他们无法觐见皇帝，因为他在高卢忙于对萨尔马提亚人的战争；但他们向利贝里乌呈递了信件。利贝里乌起初拒绝接纳他们，说他们是亚略派的人，教会不可能接纳他们共融，因为他们已拒绝了尼西亚信经。他们回答说，他们已改变心意，承认了真理，早已弃绝了不同质信经，并宣认圣子在各方面\Quote{与父相似}；而且他们认为\Quote{相似}（\OriginalTerm{相似}{homoios}）与\Quote{同质}（\OriginalTerm{同质}{homoousios}）这两个词含义完全相同。他们作了这一陈述后，利贝里乌要求他们写一份信仰告白；于是他们向他呈递了一份文件，其中包含了尼西亚信经的要义。我在此未引述那些来自士麦拿、亚细亚、彼西底、以扫利亚、旁非利亚和吕基亚的信件，因为篇幅太长；这些地方都曾召开过主教会议。由欧斯塔修带去的代表们交给利贝里乌的书面信仰表白如下：\par

\Quote{致我们的主、兄弟及同工利贝里乌：欧斯塔修、特奥斐禄和西尔瓦诺斯在主内问安。}\par

\Quote{鉴于异端者的疯狂主张，他们不断在公教会中引入绊脚石，我们切望遏止其行径，遂出面表示赞同在\Place{兰普萨库斯}、\Place{士每拿}及其他各地召开的 orthodox 主教会议所承认的教义。我们由该会议组成代表团，向阁下及所有意大利与西方主教呈上一封书信，借此宣告：我们持守并维护在蒙福的君士坦丁治下于\Place{尼西亚}举行的神圣会议中，由三百一十八位主教所确立的大公信仰，该信仰至今保持完整无缺。在此信经中，\OriginalTerm{同质}{希腊文\GreekText{ὁμοούσιος}}一词被圣善而虔诚地用以对抗\Person{亚流}的毒害教义。因此，我们连同上述我们所代表的人，亲手承认我们已持守、正在持守并愿持守同一信仰，直至终末。我们谴责\Person{亚流}及其不敬神的教义，连同他的门徒及赞同其思想者；同样谴责\Person{撒伯里乌}的异端、\OriginalTerm{圣父受苦派}{Patripassians}、马西昂派、\Person{佛提诺}派、\Person{马尔塞鲁}派、\Person{萨莫萨塔的保罗}，以及支持此类主张者；总之，一切与上述神圣信经相悖的异端，该信经由圣教父们在\Place{尼西亚}以虔诚而大公的精神颁布。但我们特别咒诅在\Place{阿里米努姆}会议中宣读的那份信经形式，因其完全与之前\Place{尼西亚}神圣会议的信经相左——君士坦丁堡的主教们因诡计和伪誓而被欺骗（因其来自色雷斯的\Place{尼斯}镇）而签署了它。我们以及我们所代表之人的信经如下：}\par

\Quote{我们相信唯一的天主，全能的圣父，造天地及一切有形无形万物的主宰；并相信一位独一的独生子天主，主耶稣基督，天主子，由圣父所生，即由圣父的性体所生：天主出自天主，光明出自光明，真天主出自真天主，生而非受造，与圣父同体；万物藉祂而成，无论是天上的还是地上的。祂为了我们人类，并为了我们的救恩，降下，降生成人，成为人；受难，第三日复活，升天，将要来临审判生者死者。我们也相信圣神。但天主的公教会和宗徒教会咒诅那些断言\InnerQuote{祂曾有一时不存}、\InnerQuote{祂在受生之前不存在}、\InnerQuote{祂是由无中生有}的人；或那些说\InnerQuote{天主子与圣父是不同\OriginalTerm{位格}{hypostasis}或性体}的人，或说\InnerQuote{祂是可变的或可能改变的}。}\par

\Quote{我，\Person{塞巴斯蒂亚}城主教\Person{欧斯塔修}，连同\Place{兰普萨库斯}、\Place{士每拿}及其他会议的代表\Person{德奥斐洛}和\Person{息尔瓦诺}，自愿亲手签署此信仰宣认。如果在此信经公布后，有人胆敢诽谤我们或派遣我们的人，愿他持着阁下圣善的信函，到您圣善所认可的正统主教面前，与我们辩论；若有任何指控被证实，犯罪者应受处罚。}\par

\Person{利伯略}凭此文件稳妥地保证了代表们，接纳他们共融，随后以这封信遣返他们：\par

\Emph{罗马主教利伯略致马其顿派主教书}\par

致我们亲爱的弟兄和同工：\Person{埃维提乌}、\Person{济利禄}、\Person{许佩里基乌}、\Person{乌拉尼乌}、\Person{赫隆}、\Person{厄尔皮狄乌}、\Person{马克西穆}、\Person{欧塞比乌}、\Person{欧卡尔皮乌}、\Person{赫奥尔塔西乌}、\Person{尼昂}、\Person{欧马提乌}、\Person{福斯提诺}、\Person{普罗克利诺}、\Person{帕西尼库}、\Person{阿塞尼乌}、\Person{塞维鲁}、\Person{狄底弥昂}、\Person{布利塔尼乌}、\Person{卡利克拉特}、\Person{达尔马提乌}、\Person{埃德西乌}、\Person{欧斯托基乌}、\Person{安布罗西}、\Person{格洛尼乌}、\Person{帕尔达利乌}、\Person{马其顿尼乌}、\Person{保罗}、\Person{马尔塞鲁}、\Person{赫拉克利乌}、\Person{亚历山大}、\Person{阿多利乌}、\Person{马尔西安}、\Person{斯忒涅鲁}、\Person{若望}、\Person{马切尔}、\Person{卡里西乌}、\Person{息尔瓦诺}、\Person{佛提诺}、\Person{安多尼}、\Person{艾图斯}、\Person{策尔苏}、\Person{欧弗拉农}、\Person{米勒西乌}、\Person{帕特里基乌}、\Person{塞维里安}、\Person{欧塞比乌}、\Person{欧莫尔皮乌}、\Person{阿塔纳西}、\Person{狄奥凡图}、\Person{梅诺多鲁}、\Person{狄奥克勒}、\Person{克里桑佩鲁}、\Person{尼昂}、\Person{欧根尼}、\Person{欧斯塔修}、\Person{卡利克拉特}、\Person{阿塞尼乌}、\Person{欧根尼}、\Person{马蒂里乌}、\Person{希拉基乌}、\Person{莱昂提乌}、\Person{菲拉格里乌}、\Person{路济乌}，以及东方所有正统主教，意大利主教\Person{利伯略}及整个西方的主教们，在主内常致问候。\par

你们的信函，亲爱的弟兄，充满了信德的光芒，由我们极为敬重的弟兄、主教欧斯塔修、西尔瓦诺和德奥菲洛转交给我们，带来了我们长久盼望的和平与合一之喜乐；这主要是因为它们表明并给我们保证，你们的意见和情感与我们卑微之见，以及与意大利和西方所有主教的观点完全一致。我们承认这是公教和宗徒的信仰，直到尼西亚会议时期，它一直保持纯洁和不动摇。你们的使节声称他们自己持守这份信经，并且令我们极为喜悦的是，他们不仅以言语，而且以书面作证，抹去了所有有害猜疑的痕迹和印记。我们认为有必要在这封信后附上他们声明的副本，以免给异端者留下任何进行新阴谋的借口，使他们可以煽动自己恶意的余烬，并照他们的习惯重新点燃纷争之火。此外，我们极为敬重的弟兄欧斯塔修、西尔瓦诺和德奥菲洛也声称，他们自己和你们的爱德，始终持守并将持守到底由318位正统主教在尼西亚所认可的信经；这信经包含完全的真理，并驳斥和推翻所有异端的蜂群。因为如此众多的主教聚集反对亚略的疯狂，不是出于他们的意愿，而是出于天主的安排，其数目相当于蒙福的亚巴郎因信德消灭成千上万的敌人时所借助的人数。\BibleRef{创 14:14} 这信仰包含在\OriginalTerm{本体}{hypostasis}和\OriginalTerm{同体}{homoousios}这两个术语中，如同坚固不可攻破的堡垒，阻挡并击退亚略派乖戾的一切攻击和徒劳的阴谋。因此，当所有西方主教在阿里米农聚集时（亚略派的狡诈将他们引到那里，以便用欺骗性的劝说，或者更确切地说，用世俗权力的胁迫，来抹去或间接撤销如此谨慎地引入信经的内容），他们的诡计毫无成效。因为几乎所有在阿里米农被诱入错误或当时受骗的人，事后都对此事采取了正确的看法；在诅咒阿里米农会议所颁布的信仰声明后，他们签署了在尼西亚颁布的公教和宗徒信经。他们与我们共融，并对亚略及其门徒的教义更加憎恶，甚至对其感到愤怒。当你们的爱德的使节看到这些无可置疑的证据时，他们将你们自己附到自己的签署上，诅咒亚略和在尼西亚认可的信经下在阿里米农所行之事，甚至连你们自己，受伪誓的欺骗，也曾被诱导签署。因此，我们认为适合致信你们的爱德，并同意你们正当的请求，尤其是因为我们从你们使节的声明中确信，东方主教们已经恢复理智，现在与西方的正统派意见一致。我们还告知你们，免得你们不知道，阿里米农会议的亵渎已被当时似乎受欺诈所骗的人诅咒，并且所有人都承认了尼西亚信经。因此，你们应当普遍宣传，使那些因暴力或诡计而使信德受损的人，现在可以从异端的黑暗中浮现到公教自由的神圣光明中。此外，在这次会议之后，他们中凡不通过弃绝并诅咒亚略的所有亵渎而吐出败坏教义毒液的人，让他们知道，他们自己，连同亚略及其门徒和其他蛇类，无论是撒伯里乌派、圣父受难派或任何其他异端的追随者，都与教会的集会分离并被逐出，因为教会不接受非法的子女。愿天主使你们坚定不移，亲爱的弟兄。\par

当欧斯塔修的支持者收到这封信后，他们前往西西里，在那里召集了一次西西里主教会议，并在他们面前公开承认\OriginalTerm{同体}{homoousion}信仰，并声明他们坚持尼西亚信经；然后也从他们那里收到一封与前一封信内容相同的信函，便返回了派遣他们的人那里。他们这一方，在收到利伯略的信后，派遣代表从一城到另一城去见\OriginalTerm{同体}{homoousion}教义的主要支持者，劝告他们同时聚集在基里基雅的塔尔索，以确认尼西亚信经，并结束随后产生的所有争论。如果亚略派主教优多克修斯（他当时对皇帝有很大影响）没有阻挠他们的计划，这很可能已经实现；因为当他得知[在塔尔索]召集会议的消息时，他变得如此恼怒，以致加剧了对他们的迫害。关于马其顿派派遣使节给利伯略并被接纳与他共融并公开承认尼西亚信经一事，萨比努斯本人在他的\WorkTitle{会议事务汇编}中作了证明。\par

% structure: p0039 source=h2 target=subsection reason=relative-heading-rank; base=h2

\subsection{第13章。恩诺米乌斯与优多克修斯分离；亚大纳削在亚历山大里亚被优多克修斯挑起骚乱，再次自愿流亡，但由于民众的呼声，皇帝召回并恢复了他的主教职位。}

大约同时，欧诺米乌斯与欧多克修斯分离，另外举行集会，因为他多次恳求其导师埃提乌斯被接纳共融，欧多克修斯却一直反对。欧多克修斯这样做并非出于本意，因为他并不拒绝与埃提乌斯相同的意见；但他屈服于自己一派中的普遍情绪，他们反对埃提乌斯，视其为异端。这是欧诺米乌斯与欧多克修斯分裂的原因，君士坦丁堡的情况便是如此。然而，亚历山大里亚的教会因近卫总长经欧多克修斯传达的谕令而动荡不安。于是，亚大纳削惧怕民众不理性的冲动，又担心自己会成为可能发生之暴行的肇因，便躲入一座祖坟中，整整四个月。但因民众出于对他的爱戴，因他不在而不安、发生骚乱，皇帝得知此事后，确认亚历山大里亚因此动荡不安，遂下诏命令允许亚大纳削安然治理教会不受干扰；这就是亚历山大里亚教会在亚大纳削去世前得以安宁的原因。亚大纳削死后，亚流派如何占据教会，我们将在叙述历史的过程中说明。\par

% structure: p0041 source=h2 target=subsection reason=relative-heading-rank; base=h2

\subsection{第14章 欧多克修斯在君士坦丁堡去世后，亚流派任命德莫斐卢斯；正统派则推举埃瓦格里乌斯为继任者。}

皇帝瓦伦斯再次离开君士坦丁堡，向安提约基雅进发；但他抵达比提尼亚的尼科美底亚时，因以下情况而受阻。亚流派教会的主教欧多克修斯，在皇帝离开该城后不久，即瓦伦提尼安和瓦伦斯第三次执政之年去世，他占有君士坦丁堡教会席位已达十九年。因此，亚流派任命德莫斐卢斯继任；但\OriginalTerm{同质}{homoousion}派认为时机已至，推举了某位持守他们原则的埃瓦格里乌斯，并由曾任安提约基雅主教的尤斯塔修斯正式祝圣他。尤斯塔修斯被约维安从流放中召回，此时私下前来君士坦丁堡，以坚定\OriginalTerm{同质}{homoousion}道理的信众。\par

% structure: p0043 source=h2 target=subsection reason=relative-heading-rank; base=h2

\subsection{第15章 皇帝放逐埃瓦格里乌斯和尤斯塔修斯。亚流派迫害正统派。}

此事完成后，亚流派重新迫害同质派；皇帝很快得知所发生的事，担心因民众骚乱而导致城市倾覆，立即从尼科美底亚派遣军队前往君士坦丁堡；下令逮捕被祝圣者及祝圣者，并将他们流放到不同的地区。于是，尤斯塔修斯被流放到色雷斯的比齐亚城；埃瓦格里乌斯被解往别处。此后，亚流派更加大胆，严重骚扰正统派，经常殴打他们、辱骂他们、监禁他们、罚款他们；简言之，他们对他们施以痛苦难忍的侵扰。受难者被诱导向皇帝申诉，以求摆脱压迫者的保护，或许能从这种压迫中获得一些缓解。但他们从这一途径可能抱有的任何平反希望，都完全落空了，因为他们不过是向苦难的始作俑者陈诉冤情而已。\par

% structure: p0045 source=h2 target=subsection reason=relative-heading-rank; base=h2

\subsection{第16章 瓦伦斯下令在船上焚烧某些司铎。弗里吉亚发生饥荒。}

圣职人员中有八十位虔诚者，其中以乌尔巴努斯、特奥多雷和墨涅德摩斯为首，前往尼科美底亚，向皇帝呈递请愿书，告知并控诉他们所遭受的虐待。皇帝满心愤怒，但在他们面前隐藏不悦，秘密命令总督莫德斯图斯逮捕这些人并处死他们。他们被处死的方式不同寻常，值得记载。总督担心若公开处决这么多人，会激起民众对他发动暴乱，便假装将这些人遣送流放。于是，当他们以极大坚毅接受自己的命运时，总督命令将他们装船，仿佛要送往各自的流放地；同时指示水手们，一旦到达海中，就点燃船只，使被处死者甚至被剥夺埋葬。这道命令被执行了：当他们到达阿斯塔西亚湾中部时，船员们点燃船只，然后躲入紧随其后的一艘小艇中逃脱。与此同时，强劲的东风刮起，燃烧的船被猛烈推动，但行驶更快，一直保持完整，直到抵达一个名为达西迪祖斯的港口，在那里，船连同被囚禁在其中的人一起被彻底烧毁。许多人声称，这一不虔的行为并未免于惩罚：因为此后立即在整个弗里吉亚发生了如此严重的饥荒，以至于大部分居民被迫暂时离开家园，一些人逃往君士坦丁堡，一些人逃往其他省份。因为君士坦丁堡尽管供养着庞大的人口，却始终富足生活必需品，各种食物从不同地区海运而来；邻近的攸克辛海也供应它所需要的一切小麦。\par

% structure: p0047 source=h2 target=subsection reason=relative-heading-rank; base=h2

\subsection{第17章 皇帝瓦伦斯在安提约基雅再次迫害\OriginalTerm{同质}{homoousion}派信徒。}

瓦伦斯皇帝对饥荒造成的灾难漠不关心，前往叙利亚的安提约基雅，并在那里逗留期间残酷迫害那些不接受亚略主义的人。因为他不仅满足于将坚持\OriginalTerm{同质}{homoousian}观点的人几乎从东方所有的教会中驱逐出去，还对他们施加了各种处罚。他甚至比以前毁灭了更多的人，将他们交付各种不同的死法，尤其是溺死在河中。\par

% structure: p0049 source=h2 target=subsection reason=relative-heading-rank; base=h2

\subsection{第十八章。埃德萨事件：虔诚市民的坚毅与一位虔诚妇女的勇气。}

但我们必须在此提及在美索不达米亚的埃德萨发生的某些情况。在那座城市里，有一座宏伟的教堂奉献给宗徒圣多默，由于该地的神圣性，宗教集会不断举行。瓦伦斯皇帝想视察这座建筑，并得知所有通常聚集在那里的人都反对他所偏爱的异端，据说他亲手掌掴了长官，因为后者未能也将他们从那里驱逐出去。长官忍受了这种羞辱后，极不情愿地被迫顺从皇帝对他们的愤怒——因为他并不希望杀害如此众多的人——他私下暗示不应有人在那里被发现。但没有人听从他的劝诫或威胁；因为第二天，他们都涌向教堂。当长官带着一支大军前往那里以满足皇帝的愤怒时，一个穷苦妇女牵着她的幼子匆匆赶路，前往教堂，冲破了长官士兵的行列。长官被激怒，命令将她带到他面前，这样对她说：\Quote{不幸的妇人！你如此混乱地跑向哪里？}她回答说：\Quote{去别人都赶去的地方。}他说：\Quote{你没有听说长官将要处死所有在那里被发现的人吗？}妇人说：\Quote{是的，因此我急忙赶去，好让我在那里被找到。}长官说：\Quote{你拖着那个孩子去哪里？}妇人回答说：\Quote{好让他也配得殉道。}长官听了这些话，猜想其他聚集在那里的人也抱有同样的决心，立即回到皇帝那里，告知他所有的人都准备为自己的信仰而死。他补充说，一次毁灭如此众多的人是不明智的，于是说服皇帝控制他的怒火。这样，埃德萨人得以幸免，没有被皇帝下令屠杀。\par

% structure: p0051 source=h2 target=subsection reason=relative-heading-rank; base=h2

\subsection{第十九章。因异教预言，瓦伦斯屠杀多人，并记录其姓名。}

此时，一个可憎的恶魔滥用了皇帝的残忍性情，诱使某些好奇之人通过招魂术调查谁将继承瓦伦斯的王位。对他们的魔法咒语，恶魔给出了不清晰且不明确的回应，但正如通常的做法，充满了歧义；因为显示了四个字母\OriginalTerm{四字母}{q, e, o, and d}，他宣告瓦伦斯继任者的名字以这些字母开头，并且是一个复合名称。皇帝得知这个神谕后，没有将此事交托给唯一能洞察未来的天主，反而违背了他声称最热切信奉的基督徒原则，处死了许多他怀疑觊觎王权的人：这样，名叫\Quote{特奥多尔}、\Quote{特奥多图斯}、\Quote{特奥多西乌斯}、\Quote{特奥杜鲁斯}等的人，都成了皇帝恐惧的牺牲品；其中还有特奥多西奥卢斯，一个非常勇敢的人，出身于西班牙的贵族家庭。因此，许多人为了避免危险，更改了自己的名字，放弃了他们幼儿时从父母那里得到的名字，因为这些名字有危险。关于此事就说到这里。\par

% structure: p0053 source=h2 target=subsection reason=relative-heading-rank; base=h2

\subsection{第二十章。阿塔纳修的去世，以及伯多禄继其位。}

必须说明，只要亚历山大里亚的主教阿塔纳修还活着，皇帝就被天主的眷顾所约束，避免骚扰亚历山大里亚和埃及：事实上，他很清楚那些依附阿塔纳修的人数众多；因此他小心谨慎，以免因亚历山大里亚人——一个易怒的民族——被煽动叛乱而使公共事务陷入危险。但阿塔纳修在经历了如此众多而严峻的为教会的斗争后，在格拉提安与普罗布斯第二次执政期间去世，在极大的危险中治理该教会四十六年。他指定伯多禄为继承人，一位虔诚而雄辩的人。\par

% structure: p0055 source=h2 target=subsection reason=relative-heading-rank; base=h2

\subsection{第二十一章。亚略派信徒获皇帝允许监禁伯多禄，并立路基乌斯为亚历山大里亚主教。}

于是，亚略派信徒因了解皇帝的宗教倾向而壮胆，再次鼓起勇气，立即将此事告知他。皇帝当时住在安提约基雅。那时，在该城主持亚略派信徒的欧佐依乌斯热切地抓住这个有利机会，请求允许他去亚历山大里亚，以便将那里的教会交给亚略派信徒路基乌斯。皇帝同意了这个请求，欧佐依乌斯带着帝国军队尽快立即前往亚历山大里亚。皇帝的司库玛格努斯也与他同去。此外，还向埃及总督帕拉狄乌斯发出帝国命令，指令他用军队协助他们。因此，他们逮捕了伯多禄，将他投入监狱；驱散了其余神职人员后，他们将路基乌斯安放在主教座位上。\par

% structure: p0057 source=h2 target=subsection reason=relative-heading-rank; base=h2

\subsection{第二十二章。萨比努斯对亚略派信徒的恶行保持沉默；伯多禄逃往罗马；在亚略派信徒煽动下屠杀隐修士。}

关于\Person{卢济乌斯}上任时发生的暴行，以及被逐者在法庭内外所受的待遇，以及有些人遭受各种酷刑，另一些人在经历这种折磨后还被流放，\Person{撒比努斯}对此丝毫不提。事实上，他本人对亚略主义半心半意，故意隐瞒他朋友们的暴行。然而，\Person{伯多禄}在逃离监狱后写信给所有教会时，揭露了这些事。因为这位（主教）设法逃出监狱后，逃到了罗马主教\Person{达玛苏斯}那里。亚略派人数虽不多，却因此控制了亚历山大各教会；不久后他们获得一项帝国法令，命令埃及总督驱逐拥护\Quote{同质}教义的人以及所有与\Person{卢济乌斯}作对的人，不仅从亚历山大驱逐，还要逐出埃及。此后，他们攻击、骚扰并严厉迫害旷野中的隐修院；武装暴徒以最凶残的方式冲向那些手无寸铁、毫不抵抗的人：无数不抵抗的受害者就这样被以难以形容的残忍方式屠杀。\par

% structure: p0059 source=h2 target=subsection reason=relative-heading-rank; base=h2

\subsection{第23章 一些献身独修生活的圣人之事迹}

既然我已提及埃及的隐修院，在此略加叙述或许是恰当的。它们或许在很早时期就已成立，但由一位虔诚之人，名叫\Person{阿蒙}，大大扩展和增建。他年轻时厌恶婚姻；但当一些亲戚力劝他不要轻视婚姻，而要娶妻时，他被说服并结了婚。在按惯例将新娘从宴客厅引到婚床，待双方朋友离去后，他拿起一本包含宗徒书信的书，向妻子朗读保禄致格林多人书，并向她解释宗徒对已婚者的劝诫。此外，他还列举许多外在考量，详细谈论婚姻交合的不便与不适，生育的阵痛，以及抚养家庭带来的烦恼与忧虑。他将这一切与贞洁的好处相对比；描述了节制生活的自由与无瑕纯洁；并断言童贞使人处于与天主最密切的关系中。通过这些及类似论据，他说服他的童贞新娘与他一起放弃世俗生活，在他们彼此有任何夫妻认识之前。下定这个决心后，他们一同退隐到\Place{尼特里亚}山，在那里的一间小屋里，短时间共同居住在一间苦修室，不顾性别的差异，因为按宗徒所说，\BibleQuote{在基督内成为一个}。但不久后，那新近且未被玷污的新娘这样对阿蒙说：\Quote{你实践贞洁，在如此狭小的居所里看一个女人，是不合适的；因此，如果你同意，我们分开操练吧。}这个协议两人都满意，于是他们分开，余生戒绝酒和油，只吃干面包，有时隔一天，有时禁食两天，有时更多。亚历山大主教\Person{亚大纳削}在他的\WorkTitle{安东尼传}中断言，与他同时代的该传主，看到阿蒙死后灵魂被天使接升。因此，许多人效仿阿蒙的生活方式，以至于\Place{尼特里亚}和\Place{司奇提斯}山逐渐充满了隐修士，记述他们生平需要一部专门的著作。然而，由于他们中有一些杰出虔诚之人，以严格的纪律和宗徒般的生活著称，说了和做了许多值得记录的事，我认为从众多事迹中精选一些细节穿插进我的历史，对读者是有益的。据说阿蒙从未见过自己赤裸，他惯常说：\Quote{隐修士连看见自己的身体暴露也不合适。}有一次他想过河，但不愿脱衣，他恳求天主使他能够过河而不必破坏决心；立刻有一位天使将他运到河对岸。另一位名叫\Person{狄第默}的隐修士，虽然已年届九十，却完全独自生活直到去世。他们中另一位\Person{阿尔塞尼}，不愿将年轻的犯过者与共融分离，只对年长者如此：\Quote{因为，当年轻人被绝罚时，他会变得刚硬；但年长者很快会感到绝罚的痛苦。}\Person{比奥尔}习惯于边走边进食。有人问他：\Quote{你为什么这样吃？}他说：\Quote{以免我显得把吃东西当作正经事，而只是顺便为之。}对另一个提出同样问题的人，他回答说：\Quote{免得连吃东西时我的心灵也感受到身体的享乐。}\Person{依西多尔}断言，他四十年来连在思想上也没有意识到罪；他从未赞同情欲或愤怒。\Person{庞博斯}是个不识字的人，他去某人那里学习一首圣咏；听了第三十八首圣咏的第一节：\BibleQuote{我说：我要谨守我的行径，免得我以口舌犯罪}，他没有留下来听第二节就走了，说：\Quote{这一节就够了，如果我能实际掌握它。}当给他那节经的人责备他六个月内没有见他时，他回答说他还未学会实践那节圣咏。过了相当长的时间，他的一位朋友问他是否掌握了那节经文，他回答说：\Quote{我十九年来几乎还没有成功做到。}有一个人把金子放在他手里让他分给穷人，请他数算给了他多少。他说：\Quote{不需要数算，需要的是心灵的诚实。}同一位庞博斯，应主教亚大纳削之愿，从旷野来到\Place{亚历山大}，在那里看见一位女演员，他哭了。在场的人问他为何哭泣，他回答说：\Quote{两个原因触动了我：一是这女人的毁灭；二是我努力讨我天主的喜悦，尚不及她努力讨淫秽人物的喜悦。}另一个人说：\Quote{不工作的隐修士应被视为与贪婪的人同等。}\Person{彼特鲁斯}精通自然哲学的许多分支，惯常频繁地讲解这一门或那一门科学的原理，但他总是以祈祷开始他的讲解。在那个时期的隐修士中，还有两位同名、大圣德的人，都叫\Person{玛加略}；一位来自\Place{上埃及}，另一位来自\Place{亚历山大}城。两人都以其苦修纪律、生活言谈的纯洁以及他们手中所行的奇迹而闻名。埃及的玛加略行了那么多治愈，驱了那么多魔鬼，以至于需要一部专门的论著来记录天主的恩宠所允许他做的一切。他对前来找他的人态度严峻，但同时能激发敬畏。亚历山大的玛加略，虽然各方面与埃及的同名人相似，但在此点上不同：他对来访者总是愉快；并以其温和的举止引导许多年轻人走向苦修生活。\Person{厄瓦格略}成了这些人的门徒，从他们那里获得了行为的哲学，而他先前只了解言辞的哲学。他在\Place{君士坦丁堡}由\Person{纳齐盎的额我略}被祝圣为执事，后来与他一同前往埃及，在那里结识了这些杰出人物，并效仿他们的行为方式，他手中所行的奇迹与他老师的一样多且重要。他还撰写了非常有价值的著作，其中一本名为\WorkTitle{隐修士}，或\WorkTitle{论主动德行}；另一本\WorkTitle{诺斯替者}，或\WorkTitle{致堪当知识者}：此书分为五十章。第三本名为\WorkTitle{反驳集}，包含从圣经中选出的对抗诱惑之灵的章节，按论题数目分为八部分。他还写了\WorkTitle{六百预言问题}，以及两首诗歌作品，一首致\WorkTitle{团体生活的隐修士}，另一首致\WorkTitle{童贞女}。任何人读了这些作品都会确信其卓越。我认为，在先前所述之后，附上他提及的关于隐修士的几件事，在此并非不合适。以下是他的话：\par

我们应当探究在我们之前的虔诚隐修士的习惯，以便通过他们的榜样纠正自己；因为毫无疑问，他们曾说过和做过许多极好的事。其中一人惯常说：\Quote{较为干燥且有节制的饮食，结合着爱德，会迅速引领隐修士进入宁静的港湾。}同一个人因受夜间幻象困扰而解救了他的一位弟兄，命令他在禁食时服务病人。当被问及为何如此规定时，他说：\Quote{除了藉着实践怜悯之心外，没有其他事物能如此有效地驱散这类困扰。}当时有一位哲学家来到公正者安东尼面前，对他说：\Quote{父亲，您没有书籍的安慰，怎能忍受呢？}安东尼回答说：\Quote{哲学家，我的书就是受造物的本性；每当我愿意阅读天主的话语时，它就在眼前。}那个\InnerQuote{蒙拣选的器皿}，\BibleRef{宗 9:15}——埃及长者玛加略——问我：为何铭记从人那里所受的伤害会削弱记忆官能的力量，而铭记魔鬼所做的伤害却不会使其受损？我犹豫着，几乎不知如何回答，便请他解释。他说：\Quote{因为前者是违反本性的情感，后者则合乎心灵的性情。}有一次，我在正午前往圣玛加略长老那里，因酷热和口渴而难以忍受，便求他给我一些水解渴。他回答说：\Quote{以荫凉为满足吧，因为许多正在陆路旅行或海上航行的人，连这荫凉也没有。}后来我与他讨论克己主题时，他说：\Quote{我儿，放心吧！二十年来，我从未吃过饱、喝过足、睡过够；我的面饼总是称过的，水是量过的，仅有的一点睡眠也是靠在墙上偷来的。} \BibleRef{厄上 4:10} 有人向一位隐修士报告了他父亲的死讯。他对那人说：\Quote{停止你的亵渎吧！我的父亲是不死的。}有一位弟兄除了一部福音书外别无所有，他便将它卖了，把所得的钱分给饥饿之人，并说出了这句值得铭记的话：\Quote{我卖了那本说\InnerQuote{变卖你所有的，分给穷人}的书。}\BibleRef{玛 19:21} 在亚历山大城外，有一处位于所谓玛利亚湖以北的岛屿，那里住着一位来自帕伦博莱的隐修士，在诺斯替者中享有盛誉。此人惯常说，隐修士的所有行为都出于以下五种原因之一：为天主、为本性、为习惯、为需要、或为劳作。同一个人还说，本性中只有一种德行，但它根据灵魂的性情呈现各种特征；正如太阳之光本身无形，却适应接受它的形状。另一位隐修士说：\Quote{我远离享乐，为要断绝愤怒的机会；因为我知道愤怒总是为享乐而争战，扰乱我心灵的安宁，使我无法获得知识。}一位年长隐修士说：\Quote{爱德不懂得储存食物或钱财。}他补充说：\Quote{我记不得曾两次被同一件事欺骗过。}这就是埃瓦格里乌斯在他的书《\WorkTitle{实践}》中所写的。在他称为《\WorkTitle{诺斯替者}》的书中，他说：\Quote{我们从公正者额我略学到：有四种德行，各有不同的特征：明智和勇敢，节制和正义。明智的职责是沉思那些神圣而理智的能力，超越语言表达，因为这些能力由智慧揭示；勇敢的职责是坚持真理，反对一切阻碍，永不转向虚假；节制的职责是从首席农夫接受种子，\BibleRef{玛 13:24} 但拒绝那要撒上另类种子的人；最后，正义的职责是根据每个人的状况和才能调整言辞，有时隐晦地陈述，有时比喻式地陈述，有时清楚地解释，以教导那些较不聪明的人。}那真理的柱石，卡帕多细亚的巴西略，曾说：\Quote{人们所教导的知识通过不断的研习和实践而完善；但来自天主恩宠的知识则通过实践正义、忍耐和仁慈而完善。}前者常常在那些尚受激情影响的人身上发展；而后者则是那些超越激情影响之人的份，他们在虔诚的时刻默观那照亮他们的心灵特殊之光。埃及的明灯，圣亚大纳削，告诉我们：\Quote{梅瑟奉命将桌子放在北\BibleRef{出 26:35}边。因此，让诺斯替者们知道什么风与他们为敌，从而勇敢地忍受一切诱惑，并乐意将滋养供给那些向他们求助的人。}特穆伊特教会的天使塞拉皮翁宣称：\Quote{心灵通过汲取灵性知识而完全净化}；\Quote{爱德治愈灵魂的炎症倾向}；\Quote{灵魂中滋生的邪恶欲望则被节制所抑制。}伟大的明师狄迪莫斯说：\Quote{你要持续操练自己默想天主眷顾和审判；并尽力记住你关于这些主题所听过的任何论说的素材，因为几乎所有人在这一点上都有欠缺。你会在受造形式的差异和宇宙的构造中，找到关于审判的推论；涵盖了那些引领我们从恶习与无知转向德行与知识的途径的讲道，则包含关于眷顾的论说。}\par

我们以为在此插入这些从埃瓦格里乌斯摘录的片段是适当的。在隐修士中还有另一位卓越的人物，名叫阿摩尼乌斯，他对世俗之事如此漠不关心，以致当他与阿塔纳修斯一同前往罗马时，他选择不去考察那座城市的任何宏伟建筑，只满足于参观伯多禄和保禄大殿。这位阿摩尼乌斯在被催促接受主教职务时，割掉了自己的右耳，以便通过残害自己的身体而使自己失去被祝圣的资格。但很久以后，当亚历山大主教德奥斐洛想要立埃瓦格里乌斯为主教时，埃瓦格里乌斯未以任何方式自残便逃脱了，后来偶然遇见阿摩尼乌斯，便开玩笑地对他说，他割掉自己的耳朵是做了错事，因为他在天主眼中使自己成了有罪的人。阿摩尼乌斯回答说：\Quote{你认为，埃瓦格里乌斯，你不会受到惩罚吗？你因自爱而割掉了自己的舌头，以避免行使那托付给你的言论恩赐。}当时在隐修院中还有许多其他令人钦佩和虔诚的人物，在此一一列举会过于冗长，而且如果我们试图描述每个人的生活以及他们因所领受的圣德而行的奇迹，我们将必然偏离我们既定的目标。如果有人想了解他们的历史，包括他们的行为、经历和教诲，以及他们如何驯服野兽，有一部专门的论著，由隐修士帕拉狄乌斯——他是埃瓦格里乌斯的门徒——所著，其中详细记述了这一切。在那部著作中，他还提到了几位妇女，她们实践了与上述男子相同的刻苦生活。埃瓦格里乌斯和帕拉狄乌斯都在瓦伦斯去世后不久兴盛。我们现在必须回到我们偏离的地方。\par

% structure: p0063 source=h2 target=subsection reason=relative-heading-rank; base=h2

\subsection{第二十四章 袭击隐修士及其首长，他们显示了奇迹般的能力。}

瓦伦斯皇帝颁布敕令，命令在亚历山大和埃及其他地方迫害正统派，随即发生了大规模的荒芜和毁灭：一些人被拖到法庭，另一些人被投入监狱，许多人受到各种折磨，实际上各种刑罚都施加在那些只求和平与安静的人身上。当这些暴行按卢修斯的心意在亚历山大施行后，尤佐依乌斯返回安提约基雅，而亚略派信徒卢西安由军队总司令率领一支相当数量的军队，立即前往埃及的隐修院，总司令本人比无情的士兵更加狂暴地攻击那一群圣人。到达那些僻静之处时，他们发现隐修士们正在从事他们惯常的操练：祈祷、治愈疾病、驱逐魔鬼。然而，他们不顾这些天主德能的非凡明证，不允许他们继续庄重的祈祷，而是用武力将他们赶出祈祷所。鲁菲努斯声称他不仅是这些残暴行为的见证人，也是受难者之一。因此，在他们身上重演了宗徒所说的那些事：\BibleRef{希 11:36} \BibleQuote{他们曾经受过戏弄、鞭打、捆绑、监禁，被石头砸死，被刀剑杀死，披着绵羊和山羊皮各处奔跑，受贫乏、患难、虐待，这世界原不配拥有他们，在旷野、山岭、山洞和地穴中漂流。}在这一切事上，他们因信德和行为以及基督的恩宠藉他们的手所行的治愈而\Quote{获得了美名}。但似乎天主的眷顾允许他们忍受这些恶事，\BibleRef{希 11:40} \BibleQuote{因为天主为我们预备了更好的事，} 好使他人通过他们的受苦获得天主的救恩，后来的事件似乎证明了这一点。因此，当这些非凡的人胜过了施加于他们的一切暴力时，卢修斯绝望地建议军事首领将隐修士们的院长放逐出境：他们是埃及的玛加略和他在亚历山大的同名者，两人因此被流放到一个没有基督徒居民的海岛，这个岛上有一座偶像庙宇，有一位祭司被居民敬奉为神。当这些圣人到达海岛时，那里的魔鬼充满了恐惧和战栗。恰在同一时间，祭司的女儿突然被魔鬼附身，开始狂暴地行动，推翻所遇到的一切；没有任何力量能制服她，她大声向天主的这些圣人们喊叫说：\BibleRef{玛 8:29} \BibleQuote{你为什么到这里来，也要从这里赶我们出去？} 于是，这些人也在那里显示了他们藉天主恩宠所领受的特殊能力：因为他们将魔鬼从少女身上驱逐出去，使她痊愈后交还给她的父亲，他们使祭司本人以及全岛居民都皈依了基督信仰。于是，他们立即打碎了偶像，将庙宇的形状改为教堂；受洗后，他们欢乐地接受了基督信仰的教义。因此，这些非凡的人因\OriginalTerm{同质}{homoousian}信仰而受迫害后，自身更受嘉许，成为他人得救的工具，并证实了真理。\par

% structure: p0065 source=h2 target=subsection reason=relative-heading-rank; base=h2

\subsection{第二十五章 论盲人狄迪莫斯}

大约在同一时期，天主又带领另一位忠信的人进入人们的视野，认为借着他在信德上作见证是相宜的：这就是\Person{狄迪穆}，一位极令人钦佩且富有口才的人，通晓他所处时代的一切学问。在很小的时候，当他几乎尚未掌握学问的初阶时，就患了眼疾，导致失明。但天主补偿了他身体视觉的丧失，赐予他更敏锐的理智。因为凡是他无法通过观看学习的，他都能通过听觉获得；因此，从小天赋卓越的他，很快就远远超过了那些拥有最敏锐视力的同龄伙伴。他以惊人的便利掌握了语法和修辞学的基本原理；进而转向哲学研究、辩证术、算术、音乐，以及他所专注的其他各种知识领域；他将这些科学分支如此珍藏于心，以至于随时准备好与那些通过阅读书籍而熟悉这些主题的人进行讨论。不仅如此，他对旧约和新约中所包含的神圣训言如此精通，以至于撰写了多部诠释它们的论著，此外还有三卷关于\WorkTitle{论圣三}的书。他还发表了关于\Person{奥利振}的\WorkTitle{论原理}一书的注释，在这些注释中，他称赞这些著作，说它们卓越非凡，而那些毁谤其作者、轻视其作品的人不过是吹毛求疵者。他说：\Quote{因为，他们缺乏足够的洞察力来理解那位非凡之人的深邃智慧。}那些希望公正评价\Person{狄迪穆}广博学识和强烈热情的人，必须仔细阅读他多样而精湛的著作。据说，在\Person{瓦伦斯}统治之前很久，当\Person{安当}从旷野来到\Place{亚历山大里亚}，为了对抗亚略派而与\Person{狄迪穆}交谈了一段时间后，察觉到此人的学识和智慧，就对他说：\Quote{狄迪穆，不要因你身体眼睛的失明而难过：因为剥夺你的不过是蚊蝇所共有的眼睛；反而要欢喜，因为你有天使所见的眼睛，借此可以辨识天主本身，并理解他的光明。}这位虔诚的\Person{安当}对\Person{狄迪穆}的这一番话是在我们所描述时代很早以前说的：事实上，\Person{狄迪穆}当时已被视为真信德的伟大堡垒，他回应亚略派，充分揭露他们诡辩性的吹毛求疵，胜利地驳斥了他们所有虚妄的微妙和欺骗性的推理。\par

% structure: p0067 source=h2 target=subsection reason=relative-heading-rank; base=h2

\subsection{第26章 \Place{该撒利亚}的\Person{巴西略}与\Place{纳西盎}的\Person{额我略}}

如今，天主眷顾在亚历山大城兴起迪迪默斯以对抗亚流派者。但为在其他城市驳斥他们，又兴起凯撒勒雅的巴西略和纳齐盎的额我略；在此处略述此二人合宜。确实，众所熟知的记忆足以作为他们名声的标记；其学识之广，于著作中清晰可见。然而，他们才华的运用对教会大有裨益，极大地维护了大公信仰，故本史之性质使我必须特别提及此二人。若有人将巴西略与额我略相比，思量各自之生活、德行与美德，则将难以判定应归卓越于谁：无论审视其行为之正直，抑或其对希腊文学与圣经之深谙，二人皆同样出众。青年时，他们同在雅典师从当时最负盛名的诡辩家希梅里奥斯与普罗埃雷西奥斯；随后，又赴叙利亚安提约基雅，就读于利巴尼奥斯之学府，于此精修修辞学至极致。既被认配作诡辩家，多友劝其从事教授雄辩之业；另有人劝其习法：然二人皆鄙弃此二途，放弃先前学业，投身隐修生活。他们从当时在安提约基雅执教者处略尝哲学之滋味，遂获取奥力振著作，从中汲取圣经正解；因当时奥力振声名远播天下。细心研读此伟人著作后，他们与亚流派者论战，明显占优。当亚流主义辩护者引同一作者以证实，如其所想，自家观点时，二人驳斥之，并明证对手全然不理解奥力振之推理。诚然，尽管当时亚流派之斗士欧诺米乌与许多他人被视为雄辩之士，但每试图与额我略及巴西略辩论时，相比之下显得无知与无学。巴西略被祝圣为执事后，由安提约基雅主教美利提乌斯从该品级擢升为卡帕多细亚凯撒勒雅主教，此乃其故乡。故他速赴彼处，恐亚流派教义已沾染本都诸省；为抵制之，他建立多所隐修院，勤勉教导百姓自家教义，坚固心神动摇者之信德。额我略被立为纳齐盎主教，此乃卡帕多细亚一小城，由其父亲管治，他追随巴西略相似之行径；遍历各城，坚固信心软弱者。尤频繁前往君士坦丁堡，藉其牧职安慰与坚固正统信徒，故不久后，因多位主教投票，他被选为君士坦丁堡教会主教。当此两位热忱奉献者之作为传至瓦伦斯皇帝耳中，他即刻命人从凯撒勒雅将巴西略解至安提约基雅；在总督法庭受审时，该官员问他\Quote{为何不接受皇帝的信仰？}巴西略大胆谴责其君所支持之信条的错误，并辩护同质之教义；总督以死威胁时，巴西略说：\Quote{但愿我能为真理脱离肉身的束缚。}总督劝其更严肃思量此事，据载巴西略说：\Quote{我今日是如何，明日也将如何；但我愿你未曾改变自己。}当时，巴西略当日被拘押。然不久后，皇帝之幼子加拉特患重病，医师已绝望；其母后多米尼加向皇帝保证，她梦中被可怖景象严重困扰，令她相信孩子之病是因虐待主教所受之惩罚。皇帝稍思后，召来巴西略，为试验其信仰，对他说：\Quote{若你所持之教义为真理，求主使我儿不死。}巴西略答：\Quote{若陛下信我所信，且教会合一，孩子必存活。}皇帝不同意此条件：巴西略说：\Quote{那么，天主的旨意将临于孩子。}巴西略言毕，皇帝命释放他；然孩子不久即亡。此乃此两位杰出教士历史之概要，二人均留下许多卓越著作，其中一些，鲁菲努斯说他已译为拉丁文。巴西略有两弟，彼得与额我略；前者效法巴西略之隐修生活；后者则效法其雄辩教学，并完成巴西略未竟之\WorkTitle{创世六日}论著。他还在君士坦丁堡发表了安提约基雅主教美利提乌斯的葬礼演说；他的许多其他演说亦仍存世。\par

% structure: p0069 source=h2 target=subsection reason=relative-heading-rank; base=h2

\subsection{第二十七章。论\Person{额我略}\OriginalTerm{显灵者}{Thaumaturgus}。}

但\Person{本都的额我略}与上述的额我略并非同一人，因名字相似及归于额我略名下的著作标题，容易使人混淆不同的人物。他是\Place{本都}的\Place{新凯撒勒雅}人，年代较上述额我略更早，因为他是\Person{奥力振}的门徒。这位额我略的名声在\Place{雅典}、\Place{贝里图斯}、整个\Place{本都}教区，几乎可说是全世界都闻名。他在\Place{雅典}的学校完成学业后，前往\Place{贝里图斯}研习民法，在那里听说\Person{奥力振}在\Place{凯撒勒雅}讲解圣经，便迅速前往；当他心窍开启，领悟到这些神圣经文的伟大之后，便告别了罗马法的进一步研习，从此与奥力振形影不离。从他那里获得了真正哲学的知识后，不久被父母召回，返回故乡；在那里，他虽仍是平信徒，却行了许多奇迹，医治病人，甚至凭书信驱魔，致使外教人因他的行为而信从，不亚于因他的讲论。殉道者\Person{庞斐禄}在为他辩护奥力振的著作中提到了此人；其中还附有额我略在与奥力振分离之际所写的一篇颂扬奥力振的告别演说。简而言之，当时有好几位额我略：第一位也是最古老的是奥力振的门徒；第二位是\Place{纳齐盎}的主教；第三位是\Person{巴西略}的兄弟；还有一位额我略，是亚略派在\Person{亚大纳削}被流放期间立为主教的。但关于这些人的谈论就到此为止。\par

% structure: p0071 source=h2 target=subsection reason=relative-heading-rank; base=h2

\subsection{第二十八章 论诺瓦杜及其追随者。夫黎基雅的诺瓦杜派随从犹太习俗更改守复活节的时间。}

大约此时，居住在\Place{夫黎基雅}的诺瓦杜派更改了庆祝复活节的日子。此事如何发生，我将在先说明他们在\Place{夫黎基雅}和\Place{帕夫拉哥尼亚}各省的教会直至今日仍保持严格纪律的原因之后，再加以陈述。\Person{诺瓦杜}是罗马教会的司铎，因\Person{科尔内略}主教接纳了在\Person{德基乌斯}皇帝迫害教会时曾献祭的信徒领受共融，而与教会分离。他为此退出后，被怀有相同意见的主教们擢升为主教，便致信所有教会，说：\Quote{不应让那些献过祭的人参与神圣奥迹；但要劝他们悔改，将赦免他们过犯的事交给天主，因为天主有能力赦免一切罪过。}收到这些信件后，各省中收到信件的各方按照各自的心意和判断行事。由于他要求不应让那些在领洗后犯有任何大罪的人领受圣事\BibleRef{若一 5:16}，这在有些人看来是残忍无情的做法；但另一些人则认为这条规矩是公正的，有利于维持纪律和促进更虔诚的生活。在这场争论激烈进行之时，\Person{科尔内略}主教的信件到达，允诺对领洗后的过失者宽大处理。既然这两个人写信互相矛盾，且各自援引天主圣言来证实自己的做法，那么正如通常所发生的那样，每个人都认同那有利于自己先前习惯和倾向的观点。那些乐于犯罪的人，因着当时给予他们的宽松，便借此机会纵情于各种邪恶行径。\Place{夫黎基雅}人似乎比其他民族更节制，很少犯发誓的罪。相反，\Place{斯基提亚}人和\Place{色雷斯}人生性非常易怒；而东方居民则沉溺于感官享乐。但\Place{帕夫拉哥尼亚}人和\Place{夫黎基雅}人对这两种恶习都不倾向；马戏和剧院的表演直到今日在他们中间也不受重视。正因如此，在我看来，这些人以及其他性格相同的人，很容易就同意了当时\Person{诺瓦杜}所写的信件。奸淫和通奸在他们看来是最大的恶行；众所周知，在地上没有哪个民族比\Place{夫黎基雅}人和\Place{帕夫拉哥尼亚}人在这一方面更能严格克制自己的情欲。我认为同样的理由也影响了那些住在西方、追随诺瓦杜的人。然而，尽管诺瓦杜为了更严格的纪律而成为分裂者，他在守复活节的时间上并未做任何改变，而是始终遵守西方教会所行的惯例。因为他们根据起初接受基督教时传下来的古老习惯，在春分后庆祝此节日。他本人后来在\Person{瓦莱里安}皇帝统治期间，在基督徒遭受迫害时殉道。但是，在\Place{夫黎基雅}那些以他命名的诺瓦杜派，大约在这一时期更改了庆祝复活节的日子，甚至在这件事上也不愿与其他基督徒共融。这是由该派别中几位不出名的主教在位于\Place{桑伽里乌斯}河源头附近的\Place{帕宗}村召开会议而实现的；因为在那里他们制定了一条教规，规定在犹太人每年守无酵节同一天庆祝。一位年长者，他是某位司铎的儿子，曾随父亲参加过这次会议，向我们提供了此事的消息。但是，君士坦丁堡诺瓦杜派的主教\Person{阿格里雅}、尼凯雅的\Person{玛克西穆}，以及尼苛默底雅和科提埃雍的主教都不在场，尽管诺瓦杜派的教务大部分由这些主教管理。诺瓦杜派教会因这次会议不久如何分裂为两派，将在适当的时候叙述；但现在我们必须注意大约同时发生在西方地区的事情。\par

% structure: p0073 source=h2 target=subsection reason=relative-heading-rank; base=h2

\subsection{第二十九章 达玛苏被祝圣为罗马主教。乌尔西努的竞争引发暴乱和丧命。}

当皇帝\Person{瓦伦提尼安}和平统治，不干涉任何教派时，\Person{达玛苏斯}在\Person{利贝里乌斯}之后接掌\Place{罗马}的主教职位；随后因以下原因引发了一场大骚乱。一位名叫\Person{乌尔西努斯}的该教会执事，在选举主教时与其他人一同被提名；由于\Person{达玛苏斯}被优先选中，这位\Person{乌尔西努斯}无法承受希望的破灭，便在教会之外举行分裂聚会，甚至诱使某些无甚名望的主教秘密祝圣了他。这次祝圣并非在教堂内，而是在一个名为\Place{西克纳宫}的隐秘场所举行，由此在民众中引起纷争；他们的分歧并非关乎任何信理或异端，而仅仅是谁该当主教。因此频繁发生冲突，以致在这场纷争中许多人丧生；许多神职人员和教友也因此受到城市长官\Person{马克西米努斯}的惩罚。这样，\Person{乌尔西努斯}当时被迫放弃他的要求，而那些有意追随他的人也被平息了下来。\par

% structure: p0075 source=h2 target=subsection reason=relative-heading-rank; base=h2

\subsection{第三十章 关于\Place{米兰}主教\Person{奥克森提乌斯}继任者的纷争。该省总督\Person{安布罗斯}前往平息骚乱，经民众一致同意并获皇帝\Person{瓦伦提尼安}批准，被选为该教会主教。}

大约在同一时间，\Place{米兰}发生了另一件值得记载的事件。\Person{奥克森提乌斯}去世后，他曾由亚略派人士祝圣为该教会主教，民众再次因选举继任者而骚动；因为有些人推举此人，另一些人则支持另一人，全城充满了争执和喧嚣。在此情况下，该省总督\Person{安布罗斯}——他也具有执政官身份——因担心民众的激动情绪引发灾难，便跑进教堂以平息骚乱。当他到达那里，民众安静下来后，他以一篇长篇而恰当的讲论压制了众人的非理性狂怒，提出他们认为是正确的动机，所有在场者突然达成一致，高喊\Quote{安布罗斯配得主教职位}，并要求祝圣他：\Quote{因为只有这样，}据称，\Quote{才能确保教会的和平，并使所有人在同一信仰和判断中重新合一。}由于民众如此的一致在在场主教们看来似乎是出于某种神圣的任命，他们立即按手在\Person{安布罗斯}身上；在为他施洗——因为他当时只是慕道者——之后，他们准备授予他主教职。但尽管\Person{安布罗斯}乐意接受了洗礼，他却极其认真地拒绝被祝圣；于是主教们将此事呈报给皇帝\Person{瓦伦提尼安}。这位皇帝视民众的一致同意为天主的作为，传话给主教们，要他们通过祝圣\Person{安布罗斯}来承行天主的旨意；宣称\Quote{他的拣选是来自天主的声音，而非人的投票。}\Person{安布罗斯}于是被祝圣；这样，彼此分裂的\Place{米兰}居民重新恢复了合一。\par

% structure: p0077 source=h2 target=subsection reason=relative-heading-rank; base=h2

\subsection{第三十一章 \Person{瓦伦提尼安}之死。}

此后，萨尔马特人入侵罗马领土，皇帝率大军进击。但蛮族了解到这次远征的可怕性质，便派使团来求和，提出某些条件。当使节被引见给皇帝，他觉得这些人不够体面，便问：\Quote{所有的萨尔马特人都像这样吗？}他们回答说，整个部族中最尊贵的人物已来到他面前，\Person{瓦伦提尼安}便极为震怒，大声喊道：\Quote{罗马帝国确实最不幸，竟在我执政之时，这样一个可鄙的蛮族，不满足于准许在其境内安全生存，竟敢拿起武器，入侵罗马领土，并发动公开战争。}他说这些话时的激烈态度和声音如此之大，以至于所有的血管都因用力而破裂，所有动脉都断裂；随之涌出大量鲜血，他便死了。此事发生在\Place{伯格提翁城堡}，时值\Person{格拉提安}第三次任执政官，与\Person{埃奎提乌斯}联合执政，十一月十七日，\Person{瓦伦提尼安}享年五十四岁，在位十三年。\Person{瓦伦提尼安}死后六天，意大利的军队在意大利城市\Place{阿钦库姆}宣布他的儿子\Person{小瓦伦提尼安}（当时尚幼）为皇帝。当此事告知另两位皇帝时，他们不悦，并非因为其中一位的兄弟和另一位的外甥被宣布为皇帝，而是因为军人擅自未经他们同意就宣布，他们本希望自己来宣布。不过，两人都批准了这一行动，于是\Person{小瓦伦提尼安}登上他父亲的宝座。这位\Person{小瓦伦提尼安}是\Person{尤斯蒂娜}所生，而\Person{尤斯蒂娜}是\Person{年长的瓦伦提尼安}在妻子\Person{塞维拉}尚在时娶的，情况如下：\Person{尤斯蒂娜}的父亲\Person{尤斯图斯}曾在\Person{君士坦提乌斯}统治下担任\Place{皮塞努姆}总督，他做了一个梦，似乎从自己右侧生出皇袍。这个梦告诉许多人后，最终传到\Person{君士坦提乌斯}耳中，他推测这是\Person{尤斯图斯}后裔将成为皇帝的预兆，便派人暗杀了他。\Person{尤斯蒂娜}因此失去父亲，仍是处女。一段时间后，她与皇帝\Person{瓦伦提尼安}的妻子\Person{塞维拉}结识，并经常与皇后交往，最后亲密到时常一起沐浴。当\Person{塞维拉}在浴中看到\Person{尤斯蒂娜}，她为这位处女的美貌所震惊，便向皇帝说起她；说：\Quote{尤斯图斯的女儿如此可爱，身材匀称，连她自己作为一个女人也完全被她迷住。}皇帝将妻子的描述记在心里，考虑如何能娶\Person{尤斯蒂娜}而不休掉\Person{塞维拉}，因为\Person{塞维拉}已为他生了\Person{格拉提安}，他不久前已立\Person{格拉提安}为奥古斯都。于是他制定了一项法律，并在所有城市颁布，允许任何人有两位合法妻子。法律颁布后，他娶了\Person{尤斯蒂娜}，与她生了\Person{小瓦伦提尼安}和三个女儿：\Person{尤斯塔}、\Person{格拉塔}和\Person{加拉}；前两个保持独身；但\Person{卡拉}后来嫁给了皇帝\Person{狄奥多西大帝}，\Person{狄奥多西}与她生了一个女儿名叫\Person{普拉西迪亚}。因为这位皇帝由前妻\Person{弗拉西拉}生了\Person{阿尔卡狄乌斯}和\Person{霍诺留斯}；不过我们将在适当的地方详述\Person{狄奥多西}和他的儿子们。\par

% structure: p0079 source=h2 target=subsection reason=relative-heading-rank; base=h2

\subsection{第三十二章 皇帝\Person{瓦伦斯}被哲学家\Person{特米斯提乌斯}的演说安抚，减轻了对基督徒的迫害。}

与此同时，\Person{瓦伦斯}居住在\Place{安提阿}，完全不受外战困扰；因为四方的蛮族都约束在自己边界内。然而，他向那些坚持\OriginalTerm{同质}{homoousian}教义的人发动了一场最残酷的战争，每天施加更严厉的惩罚，直到哲学家\Person{特米斯提乌斯}通过他的\WorkTitle{恳请演说}稍微缓和了他的严厉。在这次演说中，他告诉皇帝：\Quote{他不应对基督徒中存在的宗教判断分歧感到惊讶，因为这种差异与异教徒中流行的大量冲突意见相比是微不足道的；这些意见多达三百种以上；确实，争执是这种分歧的必然结果；但天主会因情感的多样性而更受荣耀，且其威严的伟大更受尊敬，因为认识祂并不容易。}哲学家说了这些及类似的话后，皇帝变得温和了一些，但并未完全放弃愤怒；因为虽然他不再处死教会人士，但继续将他们流放，直到这一愤怒也被以下事件所抑制。\par

% structure: p0081 source=h2 target=subsection reason=relative-heading-rank; base=h2

\subsection{第三十三章 哥特人在\Person{瓦伦斯}统治下接受基督教。}

住在大河对岸的被称为哥特人的蛮族，内部发生了内战，分裂为两派，一派由\Person{菲里提根}领导，另一派由\Person{阿塔那里克}领导。当\Person{阿塔那里克}明显占据优势时，\Person{菲里提根}向罗马人求助，恳求他们的帮助来对抗他的对手。这事被报告给\Person{瓦伦斯}皇帝，他命令驻守在\Place{色雷斯}的军队去帮助那些向他求助的蛮族，对抗他们更强大的同胞；凭借这一支援，他们在\Place{多瑙河}对岸彻底击败了\Person{阿塔那里克}，完全击溃了敌人。这成为许多蛮族皈依基督宗教的契机：因为\Person{菲里提根}为了表达对皇帝所赐恩惠的感激，接受了其恩人的宗教，并敦促他治下的人也这样做。因此，即使到现在，仍有如此多的哥特人沾染了亚略主义的错误，因为他们当时为了皇帝的缘故而选择信奉这一异端。\Person{乌尔菲拉}，当时他们的主教，发明了哥特字母，并将圣经翻译成他们的语言，着手用神圣的圣言教导这些蛮族。由于\Person{乌尔菲拉}不仅限于在\Person{菲里提根}的臣民中工作，还扩展到那些承认\Person{阿塔那里克}统治的人，\Person{阿塔那里克}认为这是对先祖宗教特权的侵犯，对宣认基督信仰的人施加严厉惩罚；因此，当时许多亚略派的哥特人成了殉道者。事实上，\Person{亚略}在试图驳斥利比亚人\Person{撒伯里}的观点时失败了，于是他偏离了真信仰，宣称天主圣子是\Quote{一个新的神？}：\BibleRef{申 32:7}。但是蛮族以更纯朴的心灵接受基督信仰，为了对基督的信德而轻视现世生命。我们就以此结束对归化哥特人的记述。\par

% structure: p0083 source=h2 target=subsection reason=relative-heading-rank; base=h2

\subsection{第34章　接纳逃亡的哥特人进入罗马领土，此事导致皇帝覆灭，并最终使罗马帝国灭亡。}

不久之后，蛮族彼此结成了友好联盟，但又被其他蛮族，即他们的邻居，称为匈人的蛮族击败；他们被赶出自己的国家，逃入罗马领土，表示愿臣服于皇帝，并执行他的一切命令。当\Person{瓦伦斯}得知此事，丝毫没有预见到后果，他命令应仁慈地接纳这些祈求者；仅此一次显出他的慈悲。因此，他将\Place{色雷斯}的某些区域划给他们，认为自己特别幸运：因为他预计今后他将拥有一支随时可用且装备精良的军队来对付一切袭击者；并希望蛮族将成为帝国边境比罗马人自己更强大的守卫者。因此，他后来忽视了通过罗马征兵来补充军队；并鄙视那些在先前战争中英勇奋战并制服敌人的老兵，他对各省居民按村庄习惯提供的民兵队伍估定了金钱价值，命令他的税吏要求每个士兵缴纳八十枚金币，尽管他此前从未减轻过公共负担。这一变化成为此后罗马帝国许多灾难的根源。\par

% structure: p0085 source=h2 target=subsection reason=relative-heading-rank; base=h2

\subsection{第35章　因与哥特人的战争，对基督徒的迫害减轻。}

蛮族占据\Place{色雷斯}后，安然享受该罗马省份，却不能适度承受好运，反而对恩人发起敌对攻击，蹂躏了整个\Place{色雷斯}及邻近地区。当这些事传到\Person{瓦伦斯}耳中，他停止了将\OriginalTerm{同质}{homoousion}派信徒流放的行动；并在极度惊恐中离开\Place{安提阿}，来到\Place{君士坦丁堡}，在那里，对正统基督徒的迫害也因同一原因而止息。同时，\Place{安提阿}的亚略派主教\Person{欧佐伊乌斯}去世，时值\Person{瓦伦斯}第五次任执政官、小\Person{瓦伦提尼安}第一次任执政官之年；\Person{多罗修}被任命接替他的位置。\par

% structure: p0087 source=h2 target=subsection reason=relative-heading-rank; base=h2

\subsection{第36章　萨拉森人在女王\Person{玛维亚}率领下接受基督信仰；虔诚的隐修士\Person{摩西}被祝圣为他们的主教。}

皇帝刚离开\Place{安提约基雅}，先前与罗马人结盟的撒拉逊人便在女王\Person{玛维亚}的率领下反叛，她的丈夫当时已去世。因此，整个东方地区当时都遭到撒拉逊人的劫掠。但某种神圣的天主眷顾以我即将描述的方式抑制了他们的狂暴。有一个名叫\Person{摩西}的撒拉逊人，在沙漠中过隐修生活，因其虔诚、信德和奇迹而极其著名。撒拉逊女王\Person{玛维亚}希望此人被立为她民族的主教，并以此条件许诺终止战争。罗马将领们认为以此条件达成的和平极为有利，便立即下令批准。于是\Person{摩西}被强行带走，从沙漠带到\Place{亚历山大里亚}，以便在那里接受主教职。但当他被呈献给当时管理该城教会的\Person{路济乌}时，他拒绝由他祝圣，并抗议说：\Quote{我确实认为自己不配此圣职；但若国家危急需要我担当，也不应由\Person{路济乌}按手在我头上，因他的手已沾满鲜血。}当\Person{路济乌}告诉他，他的责任是从他那里学习信仰原则，而不是口出侮辱之言时，\Person{摩西}回答说：\Quote{现在不是讨论信仰问题的时候；但你对弟兄们的恶行足以证明你的教义并非基督信仰。因为基督徒是\InnerQuote{不打人，不辱骂，不争斗}；因为\InnerQuote{上主的仆人不应当争斗。}\BibleRef{弟后 2:24}但你的行为通过那些被流放、被抛给野兽、被投入火焰的人向你发出控诉。我们亲眼所见之事远比我们从他人口中听闻的更令人信服。}当\Person{摩西}表达这些与类似的话时，他的朋友们带他上山，以便他从流放在那里的主教们手中接受祝圣。\Person{摩西}就这样被祝圣后，撒拉逊战争遂告终止。\Person{玛维亚}如此严格地遵守与罗马人订立的和约，甚至将她的女儿嫁给了罗马军队总司令\Person{维克多}。这就是有关撒拉逊人的事迹。\par

% structure: p0089 source=h2 target=subsection reason=relative-heading-rank; base=h2

\subsection{第37章 \Person{瓦伦斯}离开\Place{安提约基雅}后，\Place{亚历山大里亚}人驱逐\Person{路济乌}，恢复带着\Place{罗马}主教\Person{达玛稣}信函归来的\Person{伯多禄}。}

约在同一时期，皇帝\Person{瓦伦斯}一离开\Place{安提约基雅}，所有此前在各地遭受迫害的人便开始重新鼓起勇气，尤其是\Place{亚历山大里亚}的人。\Person{伯多禄}从\Place{罗马}带着\Person{达玛稣}（罗马主教）的信函回到该城，其中他确认了\OriginalTerm{同质}{homoousian}信仰，并批准了\Person{伯多禄}的祝圣。于是民众重拾信心，驱逐了\Person{路济乌}，后者立即乘船前往\Place{君士坦丁堡}。但\Person{伯多禄}在复职后不久即去世，临终前指定他的弟弟\Person{弟茂德}继承他。\par

% structure: p0091 source=h2 target=subsection reason=relative-heading-rank; base=h2

\subsection{第38章 皇帝\Person{瓦伦斯}因哥特人受到民众嘲弄；率军讨伐，在\Place{阿德里安堡}附近交战中阵亡。}

皇帝\Person{瓦伦斯}于五月三十日抵达\Place{君士坦丁堡}，时值他本人第六任执政官、\Person{小瓦伦提尼安}第二任执政官之年，他发现民众心情极为沮丧：因为蛮族已经蹂躏了\Place{色雷斯}，现在正在践踏\Place{君士坦丁堡}的近郊，附近没有足够的兵力抵抗他们。但当蛮族尝试逼近城墙时，民众变得极为不安，开始对皇帝怨声载道；指责他引来敌人，却懒散地拖延战事，而非立即出兵讨伐蛮族。此外，在赛马场举行的竞技中，众人异口同声地高呼反对皇帝忽视公共事务，急切地喊道：\Quote{给我们武器，我们自己战斗！}皇帝被这些骚动呼声激怒，于六月十一日出城；威胁说，若他返回，不仅因他们无礼的指责，还因他们先前支持篡位者\Person{普罗科皮乌}的图谋而惩罚市民；并宣称将彻底摧毁他们的城市，用犁耕平其废墟。他进军讨伐蛮族，以重大杀伤击溃他们，一直追击到\Place{色雷斯}城市\Place{阿德里安堡}，该城位于\Place{马其顿}边境。在那里再次与已经重新集结的敌人交战时，他于八月九日丧命，时值上述执政官之年，第289届奥林匹亚纪第四年。有些人断言他在退入的一个村庄中被烧死，该村庄遭蛮族袭击并纵火。但另一些人声称，他脱去皇袍冲入步兵主力之中；而当骑兵叛变拒绝交战时，步兵被蛮族包围，全军覆没。据称皇帝也在其中倒下，但因未穿皇袍而无法辨认。他死于五十岁，与兄长共同统治十三年，在兄长去世后又统治三年。因此本书包含十六年期间的事迹。\par

\SourceRecord{Translated by A.C. Zenos.；Nicene and Post-Nicene Fathers, Second Series；Vol. 2.；Edited by Philip Schaff and Henry Wace.；Buffalo, NY: Christian Literature Publishing Co.,；1890.；<http://www.newadvent.org/fathers/26014.htm>.}

% structure: p0001 source=h1 target=section reason=page-title

\section{教会史（第五卷）}

% structure: p0002 source=h2 target=subsection reason=relative-heading-rank; base=h2

\subsection{引言。}

在开始我们历史的第五卷之前，我们必须恳求那些可能阅读此论著的人，不要因我们立意撰写一部教会史，却仍在教会事务中穿插了可信的、对所述时期发生的战争的记述，而过于仓促地指责我们。我们这样做有几重原因：首先，是为了向读者呈现事实的准确陈述；其次，是唯恐读者因不断读到主教们之间的争议和彼此暗算而心生厌烦；但尤其是为了表明，每当国家事务动荡时，教会事务也仿佛因某种生命上的共鸣而陷入混乱。事实上，但凡仔细考察此事者，必会发现国家的灾祸与教会的纷扰密不可分地相连；因为他将看到，它们要么同时兴起，要么接踵而至。有时教会事务在先，随后国家动乱；有时则相反，故而我不信这种不断的交替仅是偶然，而确信它出自我们的不义；这些灾祸是作为应得的惩罚降于我们，诚如宗徒真实所说的：\BibleQuote{有些人的罪是明显的，如同先到审判案前；有些人的罪是随后跟去的。}\BibleRef{弟前 5:24}为此，我们在教会史中穿插了许多国家事务。关于君士坦丁统治时期进行的战争，我们未曾提及，因为未曾找到任何可靠的记载；但关于此后的事件，我们尽可能从当时尚存者那里收集到的、按发生顺序的信息，已作了简略的叙述。我们不断将皇帝纳入这些历史细节，因为自从他们开始信奉基督宗教以来，教会事务便有赖于他们，甚至连最重大的主教会议也曾且仍由他们召集而举行。最后，我们特别注意到了亚流派异端，因为它极大地搅扰了教会。以上所述作为序言，当已足够；现在我们继续我们的历史。\par

% structure: p0004 source=h2 target=subsection reason=relative-heading-rank; base=h2

\subsection{第一章 瓦伦斯死后，哥特人再次进攻君士坦丁堡，被市民及一些撒拉森援军击退。}

瓦伦斯皇帝以始终未查明的方式丧命后，蛮族再次逼近君士坦丁堡城墙，四面蹂躏郊区。民众愤慨，各自拿起手边可得的武器，自动出击迎敌。皇后多米尼卡从国库中向自愿出战者发放了与士兵通常所得的相同饷银。一些撒拉森人也作为同盟协助市民，他们是马维亚女王派来的；后者我们已提过。如此，民众在此次战斗中作战后，蛮族便退离城市很远。\par

% structure: p0006 source=h2 target=subsection reason=relative-heading-rank; base=h2

\subsection{第二章 格拉提安皇帝召回正统派主教，将异端者逐出教会。他立狄奥多西为帝国同僚。}

格拉提安当时与年幼的瓦伦提尼安共同拥有帝国，他谴责其叔父瓦伦斯对\OriginalTerm{正统派}{orthodox}基督徒的残酷政策，召回了被他流放的人。他还颁布法令，允许所有派别的人不加区别地安全聚集在各自教堂中；只有\OriginalTerm{欧诺米派}{Eunomians}、\OriginalTerm{福提诺派}{Photinians}和\OriginalTerm{摩尼教}{Manichæans}被排除在教堂之外。他同时意识到罗马帝国的衰落和蛮族势力的增长，看出国家需要一位勇敢而谨慎的人，便立狄奥多西为帝权同僚。这位\Person{狄奥多西}出身西班牙显贵家族，因战功卓著享有盛名，早在格拉提安选择他之前，已被普遍认为堪当皇帝。因此，在奥索尼乌斯和欧吕布里乌斯执政之年，于伊利里库姆城西尔米乌姆，一月十六日，格拉提安立他为皇帝，并与他分任对蛮族战争的管理之责。\par

% structure: p0008 source=h2 target=subsection reason=relative-heading-rank; base=h2

\subsection{第三章 当时的主要主教。}

此时，继利贝里之后治理罗马教会的是达玛稣。耶路撒冷教会仍由西里尔执掌。安提约基雅教会，如我们所述，一分为三：\OriginalTerm{亚流派}{Arians}已选多罗费为他们的主教欧佐伊的继任；其余人中，一部分在保利诺领导下，另一部分则归属于从流放中被召回的美利提。\Person{路济}虽被迫离开亚历山大里亚而缺席，却仍在该城亚流派中维持主教权威；那里的\OriginalTerm{同质论者}{Homoousians}由继伯多禄的弟茂德领导。在君士坦丁堡，继欧多克西的德默斐罗治理亚流派派系，并拥有教堂；但不愿与他共融者则另行集会。\par

% structure: p0010 source=h2 target=subsection reason=relative-heading-rank; base=h2

\subsection{第四章 曾认同\Quote{同质论}教义的马其顿派，复归前错。}

马其顿派派遣代表给\Person{利贝里乌斯}后，该派被允许与各城教会完全共融，与那些从一开始就接受\Place{尼西亚}所公布信条的人混杂在一起。但当\Person{格拉提安}皇帝的法律允许各派别在公共宗教仪式中不受限制地重新联合时，他们再次决定分离；并在\Place{叙利亚}的\Place{安提约基雅}聚集，决定再次避免使用\OriginalTerm{同质}{homoousios}一词，并且绝不与尼西亚信经的支持者共融。然而他们从这次尝试中一无所获；因为他们自己派系中的多数人对他们时而坚持一种意见、时而另一种的反复无常感到厌恶，离开了他们，从此成为那些承认\OriginalTerm{同质}{homoousion}教义者的坚定支持者。\par

% structure: p0012 source=h2 target=subsection reason=relative-heading-rank; base=h2

\subsection{第五章——与\Person{保利努斯}和\Person{梅利修斯}有关的\Place{安提约基雅}事件。}

大约此时，在\Place{叙利亚}的\Place{安提约基雅}因\Person{梅利修斯}爆发了一场严重冲突。我们已经注意到，该城主教\Person{保利努斯}因卓越的虔诚未被流放；而\Person{梅利修斯}在\Person{尤利安}恢复后，又被\Person{瓦伦斯}流放，最终在\Person{格拉提安}统治时期被召回。他回到\Place{安提约基雅}时，发现\Person{保利努斯}因年迈而非常衰弱；因此他的支持者立即竭尽全力让他与那位主教联合担任主教职务。当\Person{保利努斯}宣称\Quote{让一位由亚略派祝圣的人担任助理是违反教规的}时，民众诉诸暴力，使他在城外的一座教堂中被祝圣。此事发生后，引发了巨大骚乱；但后来民众在以下条件下得以团结。他们召集了被视为有资格候选主教的六位神职人员，其中一位是\Person{弗拉维安}。他们让所有这些人发誓，当两位主教中任何一位去世时，不努力使自己被祝圣，而是允许幸存者不受干扰地保留已故主教的教座。如此立下誓言后，民众获得了和平，不再彼此争吵。然而，路济弗尔派从其余人中分离出来，因为由亚略派祝圣的\Person{梅利修斯}被接纳为主教。在\Place{安提约基雅}教会的这种状况下，\Person{梅利修斯}不得不前往\Place{君士坦丁堡}。\par

% structure: p0014 source=h2 target=subsection reason=relative-heading-rank; base=h2

\subsection{第六章——\Person{纳西盎的格列高利}被调至\Place{君士坦丁堡}教区。\Person{狄奥多西}皇帝在战胜蛮族后于\Place{特撒洛尼迦}患病，由主教\Person{阿肖利乌斯}为其施洗。}

在许多主教的共同投票下，\Person{纳西盎的格列高利}此时从\Place{纳西盎}教区调至\Place{君士坦丁堡}教区，此事以前文所述方式发生。大约同时，皇帝\Person{格拉提安}和\Person{狄奥多西}各自战胜了蛮族。\Person{格拉提安}立即出发前往\Place{高卢}，因为阿勒曼尼人正在蹂躏那些行省；而\Person{狄奥多西}在立起胜利纪念碑后，急忙向\Place{君士坦丁堡}进发，并抵达\Place{特撒洛尼迦}。他在那里患了重病，并表示希望接受基督徒的洗礼。他自幼受祖先的基督教原则教导，并承认\OriginalTerm{同质}{homoousian}信仰。因此，随着病情加重，他越发渴望受洗，派人请来\Place{特撒洛尼迦}的主教，首先问他持何种教义观点。主教回答说：\Quote{\Person{亚略}的意见尚未侵入\Place{伊里利库姆}行省，那个异端所产生的新奇事物尚未开始侵害那些地区的教会；他们继续坚定不移地保存从起初由宗徒们传授、并在\Place{尼西亚}大公会议上确认的信仰。}于是皇帝非常喜乐地由主教\Person{阿肖利乌斯}施洗；病愈后不久，于\Person{格拉提安}第五次执政、\Person{狄奥多西}首次执政的同年十一月二十四日抵达\Place{君士坦丁堡}。\par

% structure: p0016 source=h2 target=subsection reason=relative-heading-rank; base=h2

\subsection{第七章——\Person{格列高利}发现对他的任命有不满，辞去\Place{君士坦丁堡}主教职务。皇帝命令亚略派主教\Person{德莫菲鲁斯}要么同意\OriginalTerm{同质}{Homoousion}，要么离开城市。他选择了后者。}

当时\Person{纳齐盎的额我略}在迁至\Place{君士坦丁堡}后，在城内一个小祈祷堂中举行集会，后来皇帝们在其旁建造了一座宏伟的教堂，命名为\Emph{阿纳斯塔西亚}。但额我略在口才和虔诚上远超他同时代的人，他知道有人因他是外地人而对他的升迁有所不满，在表达了对皇帝到来的喜悦之后，便辞去了\Place{君士坦丁堡}的主教职务。当皇帝发现教会处于这种状况时，他开始考虑如何实现和平、达成联合、并扩大教会。他立即向领导亚略派的\Person{德摩菲禄}表达了他的愿望，询问他是否愿意同意\WorkTitle{尼西亚信经}，从而重新联合人民，建立和平。\Person{德摩菲禄}拒绝接受这一提议后，皇帝对他说：\Quote{既然你拒绝和平与和谐，我命令你离开教堂。}\Person{德摩菲禄}听到这话，心中权衡与强大力量对抗的困难，便召集他的追随者在教堂中，站在他们中间说：\Quote{弟兄们，福音上记载\BibleRef{玛 10:23}：\InnerQuote{几时他们在这城迫害你们，你们就逃往另一城去。}因此，既然皇帝需要这些教堂，请注意，我们今后将在城外举行集会。}说完这话他就离开了，但并未正确理解福音作者那句话的意思，因为神圣谕言的真正含义是，那些想要避开今世道路的人必须寻求天上的\Place{耶路撒冷}。于是他出到城门外，在那里以后举行集会。\Person{路爵}也随他出去，他先前从\Place{亚历山大}被驱逐，如我们之前所述，逃到了\Place{君士坦丁堡}并住在那里。因此，亚略派在占据教堂四十年之后，由于反对\Person{德奥多西}皇帝提出的和平，在\Person{格拉提安}第五次执政、\Person{德奥多西·奥古斯都}第一次执政之年，十一月二十六日被赶出城。坚持\OriginalTerm{同质}{homoousian}信仰的人就这样重新获得了教堂。\par

% structure: p0018 source=h2 target=subsection reason=relative-heading-rank; base=h2

\subsection{第八章 由一百五十位主教组成的会议在君士坦丁堡召开。所颁布的法令。奈克塔略的祝圣。}

皇帝没有拖延，立即召集了他自己信仰的主教会议，以便确立\WorkTitle{尼西亚信经}，并任命\Place{君士坦丁堡}的主教；因为他并非没有希望将马其顿派争取到自己的观点，就也邀请了该派的首领出席。因此，此次与会的同质派有：来自\Place{亚历山大}的\Person{弟茂德}，来自\Place{耶路撒冷}的\Person{济利禄}——他当时已撤回先前意见，承认\Emph{\OriginalTerm{同质}{homoousion}}教义；来自\Place{安提约基雅}的\Person{默里济乌斯}，他先前已到达那里协助\Person{额我略}就职；来自\Place{得撒洛尼}的\Person{阿斯科利乌斯}，以及其他许多人，总数共一百五十位。马其顿派的首领是\Place{基齐库斯}的\Person{厄肋乌西乌斯}和\Place{兰普萨库斯}的\Person{玛尔奇安}；他们连同其余的人，大多来自\Place{赫勒斯滂}各城，共三十六人。会议在五月份，\Person{欧恰里乌斯}和\Person{厄瓦格利乌斯}执政年，开始举行。皇帝与持同样意见的主教们一道，竭尽全力争取\Person{厄肋乌西乌斯}及其追随者到自己的方面。皇帝提醒他们，曾由\Person{欧斯塔提乌斯}派遣代表团前往当时\Place{罗马}主教\Person{利贝里}那里；他们不久前曾自愿与正统派完全共融；并向他们指出他们行为的前后不一和摇摆不定——如今竟试图推翻他们曾经承认并与公教徒一致赞同的信仰。但他们对劝诫和责备都置之不理，宁愿坚持亚略派教义，也不赞同同质（homoousian）教义。发表这一声明后，他们离开了\Place{君士坦丁堡}；此外，他们写信给各城的同党，嘱咐他们决不可与\WorkTitle{尼西亚会议}的信经协调一致。另一派的主教留在\Place{君士坦丁堡}，商讨祝圣主教的事；因为\Person{额我略}，如前所述，已辞去该教座，准备返回\Place{纳齐盎}。当时有一个人名叫\Person{奈克塔略}，出身元老家族，性情温和良善，整个生活方式令人钦佩，尽管那时担任代行执政官之职。民众抓住他，选举他为主教，并由在场的一百五十位主教祝圣。这些主教还颁布了一项法令，规定\Quote{\Place{君士坦丁堡}的主教在\Place{罗马}主教之后应享有次等的尊荣，因为该城是新的\Place{罗马}。}他们再次确认了\WorkTitle{尼西亚信经}。此外，设立了宗主教，并划分了教区，以便没有一位主教能超越自己的教区对其他教会行使管辖权：因为之前由于迫害，这种事常常被随意进行。因此，\Person{奈克塔略}获得了大城和\Place{色雷斯}。\Person{赫拉狄乌斯}，\Person{巴西略}在\Place{卡帕多细亚}的\Place{凯撒勒雅}主教区的继任者，与\Person{巴西略}的兄弟、\Place{卡帕多细亚}的\Place{尼撒}主教\Person{尼撒的额我略}，以及\Place{亚美尼亚}的\Place{默利提纳}主教\Person{欧特雷乌斯}，共同获得了\Place{本都}教区的宗主教区。\Place{依科尼雍}的\Person{盎博罗削}和\Place{丕息狄雅}的\Place{安提约基雅}的\Person{奥普提姆斯}被分配了\Place{亚细亚}教区。整个\Place{埃及}教会的监督权委托给了\Place{亚历山大}的\Person{弟茂德}。\Place{劳狄刻雅}的\Person{贝拉基}和\Place{塔尔索}的\Person{狄奥多若}承担了东方教会的管理职责；但这并不损害保留给\Place{安提约基雅}教会并授予当时在场的\Person{默里济乌斯}的尊荣特权。他们还规定，鉴于需要，各省的教会事务应由该省的主教会议处理。这些安排得到皇帝的认可。这就是本次会议的结果。\par

% structure: p0020 source=h2 target=subsection reason=relative-heading-rank; base=h2

\subsection{第九章 \Place{君士坦丁堡}主教\Person{保禄}的遗体被隆重地从流放地迁回。\Person{默肋提}之死。}

皇帝当时下令将\Person{保禄}主教的遗体从\Place{安奇拉}城迁走。这位\Person{保禄}主教是禁卫军长官\Person{斐理伯}在\Person{马其顿尼}的唆使下流放，并下令在\Place{亚美尼亚}的一个小镇\Place{库库苏斯}将其勒死的，如我先前所述。因此，皇帝极其恭敬地接收了遗骸，并将其安放在如今以\Person{保禄}命名的教堂里。该教堂先前由马其顿派占据，当时他们与亚略派保持分离，但因拒绝采纳皇帝的观点而被皇帝驱逐。大约在此期间，\Place{安提约基雅}主教\Person{默肋提}患病去世。\Person{巴西略}的兄弟\Person{额我略}为他致了葬礼悼词。已故主教的遗体被他的朋友们运往\Place{安提约基雅}，那里那些认同\Person{默肋提}利益的人再次拒绝服从\Person{保利诺}，反而让\Person{弗拉维安}取代\Person{默肋提}，民众又开始争吵。于是，\Place{安提约基雅}教会再次分裂为敌对派别，并非基于信仰上的任何差异，而仅仅是对主教人选的不同偏好。\par

% structure: p0022 source=h2 target=subsection reason=relative-heading-rank; base=h2

\subsection{第十章 皇帝下令召集所有不同派别的大会。\Person{阿尔卡狄乌斯}被宣布为奥古斯都。\OriginalTerm{诺瓦替安派}{Novatians}被允许在\Place{君士坦丁堡}城内举行集会；其他异端者被驱逐。}

在其他城市也发生了大骚动，因为亚略派信徒被逐出教堂。但我无法充分赞叹皇帝在此危急关头所表现出的审慎。因为他尽可能不愿使城市充满纷扰，于是不久之后，他便召集了一次各教派的全体会议，认为通过各主教之间的讨论，他们相互间的分歧或许可以化解，并达成一致。我相信，皇帝的这一意图正是当时他的事务如此兴旺的原因。事实上，由于天主上智的特殊安排，野蛮民族都臣服于他：其中，哥特人的国王\Person{阿塔那里克}自愿向他降服，连同他所有的子民，不久后便在\Place{君士坦丁堡}去世。在此关头，皇帝于一月十六日，在\Person{梅罗保德}和\Person{萨图尼鲁斯}第二次执政官任期内，宣布他的儿子\Person{阿卡迪乌斯}为奥古斯都。不久之后，在同年六月，同一执政官任期内，各教派的主教从各地抵达；于是皇帝召见主教\Person{内克塔里乌斯}，与他商议如何以最佳方式使基督信仰摆脱纷争，并使教会归于合一。\Quote{争论的问题，}他说，\Quote{应当公平讨论，以便通过发现并消除不和的根源，达成普遍的一致。}听到这个提议，\Person{内克塔里乌斯}感到不安，便告知了诺瓦替派的主教\Person{阿格利乌}，因为\Person{阿格利乌}在信仰问题上与他持相同意见。此人虽然十分虔诚，却完全不能胜任教义问题的辩论；因此他提议将此事委托给他的诵读员\Person{西西尼乌斯}，认为他是主持辩论的合适人选。\Person{西西尼乌斯}不仅博学，且经验丰富，精通圣经诠释和哲学原理。他深信辩论非但不能治愈分裂，反而通常会滋生更顽固的异端，因此他深知古代作者从未将存在的开端归于天主子，而是认为祂与圣父同永，于是向\Person{内克塔里乌斯}提出以下建议：他们应避免辩证法的论战，而提出古代作者的见证作为真理的证据。\Quote{让皇帝询问各教派的首领，}他说，\Quote{他们是否尊重那些在分裂扰乱教会之前活跃的古代作者，还是将他们视为与基督信仰疏远而予以拒绝？如果他们拒绝这些作者的权威，那么就叫他们也谴责他们：如果他们胆敢采取这一步，他们自己就会立刻被民众赶走，这样真理将明显得胜。但另一方面，如果他们不愿弃绝教父们，那么我们就应当出示他们的著作，这些著作将充分证实我们的观点。}\Person{内克塔里乌斯}听了\Person{西西尼乌斯}的话，急忙前往皇宫，向皇帝报告了向他提出的计划；皇帝立刻看出了其中的智慧与恰当，便以极大的审慎予以执行。因为他不暴露自己的目的，只是简单地问异端首领们，他们是否尊重并接受那些在教会分裂之前生活的导师的教训？他们没有否认，反而回答说他们极其尊崇他们，视他们为师长；皇帝于是又追问他们，是否会将这些导师视为基督教义的可靠见证而服从他们？面对这个问题，各派的首领连同他们的逻辑辩手——因为许多人已准备好进行诡辩——感到极为尴尬。因为他们中产生了分歧：一些人认可皇帝提议的合理性，另一些人则退缩，意识到这对他们极为不利：因此所有人在对待古代著作的态度上各不相同，他们彼此之间无法再达成一致，不仅与其他教派意见相左，甚至同派之间也彼此不同。因此，一致的恶意，如同古时巨人的舌头，被混乱了，他们的邪恶之塔被倾覆。皇帝从他们的混乱中看出，他们唯一的信心在于巧辩，他们害怕诉诸教父们的诠释，于是采取了另一种方法：他命令每个教派以书面形式陈述自己特有的教义。于是，他们中被认为最有技巧的那些人，草拟了一份各自信经的陈述，措辞极为谨慎；指定了一日，为此选出的主教出现在皇宫。\Person{内克塔里乌斯}和\Person{阿格利乌}作为\OriginalTerm{同质}{homoousian}信仰的辩护者出现；\Person{德摩菲卢斯}支持亚略派的教义；\Person{欧诺弥乌斯}本人代表欧诺弥派；\Place{基齐库斯}的主教\Person{厄琉息乌}代表了被称为马其顿派的人的意见。皇帝礼貌地接待了他们所有人；收到每人书面的信仰声明后，他独自闭门，热切祈祷，求天主帮助他努力辨明真理。然后，他仔细审阅了每份提交给他的声明，谴责了其余所有声明，因为它们引入了对圣三的分割，只批准了包含\Emph{\OriginalTerm{同质}{homoousion}}教义的那一份。这个决定使诺瓦替派再次兴旺，并在城内举行集会：因为皇帝对他们的信仰声明与他所采纳的一致感到欣喜，颁布了一项法律，确保他们和平拥有自己的教堂建筑，并授予他们的教会与他特别认可的那些教会同等的特权。但其他教派的主教们，由于他们彼此意见不合，甚至被自己的追随者轻视和指责：因此他们满心困惑和恼怒，离开了，给他们的拥护者写了安慰信，劝诫他们不要烦恼，因为许多人离开了他们，投向了\OriginalTerm{同质}{homoousian}派；因为他们说\BibleQuote{蒙召的人多，被选的人少}\BibleRef{玛 20:16}——这句话他们从未在因暴力和恐怖而使民众多数站在他们一边时使用过。尽管如此，正统信徒并非完全免于不安；因为\Place{安提约基雅}教会的事务在出席大公会议的人中引起了分裂。\Place{埃及}、\Place{阿拉伯}和\Place{塞浦路斯}的主教们联合反对\Person{弗拉维安}，坚持将他逐出\Place{安提约基雅}；但\Place{巴勒斯坦}、\Place{腓尼基}和\Place{叙利亚}的主教们同样热心地为他争辩。这场争论的结果如何，我将在适当的地方叙述。\par

% structure: p0024 source=h2 target=subsection reason=relative-heading-rank; base=h2

\subsection{第十一章　僭主玛克西慕斯设奸谋杀害皇帝格拉提安；尤斯蒂娜因此恐惧，停止迫害盎博罗削。}

几乎在举行这些君士坦丁堡会议的同时，西方发生了以下事件。\Person{玛克西慕斯}从不列颠岛起兵反叛罗马帝国，袭击了当时因与阿勒曼尼人作战而疲惫不堪的\Person{格拉提安}。在意大利，\Person{瓦伦提尼安}尚年幼，\Person{普罗布斯}（一位曾任执政官的人）执掌政务，当时担任禁卫军长官。年轻皇帝的母亲\Person{尤斯蒂娜}怀有阿里乌派思想，在其丈夫在世时未能加害于同质派；但她趁儿子年幼前往米兰时，对主教\Person{盎博罗削}表现出极大敌意，并下令将他流放。由于民众对\Person{盎博罗削}怀有极其深厚的感情，因而抵抗那些奉命将他驱逐出境的人。此时传来消息说，\Person{格拉提安}已被僭主\Person{玛克西慕斯}设奸谋杀害。事实上，\Person{玛克西慕斯}麾下将领\Person{安德拉加提乌斯}藏在一顶形似卧榻的轿子里，由骡子抬着，并命令他的卫兵在他面前散布谣言，说轿中载有皇帝\Person{格拉提安}的妻子。他们在法兰西的里昂城附近与皇帝相遇，当时皇帝刚过河；皇帝以为那是自己的妻子，未怀疑有诈，便如盲人堕入坑中一般落入敌手；因为\Person{安德拉加提乌斯}突然从轿中跃出，将他杀死。\Person{格拉提安}遂在\Person{梅罗高德斯}与\Person{撒图尔尼努斯}任执政官之年遇害，时年二十四岁，在位十五年。此事发生后，皇后\Person{尤斯蒂娜}对\Person{盎博罗削}的怒气得以平息。后来\Person{瓦伦提尼安}尽管极不情愿，但因形势所迫，接纳\Person{玛克西慕斯}为帝国同僚。\Person{普罗布斯}对\Person{玛克西慕斯}的权势感到惊恐，决定退往东方地区；于是离开意大利，前往伊利里库姆，定居于马其顿城市得撒洛尼。\par

% structure: p0026 source=h2 target=subsection reason=relative-heading-rank; base=h2

\subsection{第十二章　皇帝德奥多西备战对付玛克西慕斯时，其子欧诺留出生；随后他前往米兰以迎战僭主。}

然而，皇帝\Person{德奥多西}极为忧虑，并征召了一支强大的军队对付僭主，唯恐他图谋杀害年幼的\Person{瓦伦提尼安}。在筹备期间，波斯使团抵达，请求与皇帝议和。与此同时，皇后\Person{弗拉基拉}为他生下一子，取名\Person{欧诺留}，时在\Person{里科梅卢斯}与\Person{克莱阿尔库斯}任执政官之年的九月九日。同一年及稍早时，诺瓦提安派主教\Person{阿格利乌斯}去世。次年，当\Person{阿卡狄}\Person{奥古斯都}与\Person{包东}共任首届执政官时，亚历山大里亚主教\Person{弟茂德}去世，\Person{德奥斐卢斯}继任主教之职。约一年后，阿里乌派主教\Person{德莫斐卢斯}离世，阿里乌派人从色雷斯召来他们异端的一位首领\Person{马里努斯}，将主教职位托付给他；但\Person{马里努斯}并未久居此位，因为在他领导下，该派分裂为两派，正如我们以后将要说明的；因为他们邀请\Person{多若瑟}从叙利亚的安提约基雅前来，并立他为自己的主教。与此同时，皇帝\Person{德奥多西}前往进行对抗\Person{玛克西慕斯}的战争，将儿子\Person{阿卡狄}以皇权留于君士坦丁堡。因此，他抵达得撒洛尼时，发现\Person{瓦伦提尼安}及其随从十分焦虑，因为他们被迫承认了僭主为皇帝。然而\Person{德奥多西}并未公开表达自己的看法；他既未拒绝也未接受\Person{玛克西慕斯}的使团；只是不能容忍僭称帝号者统治罗马帝国的暴政，于是迅速集结兵力，进军至米兰，僭主已先期到达那里。\par

% structure: p0028 source=h2 target=subsection reason=relative-heading-rank; base=h2

\subsection{第十三章　阿里乌派在君士坦丁堡引发骚乱。}

当皇帝忙于军事远征之时，阿里乌派用以下手段在君士坦丁堡引起了极大的骚乱。人们喜欢编造关于自己所不知之事的说法；一旦有机会，他们便把关于自己所期望之事的小道消息极度扩大，总是喜好变化。这在当时君士坦丁堡表现得很突出：每个人都按自己心血来潮，编造关于远方战争的新闻，并总以最灾难性的结果为前提；而且在战争尚未开始时，他们便以如同亲临现场那般自信的语气谈论自己一无所知的事。于是有人肯定地说：\Quote{僭主已击败了皇帝的军队}，甚至具体说明了双方阵亡人数；还说\Quote{皇帝本人几乎落入僭主之手}。然后，阿里乌派因那些曾受他们迫害的人获得了城内教堂而感到极度愤慨，便开始在这些谣言中添油加醋。但由于这样的故事日复一日地流传，夸大其词，最终连农民本身也相信了——那些凭传闻散布谣言的人向谎言的制造者证实，他们从后者那里听到的消息已在别处得到了充分证实——这时阿里乌派便敢于采取暴力行为，其中一项暴行是纵火焚烧主教\Person{内克塔里乌斯}的住所。这事发生在\Person{德奥多西}\Person{奥古斯都}与\Person{西内吉乌斯}共任的第二任执政官期间。\par

% structure: p0030 source=h2 target=subsection reason=relative-heading-rank; base=h2

\subsection{第十四章　僭主玛克西慕斯的覆灭与死亡。}

当皇帝向篡位者进军时，关于他所作可怕准备的传闻使\OriginalTerm{马克西穆斯}{Maximus}的军队如此惊恐，以致他们不为他作战，反而将他捆绑交给皇帝；皇帝于同年执政官任期下，在八月二十七日将他处死。亲手杀死\OriginalTerm{格拉提安}{Gratian}的\OriginalTerm{安德拉加提乌斯}{Andragathius}得知马克西穆斯的命运后，纵身投入附近的河流淹死。随后，得胜的皇帝们由\OriginalTerm{狄奥多西}{Theodosius}之子、尚是幼童的\OriginalTerm{霍诺留}{Honorius}陪同，进入罗马城；霍诺留是在马克西穆斯被击败后，其父立即从\OriginalTerm{君士坦丁堡}{Constantinople}召来的。他们在\OriginalTerm{罗马}{Rome}停留，举行凯旋庆典；在此期间，狄奥多西皇帝对\OriginalTerm{叙马库斯}{Symmachus}——一位曾任执政官并担任罗马元老院首领的人——表现出了非凡的宽仁。这位叙马库斯以辩才著称，他的许多演说辞至今仍存，以拉丁语写成；但因他曾为马克西穆斯撰写颂词并当众宣读，事后被控叛国罪；为免死刑，他躲进教堂寻求庇护。皇帝出于对宗教的敬意，不仅尊敬自己教会的众主教，也善待持'同质'信经（\OriginalTerm{同质}{homoousian}）的\OriginalTerm{诺瓦替雅奴派}{Novatians}主教；因此，为了取悦代表叙马库斯求情的罗马诺瓦替雅奴派主教\OriginalTerm{莱昂提乌斯}{Leontius}，皇帝仁慈地赦免了叙马库斯的罪行。叙马库斯获赦后，致信狄奥多西皇帝为自己辩护。如此，这场起初看来十分可怕的战争便迅速结束了。\par

% structure: p0032 source=h2 target=subsection reason=relative-heading-rank; base=h2

\subsection{第15章 论安提约基雅主教\OriginalTerm{弗拉维亚诺}{Flavian}。}

大约同一时期，叙利亚的\OriginalTerm{安提约基雅}{Antioch}发生了以下事件。\OriginalTerm{保利诺}{Paulinus}去世后，原受他管辖的人们拒绝服从弗拉维亚诺的权威，却另立\OriginalTerm{埃瓦格里乌斯}{Evagrius}为他们一方的主教。埃瓦格里乌斯祝圣后不久便去世了，由于弗拉维亚诺从中干预，无人接替他；然而，那些因弗拉维亚诺违背誓言而厌恶他的人仍自行集会。与此同时，弗拉维亚诺'不遗余力'——如谚语所说——也要将这些人置于控制之下；不久后他果然做到了，因为他平息了当时亚历山大里亚主教\OriginalTerm{泰奥菲罗}{Theophilus}的怒气，并通过泰奥菲罗的调解也安抚了罗马主教\OriginalTerm{达马苏斯}{Damasus}。因为这两位都对弗拉维亚诺极为不满，既因他所犯的伪誓罪，也因他在原本合一的人中引起的分裂。于是泰奥菲罗被安抚后，派遣司铎\OriginalTerm{依西多禄}{Isidore}前往\OriginalTerm{罗马}{Rome}，说服仍怀怨的达马苏斯；向他指出，为了促成人民之间的和睦，理当忽略弗拉维亚诺以往的不当行为。这样，弗拉维亚诺恢复了共融，安提约基雅的人们不久后也接受了他所带来的合一。安提约基雅的此事就此了结。但该城的亚略人被逐出教堂后，习惯在郊区举行集会。与此同时，耶路撒冷主教\OriginalTerm{济利禄}{Cyril}大约在此时去世，由\OriginalTerm{若望}{John}继任。\par

% structure: p0034 source=h2 target=subsection reason=relative-heading-rank; base=h2

\subsection{第16章 \OriginalTerm{亚历山大里亚}{Alexandria}偶像庙宇的拆毁，及随之而来的外教徒与基督徒之间的冲突。}

应亚历山大主教狄奥斐卢斯之请，皇帝此时下令拆毁该城异教神庙，并命在狄奥斐卢斯指导下执行。狄奥斐卢斯抓住此机，竭力羞辱异教秘仪。他首先使人清理密特拉神庙，公开展示其血腥秘仪的标记。随后他摧毁了塞拉比斯神庙，并公开嘲弄密特拉神庙的血腥仪式；他还展示塞拉比斯神庙充满了荒诞的迷信，并将普里阿普斯的阳物像抬过广场中央。亚历山大的异教徒，尤其是哲学教授们，因这一曝光而怒不可遏，其报复之凶残超过了以前的暴行：因为他们同声一气，按约定信号，猛烈冲向基督徒，杀害所有能抓住的人。基督徒也试图抵抗袭击者，于是祸害更加扩大。这场殊死混战持续到杀得尽兴才告结束。事后发现异教徒被杀者甚少，而基督徒死伤众多；双方受伤者几乎不计其数。异教徒因所作所为而恐惧，担心皇帝不悦。他们随心所欲行事，以流血浇灭了勇气后，有的逃向这里，有的逃向那里，许多离开亚历山大，分散到各城。其中有两位文法家赫拉狄乌斯与阿摩尼乌斯，我年轻时在君士坦丁堡曾受教于他们。据说赫拉狄乌斯是朱庇特的司祭，阿摩尼乌斯是赛米乌斯的司祭。这场骚乱平息后，亚历山大总督与埃及军队统帅协助狄奥斐卢斯拆毁异教神庙。于是这些神庙被夷为平地，其神像被熔铸成锅及其他便利器皿，供亚历山大教会使用；因为皇帝曾指示狄奥斐卢斯将其分给穷人。所有神像都被打碎，唯独一座前述神祇的雕像，狄奥斐卢斯保留下来，立于公共场所；他说：\Quote{免得日后异教徒否认他们曾崇拜过这样的神。}此举尤其令文法家阿摩尼乌斯大为不快，据我所知，他常说：\Quote{外邦宗教遭到严重侮辱，因为这一尊雕像没有被熔化，反而被保留下来，以嘲笑那宗教。}然而赫拉狄乌斯曾在一些人面前夸口说，他在那次猛烈攻击中亲手杀死了九个人。当时在亚历山大的情况便是如此。\par

% structure: p0036 source=h2 target=subsection reason=relative-heading-rank; base=h2

\subsection{第十七章　论在塞拉比斯神庙中发现的象形文字。}

当塞拉比斯神庙被拆毁并暴露时，其中发现刻在石头上的某些文字，他们称之为象形文字，形状为十字架。基督徒与异教徒看见后，各自将其纳入并应用于自己的宗教：因为断言十字架是基督救难标记的基督徒，声称这文字专属他们；而异教徒则声称它可能同时属于基督与塞拉比斯；\InnerQuote{因为，}他们说，\InnerQuote{它对基督徒意味着一种东西，对异教徒又意味着另一种。}当双方为此争论时，一些改信基督教的异教徒，通晓这些象形文字，解释十字架的形状，说它意味着\Quote{未来的生命}。基督徒欣然抓住这一点，认为对其宗教极为有利。但后来破译了其他象形文字，其中包含预言：\Quote{当十字架出现时——即指\InnerQuote{未来的生命}——塞拉比斯神庙必将被毁}，于是大批异教徒接受了基督教，承认自己的罪并接受了洗礼。关于发现这个十字架形状的符号，我所听到的报告即是如此。但我不能想象埃及司祭在刻十字架图形时预知了关于基督的事。因为\Quote{我们的救主降临世界，是历世历代所隐藏的奥秘}，如宗徒所宣告；连邪恶之首魔鬼也一无所知，其仆役埃及司祭更不可能知晓；但天主眷顾无疑旨意，在这符号的探究中，会发生类似当初保禄宣讲时发生的事。他因圣神而智慧，对雅典人采用了类似的方法，\BibleRef{宗 17:23}，读了一座祭台上的题词，将其运用于自己的讲论，使许多人归信。除非有人要说，天主圣言曾在埃及司祭内工作，如同在巴郎\EditorialNote{户 24}与盖法身上一样；\BibleRef{若 11:51}，因这两人不由自主地预言了善事。关于此事，就此足矣。\par

% structure: p0038 source=h2 target=subsection reason=relative-heading-rank; base=h2

\subsection{第十八章　狄奥多西皇帝改革罗马的弊端。}

皇帝狄奥多西在意大利短暂停留期间，通过赏赐与废除两种方式，为罗马城带来了最大的恩惠。他的赏赐确实非常慷慨；并且他清除了城中两个最恶劣的弊政。其一如下：在古罗马，往昔曾建有规模宏大的建筑，用以制作面包分发给民众。负责管理这些建筑的人，拉丁语称为\OriginalTerm{曼契佩斯}{Mancipes}，日久天长，竟将它们变成了盗贼的巢穴。由于这些建筑中的烤房设在地下，他们便在每座烤房旁开设酒馆，豢养娼妓；以此诱捕许多前去歇息或满足淫欲之人，因为通过某种机械机关，他们能将人从酒馆摔入下方的烤房。这种勾当主要针对外乡人；如此被绑架的人被迫在烤房中劳作，其中许多人被囚禁至老，不得外出，使其亲友误以为他们已经去世。恰巧狄奥多西皇帝的一名士兵落入了这个陷阱；他被关在烤房中，无法出去，便拔出随身携带的匕首，杀死了阻拦他的人；其余人恐惧万分，便放他逃脱。皇帝得知此事后，惩罚了那些\OriginalTerm{曼契佩斯}{Mancipes}，并下令拆毁这些不法之徒的巢穴。这是皇帝为皇城清除的丑恶弊政之一；另一件性质如下：当一名女子被控通奸时，他们处罚犯法者的方式非但不是纠正，反而加重其罪行。他们将女子关进狭小的妓院，迫使其以极其令人作呕的方式卖淫；并在这污秽行为发生时摇响小铃，使路人皆知内中正在发生何事。这无疑是为了在公众舆论中给罪行打上更大的耻辱烙印。皇帝一旦得知这一不雅习俗，绝不容忍；他下令拆毁这些称为\OriginalTerm{西斯特拉}{Sistra}的惩罚性卖淫场所，并制定其他法律惩罚通奸妇女。就这样，狄奥多西皇帝使城市摆脱了两种最可耻的弊政；当他满意地安排完其他所有事务后，便将瓦伦提尼安皇帝留在罗马，自己与儿子何诺留返回君士坦丁堡，于塔提安和叙马库斯任执政官之年的11月10日进入该城。\par

% structure: p0040 source=h2 target=subsection reason=relative-heading-rank; base=h2

\subsection{第十九章　论补赎司铎的职务及其废除}

此时，人们认为有必要废除那些在教会中负责补赎事宜的司铎的职务；此事缘由如下。当诺瓦替安派因不愿与德西乌迫害期间背教者共融而脱离教会时，主教们便在教会法规中增设了一位补赎司铎，以便那些在领洗后犯罪的人向这位指定的司铎告明他们的罪过。这种纪律在其他所有异端派别中仍保留为他们异端制度的一部分；只有同质派以及持有相同教义观点的诺瓦替安派废除了它。后者其实从未允许设立此职；而如今掌管教会的同质派，在维持此职相当长一段时间后，于奈克塔留时代，因君士坦丁堡教会发生的一件事而废除了它，事件如下：一位贵族妇女来到补赎司铎处，全面告明了她领洗后所犯的罪过；司铎命令她持续守斋祈祷，以便连同承认过错，她也要显出与悔改相称的作为。过了一段时间，这位夫人再次前来，告明她犯了另一桩罪行，教会的一位执事曾与她同寝。此事得到证实后，执事被逐出教会；但民众非常愤怒，不仅因所发生的事而反感，还因这事给教会带来了丑闻和贬抑。因此，当神职人员遭受嘲讽和指责时，一位原籍亚历山大里亚的教会司铎\OriginalTerm{犹得蒙}{Eudæmon}劝说主教\OriginalTerm{奈克塔留}{Nectarius}废除补赎司铎的职务，而让人人在领受神圣奥迹一事上各凭良心；因为依他之见，唯有这样教会才能免受毁谤。听了犹得蒙对此事的解释后，我冒昧地将其写入本书；因为正如我常说的，我不遗余力地从最知情者那里获取真实记载，并审慎查考每项报告，以免陈述不实之事。当犹得蒙首次叙述此事时，我对他说的是：\Quote{司铎啊，你的建议对教会是否有益，天主知道；但我看到它取消了互相责备过失的途径，并且使我们不能遵行宗徒的这句训诫：\BibleQuote{不要与黑暗无益的工作有分，倒要指摘它们。}}关于此事，论述至此即可。\par

% structure: p0042 source=h2 target=subsection reason=relative-heading-rank; base=h2

\subsection{第二十章　亚流派及其他异端中的分裂}

此外，我认为有必要不略过其他宗教团体的行径，即\OriginalTerm{亚略派}{Arians}、\OriginalTerm{诺瓦替安派}{Novatians}，以及那些从玛塞督尼（\Person{Macedonius}）和欧诺米（\Person{Eunomius}）得名的派别。因为教会一旦分裂，便不止于那次分裂，而分离者又因最轻微、最琐碎的借口彼此不和。其方式、时间以及他们挑起彼此纷争的原因，我们将在行文中陈述。但于此应注意到，皇帝德奥多西（\Person{Theodosius}）除了欧诺米之外，并未迫害他们中的任何人；然而，由于后者在君士坦丁堡的私宅中集会，宣读他所撰写的著作，以其教义败坏了许多人，皇帝便下令将他流放。对于其他异端者，他没有干预任何人，也没有强迫他们与自己共融；相反地，他允许所有人在自己的集会处聚集，并就基督信仰之点持有自己的意见。至于其他人，准许他们在城外建造教堂；但鉴于诺瓦替安派在信仰上与他持完全相同的观点，他下令他们应如我先前所述，在城内教堂中不受干扰地继续存在。关于这些，我认为在此处稍作进一步记述是合宜的，因此将追溯其历史中的若干情况。\par

% structure: p0044 source=h2 target=subsection reason=relative-heading-rank; base=h2

\subsection{第二十一章 诺瓦替安派特有的分裂}

君士坦丁堡的诺瓦替安教会，阿格里乌斯（\Person{Agelius}）担任主教四十年，即从君士坦丁在位至德奥多西皇帝第六年，如我先前在某处所陈述。他见自己大限将至，便祝圣西辛尼乌（\Person{Sisinnius}）接任主教职。此人原是阿格里乌斯所辖教会的一位司铎，口才极佳，曾与尤利安皇帝同时受教于玛克西穆（\Person{Maximus}）学习哲学。然而，诺瓦替安教友对此任命不满，宁愿他祝圣玛尔西安（\Person{Marcian}）——一位德高望重之人，因其影响，该派在瓦伦斯（\Person{Valens}）统治期间未受骚扰。故而，阿格里乌斯为平息民意，也按手在玛尔西安身上。他病情稍有好转后，前往教会，主动对会众说：\Quote{我死后，让玛尔西安作你们的主教；玛尔西安之后，让西辛尼乌作主教。}他此话后不久便去世了。于是，玛尔西安被立为诺瓦替安派主教，该派教会也因如下原因产生分裂。玛尔西安擢升一位改宗犹太教徒——名为撒巴提乌（\Person{Sabbatius}）——为司铎，但此人仍保留许多犹太偏见；并且他极其渴望成为主教。因此，他私下争取了两位司铎——德奥克提斯都（\Person{Theoctistus}）和玛卡里乌（\Person{Macarius}）——的支持，并决心维护诺瓦替安派在瓦伦斯时代于夫里家的帕祖姆村（\Place{Pazum}）关于逾越节庆期所作的革新，此事我已提及。最初，他以更苦修禁欲为借口，私下离开教会，声称\Quote{他因怀疑某些人不配领受圣事而忧伤}。然而，不久便发现他的目的是要举行分开的集会。玛尔西安得知此事后，痛责自己错误，竟祝圣如此热衷虚荣之人为司铎；并常说：\Quote{他宁可按手在荆棘上，也不应按手在撒巴提乌头上。}为制止其行为，玛尔西安促使在彼提尼雅的赫肋诺波利（\Place{Helenopolis}）附近的商业市镇安加鲁姆（\Place{Angarum}）召开诺瓦替安派主教会议。聚集于此的与会者传讯撒巴提乌，要求他解释不满的原因。他声称自己因逾越节庆期上的分歧而困扰，认为应依照犹太人的习俗，并遵循帕祖姆会议所订的规定来守节。与会主教看出此说纯属掩饰其寻求主教之位的借口，便强迫他发誓永不接受主教职。他发誓后，主教们便通过了一条关于此庆期的法规，称之为\OriginalTerm{两可法规}{Indifferent Canon}，宣布\Quote{在此事上的分歧并不构成脱离教会的充分理由；帕祖姆会议并未损害大公教会法规。即使最接近宗徒时代的古人对此庆期的遵守有差异，也未妨碍他们的彼此共融，也未引起任何不和。此外，帝国罗马的诺瓦替安派从未遵从犹太习俗，总是在春分后守逾越节；但他们并未与在不同日子庆祝的同信仰者分离。}基于这些及许多类似考量，他们制定了上述关于逾越节的\Quote{两可法规}，允许每个人若愿意，可自由保持自己偏爱的习俗；就共融而言，这并无差别；即使庆祝方式不同，他们在教会中仍应和睦。此规则确定后，撒巴提乌既受誓言约束，每当逾越节庆典时间存在差异时，便私下先行守斋，彻夜祈祷，庆祝逾越节安息日；次日则前往教会，与其他会众一同领受圣事。他如此行多年，以致无法瞒过众人；一些较为无知者，尤其是夫里家人和迦拉达人，效仿他，自以为可因此称义，便私下照他的方式守逾越节。但后来撒巴提乌不顾放弃主教职位的誓言，举行分裂集会，并被其追随者立为主教，如我们日后将述。\par

% structure: p0046 source=h2 target=subsection reason=relative-heading-rank; base=h2

\subsection{第二十二章 作者关于庆祝复活节、圣洗、斋戒、婚姻、圣体及其他教会礼仪的观点}

由于我们已经触及这个话题，我认为就复活节说几句话并非不合情理。在我看来，无论是古人还是那些自称追随犹太人的现代人，都没有任何合理的理由就这个问题如此固执地争论。因为他们没有考虑到这样一个事实：当犹太教转变为基督教时，遵守梅瑟法律和礼仪预像的义务就终止了。这件事的证据是明显的；因为基督的法律没有允许基督徒模仿犹太人。相反，宗徒明确禁止这样做；不仅拒绝割礼，而且不鼓励关于节日的争论。他在写给迦拉达人的书信中写道：\BibleQuote{你们愿意在法律之下的，请告诉我：你们没有听见法律吗？}\BibleRef{迦 4:21}他继续他的论证，指出犹太人如同仆人般受束缚，但那些来到基督面前的人却\BibleQuote{被召叫进入子女的自由}。\BibleRef{迦 5:13}此外，他劝戒他们绝不可重视\BibleQuote{日子、月份和年份}。\BibleRef{迦 4:10}此外，他在写给哥罗森人的书信中明确宣告，这样的遵守不过是影子：因此他说：\BibleQuote{不可让人在饮食上，或在节期、月朔、安息日方面评断你们；这些不过是将来之事的影子。}\BibleRef{哥 2:16}同样的真理也在他写给希伯来人的书信中得到证实：\BibleQuote{司祭职既已改变，法律也必须改变。}\BibleRef{希 7:12}因此，无论是宗徒还是福音书，都没有在任何地方将\BibleQuote{奴役的轭}\BibleRef{迦 5:1}强加给接受真理的人；而是将复活节及其他一切节期留给蒙受恩宠者以感恩之心来尊崇。因此，既然人们喜爱节期，因为它们提供劳作后的休息：各处的人便按照自己的喜好，依普遍的习惯来纪念救赎的苦难。救主和他的宗徒没有以任何法律命令我们遵守这个节期；福音书和宗徒也没有像梅瑟法律对待犹太人那样，因忽略它而用任何惩罚、刑罚或咒诅来威胁我们。仅仅是为了历史的准确性，并为了责备犹太人，因为他们在自己的节期上以血玷污了自己，福音书才记载我们的救主在\BibleQuote{无酵节}期间受难。宗徒的目的不是规定节期，而是教导正义的生活和虔诚。而且在我看来，正如许多其他习俗因习惯而在各地建立一样，复活节的庆祝也因各民族的个别特点而在各处实行，因为宗徒中没有人在此问题上立法。而且实际上，这种遵守并非源于立法，而是作为一种习俗，事实本身表明了这一点。在小亚细亚，大多数人遵守月之第十四日，不顾及安息日；但他们从未与那样做的人分离，直到罗马主教\Person{维克托}，受过于热心的影响，对亚细亚的十四日派发出了绝罚的判决。因此，法兰西里昂的主教\Person{依勒内}也写信严厉谴责\Person{维克托}的过度热情；告诉他，尽管古人在庆祝复活节上有所不同，但他们并未停止互相共融。还有，\Place{士每拿}主教\Person{波利卡普}，后来在\Person{戈尔狄安}统治下殉道，他虽按故土\Place{士每拿}的惯例在月之第十四日庆祝复活节，但仍继续与罗马主教\Person{阿尼塞特}共融，正如\Person{欧瑟伯}在他的\WorkTitle{教会史}第五卷中所证明的那样。因此，当小亚细亚的一些人遵守上述日期时，东方其他人却确实在安息日庆祝，但在月份上有所不同。前者认为应该跟随犹太人，尽管他们并不精确；后者则在春分后庆祝复活节，拒绝与犹太人一同庆祝；他们说：\Quote{因为应该在太阳进入白羊座时庆祝，即在安提约基雅人称为赞提克斯月、罗马人称为四月的那月。}他们声称，在此做法上，他们并非遵循几乎在一切事上都有错误的现代犹太人，而是遵循古人，并遵循\Person{约瑟夫}在他的\WorkTitle{犹太古史}第三卷中所记载的。因此，这些人彼此有分歧。但是西方所有其他基督徒，直到大洋之滨，都按照非常古老的传统在春分后庆祝复活节。事实上，这样行的人在这个问题上从未有过分歧。有些人声称\Person{君士坦丁}时期的会议改变了这个节期，这是不真实的：因为\Person{君士坦丁}本人写信给那些对此有不同意见的人，建议他们既然人数稀少，可以与多数弟兄达成一致。他的信完整地收录在\Person{欧瑟伯}\WorkTitle{君士坦丁传}第三卷中；其中有关复活节的段落如下：\par

\Quote{全世界西方、南方和北方的所有教会，以及东方的一些地方，都遵守一个合宜的秩序。因此，所有人在此刻都判定为正确，并且我本人也保证你们审慎的心会同意，即在\Place{罗马}城、整个\Place{意大利}、\Place{阿非利加}、整个\Place{埃及}、\Place{西班牙}、\Place{法兰西}、\Place{不列颠}、\Place{利比亚}、整个\Place{希腊}、\Place{亚细亚}和\Place{本都}教区以及\Place{基里基雅}所一致遵守的，你们的智慧也乐于接受；不仅要考虑到上述地方的教会数量更多，而且要考虑到，既然在最有理性的事上应有普遍一致，我们就应当与背信的犹太人没有任何共同之处。}\par

皇帝的書信大意如此。此外，十四日派聲稱，遵守十四日這個規矩是由若望宗徒傳給他們的；而羅馬人與西部地區的人則向我們保證，他們的習俗源於宗徒伯多祿和保祿。然而，這兩派都無法出示任何書面證據來證實他們的說法。不過，我推斷，各地慶祝復活節的時間不同是取決於習俗，因為信仰相同的人，在習俗問題上卻彼此有分歧。在此提及教會中禮儀的多樣性，或許並非不合時宜。人們會發現，復活節前的齋戒在不同民族中有不同的實行方式。羅馬人在復活節前連續齋戒三週，但星期六和星期日除外。在伊利里亞、整個希臘和亞歷山大港，人們齋戒六週，稱之為\Quote{四十天齋戒}。另一些人從復活節前第七週開始齋戒，只間歇性地齋戒三天或五天，卻也稱那段時間為\Quote{四十天齋戒}。我確實感到驚訝，他們在齋戒天數上如此不同，卻都用同一個名稱；有些人給出一個理由，另一些人又給出另一個理由，各隨己見。人們也能看到，不僅在齋戒天數上，而且在禁食的方式上也有分歧。有些人完全禁絕有生命的東西；另一些人只吃魚這類活物；許多人連同魚還吃禽鳥，說根據梅瑟\BibleRef{創 1:20}，這些也是由水造成的。有些人禁食蛋類和各種水果；另一些人只吃乾麵包；還有人連乾麵包也不吃；又有人齋戒到第九時辰，之後不加區別地吃任何食物。在各民族中還有其他習俗，為此有無數的理由被提出。然而，既然沒有人能出示書面命令作為權威，顯然宗徒們在這事上讓各人自由選擇，為使每人能出於自願而非強制或必要地行善。教會在齋戒問題上就是如此分歧。在宗教集會方面差異也不小。雖然全世界幾乎所有教會都在每週的安息日舉行神聖奧跡，但亞歷山大港和羅馬的基督徒因某些古老傳統，已停止這樣做。亞歷山大港附近的埃及人和底比斯地區的居民在安息日舉行宗教集會，但並不以基督徒通常的方式領受奧跡：因為他們在吃飽、用各種食物滿足自己之後，在傍晚獻上祭品時才領受奧跡。又在亞歷山大港，在聖週的星期三和耶穌受難日，他們閱讀聖經，教師們解釋經文；集會中舉行所有常規禮儀，但舉行奧跡除外。亞歷山大港的這個慣例非常古老，因為奧力振似乎最常在那些日子在教會中教導。他是一位精通聖經的博學教師，看出梅瑟法律的\BibleQuote{無能？}\BibleRef{羅 8:3}因字面解釋而削弱，便給予屬靈的解釋；他聲稱從來只有一個真正的復活節，就是救主懸在十字架上時所慶祝的：因為那時他戰勝了敵對的勢力，並以此作為樹立起來對抗魔鬼的戰利品。在同一座亞歷山大城，讀經員和歌詠員不分望教者和信友都可被選任；而在其他所有教會中，只有信友才能被晉升到這些職務。我自己也曾聽說過帖撒利亞的另一種習俗。如果該地的一位神職人員在領受聖職後，與他在晉鐸前合法娶的妻子同寢，就會被革職。在東方，實際上所有神職人員，甚至主教們，都戒絕與妻子往來；但這是他們自願的，並非出於任何法律的強制；因為在他們當中曾有許多主教，在任主教期間與合法妻子生了孩子。據說，帖撒利亞流行的習俗的創始人是該地特黎卡的主教赫略多魯；他年輕時寫了一部愛情小說，題為《\WorkTitle{埃塞俄比亞奇遇記}》。同樣的習俗也在得撒洛尼、馬其頓和希臘盛行。我還知道帖撒利亞的另一個特點，就是他們只在復活節的日子施行聖洗；結果許多人未領受聖洗就去世了。在敘利亞的安提約基雅，教堂的方位是倒轉的，祭台不朝東，而朝西。然而在希臘、耶路撒冷和帖撒利亞，人們在點燈後立即祈禱，與君士坦丁堡的諾窪天派一樣。在凱撒勒雅、卡帕多細亞和塞浦路斯，長老和主教們在點燈後於晚上講解聖經。赫勒斯滂的諾窪天派祈禱的方式不完全與君士坦丁堡的相同；但大多數方面，他們的慣例與主流教會相似。總之，在所有派別中，無法找到兩個在祈禱禮儀上完全一致的教會。在亞歷山大港，不允許任何長老公開講道；這項規定是在亞略擾亂該教會之後制定的。在羅馬，他們每個星期六都齋戒。在卡帕多細亞的凱撒勒雅，他們像諾窪天派一樣，拒絕那些在受洗後犯罪的人共融。同樣的紀律也由赫勒斯滂的馬其頓派和亞細亞的十四日派實行。在弗裡吉亞的諾窪天派不接受再婚的人；但君士坦丁堡的諾窪天派既不公開接納，也不公開拒絕，而在西部地區則公開接納。我認為，這種多樣性是由於各時代管理教會的主教們造成的；那些接受了這些不同禮儀和習俗的人，將它們作為法律傳給後代。然而，要完整列出每個城市和國家中使用的所有各種習俗和禮儀，是困難的——甚至是不可能的；但我們所舉的例子足以說明，復活節慶典從某些古老先例開始，就在每個省份以不同方式慶祝。因此，那些聲稱在尼西亞大公會議中更改了復活節時間的人是在胡說；因為在那裡聚會的主教們努力使最初持不同意見的少數人與其他人統一起來。事實上，即使在宗徒時代的教會中，由於這類主題也已經存在許多分歧，這甚至對宗徒們自己也不是秘密，正如《\BibleBook{宗徒}{大事录}》所見證。因為當他們得知因外邦人的爭論而在信徒中發生了騷亂時，他們聚集在一起，頒布了一條神聖法律，以書信的形式發出。藉此法令，他們使基督徒擺脫形式規條的束縛和對此類事物的無謂爭論；他們教導了真正虔誠的道路，只規定有助於達到此目的的事。我將在此轉錄的這封書信本身，記載於《\BibleBook{宗徒}{大事录}》中。\par

\Quote{宗徒和长老与弟兄们，向在安提约基雅、叙利亚和基里基雅的外邦弟兄们请安。我们听说有些从我们这里出去的人，用言语搅乱了你们，颠覆了你们的灵魂，说你们必须受割礼并遵守法律；其实我们没有吩咐他们这样。我们同心合意，认为该拣选几个人，同我们亲爱的巴尔纳伯和保禄，打发到你们那里去；这二人曾为我们主耶稣基督的名，不顾性命。所以我们打发犹达和息拉，他们也要亲口告诉你们同样的事。因为圣神和我们定意，不将更多的重担加给你们，唯有这几件事是必要的：就是禁戒祭偶像的物，和血，和勒死的牲畜，并禁戒奸淫。你们若保守自己，远离这些，就必做得好。愿你们平安。}\par

这些事确实中悦天主：因为信上明明说：\Quote{圣神定意，不将更多的重担加给你们，唯有这几件事是必要的。}然而有些人却不理会这些训诫，以为一切奸淫都是无关紧要的事；却为守节期争论不休，好像性命攸关，如此违背天主的命令，为自己立法，使宗徒的法令无效；他们也不觉察自己正行相反天主所悦纳的事。我们很容易将关于逾越节的论述加以引申，并证明犹太人在举行逾越节庆典的时间或方式上并没有任何精确的规矩；而且撒玛黎雅人，虽是犹太人的一支，总是在春分之后举行这庆节。但这题目需要一部专门详尽的论著；因此我只再补充一点：那些如此效法犹太人、如此执著于精确遵守预像的人，不应在任何事上偏离他们。因为如果他们选择如此精确，那么不但要遵守日子和月份，而且要遵守其他一切，就是基督（他\InnerQuote{生于法律之下}\BibleRef{迦 4:4}）以犹太人方式所做的一切，或是他无端被他们迫害的一切，或是他为众人益处行预表性的事。他上了一只船并教导。他吩咐在楼上预备逾越节。他命令解开拴着的驴。他提出一个拿着水罐的人作为指示，催促他们预备逾越节。［他做了］无数其他这类记载于福音中的事。然而那些以为自己因守这节庆而成义的人，却认为在肉身层面遵守任何这些事都是荒谬的。因为没有医生曾梦想从船上讲道——没有人认为必须上楼去那里守逾越节——他们从不拴住又解开驴——最后，没有人吩咐别人拿水罐，以便应验这些象征。他们很公正地认为这类事带有更多犹太主义的味道：因为犹太人更关心外在的典礼，而非内心的服从；因此他们落在咒诅之下，因为他们不辨识梅瑟法律的属灵意义，只停留在它的预像和阴影中。那些偏袒犹太人的人承认这些事的寓意；然而他们却对遵守日子和月份进行死战，而不将类似的意义应用于其上；如此他们必然与犹太人一同被定罪。\par

但我认为关于这些事已说得够了。现在让我们回到先前所论的主题，即教会一旦分裂，并未停留在那分裂中，而是那些分离的人又彼此分裂，因最微不足道的缘由而起纷争。诺瓦提安派，如我已说过，因逾越节庆节而彼此分裂，争论并不限于一点。因为在不同省份，有人持一种看法，有人持另一种，不仅关于月份，也关于星期几和其他不重要的事情；有些地方因此举行分离的集会，有些地方则彼此共融联合。\par

% structure: p0053 source=h2 target=subsection reason=relative-heading-rank; base=h2

\subsection{第二十三章 君士坦丁堡亚流派中的进一步分歧。撒提良派。}

但亚略派人士也因这事起了纷争。他们中间每日争论不休的问题，引导他们提出极其荒谬的主张。因为教会一向相信天主是圣子、圣言的圣父，于是有人问：在圣子存在之前，天主能否被称为\Quote{圣父}？他们断言天主的圣言不是由圣父所生，而是从\Quote{无}中受造的，因此在主要关键上陷入谬误，就理所当然地在一个纯粹的名号上陷于荒谬的挑剔。因此，他们从安提约基雅请来多罗特乌斯，他主张：在圣子存在之前，天主既不是、也不能被称为\Quote{圣父}。但他们在多罗特乌斯之前已从色雷斯召来了马利努斯，马利努斯因对手受到更多尊敬而恼羞成怒，便起来维护相反的意见。因此他们当中发生了分裂，并因这术语而分为两派，各自举行集会。多罗特乌斯一派保留了原来的集会场所；而马利努斯的追随者则另建祈祷所，并声称：即使圣子不存在时，圣父始终是圣父。亚略派的这一分支被称为\OriginalTerm{蛋糕贩派}{Psathyrians}，因为这一意见最热心的捍卫者之一特奥克提斯图斯是叙利亚人，以卖蛋糕（\OriginalTerm{蛋糕贩}{Psathyropola}）为生。哥特人的主教塞勒纳斯采纳了这一派的观点，他是个混血儿：父亲是哥特人，母亲是弗里吉亚人，因此他在教会中能流利地使用这两种语言教导。但这一派不久又内部分裂，因为马利努斯与他自己提升为厄弗所主教的阿加丕乌斯意见不合。他们争论的不是任何宗教信仰问题，而是心胸狭隘地争优先地位，哥特人支持阿加丕乌斯。因此，他们管辖下的许多神职人员厌恶这二人的虚荣之争，便抛弃了他们两人，归属了\Quote{同质}信仰。亚略派人士就这样分裂了三十五年，后在狄奥多西二世执政时期，在军队总司令普林塔任执政官之年重新联合，普林塔本人属于蛋糕贩派；他们被说服停止争论。随后他们通过一项决议，赋予其法律效力：导致他们分离的问题永远不得再提出。但这和解仅止于君士坦丁堡；在其他城市，凡有这两派中任何一派之处，他们仍坚持先前的分裂。关于亚略派的分裂，就说这么多。\par

% structure: p0055 source=h2 target=subsection reason=relative-heading-rank; base=h2

\subsection{第24章 欧诺米派分裂为数派}

但欧诺米的追随者也未能免于纷争。因为欧诺米本人早在此之前就已与祝圣他为基齐库斯主教的欧多克西乌斯分离，原因是欧多克西乌斯拒绝让被驱逐的他的老师埃提乌斯恢复共融。但以他命名的那些人后来也分裂为数派。首先，卡帕多西亚人德奥弗洛尼乌斯，曾受教于欧诺米的辩论技巧，并略知亚里士多德的\WorkTitle{范畴篇}和\WorkTitle{解释篇}，他撰写了题为\Quote{论心智的操练}的论文。然而，他招致了本派的谴责，被视为背教者而被驱逐。此后他离开他们独自举行集会，留下了一个以他命名的异端。此外，在君士坦丁堡，一位名叫欧提基乌斯的人因某种荒谬的争论脱离欧诺米派，至今仍单独举行集会。德奥弗洛尼乌斯的追随者被称为\OriginalTerm{欧诺米-德奥弗洛尼派}{Eunomiotheophronians}，而欧提基乌斯的追随者被称为\OriginalTerm{欧诺米-欧提基派}{Eunomieutychians}。我认为他们争论的那些无意义的术语不值得在此史中记载，以免我涉及与主题无关的事情。我只指出，他们篡改了圣洗：因为他们不因天主圣三之名施洗，而是施洗归于基督之死。马其顿派内也曾一度分裂，当时司铎欧特罗丕乌斯举行单独的集会，而卡尔特留斯与他意见不合。可能在其他城市也有这些派别中衍生出来的小派。但我住在君士坦丁堡，在那里出生和受教育，所以更详细地描述该城发生的事；因为我亲眼目睹了其中一些事件，而且那里发生的事极为重要，因而更值得记载。但应当明白，我在这里叙述的事情发生在不同时期，并非同时。如果有人想了解这些不同派别的名称，他很容易满足，只需阅读塞浦路斯主教埃皮法尼乌斯所著的\WorkTitle{锚}一书；但我满足于已经叙述的内容。国事又因我现在要提及的原因而陷入动荡。\par

% structure: p0057 source=h2 target=subsection reason=relative-heading-rank; base=h2

\subsection{第25章 僭主欧根尼谋害小瓦伦提尼安。狄奥多西击败他。}

在西方地区有一位名叫\Person{欧根尼乌斯}的文法教师，他教授拉丁语一段时间后，离开学校，被任命在宫廷服务，成为皇帝的首席秘书。他颇有口才，因此受到比他人更高的礼遇，但他无法以节制的心态承受自己的好运。他勾结了\Person{阿波加斯特}——一位\Place{小加拉提亚}人，当时指挥一支军队，此人性格粗暴且嗜血成性——决定篡夺最高权力。于是二人合谋杀害皇帝\Person{瓦伦提尼安}，收买了寝宫中的宦官。这些宦官在得到诱人的升迁许诺后，趁皇帝熟睡时将其勒死。\Person{欧根尼乌斯}立即在帝国西部篡取了最高权力，其行为正如人们预料的一个篡位者。皇帝\Person{狄奥多西}得知此事后，极为痛苦，因为他击败\Person{马克西姆斯}只是为新的麻烦铺平了道路。于是他集结军队，于1月10日，在他与\Person{阿布恩丹提乌斯}共同担任执政官的第三个任期内，立其子\Person{霍诺留}为\OriginalTerm{奥古斯都}{Augustus}，再次匆忙向西方出发，将两个儿子留在\Place{君士坦丁堡}执掌帝国权力。当他向\Person{欧根尼乌斯}进军时，\Place{多瑙河}以外的众多蛮族自愿效劳，随他出征。经过快速行军，他率领大军到达\Place{高卢}，\Person{欧根尼乌斯}也率领大量军队在那里等候他。于是双方在\Place{弗里吉杜斯河}附近交战，该河\EditorialNote{约三十六英里}远\EditorialNote{距阿奎莱亚}。在那场战斗中，罗马人与其同胞作战的部分，战况难分胜负；但在皇帝\Person{狄奥多西}的蛮族援军参战的部分，\Person{欧根尼乌斯}的军队大占优势。皇帝看到蛮族士兵大量死亡，痛苦地伏地祈求天主的帮助；他的请求没有落空。他的主要军官\Person{巴库里乌斯}在突然而异常的热忱激励下，率领前锋冲向蛮族受困之处，突破敌军阵线，使那些不久前还在追击的敌人溃逃。还发生了另一件奇妙的事：突然刮起一阵强风，将\Person{欧根尼乌斯}士兵投掷的标枪吹回他们自己身上，同时将帝国军队投掷的标枪以更大的力量吹向对手。皇帝的祈祷如此有效。战局如此扭转后，篡位者匍匐在皇帝脚下，请求饶命；但他正跪在皇帝脚下哀求时，被士兵于9月6日斩首，当时是\Person{阿卡迪乌斯}第三次担任执政官、\Person{霍诺留}第二次担任执政官。\Person{阿波加斯特}是这场大祸的主要起因，他在战后继续逃亡两天，见逃脱无望，便用他自己的剑自尽了。\par

% structure: p0059 source=h2 target=subsection reason=relative-heading-rank; base=h2

\subsection{第二十六章 年长的\Person{狄奥多西}患病与去世}

皇帝\Person{狄奥多西}因这场战争带来的焦虑和疲劳而病倒；他相信自己得的病将是致命的，便更加关心国家大事而非自己的性命，考虑君主的去世后人民往往遭遇多大的灾难。他急忙从\Place{君士坦丁堡}召来儿子\Person{霍诺留}，主要目的是安排帝国西部的事务。他儿子到达\Place{米兰}后，他似乎稍有好转，并下令为庆祝胜利举行赛马场竞技。用餐前他状况尚好，观看了比赛；但用餐后他突然病重无法返回，派儿子代他主持。夜晚来临时，他去世了，时值1月17日，\Person{奥利布里乌斯}和\Person{普罗布斯}担任执政官期间。这是第二百九十四届\OriginalTerm{奥林匹亚德}{Olympiad}的第一年。皇帝\Person{狄奥多西}活了六十岁，在位十六年。因此本卷涵盖十六年又八个月的事务。\par

\SourceRecord{Translated by A.C. Zenos.；Nicene and Post-Nicene Fathers, Second Series；Vol. 2.；Edited by Philip Schaff and Henry Wace.；Buffalo, NY: Christian Literature Publishing Co.,；1890.；<http://www.newadvent.org/fathers/26015.htm>.}

% structure: p0001 source=h1 target=section reason=page-title

\section{\WorkTitle{教会史（第六卷）}}

% structure: p0002 source=h2 target=subsection reason=relative-heading-rank; base=h2

\subsection{引言。}

哦，天主的人\Person{特奥多勒}，你交付给我们的任务，我们已在前五卷中完成；在那五卷里，我们尽己所能，从\Person{君士坦丁}时代起记述了教会历史。但请注意，我们绝未刻意雕琢文体；因为我们以为，若过分讲究辞藻优雅，反倒可能妨碍既定目标。即便假设我们的目的仍能达成，我们也完全无法运用古代历史学家似乎大量使用的那种自由裁量权——他们中的每一位都自以为可以随意增减事实。此外，精炼的文体对大众和目不识丁者毫无裨益，他们只想知道事实，而非欣赏文辞之美。因此，为使我的作品对两类读者——一方面是有学问的人，因为无论怎样精炼的语言都无法令他们满意，无法将其与古代作家的宏富文采相提并论；另一方面是无学问的人，因为若事实被辞藻的炫耀所遮蔽，他们便无法理解——都不致无益，我们有意采用了一种平实清晰的文体，毫无刻意追求崇高的矫饰。\par

然而，当我们开始第六卷时，必须预先说明：我们承担详述本时代事件的任务，担心会提出许多令不少人不悦的事情；或是因为，正如谚语所言\Quote{真理是苦的}；因为未能以赞词提及某些人所喜爱之人的名字；或是因为我们没有夸大他们的行为。我们教会的热心者会谴责我们未称主教们为\Quote{至亲于天主者}、\Quote{至圣者}等等。另一些人会争论，因为我们未将\Quote{至神圣者}和\Quote{主}等称号加于皇帝，也未使用通常赋予他们的其他尊称。但我能轻易从古代作者的见证中证明，在他们当中，仆人惯常仅以名字称呼主人，而不提及他的尊贵或头衔，因事务繁忙之故。同样，我将服从历史法则，它要求简明忠实的叙述，不为任何面纱所遮蔽。我将准确记录我亲眼所见或从实际目击者那里确知之事，通过证人对同一事件的一致见证以及我所能掌控的一切手段来检验真相。查明真相的过程确实艰辛，因为许多不同的人给出不同的说法，有些人声称是目击者，另一些人则自称比任何人都更熟悉这些事情。\par

% structure: p0005 source=h2 target=subsection reason=relative-heading-rank; base=h2

\subsection{第一章 论\Person{狄奥多西}之死，其两子分治帝国；\Person{鲁菲努斯}在\Person{阿卡狄乌斯}脚下被杀。}

皇帝\Person{狄奥多西}驾崩后，时在\Person{奥利布里乌斯}与\Person{普罗比努斯}执政之年，即一月十七日，其两子接管罗马帝国。于是\Person{阿卡狄乌斯}治理东方，\Person{霍诺里乌斯}治理西方。其时\Person{达玛苏斯}为帝国罗马教会主教，\Person{特奥菲卢斯}为亚历山大里亚主教，\Person{若翰}为耶路撒冷主教，\Person{弗拉维亚努斯}为安提约基雅主教；而君士坦丁堡（或称新罗马）的主教席位则由\Person{内克塔里乌斯}担任，正如我们在前卷所述。皇帝\Person{狄奥多西}的遗体于同年十一月八日运抵君士坦丁堡，并由其子\Person{阿卡狄乌斯}以通常的葬礼隆重安葬。不久之后，同月二十八日，曾在皇帝\Person{狄奥多西}麾下参加讨伐僭主之战的军队也到达。因此，当皇帝\Person{阿卡狄乌斯}依惯例在城门外迎接军队时，士兵们杀死了近卫军长官\Person{鲁菲努斯}。因为有人怀疑他觊觎皇位，且据传他曾邀请蛮族\OriginalTerm{匈人}{Huns}进入罗马领土，后者已蹂躏\Place{亚美尼亚}，并正在其他东方行省进行劫掠。在\Person{鲁菲努斯}被杀的当天，诺瓦提安派主教\Person{马尔西安}去世，由我们已提及的\Person{西西尼乌斯}继任主教。\par

% structure: p0007 source=h2 target=subsection reason=relative-heading-rank; base=h2

\subsection{第二章 \Person{内克塔里乌斯}之死与\Person{若翰}的祝圣。}

不久之后，\Place{君士坦丁堡}主教\Person{聂克塔留斯}也在\Person{凯撒留斯}和\Person{阿提库斯}执政时期，于九月二十七日去世。一场关于继任者的争夺随即展开，有人推举这人，有人推举那人；但最终决定派人去请\Place{安提约基雅}教会的一位司铎\Person{若望}，因为据报告他非常善于教导，且口才出众。因此，在圣职人员和教友的一致同意下，他不久后被\Person{阿尔卡狄乌斯}皇帝召到\Place{君士坦丁堡}：为了使祝圣更具权威和庄严，邀请了几位主教出席，其中也包括\Place{亚历山大里亚}主教\Person{德奥斐洛}。此人竭尽所能诋毁\Person{若望}的声誉，因为他想提拔自己教会的一位司铎\Person{依西多禄}到该教座，他非常依恋\Person{依西多禄}，因为\Person{依西多禄}曾为他的利益承担了一件非常棘手和危险的事务。我现在必须说明这是什么事。当\Person{德奥多西}皇帝准备进攻篡位者\Person{马克西穆斯}时，\Person{德奥多西}派\Person{依西多禄}带着礼物和双重信件，命令他将礼物和适当的信件交给最终获胜者。\Person{依西多禄}遵照指示抵达\Place{罗马}，在那里等待战争的结果。但这件事情不久就泄露了：因为陪同他的一个读经员私下扣留了信件；于是\Person{依西多禄}惊恐万分地返回了\Place{亚历山大里亚}。这就是\Person{德奥斐洛}如此偏爱\Person{依西多禄}的原因。然而宫廷却偏爱\Person{若望}：并且由于许多人重新提出对\Person{德奥斐洛}的指控，并准备了各种控诉的备忘录提交给当时聚集的主教们，皇室寝宫总管\Person{欧特罗皮乌斯}收集了这些文件，并拿给\Person{德奥斐洛}看，告诉他\Quote{要么祝圣若望，要么接受对你提出的指控的审判。}\Person{德奥斐洛}被这一选择吓住了，同意祝圣若望。于是\Person{若望}在二月二十六日接受了主教职，当时\Person{霍诺留}皇帝在\Place{罗马}以公共竞技庆祝，而\Person{欧提基安}（当时在\Place{君士坦丁堡}任近卫总长）的执政年。但是，既然此人以他所留下的著作和他所陷入的许多麻烦而闻名，我认为我不应该对他的事情默不作声，而是尽可能简明地叙述他来自何方，出身如何；还有他升任主教的具体情况，以及他后来被贬黜的方式；最后叙述他死后如何比生前更受尊敬。\par

% structure: p0009 source=h2 target=subsection reason=relative-heading-rank; base=h2

\subsection{第三章 \Place{君士坦丁堡}主教\Person{若望}的出生与教育}

\Person{若望}是\Place{叙利亚-科莱的安提约基雅}本地人，父亲\Person{瑟孔多斯}，母亲\Person{安图萨}，是该国贵胄之子。他跟随修辞学家\Person{利巴尼乌斯}学习修辞学，跟随哲学家\Person{安德拉加提乌斯}学习哲学。他正要开始从事民法实践，但反思那些献身于法庭实践之人的不安与不义行为，便转向了平静的生活方式，效法\Person{埃瓦格里乌斯}的榜样。\Person{埃瓦格里乌斯}本人曾在同样的老师门下受教育，并且在此前一段时间已退隐过私人生活。于是他脱下法袍，专心阅读圣经，非常勤奋地前往教堂。他还说服了与他一起在修辞学家\Person{利巴尼乌斯}门下求学的同学\Person{德奥多若}和\Person{马克西穆斯}，放弃以谋利为首要目标的职业，拥抱更简朴的生活。这两人中，\Person{德奥多若}后来成为\Place{基里基雅的摩普绥提亚}主教，\Person{马克西穆斯}成为\Place{伊扫里亚的塞琉西亚}主教。那时，他们热心追求成全，在\Person{狄奥多若}和\Person{卡尔特利乌斯}的指导下开始隐修生活，这两人当时主持一座隐修院。前者后来被提升为\Place{塔尔索斯}主教，并撰写了许多论文，在其中他局限于圣经的字面意义，避免神秘解释。关于这些人就说到这里。当时\Person{若望}与\Person{巴西略}交往最为密切，\Person{巴西略}当时由\Person{默肋提乌斯}立为执事，后来被祝圣为\Place{卡帕多细亚的凯撒利亚}主教。于是，\Person{策诺}主教从\Place{耶路撒冷}回来后，任命他为\Place{安提约基雅}教会的读经员。在他担任读经员期间，他撰写了\WorkTitle{驳犹太人}一书。\Person{默肋提乌斯}不久后授予他执事职，他创作了\WorkTitle{论司祭职}和\WorkTitle{驳斯塔吉利乌斯}等作品；此外还有\WorkTitle{论天主本性的不可理解性}和\WorkTitle{论与教士同居的妇女}。后来，\Person{默肋提乌斯}在\Place{君士坦丁堡}去世——因为他为了\Person{纳西盎的额我略}的祝圣去了那里——\Person{若望}离开了默肋提乌斯派，没有与\Person{保利努斯}共融，隐居了整整三年。后来，\Person{保利努斯}去世后，他被\Person{保利努斯}的继任者\Person{埃瓦格里乌斯}祝圣为司铎。以上是\Person{若望}被召任主教职之前生涯的简述。据说，由于他对节制的热情，他严肃而严厉；他的一位早期朋友曾说\Quote{他年轻时表现出易怒的倾向，而非谦逊。}由于他生活的正直，他对未来无忧无虑，而他的纯朴性格使他坦率而真诚；然而他允许自己的言论自由冒犯了很多人。在公开教导中，他在改正听者的道德方面大有功效；但在私下交谈中，不了解他的人常常认为他傲慢自大。\par

% structure: p0011 source=h2 target=subsection reason=relative-heading-rank; base=h2

\subsection{第四章 论执事\Person{塞拉皮昂}，\Person{若望}因他而遭其圣职人员憎恶}

\Person{若望}在性情和举止上如此，又被擢升为主教，遂以过分倨傲的态度对待其神职人员，意图纠正属下神职人员的道德。如此触怒了这些教会人士，他们便厌恶他；许多人因他暴躁而疏远他，另一些人则成为他的死敌。他身边的执事\Person{塞拉皮翁}更使他与众人疏远；一次在全体现职神职人员面前，他大声对主教喊道：\Quote{主上，除非你用棍子把他们全都赶走，否则你永远无法治理这些人。}这番话激起众人对主教的普遍敌意；主教不久也将其中许多人逐出教会，或为此故，或为彼故。正如凡掌权者采取此类激烈手段时常见的情形，那些被他驱逐的人结成团体，在民众面前攻击他。使这些指控更显可信的重大因素是：主教不愿与任何人一同进食，从不接受宴请。为此反对他的阴谋四处蔓延。他为何不与别人共食，无人确知；但有人为他的行为辩护，说他肠胃极弱、消化不佳，不得不谨慎饮食，因此独食；另有人则认为这是出于他严格的习惯性节制。无论真实动机如何，这一事实本身为他的诽谤者提供了不少口实。然而民众仍因他在教会中有价值的讲道而爱戴尊敬他，那些图谋诋毁他的人反而自取其辱。他的讲道辞——无论自己发表的那些，还是速记员在他讲道时记录的那些——何等雄辩、有说服力且打动人，我们又何必赘述？凡愿对其有充分认识者，须自行阅读，并将由此获得愉悦与裨益。\par

% structure: p0013 source=h2 target=subsection reason=relative-heading-rank; base=h2

\subsection{第五章  \Person{若望}招致许多显贵之人的不悦。论宦官\Person{欧特罗皮乌斯}。}

当\Person{若望}仅与神职人员冲突时，针对他的阴谋完全无效；但当他开始以过度激烈的方式斥责许多在职官员时，不受欢迎的潮流便以远为强大的势头向他涌来。于是，许多贬损他的故事被传开，且大多有热衷听信之人。当时他针对\Person{欧特罗皮乌斯}的一篇演说更使这种日渐增长的偏见大大增加。因为\Person{欧特罗皮乌斯}是皇帝寝宫的首席宦官，也是所有宦官中第一位荣任执政官者。他欲报复某些躲入教堂的人，便怂恿皇帝立法剥夺罪犯的避难特权，授权逮捕那些寻求圣殿庇护的人。但立法者几乎立刻为此受到惩罚；法律颁布不久，\Person{欧特罗皮乌斯}本人失宠于皇帝，逃到教堂寻求保护。主教见\Person{欧特罗皮乌斯}战战兢兢地躲在祭台桌下，便登上他惯常向民众讲话以便更清晰被听到的讲坛，对他发表了抨击之辞；因此，在某些人看来，他似乎引起了更大的不满，因为他不仅拒绝怜悯不幸者，而且在残忍之上又增添了侮辱。然而，根据皇帝的命令，\Person{欧特罗皮乌斯}虽身居执政官之职，却因所犯某些罪行而被斩首，其名从执政官名单中除去，只保留其同僚\Person{特奥多雷}为当年在任者。据说\Person{若望}后来也对当时担任军队统帅的哥特人\Person{盖纳斯}使用了同样的直言不讳；他以其特有的粗鲁对待\Person{盖纳斯}，因为\Person{盖纳斯}曾请求皇帝将城中一座教堂分给与他持相同观点的亚略派。此外，他还以同样不拘礼节的自由责备了许多其他高阶人士，因而树立了许多强大的敌人。因此，\Place{亚历山大}的主教\Person{特奥斐洛}在\Person{若望}祝圣后立即图谋推翻他，并与身边的朋友秘密策划，还与远方的人通信策划。因为让\Person{特奥斐洛}介意的，与其说是\Person{若望}任意鞭挞他所厌恶之事的胆量，不如说是他自己未能将宠爱的司铎\Person{伊西多尔}置于\Place{君士坦丁堡}主教之位。当时\Person{若望}主教的处境即是如此；他的主教任期一开始便有祸患威胁。但我们将在后文中更详细地叙述这些事。\par

% structure: p0015 source=h2 target=subsection reason=relative-heading-rank; base=h2

\subsection{第六章  哥特人\Person{盖纳斯}图谋篡夺最高权力；使\Place{君士坦丁堡}充满混乱后，被杀。}

我现在要叙述当时发生的一些值得纪念的情况，从中可以看出天主眷顾如何通过非凡的工具介入，以保全城市和罗马帝国免于极大的危险。盖纳斯出身蛮族，但成为罗马臣民后，服役于军队，并逐步升迁，最终被任命为罗马骑兵和步兵的总司令。当他获得这一高位后，他忘记了自己的身份和关系，无法自制，另一方面，正如俗语所说\Quote{不遗余力}，以图控制罗马政府。为此，他派人从哥特人中召来同族，并把军队中的主要指挥权交给了他的亲属。后来，当他的一个亲属\Person{特里比吉尔杜斯}——他在\Place{弗里吉亚}统率军队——在\Person{盖纳斯}的煽动下公开反叛，使\Place{弗里吉亚}人民陷入混乱和惊恐时，他设法让自己被委派去管理那个动乱的省份。当时，皇帝\Person{阿卡狄乌斯}没有怀疑任何危害，把这些事务的职责交给了他。因此\Person{盖纳斯}立即率领大量蛮族哥特人出发，表面上是对\Person{特里比吉尔杜斯}进行征讨，但实际意图是建立他自己不义的统治。到达\Place{弗里吉亚}后，他开始颠覆一切。因此，罗马人的事务立刻陷入极大的恐慌，不仅因为\Person{盖纳斯}指挥的庞大蛮族军队，也因为东方最富饶的地区面临被蹂躏的威胁。在这一紧急关头，皇帝行事非常谨慎，试图通过交涉来阻止蛮族的行动：于是他派遣一个使团，指示他们用各种让步来暂时安抚他。\Person{盖纳斯}要求将\Person{萨图尔尼努斯}和\Person{奥雷利安}——两位元老院中最显要的人物，具有执政官身份，他知道他们反对他的企图——交给他，皇帝极不情愿地屈从于危机的紧急情况；而这两个人，准备为公共利益而死，高贵地服从了皇帝的处置。于是他们前去会见这个蛮族，地点在离\Place{迦克墩}不远的一个赛马场，决心忍受他可能施加的任何惩罚；然而他们并没有遭受伤害。这个篡位者假装不满，前进到\Place{迦克墩}，皇帝\Person{阿卡狄乌斯}也去那里见他。两人随后进入了存放殉道者\Person{尤菲米娅}遗体的圣堂，并在那里互相宣誓，承诺彼此不图谋反对对方。皇帝确实遵守了他的承诺，因为他虔诚地尊重誓言，并因此为天主所爱。但\Person{盖纳斯}很快就背弃了誓言，并没有偏离他最初的目的；相反，他致力于屠杀、掠夺和纵火，不仅针对\Place{君士坦丁堡}，也针对整个\Place{罗马}帝国，如果他能设法实现的话。因此，城市被蛮族淹没，居民沦落到与俘虏相当的地步。此外，城市的危险如此之大，以至于一颗从未见过的、巨大无比、从天到地的彗星预兆了它。\Person{盖纳斯}首先最无耻地试图夺取商店中公开陈列出售的银器：但当店主们事先得到消息，得知他的意图，因而避免将其摆放在柜台上时，他的想法转向另一个目标，即派遣一大群蛮族在夜间焚烧皇宫。那时，确实清楚地表明天主对这座城市有眷顾：因为一大群天使以全副武装的巨人形象出现在叛军面前，蛮族以为他们是庞大的英勇军队，惊恐地转身离去。当这件事报告给\Person{盖纳斯}时，他觉得完全不可信——因为他知道罗马军队的大部分都在远处，分散驻防在东方的城市——于是他在第二天晚上并随后多次派遣其他人。由于他们总是带着同样的报告返回——因为天主的天使总是以同样的形象出现——他带领一大群人前来，最终亲自目睹了这一奇迹。于是，他以为他所看见的确实是一队士兵，白天隐藏起来，夜间挫败他的计划，他停止了尝试，并采取了另一个他认为会对罗马人不利的决定；但结果证明这对他们大大有利。他假装被魔鬼附身，假装去祈祷，前往离城七英里的\Emph{圣若望宗徒圣堂}。与他同行的蛮族携带了武器，将武器藏在酒桶和其他伪装覆盖物中。当守卫城门的士兵发现这些武器，不让他们通过时，蛮族拔出剑杀死了他们。城中因此产生了可怕的骚动，死亡似乎威胁着每个人；尽管如此，当时城市仍然安全，因为城门处处得到很好的防守。皇帝以适时的智慧宣布\Person{盖纳斯}为公敌，并下令杀死所有仍被困在城内的蛮族。于是在城门守卫被杀后的一天，罗马人在城内的哥特人圣堂附近攻击了蛮族——因为那些留在城中的蛮族都躲到了那里——在消灭了大量蛮族后，他们放火烧了圣堂，将其夷为平地。\Person{盖纳斯}得知他那些未能逃出城的人被屠杀，并意识到他所有诡计的失败，便离开了圣若望圣堂，迅速向\Place{色雷斯}前进。到达\Place{克索尼苏斯}后，他试图从那里渡过并攻占\Place{兰普萨库斯}，以便从那个地方成为东方地区的主人。由于皇帝立即派遣了军队从陆路和海路追击，另一个天主眷顾的奇妙干预发生了。因为当蛮族没有船只，匆忙扎成木筏并试图渡河时，突然罗马舰队出现，西风开始猛烈吹刮。这为罗马人提供了便利的通道；但蛮族连同他们的马匹，在狂风中于脆弱的木筏上颠簸，最终被波涛吞没；他们中许多人也为罗马人所杀。就这样，在渡河过程中，大量蛮族丧生；但\Person{盖纳斯}从那里逃走后逃往\Place{色雷斯}，在那里遇到另一支罗马军队，他和随同他的蛮族都被杀。关于\Person{盖纳斯}的简短叙述到此为止。\par

那些想要了解那场战争更多细节的人，应该阅读\Person{欧瑟比乌斯·斯科拉斯提库斯}的\WorkTitle{盖尼亚}，他当时是智者\Person{特罗伊鲁斯}的学生；他亲眼目睹了那场战争，并以一部由四卷组成的英雄史诗记述了其事；由于所提及的事件刚刚发生不久，他为自己赢得了极大的声誉。诗人\Person{阿摩尼乌斯}最近也以诗句形式创作了关于相同事件的另一篇描述，他在\Person{小德奥多西}与\Person{福斯图斯}共同担任执政官的第十六任期内，在皇帝面前朗诵了该作品。\par

这场战争在\Person{斯提利科}和\Person{奥勒良}担任执政官期间结束。次年，由\Person{弗拉维图斯}（也是哥特裔）主持执政官庆典，他受到罗马人的尊敬，并向他们表现出极大的忠诚与依恋，在这场战争中提供了重要的服务。为此，他获得了执政官的尊荣。同年四月十日，皇帝\Person{阿卡狄俄斯}得一子，即良善的\Person{德奥多西}。\par

然而，当国家事务如此动荡不安时，教会的显贵们丝毫没有停止他们彼此之间可耻的密谋，这给基督信仰带来了极大的耻辱；因为在此期间，神职人员煽动彼此间的骚乱。这祸害的根源起源于\Place{埃及}，方式如下。\par

% structure: p0020 source=h2 target=subsection reason=relative-heading-rank; base=h2

\subsection{第七章. \Place{亚历山大}的主教\Person{德奥斐洛}与旷野隐修士之间的分歧。\Person{奥力振}著作的谴责。}

在此之前不久，有人提出了这样一个问题：天主是否是有形体的存在，具有人的形象？或者他是无形体的，没有人的或一般所说的任何其他形状？由此问题，在为数众多的人中引起了纷争和争论，一些人赞成这一观点，另一些人支持相反的观点。许多比较单纯的隐修士断言天主是有形体的，且具有人形；但大多数人谴责他们的判断，主张天主是无形体的，没有任何形式。亚历山大里亚主教德奥斐洛完全同意后一种观点，以致在教会中当着全体人民的面，他痛斥那些将人形归于天主的人，明确教导说神体完全是无形体的。\newline{}当埃及隐修士得知此事后，他们离开隐修院来到亚历山大里亚，在那里对主教发动骚乱，指控他不虔敬，并威胁要处死他。德奥斐洛意识到自己的危险，经过一番思考，他采取了这样的权宜之计来摆脱死亡的威胁。他走到隐修士们面前，用和解的语气对他们说：\Quote{我看见你们，就是看见了天主的容貌。} 这句话平息了这些人的怒气，他们回答说：\Quote{如果你真的承认天主的容貌像我们一样，就请诅咒奥力振的书；因为有些人从他的书中引证论据来反对我们的意见。如果你不这样做，就等着被我们当作不虔敬的人和天主的敌人来对待吧。} \Quote{至于我，}德奥斐洛说，\Quote{我很乐意按你们的要求去做。不要对我生气，因为我自己也不赞同奥力振的作品，并认为那些支持它们的人应受责备。} 这样，他成功安抚并遣走了那些隐修士；若非紧接着发生的另一件事，这场争论很可能就此平息。\newline{}在埃及的隐修院中，有四位虔诚的监督者，名叫狄奥斯哥鲁、阿摩纽、欧瑟伯和欧提米；他们是兄弟，因其身高而被称作\OriginalTerm{高修士}{Tall Monks}。他们更因其生活的圣洁和学识的广博而著称，因此在亚历山大里亚享有很高的声望。特别是该城的主教德奥斐洛，极其敬爱他们，以至于不顾狄奥斯哥鲁的意愿，强行将他从隐修中拉出，立他为赫尔莫波利斯主教。他恳求另外两人继续与他同住，费尽力气才说服他们；他运用主教权威达到了目的：于是授予他们圣职，委托他们管理教会事务。他们被形势所迫，成功地履行了加给他们的职责；然而他们并不满意，因为他们无法追求哲学研究和隐修操练。随着时间的推移，他们觉得自己的灵性受到损害，看到主教耽于牟利、贪婪地聚敛财富，正如俗话所说，\Quote{为了利益不择手段}，他们便拒绝再与他同住，声称他们喜爱独处，远比住在城里更喜欢。起初不知他们离去的真实动机时，他恳切地请求他们留下；但当他发现他们不满他的行为时，便极为恼怒，威胁要给他们种种伤害。但他们不理会他的威胁，退入了旷野。于是，德奥斐洛显然性情急躁且恶毒，对他们大加指责，并用尽一切办法竭力伤害他们。他还对他们的兄弟、赫尔莫波利斯主教狄奥斯哥鲁产生了厌恶。此外，他对隐修士们对狄奥斯哥鲁的尊敬和崇敬感到极为恼火。\newline{}但他知道，除非他能煽动隐修士们对这些人的敌意，否则无法伤害他们，于是用了这个计谋。他很清楚，这些人经常与他进行神学讨论时，坚持主张神性是无形体的，绝无人形，因为这样的构成必然伴随着人类情感的附属。这一点已由古代作家、尤其是奥力振所证明。然而，德奥斐洛虽然自己对于神性也持有完全相同的看法，但为了满足报复心理，不惜歪曲他和他们曾正确教导的真理：他欺骗了大多数隐修士——这些人是真诚但\BibleQuote{拙于言辞}\BibleRef{格后 11:6}的人，其中大部分完全是文盲。他向旷野中的隐修院发送信件，劝告他们不要听从狄奥斯哥鲁和他的兄弟，因为他们断言天主没有身体。\Quote{然而，}他说，\Quote{按照圣经，天主像人一样有眼、耳、手、脚；但狄奥斯哥鲁的同党，作为奥力振的追随者，引入了亵渎性的教条，说天主没有眼、耳、脚、手。} 他用这种诡辩利用了这些隐修士的单纯，从而在他们中间激起了激烈的纷争。\newline{}确实，有教养的人不被这种似是而非的言论所迷惑，因此仍然追随狄奥斯哥鲁和奥力振；但更无知的人（在数量上远超他人）在缺乏知识的狂热中被点燃，立刻对弟兄们发出攻击。这样一分裂，双方都指责对方不虔敬；一些听从德奥斐洛的人称他们的弟兄为\OriginalTerm{奥力振派}{Origenists}和\Quote{不虔敬的人}，而另一些人则称那些被德奥斐洛说服的人为\OriginalTerm{人形派}{Anthropomorphitæ}。为此爆发了激烈的争吵和隐修士之间无法熄灭的战争。德奥斐洛得知其计谋成功后，便带着一群人前往隐修院所在的尼特里亚，武装隐修士反对狄奥斯哥鲁和他的兄弟；这些人生命垂危，好不容易才逃脱。\par

当这些事在埃及进行时，君士坦丁堡主教若望对此一无所知，但他的口才越发彰显，其讲道也日益著名。此外，他首次扩充了夜间赞美诗中的祈祷，原因如下。\par

% structure: p0023 source=h2 target=subsection reason=relative-heading-rank; base=h2

\subsection{第8章——亚略派信徒与\Quote{同质者}的支持者举行夜间集会并咏唱对唱圣歌，此类作品被认为出自依纳爵（别号德奥福禄）。两派之间的冲突。}

如前所述，亚略派信徒在城外聚会。每逢节日——即每周的星期六和主日，即教会通常举行聚会的日子——他们聚集在城门内的公共广场上，咏唱符合亚略异端的应答诗句。他们这样做几乎彻夜；清晨，他们唱着同样的歌曲（他们称之为\Quote{应答歌}），游行穿过城中，然后出城前往他们的聚会地点。但他们并未停止使用侮辱性言辞攻击同质论派，常常唱起这样的歌词：\Quote{那些说三物只是一权的人在哪里？}——若望担心一些较为单纯的人会被这类赞美诗引诱离开教会，便安排自己的一些人也参与咏唱夜间赞美诗，以削弱亚略派信徒的努力，并巩固自己一方对信仰的宣认。若望的设计似乎出于好意，却引起了骚乱和危险。因为当同质论派以更壮观的方式咏唱夜间赞美诗时——若望为他们发明了银制十字架，上面点燃着蜡烛，由皇后欧多克西亚提供费用——人数众多且充满嫉妒的亚略派信徒决意以绝望而暴乱的攻击来报复对手。由于他们仍记得自己不久前的统治地位，他们充满自信，轻视对手。于是，在一个夜晚，他们立即投入战斗；皇后的一位太监布里索，当时正带领这些赞美诗的领唱者，被一块石头击中额头受伤，双方也有一些平民被杀。皇帝得知后大怒，禁止亚略派信徒再在公开场合咏唱他们的赞美诗。这就是当时发生的事件。\par

然而，我们现在必须提及教会中应答歌唱这一习俗的起源。安提约基雅的第三位主教依纳爵（从宗徒伯多禄算起，他曾与宗徒们亲身交往）看见天使交替咏唱赞颂天主圣三的异象。于是他将所见到的歌唱方式引入安提约基雅教会，并由此通过传统传递到所有其他教会。这就是我们对于这些应答赞美诗的记述。\par

% structure: p0026 source=h2 target=subsection reason=relative-heading-rank; base=h2

\subsection{第9章——德奥斐洛与伯多禄之间的争论导致前者企图罢免君士坦丁堡主教若望。}

此后不久，旷野隐修士们与狄奥斯科若及其弟兄们来到君士坦丁堡。与他们同来的还有依西多尔，他曾是德奥斐洛主教最亲密的朋友，但如今成了他最大的敌人，原因如下：当时亚历山大教会有一位名叫伯多禄的总司铎；德奥斐洛对此人不满，决定将他逐出教会，并以此为借口指控他让一位摩尼教派女子参与神圣奥迹，而未先迫使她弃绝摩尼教异端。当伯多禄辩护说，不仅该女子的错误先前已被弃绝，而且德奥斐洛本人也曾批准她领受圣体圣事时，德奥斐洛感到愤慨，仿佛受了严重诽谤；于是他声言自己对整个事件一无所知。因此伯多禄传唤依西多尔作证，证明主教对这名女子的事实知情。当时依西多尔恰好在罗马，奉德奥斐洛之命前往帝都主教达玛稣处，为促成他与安提约基雅主教弗拉维雅诺之间的和解；因为默勒提派的追随者厌恶弗拉维雅诺的伪誓而与他分离，正如我们之前所述。当依西多尔从罗马返回，并被伯多禄召为证人时，他宣誓声明该女子是经主教同意被接纳的，而且他本人曾为她施行圣事。德奥斐洛闻讯大怒，愤怒地将两人都驱逐了。这就是依西多尔与狄奥斯科若及其弟兄们前往君士坦丁堡的原因，为将德奥斐洛对他们的不义和暴力提交皇帝和主教若望审理。若望得知事实后，热情接待了这些人，并未在祈祷中拒绝与他们共融，但推迟了他们领受神圣奥迹，直到他们的事务得到审查。正在此时，一个虚假传闻传到德奥斐洛耳中，说若望已准许他们参与奥迹，并准备帮助他们；于是他决心不仅要报复依西多尔和狄奥斯科若，而且若有可能，还要将若望逐出主教之位。为此，他写信给各城的主教，隐藏真实动机，表面上仅谴责奥力振的著作；而他的前任亚大纳削曾用这些著作来确证自己的信仰，在反对亚略派信徒的演讲中屡次援引奥力振著作的见证和权威。\par

% structure: p0028 source=h2 target=subsection reason=relative-heading-rank; base=h2

\subsection{第10章——塞浦路斯主教厄丕法尼召集主教会议谴责奥力振的著作。}

此外，他重新与塞浦路斯君士坦提亚主教\Person{厄丕法尼}恢复友谊，此前二人曾有分歧。因为\Person{德奥斐禄}曾指控\Person{厄丕法尼}对天主持有低下的观念，认为天主具有人形。尽管\Person{德奥斐禄}实际上并未改变看法，并曾谴责那些认为神性具有人形的人，但因对他人的仇恨，他公开否认了自己的信念；现在他声称与\Person{厄丕法尼}友好，仿佛他已改变心意，并与\Person{厄丕法尼}在天主的看法上一致。然后他设法使\Person{厄丕法尼}致信召集塞浦路斯的主教会议，以谴责\Person{奥力振}的著作。\Person{厄丕法尼}因非凡的虔诚，心地单纯，作风简朴，易受\Person{德奥斐禄}书信的影响：于是他在该岛召集了主教会议，并在会上禁止阅读\Person{奥力振}的作品。他还写信给\Person{若望}，劝他放弃研究\Person{奥力振}的著作，并召集会议颁布同样的禁令。当\Person{德奥斐禄}这样欺骗了以虔诚闻名的\Person{厄丕法尼}后，见计划如愿进展，便更加自信，自己也召集了许多主教。在那次会议上，他效法\Person{厄丕法尼}的做法，促使对已去世近两百年的\Person{奥力振}的著作作出了类似的谴责判决：但这并非他的首要目标，他的真正目的在于报复\Person{狄约斯哥洛}及其弟兄。\Person{若望}对\Person{厄丕法尼}和\Person{德奥斐禄}的来函不甚在意，专心教导各教会；他作为宣道者日益兴盛，却对针对他的阴谋不予理会。然而，当人人都清楚\Person{德奥斐禄}正试图剥夺\Person{若望}的主教职位时，所有对\Person{若望}怀有恶意的人便联合起来诋毁他。于是，许多圣职人员、许多担任公职者以及朝廷中极具影响力的人，认为找到了报复\Person{若望}的机会，便竭力促成在\Place{君士坦丁堡}召开大公会议，一方面发送书信，另一方面向各地派遣信使。\newline{}\par

% structure: p0030 source=h2 target=subsection reason=relative-heading-rank; base=h2

\subsection{第十一章。论\Person{塞味黎安}与\Person{安提约古}：他们与\Person{若望}的分歧。}

对\Person{金口若望}的憎恶因另一事件而显著增加，情况如下：当时有两位主教，叙利亚人出身，名为\Person{塞味黎安}和\Person{安提约古}；塞味黎安管理叙利亚的\Place{加巴拉}教会，安提约古管理腓尼基的\Place{托勒买}教会。两人都以口才著称；但塞味黎安虽为饱学之士，却不能完美地使用希腊语；因此他在说希腊语时暴露出他的叙利亚出身。安提约古先来到\Place{君士坦丁堡}，在教会中热忱而能干地讲道了一段时间，因而积累了一大笔钱，然后回到自己的教会。塞味黎安听说安提约古因访问君士坦丁堡而发了财，便决定效法他的榜样。他为此锻炼自己，撰写了许多讲道稿，然后动身前往君士坦丁堡。他受到若望极为友善的接待，在某种程度上，他安抚和奉承若望，他自己也受到若望同样程度的喜爱和尊敬；同时，他的讲道为他赢得了极大的声誉，以致引起了许多显要人士甚至皇帝本人的注意。恰在此时，\Place{厄弗所}的主教去世，若望必须前往厄弗所，以便祝圣一位继任者。到达该城后，由于民众在选举上意见分歧，有人推举这人，有人推举那人，若望看出双方都处于争执状态，不愿听从他的劝告，便决定不费周折地结束他们的争端，将自己的一位执事、塞浦路斯裔的\Person{赫拉克利德}擢升为主教。于是双方停止了争斗，获得了和平。由于这次在厄弗所的耽搁延长了，塞味黎安继续在君士坦丁堡讲道，日益得到听众的喜爱。若望并非不知情，因为有人立即将所发生的一切报告给他，我们前面提到过的\Person{色拉皮雍}向他传递消息，并声称教会正被塞味黎安扰乱；于是主教被激起了嫉妒之心。因此，在处理了其他一些事务后，若望剥夺了许多\OriginalTerm{诺洼天派}{Novatians}和\OriginalTerm{十四日派}{Quartodecimans}的教会，然后返回君士坦丁堡。他重新亲自照管自己特别管辖下的教会。但是没有人能忍受色拉皮雍的傲慢；因为他赢得了若望无限的信任和敬重，以至于得意忘形，对每个人都轻蔑以待。也正因如此，对主教的敌意更加炽烈。有一次，当塞味黎安从他身边经过时，色拉皮雍未向他致以主教应有的敬意，而是继续坐着（没有起立），清楚地表明他多么不在乎塞味黎安的出现。塞味黎安无法耐心容忍这种（所谓的）粗鲁和轻蔑，便大声对在场的人说：\Quote{如果色拉皮雍作为基督徒而死，基督就没有降生成人。}色拉皮雍借此言论公开煽动金口若望对塞味黎安的敌意：他删去了句子中的条件从句\Quote{如果色拉皮雍作为基督徒而死}，而声称塞味黎安作出了断言\Quote{基督没有降生成人}，并带来了他自己一方的几位证人支持这一指控。但皇后\Person{欧多克霞}得知此事后，严厉斥责了若望，并命令立即将塞味黎安从\Place{比提尼亚}的\Place{加采东}召回。他立即返回；但若望不愿与他有任何来往，也不听任何人的劝说，直到最后皇后欧多克霞本人在称为\Emph{宗徒教堂}的教堂里，将她儿子\Person{德奥多西}（如今幸运地统治着帝国，但当时还是个婴儿）放在若望膝前，再三以她年幼的儿子小王储的名义恳求若望，才好不容易使他与塞味黎安和好。这样，他们表面上和好了；但彼此之间仍然怀有怨恨。这就是[\Person{若望}]对塞味黎安敌意的起因。\par

% structure: p0032 source=h2 target=subsection reason=relative-heading-rank; base=h2

\subsection{第十二章 \Person{厄丕法尼}为了取悦\Person{德奥菲禄}，在\Place{君士坦丁堡}未经\Person{若望}许可执行祝圣。}

不久之后，在\Person{德奥菲禄}的建议下，\Person{厄丕法尼}主教再次从\Place{塞浦路斯}来到\Place{君士坦丁堡}；他还带来了一份主教会议法令的副本，该法令并未革除\Person{奥利振}本人的教籍，但谴责了他的著作。到达距城七英里的\Emph{圣若望教堂}时，他下船，在那里举行了礼仪；然后祝圣了一位执事后，再次进城。为了讨好德奥菲禄，他谢绝了若望的好意，租住在一间私宅。随后，他召集当时在都城的那些主教，出示了他那份谴责奥利振著作的主教会议法令副本，并在他们面前宣读；他无法为这一判决提出任何理由，只能说德奥菲禄和他自己认为应该拒绝它们。有些人出于对厄丕法尼的尊敬，签署了法令；但许多人拒绝签署，其中包括\Place{斯基泰}主教\Person{德奥提摩}，他对厄丕法尼说：\Quote{我不愿意侮辱很久以前就虔诚结束生命者的记忆；也不敢犯下如此不敬的行为，即谴责我们的前人未曾拒绝的东西；尤其是我知道奥利振的著作中并不含有邪恶的教义。}说完这话，他拿出该作者的一部著作，读了其中的几段，表明所提出的观点与正统信仰完全一致。他接着补充说：\Quote{那些诋毁这些著作的人，在无意中污辱了这些原则所源自的圣经。}这就是德奥提摩——一位以虔诚和正直生活著称的主教——对厄丕法尼的回答。\par

% structure: p0034 source=h2 target=subsection reason=relative-heading-rank; base=h2

\subsection{第十三章 作者为\Person{奥利振}辩护}

然而，由于吹毛求疵的中伤者蒙蔽了许多人，并成功地阻止他们阅读\Person{奥力振}的作品，仿佛他是亵渎性的作者，我认为现在对他略加评论并非不合时宜。无价值的人物，以及那些自身无能获得卓越声誉之人，往往通过贬低胜过他们的人来引人注目。首先是里基亚地区\Place{奥林彼斯}城的主教\Person{美多迪乌斯}患有此症；其次是曾在\Place{安提约基雅}短暂管理教会的\Person{欧斯塔提乌斯}；然后是\Person{阿波利纳里斯}；最后是\Person{特奥菲鲁斯}。这四位诋毁者诽谤了\Person{奥力振}，但并非出于相同理由：一人找到一种指控他的原因，另一人找到另一种；因此每个人都表明，凡他未提出异议的，他都完全接受了。因为既然有人特别攻击某一观点，另一人批评另一观点，显然每个人都承认了他未攻击的内容为真，对他未进攻的部分给予默许的赞同。\Person{美多迪乌斯}确实，在多次痛斥\Person{奥力振}之后，后来仿佛收回了他之前所说的一切，在一部名为\WorkTitle{Xenon}的对话录中表达了对他的钦佩。但我断言，从这些人的谴责中，\Person{奥力振}获得了更高的赞誉。因为那些探求了他们认为在他身上该受斥责的一切、却从未指控他持有关于至圣天主圣三的错误观点的人，由此最明确地表明为他正统的虔诚作证：通过不在这一点上责备他，他们用自己的证词赞扬了他。但是，\OriginalTerm{同质}{consubstantiality}论的捍卫者\Person{亚大纳西}，在其\WorkTitle{驳亚略人论集}中不断引用这位作者作为他自己信仰的见证，将他的话与自己的话交织在一起，并说：\Quote{最值得钦佩且勤勉的\Person{奥力振}，以他自己的证词确认了我们关于天主圣子的教义，肯定他与圣父\OriginalTerm{同永恒}{co-eternal}。}因此，那些将耻辱加于\Person{奥力振}的人，忽视了他们的咒骂同时落在赞扬\Person{奥力振}的\Person{亚大纳西}身上。为\Person{奥力振}辩护至此已足够；我们现在回到历史的进程。\par

% structure: p0036 source=h2 target=subsection reason=relative-heading-rank; base=h2

\subsection{第14章。\Person{厄丕法尼乌斯}被邀与\Person{若望}会面；拒绝后，他因违反教规的行为受到告诫；为此惊恐，他离开了\Place{君士坦丁堡}。}

\Person{若望}并未因\Person{厄丕法尼乌斯}违反教规在其教会内举行了祝圣而生气，反而邀请他留在主教府。但他回答说，除非\Person{若望}驱逐\Person{狄奥斯柯鲁斯}及其弟兄们出城，并亲笔签署谴责\Person{奥力振}著作的判决书，否则他既不会留下，也不会与他一同祈祷。现在\Person{若望}推迟了这些事，说在普世会议审议之前不应草率行事，\Person{若望}的对手们便引导\Person{厄丕法尼乌斯}采取另一种方式。因为他们策划，在名为\Place{宗徒}的教堂里举行集会时，\Person{厄丕法尼乌斯}站出来，在全体民众面前谴责了\Person{奥力振}的著作，绝罚了\Person{狄奥斯柯鲁斯}及其追随者，并指控\Person{若望}纵容他们。这些事报告给了\Person{若望}；于是第二天，当\Person{厄丕法尼乌斯}刚进入教堂时，\Person{若望}就向他发送了以下讯息：\par

\Quote{厄丕法尼乌斯，你做了许多违反教规的事。首先，你在属我管辖的教会内进行了祝圣；然后，未经我任命，你自行在这些教会中行使职务。此外，以前我邀请你来此处时，你拒绝前来，现在你却自行其是。因此你要当心，免得在民众中激起骚乱时，你自己也因而遭遇危险。}\par

\Person{厄丕法尼乌斯}听到这些告诫后惊慌起来，离开了教堂；在指控了\Person{若望}许多事后，他启程返回\Place{塞浦路斯}。有人说，他临行前对\Person{若望}说：\Quote{我希望你不会以主教身份死去}；\Person{若望}回答说：\Quote{不要指望回到你自己的国家。}我不能确定那些向我报告这些事的人说的是否真实；但尽管如此，两人的结局与上述预言相符。因为\Person{厄丕法尼乌斯}没有到达\Place{塞浦路斯}，在航行途中死于船上；而\Person{若望}在不久之后被逐出了他的教座，正如我们将在继续叙述中所示。\par

% structure: p0040 source=h2 target=subsection reason=relative-heading-rank; base=h2

\subsection{第15章。\Person{若望}因指责妇女而被\Place{迦克墩}会议逐出教会。}

当\Person{圣厄丕法尼}离去后，有人告知\Person{圣若望·金口}，说\Person{欧多克西亚}皇后曾煽动\Person{圣厄丕法尼}反对他。\Person{圣若望}性情如火，口才敏捷，不久后便在公开讲道中抨击女性。民众立刻认为这是影射皇后，于是这番话被心怀恶意者抓住，报告给当权者。皇后得知后，立即向丈夫诉苦，称对她的侮辱同样是对他的侮辱。皇帝于是授权\Person{德奥斐禄}立即召集一次主教会议反对\Person{圣若望}；\Person{塞维黎安}也协力推动此事，因为他仍对\Person{圣若望}怀恨在心。不久后\Person{德奥斐禄}抵达，他引领了数位来自不同城市的主教同行；不过这些主教也是奉皇帝之命召集的。\Person{圣若望}在前往\Place{厄弗所}并祝圣\Person{赫辣克利德}时，曾罢免了\Place{亚洲}的许多主教。因此他们预先在\Place{彼提尼雅}的\Place{加采东}集会。当时\Place{加采东}的主教是\Person{奇黎诺}，一位埃及人，他向主教们说了许多贬损\Person{圣若望}的话，斥责他为\Quote{不敬者}、\Quote{傲慢者}、\Quote{铁石心肠者}。他们对此非常满意。但\Place{美索不达米亚}的主教\Person{玛路塔}无意中踩了\Person{奇黎诺}的脚，\Person{奇黎诺}剧痛难忍，未能与其余人一同渡海前往\Place{君士坦丁堡}，只得留在\Place{加采东}。其余人渡过了海。\Person{德奥斐禄}如此公开表明对\Person{圣若望}的敌意，以致没有一个神职人员出去迎接他，或向他致以任何敬意；但碰巧有些\Place{亚历山大港}的水手在场——因为那时运粮船正在那里——以欢呼声迎接了他。他推辞不进入教堂，而住进了一座名为\Quote{普拉基迪亚}的皇家庄园。于是，针对\Person{圣若望}的控告如潮水般涌来；因为此时已不再提及\Person{奥力振}，所有人都专注于提出各种指控，其中许多是荒谬可笑的。初步事务就这样安排妥当后，主教们在\Place{加采东}的一个郊区集会，那个地方叫\Place{橡树}，并立即传唤\Person{圣若望}来回答对他的指控。他也传唤了执事\Person{塞拉丕翁}；阉人长老\Person{提格黎斯}和诵读员\Person{保禄}也被要求与他一同到场，因为这些人被列在控告中，被视为他的共犯。\Person{圣若望}向传唤他的人提出异议，认为他们是他的敌人，因此拒绝出席，并要求召开一次大公会议；但他们立即连续四次重复传唤；他坚持拒绝以他们为审判者，总是给出同样的回答，于是他们判决他有罪，罢免了他，除拒绝服从传唤外，没有为罢免提出任何其他理由。这一决定在傍晚宣布，激起了民众极为危险的骚动；他们彻夜守护，绝不允许将他从教堂带走，并高喊他的案件应由更庞大的集会裁决。然而，皇帝的敕令命令立即驱逐他，流放他处；\Person{圣若望}得知后，在判决后的第三天中午，趁民众不知情时自动投降：因为他担心因自己而发生任何暴动，于是被带走了。\par

% structure: p0042 source=h2 target=subsection reason=relative-heading-rank; base=h2

\subsection{第十六章 因\Person{圣若望·金口}被流放而起的骚动。他被召回。}

民众变得难以忍受地骚动；如同此类情况中常见的那样，许多先前对\Person{圣若望}怀有敌意的人，此时将敌意转为同情，并称他们不久前还希望他被罢免的那个人，是被诽谤的。因此，许多人开始同时反对皇帝和主教会议；但阴谋的起源他们更特别指向\Person{德奥斐禄}。因为他的欺诈行为再也无法隐藏，许多其他迹象暴露了这一点，尤其是他在\Person{圣若望}被罢免后立即与\Person{狄奥斯克洛}以及被称为\Quote{长修士}的人共融一事。但\Person{塞维黎安}在教堂讲道时，认为这是抨击\Person{圣若望}的合适时机，他说：\Quote{即使\Person{圣若望}没有因其他事而被定罪，他举止的傲慢也足以构成罢免他的罪行。人确实会得到其他所有罪的赦免，但\BibleQuote{天主拒绝骄傲的人}，正如圣经教导我们的。}这些指责使民众更加反对；于是皇帝下令立即召回他。于是，皇后的一名太监\Person{布里索}被派去追他，在\Place{普莱奈屯}——一个面对\Place{尼科美底亚}的商业城镇——找到他，并带他返回\Place{君士坦丁堡}。\Person{圣若望}被召回后，拒绝进城，宣称除非他的清白得到更高法庭的承认，否则他决不进城。于是他留在了一个名为\Place{玛利亚奈}的郊区。他在那里拖延时，骚动加剧，民众对统治者说出非常愤怒和辱骂的话；为了平息他们的愤怒，\Person{圣若望}被迫继续前行。在路上，一大群人怀着敬意和荣耀，立即引他进入教堂；在那里，他们恳求他坐到主教座位上，赐予他们惯常的祝福。他试图推辞，说\Quote{这应当由审判我的法官下令，必须由那些定我罪的人首先撤销他们的判决}，但民众却更加热切地希望看到他复职，并再次听他对他们讲话。民众最终说服他重新坐下，按惯例为他们祈求和平；之后，在同样的压力下，他向他们讲道。\Person{圣若望}的这一顺从给他的对手提供了另一个指控的理由；但那时他们没有就此采取行动。\par

% structure: p0044 source=h2 target=subsection reason=relative-heading-rank; base=h2

\subsection{第十七章 君士坦丁堡人与亚历山大里亚人之间关于\Person{赫拉克利德斯}的冲突；\Person{塞奥菲鲁斯}及其同党主教的逃亡。}

于是，首先，\Person{塞奥菲鲁斯}试图调查\Person{赫拉克利德斯}的祝圣案件，以便若有可能，借此找到再次罢免\Person{若望}的借口。\Person{赫拉克利德斯}并未出席此次审查。尽管如此，他仍被缺席审判，罪名是不公正地殴打了一些人，随后用锁链拖着他们穿过\Place{厄弗所}城中心。当\Person{若望}及其追随者抗议对缺席者判决的不公正时，\Place{亚历山大里亚}人则主张，即使\Person{赫拉克利德斯}不在场，也应听取控告者的陈述。因此，\Place{亚历山大里亚}人与\Place{君士坦丁堡}人之间爆发了激烈的争执，并引发了骚乱，许多人受伤，一些人丧生。\Person{塞奥菲鲁斯}看到事态发展，便不告而别，逃往\Place{亚历山大里亚}；其他主教，除了少数支持\Person{若望}者，也效仿他，返回各自教区。这些事件之后，\Person{塞奥菲鲁斯}在众人眼中声名狼藉；但他继续无耻地阅读\Person{奥力振}著作的做法，更使人们对他的憎恶大大增加。当有人问他为何如此支持他曾公开谴责的东西时，他回答说：\Quote{\Person{奥力振}的书就像一片草地，点缀着各种花朵。因此，如果我偶然在其中找到一朵美丽的，我便采摘它；但在我看来有刺的东西，我便跨过去，免得被刺。}然而\Person{塞奥菲鲁斯}回答时，没有思想智者\Person{撒罗满}的箴言\BibleRef{训 12:11}：\Quote{明智人的话如同刺棒}；那些被其中训诫刺痛的人，不该反踢它们。因此，\Person{塞奥菲鲁斯}被所有人轻视。\Person{迪奥斯库鲁斯}，\Place{赫尔莫波利斯}主教，被称为\Quote{高修士}者之一，在\Person{塞奥菲鲁斯}逃亡后不久去世，并受到隆重葬礼，被安葬在\Quote{橡树}教堂，那里曾因\Person{若望}的事而召开过会议。与此同时，\Person{若望}勤奋地致力于讲道。他祝圣\Person{塞拉皮翁}为\Place{色雷斯}的\Place{赫拉克利亚}主教，正是因他而引发了对\Person{若望}的憎恨。不久之后，发生了以下事件。\par

% structure: p0046 source=h2 target=subsection reason=relative-heading-rank; base=h2

\subsection{第十八章 关于\Person{欧多克西亚}的银像；因此\Person{若望}第二次被流放。}

此时，皇后\Person{欧多克西亚}的一座银像，身披长袍，被竖立在斑岩柱上，柱下有高大的基座支撑。这雕像既不靠近也不远离名为\Place{索菲亚}的教堂，而是相隔半条街宽。在这雕像处习惯举行公开游戏；\Person{若望}认为这是对教堂的侮辱，他恢复了往常的直言快语，用口舌攻击那些容忍这些游戏的人。然而，用恳求请愿的方式劝导当局停止游戏本是合宜的，他却没有这样做，反而使用辱骂性的语言嘲笑那些下令举行这类活动的人。皇后再次将他的话视为对她本人的明显蔑视：于是她设法召集另一个主教会议来反对他。当\Person{若望}得知此事，他在教堂发表了那篇著名的演讲，开头的话是：\Quote{\Person{黑落狄雅}再次狂怒；她再次不安；她再次跳舞；再次渴望在盘子里得到\Person{若翰}的头。}这当然更加激怒了皇后。不久之后，下列主教到达：亚细亚省\Place{安奇拉}的主教\Person{良提乌}、丕息狄雅省\Place{劳狄刻雅}的主教\Person{阿摩尼乌}、色雷斯省\Place{斐理伯}的主教\Person{布里索}、叙利亚省\Place{贝洛雅}的主教\Person{阿卡基乌}，以及其他几位。\Person{若望}无畏地出现在他们面前，要求调查针对他的指控。但当我们救主的诞辰周年纪念来临，皇帝没有像往常一样前往教堂，而是派人给\Person{金口若望}带信，大意是除非他澄清自己所面临指控的罪行，否则皇帝不会与他一同领受共融。由于\Person{若望}保持着勇敢和炽热的态度，而他的控告者似乎变得胆怯，在场的众主教撇开所有其他事情，说他们只局限于这一考虑：他是在被废黜之后，未经教会议会授权，自行又坐回主教座位。当他声称与他共融的六十五位主教已经恢复了他的职务时，良提乌的党羽反对说：\Quote{在主教会议中，有更多的人投票反对你，\Person{若望}。}但虽然\Person{若望}当时争辩说，这是亚略派的教规，而非大公教会的教规，因此对他无效——因为这是在\Place{安提约基雅}为反对\Person{亚太纳削}而召开的会议中制定的，目的是为了推翻\OriginalTerm{同质}{consubstantiality}教义——但众主教不愿听取他的辩护，立刻判他有罪，没有考虑到使用这条教规就等于认可了\Person{亚太纳削}本人被废黜。这个判决在复活节前不久宣布；皇帝于是派人告诉\Person{若望}，他不能前往教堂，因为两次主教会议已经判他有罪。因此，\Person{金口若望}被禁止发言，不再前往教堂；但属于他派别的人在被称为君士坦提安的公共浴场庆祝复活节，此后离开了教会。其中有众多主教和司铎，以及教会其他等级的人，他们从那时起在不同地方分开举行集会，并因他得名\OriginalTerm{若望派}{Johannites}。两个月期间，\Person{若望}避免公开露面；此后，皇帝的一道谕令将他放逐。于是\Person{若望}被迫被带往放逐地，在他离开的那一天，一些若望派的人放火烧了教堂，由于强劲的东风，火势蔓延到元老院。这场火灾发生在六月二十日，正值\Person{霍诺留}第六次执政官任期，他与\Person{阿里斯泰内图斯}共同担任执政官。\Place{君士坦丁堡}总督\Person{奥普塔图斯}，宗教上为外教人，且憎恨基督徒，对\Person{若望}的朋友施加的严酷手段，以及他因为这次纵火行为处死了其中许多人，我相信，我应当略过不提。\par

% structure: p0048 source=h2 target=subsection reason=relative-heading-rank; base=h2

\subsection{\Person{阿尔萨修斯}被任命为\Person{若望}的继任者。\Place{卡尔西顿}主教\Person{奇里努斯}身体不适。}

几天之后，\Person{阿尔萨修斯}被祝圣为\Place{君士坦丁堡}主教；他是\Person{聂克塔略}的兄弟，聂克塔略在\Person{若望}之前非常能干地管理了这个教区，但阿尔萨修斯当时年事已高，年逾八十。当他非常温和和平地管理主教职务时，\Place{卡尔西顿}主教\Person{奇里努斯}，\Place{美索不达米亚}主教\Person{马鲁塔斯}无意中踩了他的脚，这意外使他受到严重影响，导致坏疽，不得不截去他的脚。这次截肢并非只进行一次，而是需要多次重复：因为受伤的肢体被切除后，病患蔓延至全身，另一只脚也受到疾病影响，不得不接受同样的手术。我提及此事，是因为许多人断言，他所受的痛苦是对他诽谤\Person{若望}的惩罚，他常称\Person{若望}为傲慢和无情，正如我已说过的。此外，在上述执政官任期内的九月三十日，\Place{君士坦丁堡}及其郊区发生了一场异常的大冰雹，冰雹尺寸巨大，这也被说成是神圣义怒的表达，是因为\Person{金口若望}不公正地被废黜：而皇后的死亡使这些传言更加可信，因为她在冰雹袭击后四天去世。然而，另一些人断言，\Person{若望}被废黜是罪有应得，因为他在\Place{亚细亚}和\Place{吕底亚}行使暴力，剥夺了\OriginalTerm{诺瓦替安派}{Novatians}和\OriginalTerm{十四日派}{Quartodecimans}的许多教堂，当时他前往\Place{厄弗所}并祝圣了\Person{赫拉克利德斯}。但是，\Person{若望}被废黜是否公正，如他的敌人所说；或\Person{奇里努斯}是否因他诽谤辱骂而受惩罚；冰雹是否下降，或皇后是否因\Person{若望}而死，或者这些事是出于其他原因，或是这些原因与其他原因共同作用，只有天主知道，祂是洞悉秘密者，是真理的公正审判者。我只是记录那时流传的说法。\par

% structure: p0050 source=h2 target=subsection reason=relative-heading-rank; base=h2

\subsection{\Person{阿尔萨修斯}之死与\Person{阿提库斯}的任命。}

但是\Person{阿尔萨西乌斯}继任主教后不久便去世了；他在下一个执政官时期，即\Person{斯提里科}第二次执政、\Person{安特弥乌斯}第一次执政期间，于11月11日去世。由于主教职位变得令人向往，许多人觊觎这个空缺的教座，因此在选举继任者之前过去了很长时间；但最终在随后的执政官时期，即\Person{阿卡狄乌斯}第六次执政、\Person{普罗布斯}第一次执政期间，一位虔诚的人，名叫\Person{阿提库斯}，被提升为主教。他是\Place{亚美尼亚}的\Place{塞巴斯蒂亚}本地人，自幼便过着苦修生活；此外，除了有适度的学识外，他还拥有大量的自然审慎。但我将在稍后更详细地谈论他。\par

% structure: p0052 source=h2 target=subsection reason=relative-heading-rank; base=h2

\subsection{第二十一章 \Person{若望}死于流放}

\Person{若望}被流放，于9月14日死在\Place{攸克辛}的\Place{科马纳}，当时是下一个执政官时期，即\Person{霍诺留}第七次执政、\Person{狄奥多西}第二次执政。正如我们之前所观察到的，他是一个因对节制的热忱而倾向于愤怒而非宽容的人；他个人的圣德品格使他纵容一种对他人来说无法容忍的言论自由。确实，对我来说最不可理解的是，他怀着如此炽热的自我克制和品行无瑕的热忱，却在讲道中似乎教导一种松散的节制观。因为主教会议只接受那些在领洗后犯罪的人悔改一次；他却毫不犹豫地说：\Quote{来接近吧，即使你已经悔改了一千次。}对于这个教导，甚至他的许多朋友都指责他，尤其是\OriginalTerm{诺瓦提安派}{Novatian}的主教\Person{西西尼乌斯}；他写了一本书谴责\Person{金口}的上述言论，并严厉斥责他。但这发生在很久以前。\par

% structure: p0054 source=h2 target=subsection reason=relative-heading-rank; base=h2

\subsection{第二十二章 论\OriginalTerm{诺瓦提安派}{Novatian}的主教\Person{西西尼乌斯}。他应对的机敏}

我认为在此处对\Person{西西尼乌斯}作一些记述并非不合适。正如我经常所说，他是一个非常雄辩的人，并且精通哲学。但他尤其修习了逻辑学，并精通圣经的解释；以至于异端者\Person{欧诺米乌斯}常常回避他推理所展现的敏锐。至于他的饮食，他并不简朴；因为他虽然实行最严格的节制，但他的餐桌总是丰盛地摆设。他还习惯于放纵自己穿白色衣服，每天在公共浴池洗两次澡。当有人问他\Quote{为何你一个主教每天洗两次澡？}他回答说：\Quote{因为洗三次不方便。}有一天，他出于礼貌去拜访\Person{阿尔萨西乌斯}主教，被那位主教的一个朋友问道：\Quote{为什么你穿一件如此不适合主教的衣服？哪里写着神职人员应穿白色？}他说：\Quote{你先告诉我，哪里写着主教应穿黑色？}当询问者不知道如何回答这个反问时，\Person{西西尼乌斯}反驳说：\Quote{你不能证明司铎应穿黑色。但\Person{撒罗满}是我的权威，他的劝言是：\BibleQuote{愿你的衣服常白。}\BibleRef{训 9:8}而且我们的救主在福音中显现时穿着白衣；此外，他向宗徒们显示了\Person{梅瑟}和\Person{厄里亚}，也都穿着白衣。}他对这些及其他问题的迅速回答引起了在场之人的钦佩。再次，当\Place{小加拉提亚}的\Place{安卡拉}主教\Person{良提乌斯}（他曾从\OriginalTerm{诺瓦提安派}{Novatian}手中夺走一座教堂）访问\Place{君士坦丁堡}时，\Person{西西尼乌斯}去找他，恳求他归还教堂。但他粗暴地接待了他，说：\Quote{你们\OriginalTerm{诺瓦提安派}{Novatian}不应拥有教堂；因为你们剥夺了悔改，并且拒绝天主的仁慈。}当\Person{良提乌斯}说出这些以及许多其他对\OriginalTerm{诺瓦提安派}{Novatian}的辱骂时，\Person{西西尼乌斯}回答说：\Quote{没有人比我更真诚地悔改。}当\Person{良提乌斯}问他\Quote{你为何悔改？}他说：\Quote{因为我来看你。}有一次，\Person{若望}主教与他争论，说：\Quote{这城不能有两个主教。}\Person{西西尼乌斯}说：\Quote{也没有。}\Person{若望}被这个回答激怒，说：\Quote{你看，你假装只有你一个人是主教。}\Person{西西尼乌斯}反驳说：\Quote{我没有那样说；只是在你眼中我不是主教，而我在别人眼中是。}\Person{若望}对这个回答更加恼火，说：\Quote{我要阻止你讲道；因为你是个异端者。}\Person{西西尼乌斯}幽默地回答说：\Quote{如果你能让我免除如此艰巨的职责，我会给你奖赏。}\Person{若望}被这个回答稍微缓和了一点，说：\Quote{如果你觉得讲道这么麻烦，我不会让你停止讲道。}\Person{西西尼乌斯}就是这样风趣，如此善于应对；但详细谈论他的俏皮话会令人厌烦。因此，通过几个例子我们已经说明了他是什么样的人，认为这些已足够。我只补充说，他以博学著称，因此所有继任的主教都爱戴并尊敬他；不仅他们，而且元老院的领导成员也都尊重和钦佩他。他写了许多著作；但它们过于追求文辞优雅，并大量掺杂诗意的表达。因此，他作为演说家比作为作家更受赞赏；因为他的面容和声音，以及他的形体和神态，都带有尊严，他的一举一动都很优雅。由于这些特征，他受到所有派别的爱戴，并且特别受到\Person{阿提库斯}主教的青睐。但我必须结束对\Person{西西尼乌斯}的简短记述。\par

% structure: p0056 source=h2 target=subsection reason=relative-heading-rank; base=h2

\subsection{第二十三章 \Person{阿卡狄乌斯}皇帝驾崩}

\Person{若望}死后不久，\Person{阿卡狄乌斯}皇帝也去世了。这位君王性情温和宽厚，晚年时因下述事迹被认为深受天主宠爱。\Place{君士坦丁堡}有一座巨大的府邸，称为\Place{卡里亚}；因为它的庭院中有一棵胡桃树，据说\Person{阿卡西乌}正是在这棵树上被吊死而殉道；因此，其附近建了一座小堂，\Person{阿卡狄乌斯}皇帝有一天认为应当去参观，并在那里祈祷后离去。所有住在这座小堂附近的人都蜂拥而出观看皇帝；有些人从那座府邸出来，尽力占据街道以便更清楚地看到他们的君主及其随从，另一些人则跟在他的队伍后面，直到所有住在府邸里的人，包括妇女和儿童，全都走了出来。这座巨大的建筑刚空无一人（其建筑完全环绕着教堂），整座建筑便轰然倒塌。于是响起一片惊叫声，随后是赞叹的欢呼，因为人们相信皇帝的祈祷拯救了如此众多的人免于毁灭。此事就这样发生了。五月一日，\Person{阿卡狄乌斯}去世，留下他年仅八岁的儿子\Person{狄奥多西}，时值\Person{巴苏斯}与\Person{菲利普}任执政官，第297届奥林匹克运动会第二年。他与父亲\Person{狄奥多西}共同统治了十三年，在其父去世后又独自统治了十四年，享年三十一岁。本书涵盖的时间跨度为十二年零六个月。\par

\SourceRecord{Translated by A.C. Zenos.；Nicene and Post-Nicene Fathers, Second Series；Vol. 2.；Edited by Philip Schaff and Henry Wace.；Buffalo, NY: Christian Literature Publishing Co.,；1890.；<http://www.newadvent.org/fathers/26016.htm>.}

% structure: p0001 source=h1 target=section reason=page-title

\section{教会史（第七卷）}

% structure: p0002 source=h2 target=subsection reason=relative-heading-rank; base=h2

\subsection{第一章 近卫总长安特弥乌斯代为年轻的狄奥多西管理东方政务。}

阿卡狄乌斯于巴苏斯与菲利普执政期间五月一日去世后，其弟霍诺留仍治理帝国西部，而东方行政权则落于其子、年仅八岁的狄奥多西二世之手。因此，公共事务交由近卫总长安特弥乌斯管理，他是那位菲利普的孙子；该菲利普在君士坦提乌斯统治时期，曾将保禄逐出君士坦丁堡教区，并扶植马其顿尼乌斯取而代之。安特弥乌斯下令用高墙环绕君士坦丁堡。他被视为——且确实是——当时最审慎之人，凡事鲜有鲁莽，常与最明智之友商议政务，尤其与智术师特洛伊鲁斯共商，后者既精于哲学造诣，在政治智慧上亦与安特弥乌斯相当。因此，几乎所有事务都得到特洛伊鲁斯的赞同。\par

% structure: p0004 source=h2 target=subsection reason=relative-heading-rank; base=h2

\subsection{第二章 君士坦丁堡主教阿提库斯的品格与行为。}

当皇帝狄奥多西八岁时，阿提库斯已担任君士坦丁堡教会主教第三年，此人如前所述，以学识、虔诚与明辨著称，故其治下教会极为兴盛。他不仅团结了\Quote{信德之家}，见：\BibleRef{迦 6:10}，而且以其智慧令异端者钦佩；他绝不愿骚扰他们，但若有时不得不以威吓震慑他们，随后便对他们显出温和与宽仁。他亦不忽视学问，勤读古人著作，常彻夜不倦，故不为哲学家的推理与智术师的诡辩所惑。此外，他言谈可亲、引人入胜，随时同情受苦者；简言之，用宗徒的话说，\Quote{对一切人，我就成为一切}，见：\BibleRef{格前 9:22}。早年任司铎时，他惯于将撰好的讲道词背下，然后在教堂背诵；但因勤勉练习，他获得了自信，能即席教导且口若悬河。然而，他的讲道并不常得听众热烈喝采，也不值得记录下来。关于他的才华、学识与品行，兹述如上。现在必须叙述他任期内值得记录之事。\par

% structure: p0006 source=h2 target=subsection reason=relative-heading-rank; base=h2

\subsection{第三章 论西纳达主教狄奥多西与阿加佩图斯。}

有一位名叫狄奥多西的主教，管辖弗里吉亚·帕卡塔的\Place{西纳达}；他激烈迫害该省的异端者——人数甚多——尤其是马其顿派；不仅将他们逐出城市，乃至逐出乡间。他如此行并非出于正统教会的先例，亦非为传播真信仰，而是为贪图不义之财，受贪婪驱使，藉从异端者榨取钱财以聚敛。为此，他对马其顿派用尽各种手段，武装圣职人员，对他们施以无数计谋，且毫不迟疑地将他们交给世俗法庭。他尤其困扰他们的主教，名叫阿加佩图斯。因觉该省总督权力不足，未能按他所愿惩罚异端者，他便前往君士坦丁堡，向近卫总管请求更严厉的法令。当狄奥多西为此事外出时，阿加佩图斯——如前所述，领导马其顿派——想出一个明智审慎的结论。他与圣职人员商议，召集所有受他引导的民众，劝勉他们接受\OriginalTerm{同质}{homoousian}信仰。众人同意此议后，他立即前往教堂，不仅有他自己的人跟随，更有全体民众。他在那里祈祷后，坐上狄奥多西惯用的主教座位，从此宣讲同质道理，重新团结民众，并控制了西纳达教区的各教堂。不久后，狄奥多西回到\Place{西纳达}，带着从近卫总管处获得的更大权力，对已发生之事一无所知，便径直走进教堂。他立即被一致驱逐，于是再次前往君士坦丁堡。到达后，他向主教阿提库斯诉说自己所受的待遇及被剥夺主教职位的经过。阿提库斯看出这次运动对教会有益，便尽力安慰狄奥多西，劝他以知足之心安度退隐生活，为公益牺牲私利。然后他致信阿加佩图斯，授权他保留主教职位，并命他无需因狄奥多西的抱怨而忧惧被骚扰。\par

% structure: p0008 source=h2 target=subsection reason=relative-heading-rank; base=h2

\subsection{第四章 阿提库斯在圣洗中治愈一名瘫痪的犹太人。}

这是阿提库斯管理期间教会状况的一项重要改善。这个时代也未缺乏奇迹与治愈的证明。有一个犹太人，患瘫痪病，卧床多年；各种医术和犹太同胞的祈祷都试过了，但毫无效果，他最终求助于基督教的圣洗，相信这是唯一真正的救治方法。当主教阿提库斯得知他的意愿后，就教导他基督教真理的初步知识，并向他说教要寄望于基督，然后吩咐将他连床抬到洗礼池。那瘫痪的犹太人怀着真诚的信德领受圣洗，一从洗礼池中出来，就发现自己完全痊愈了，此后一直身体健康。基督恩赐这奇迹甚至在我们这个时代显现；它的名声使许多外教人相信并领受了圣洗。但犹太人虽热切\InnerQuote{追求神迹}\BibleRef{格前 1:22}，就连实际发生的神迹也未诱使他们接受信仰。基督就是这样赐福给人类。\par

% structure: p0010 source=h2 target=subsection reason=relative-heading-rank; base=h2

\subsection{第五章 司铎萨巴提乌斯，原为犹太人，自\OriginalTerm{诺瓦替派}{Novatians}分离。}

然而，许多人并不重视这些事，反而顺从自己的败坏；因为不仅这奇迹之后犹太人仍不信，而且其他喜欢追随他们的人也表现出与他们相似的观点。其中就有萨巴提乌斯，如前所述；他不满足于已获得的司铎尊位，从一开始就觊觎主教之职，便以守犹太人的逾越节为借口，自诺瓦替派教会分离。因此，他离开自己的主教西西尼乌斯，在名为色罗洛弗斯（阿尔卡狄乌斯广场现今所在之处）的地方举行分裂集会，并胆敢做出应受最重惩罚的行为。有一天，在一次集会中，他读到福音中那段话说：\BibleQuote{那时正是犹太人的逾越节，} 接着添加了从未写过或听过的话：\Quote{凡在无酵节以外庆祝逾越节的，应受诅咒。} 这些话在民众中传开，诺瓦替派较单纯的平信徒被这诡计所骗，纷纷投奔他。但他伪造的谎言对他毫无益处；因为他的伪造导致了最灾难性的后果。不久之后，他提前守了这节期，早于基督教的复活节；许多人照常蜂拥到他那里。当他们按惯常的守夜礼仪过夜时，一种仿佛由恶灵引起的恐慌降临到他们身上，好像他们的主教西西尼乌斯正带着许多人来攻击他们。由于在这种情况下可以预见的慌乱，以及他们夜间被困在狭窄的地方，互相践踏，以致七十多人被踩死。为此，许多人离开了萨巴提乌斯；但仍有一些人坚持他的愚昧偏见，留在他那里。萨巴提乌斯如何违背誓言，后来设法使自己受祝圣为主教，我们将在后面叙述。\par

% structure: p0012 source=h2 target=subsection reason=relative-heading-rank; base=h2

\subsection{第六章 此时\OriginalTerm{亚略派}{Arianism}的领袖们。}

\Person{多罗塞乌斯}曾是亚略派的主教，如前所述，他被该派从安提约基雅调任至君士坦丁堡，享年一百一十九岁，于霍诺里乌斯第七次执政、狄奥多西·奥古斯都第二次执政那年的11月6日去世。在他之后，\Person{巴尔巴斯}主持亚略派；在他任内，亚略派因拥有两位非常雄辩的成员而受惠，这两人都有司铎品级，一位名叫\Person{蒂莫提}，另一位叫\Person{乔治}。乔治擅长希腊文学；蒂莫提则精通圣经。乔治确实手中常拿着亚里士多德和柏拉图的著作；蒂莫提则从奥力振获得灵感；他在公开讲解圣经时，也表现出对希伯来语相当熟悉。蒂莫提先前曾认同\OriginalTerm{普萨提里派}{Psathyrians}；乔治则由巴尔巴斯祝圣。我本人曾与蒂莫提交谈，对他回答最困难问题、阐明神谕中最隐晦段落的敏捷深感惊讶；他还总是引用奥力振作为无可置疑的权威来证实自己的言论。但令我惊讶的是，这两人竟继续坚持亚略派的异端；一人如此精通柏拉图，另一人则时常引用奥力振。因为柏拉图并未说第二因和第三因（他通常这样称呼）有存在的开端；而奥力振到处承认圣子与圣父同等永存。然而，虽然他们仍与自己的教会保持联系，但他们无意中改善了亚略派，用他们的教导取代了许多亚略的亵渎言论。但关于这些人已说得够多了。诺瓦替派的主教西西尼乌斯在同一执政期内去世，\Person{克里桑图斯}被祝圣接替他，关于后者我们稍后还要叙述。\par

% structure: p0014 source=h2 target=subsection reason=relative-heading-rank; base=h2

\subsection{第七章 \Person{济利禄}接替\Person{狄奥菲鲁斯}为亚历山大里亚主教。}

\Place{亚历山大里亚}主教\Person{德奥斐洛}陷入昏睡状态，于\Person{奥诺留}第九任执政、\Person{德奥多西}第五任执政之年十月十五日去世。随后立即爆发了关于继任者的激烈争论：一些人试图将总执事\Person{提摩太}推上主教席位，另一些人则希望由\Person{德奥斐洛}的侄子\Person{济利禄}出任。为此民众中发生了骚动，\Person{阿班提乌斯}（\Place{埃及}驻军司令）支持\Person{提摩太}。然而，\Person{济利禄}的拥护者取得了胜利。于是，\Person{德奥斐洛}死后第三天，\Person{济利禄}便登上了主教之位，其权力甚至超过了\Person{德奥斐洛}以往所行使的。因为从那时起，\Place{亚历山大里亚}主教职超出了司铎职能的界限，开始掌管世俗事务。\Person{济利禄}随即关闭了\Place{亚历山大里亚}的诺瓦提安派教堂，没收了他们所有祝圣过的器皿和装饰品，然后剥夺了他们的主教\Person{德奥彭普托斯}的一切。\par

% structure: p0016 source=h2 target=subsection reason=relative-heading-rank; base=h2

\subsection{第八章 \Place{美索不达米亚}主教\Person{玛鲁塔}在\Place{波斯}人中传播基督信仰}

大约在同一时期，由于以下原因，基督信仰在\Place{波斯}得到了传播。\Place{波斯}与\Place{罗马帝国}的君主之间经常互派使节，此类往来从未间断。当时，因时势所迫，罗马皇帝认为应当派遣此前曾提及的\Place{美索不达米亚}主教\Person{玛鲁塔}出使\Place{波斯}王廷。\Place{波斯}王发现此人极有圣德，便以殊礼相待，视他如同真正蒙天主所爱之人。这引起了\Person{玛吉}（在\Place{波斯}君主面前颇有势力）的嫉妒，因为他们害怕\Person{玛鲁塔}会说服国王信奉基督信仰。原来，\Person{玛鲁塔}曾藉祈祷治好了国王长期受其困扰的剧烈头痛，而\Person{玛吉}却无力缓解。于是，巫师们便使用了这样的欺骗手段：由于\Place{波斯}人崇拜火，国王习惯在一座特定的建筑内向常年不熄的火致敬，他们便在一人藏在圣坛之下，命令他在国王惯常敬拜的时刻喊出下面的话：\Quote{国王当被驱逐，因为他犯了不敬之罪，竟以为一位基督徒司铎为神所爱。}\Person{伊斯迪格德斯}——这是国王的名字——听见这话，便决定遣走\Person{玛鲁塔}，尽管他心中对他仍怀敬意。但\Person{玛鲁塔}真是一位爱天主的人，他藉恳切祈祷识破了\Person{玛吉}的计谋。于是他走到国王面前，这样说道：\Quote{国王啊，不要受骗，}他说，\Quote{但当你再次进入那建筑、听见同样的声音时，请探查地下，就会发现其中的欺诈。因为说话的不是火，而是人的诡计。}国王采纳了\Person{玛鲁塔}的建议，照常去了那间常年燃火的屋子。当他再次听见同样的声音时，便命人挖开圣坛；于是那个冒充神明发言的骗子被发现了。国王对这种欺骗行为感到愤慨，下令将\Person{玛吉}一族处以什一之刑。这事之后，他允许\Person{玛鲁塔}随心所欲地建造教堂；从此基督宗教在\Place{波斯}人中传播开来。随后\Person{玛鲁塔}被召回，去了\Place{君士坦丁堡}；但不久之后，他再次被派往\Place{波斯}宫廷为使节。\Person{玛吉}又设计谋，想尽一切办法阻止国王接见他。他们其中一计是在国王惯常经过的地方制造一种极其恶臭的气味，然后嫁祸于基督徒。然而国王早已对\Person{玛吉}有所怀疑，便非常仔细、严密地调查了此事；再次查出了制造恶臭的人。于是惩罚了其中数人，并对\Person{玛鲁塔}更加尊敬。他对\Place{罗马}这个民族极为看重，非常珍视他们的善意。他几乎要亲自接受基督信仰，因为\Person{玛鲁塔}与\Place{波斯}主教\Person{阿布达斯}共同以另一项经验证明了其大能：二人藉着多次禁食和祈祷，赶走了附着在王子身上的魔鬼。但\Person{伊斯迪格德斯}的死阻止了他公开宣认基督信仰。于是王国传给了他的儿子\Person{瓦拉兰尼斯}，在其统治期间，\Place{罗马}与\Place{波斯}之间的条约被打破，此事我们稍后将予以叙述。\par

% structure: p0018 source=h2 target=subsection reason=relative-heading-rank; base=h2

\subsection{第九章 \Place{安提约基雅}和\Place{罗马}的主教}

在此期间，\Person{弗拉维安}去世后，\Person{波菲里}接受了\Place{安提约基雅}的主教职；之后，\Person{亚历山大}被立为那教会的主教。但在\Place{罗马}，\Person{达玛苏}担任该主教职十八年后，\Person{西利西乌}继任；\Person{西利西乌}在位十五年后，\Person{阿纳斯塔西乌}治理教会三年；\Person{阿纳斯塔西乌}之后，\Person{依诺增爵}被擢升为同一教座。他是在\Place{罗马}第一位迫害诺瓦提安派的人，夺取了他们许多教堂。\par

% structure: p0020 source=h2 target=subsection reason=relative-heading-rank; base=h2

\subsection{第十章 \Person{阿拉里克}攻陷并洗劫\Place{罗马}}

大约在同一时期，罗马被蛮族攻陷。原来有一个名叫\Person{阿拉里克}的蛮族人，他曾是罗马人的盟友，与\Person{德奥多西}皇帝一同在对抗篡位者\Person{尤金尼乌斯}的战争中担任盟友，因此被授予罗马的荣誉，但他无法承受这份好运。他无意称帝，而是离开\Place{君士坦丁堡}前往西方，抵达\Place{伊利里库姆}后立即蹂躏了整个地区。然而，在他行军时，色萨利人在\Place{佩纽斯河}河口抵挡了他，那里有一条翻越\Place{品都斯山}通往\Place{伊庇鲁斯}的\Place{尼科波利斯}的隘口；双方交战，色萨利人杀死了他约三千名士兵。此后，他身边的蛮族摧毁了一切，最终攻陷了\Place{罗马}本身，他们洗劫了该城，焚烧了城内最宏伟的建筑及其他令人赞叹的艺术品。金钱和贵重物品被掠夺瓜分。许多元老院要员被以各种借口处死。此外，\Person{阿拉里克}为了嘲弄皇权，册立了一个叫\Person{阿塔鲁斯}的人为皇帝，他命令此人在一日之内享有君主的一切标志，次日却穿着奴隶的服装示众。完成这些之后，他急忙撤退，因为有消息传来说\Person{德奥多西}皇帝已派兵来攻打他。这消息并非虚假；帝国军队确实在途中；但\Person{阿拉里克}没有等待谣言成真，便拔营逃跑了。据说，在他向\Place{罗马}进军时，一位虔诚的隐修士劝诫他不要乐此不疲地施行如此暴行，不要再以杀戮和流血为乐。\Person{阿拉里克}回答说：\Quote{我并非自愿走这条路；但有一种不可抗拒的力量每日催逼着我，说：\InnerQuote{前往\Place{罗马}，摧毁那座城。}}此人便是如此行径。\par

% structure: p0022 source=h2 target=subsection reason=relative-heading-rank; base=h2

\subsection{第十一章 罗马主教}

在\Person{依诺森}之后，\Person{佐西默}治理罗马教会两年；此后\Person{波尼法爵}主持了三年。\Person{塞莱斯定}接续了他。这位\Person{塞莱斯定}也剥夺了罗马诺瓦提安派的教堂，并迫使他们的主教\Person{鲁斯提库拉}在私人住宅中秘密聚会。在此之前，诺瓦提安派在罗马极其兴盛，拥有许多教堂，会众众多。但嫉妒也攻击了他们，一旦罗马主教职，如同亚历山大里亚的，扩展到教会管辖范围之外，堕落成现今的世俗统治状态。因为从此以后，主教们甚至不容许与他们信仰一致的人享有平安集会的特权，反而剥夺他们的一切，只称赞他们在信仰上的一致。君士坦丁堡的主教们则避免了这种行径；因为他们不仅容忍他们，允许他们在城内集会（如我已说过的），还以各种基督徒的关怀对待他们。\par

% structure: p0024 source=h2 target=subsection reason=relative-heading-rank; base=h2

\subsection{第十二章 论君士坦丁堡诺瓦提安派主教克里桑特}

\Person{西西尼乌斯}死后，\Person{克里桑特}被迫接受了主教职务。他是\Person{西西尼乌斯}的前任\Person{马尔西安}的儿子，早年曾在宫中担任军职，后来在\Person{德奥多西大帝}手下被任命为\Place{意大利}总督，之后又担任不列颠群岛的总督，在两种职务中都赢得了极高的赞誉。晚年回到\Place{君士坦丁堡}，热切希望被任命为该城的行政长官，却违背自己的意愿被立为诺瓦提安派主教。因为\Person{西西尼乌斯}临终时曾提到他是最适合坐镇该教座的人，民众视此宣告为法律，立即寻求给他祝圣。当\Person{克里桑特}试图避免这强加于他的尊位时，\Person{撒巴提乌斯}以为时机已到，可以夺取教会大权，他不顾自己曾起过的誓言，在几个微不足道的主教手中为自己举行了祝圣。其中有一个叫\Person{赫尔莫根}的，曾被\Person{撒巴提乌斯}本人因他亵渎性的著作而诅咒开除教籍。但\Person{撒巴提乌斯}的这一背誓行为对他毫无益处：因为民众厌恶他的喧闹，尽力寻找\Person{克里桑特}的藏身之处；发现他隐居在\Place{比提尼亚}后，便强行将他带回，授予他主教职。他是一个谦逊和明智无与伦比的人；因此他建立并扩大了君士坦丁堡的诺瓦提安派教会。此外，他是第一个用自己的私人财产向穷人分发金子的人。而且，除了每个主日两个祝圣过的饼之外，他不从教会收取任何东西。他如此急切地促进自己教会的利益，以致从\Person{特罗伊鲁斯}的学派中招来了当时最杰出的演说家\Person{阿布拉比乌斯}，并祝圣他为司铎；他的讲道辞流传甚广，极其优雅且切中肯綮。但\Person{阿布拉比乌斯}后来被提升为\Place{尼西亚}诺瓦提安派教会的主教，同时也在那里教授修辞学。\par

% structure: p0026 source=h2 target=subsection reason=relative-heading-rank; base=h2

\subsection{第十三章 亚历山大里亚基督徒与犹太人的冲突：主教济利禄与行政长官奥勒斯提斯之间的决裂}

约在同一时期，犹太人居民被主教\Person{济利禄}逐出\Place{亚历山大}，原因如下。\newline{}亚历山大民众比任何其他民族更喜爱骚乱；一旦找到借口，就会爆发最不可容忍的暴行；因为他们的骚动从不以流血告终。\newline{}当下发生了一场民众骚动，并非出于任何重要原因，而是源于一种在几乎所有城市都流行的恶习，即对舞蹈表演的迷恋。\newline{}由于犹太人在安息日放下事务，不是去聆听法律，而是消磨于戏剧娱乐，舞者通常在当日吸引大量人群，几乎总是引发混乱。\newline{}尽管亚历山大总督在一定程度上控制了这种情况，犹太人却继续反对这些措施。\newline{}尽管他们一向敌视基督徒，但出于对舞者的缘故，他们对基督徒的反对更加激烈。\newline{}因此，当总督\Person{奥勒斯特}在剧场发布谕令——他们习惯这样称呼公告——以规范演出时，主教\Person{济利禄}的一些同党在场，欲知即将发布的命令内容。\newline{}其中有位\Person{希拉克斯}，是初级文学教师，也是主教\Person{济利禄}讲道极为热忱的听众，因急于鼓掌而引人注目。\newline{}犹太人见此人出现在剧场，立即喊道，他来这里无非是要在民众中煽动叛乱。\newline{}奥勒斯特长久以来忌惮主教们日益增长的权力，因为他们侵犯了皇帝所任命官员的管辖权，尤其是济利禄企图监视他的行动；因此他下令逮捕\Person{希拉克斯}，并在剧场公开施以酷刑。\newline{}济利禄得知此事后，召见犹太人的首领，威胁说若不停止对基督徒的骚扰，将给予最严厉的惩罚。\newline{}犹太民众听到这些威胁后，非但没有收敛暴力，反而更加狂怒，并密谋消灭基督徒；其中一项阴谋极为极端，最终导致他们被全部逐出\Place{亚历山大}；我现在就来叙述此事。\newline{}他们约定每人手指上戴一枚用棕榈枝皮制成的戒指，以便相互识别，然后决定夜间袭击基督徒。\newline{}于是他们派人到街上呼喊，说亚历山大教堂着火了。\newline{}许多基督徒闻讯，从四面八方急忙跑出，焦虑万分地要去救他们的教堂。犹太人立刻扑上去杀死他们，凭借戒指轻易识别出彼此。\newline{}天亮后，这场暴行的肇事者无法隐藏；济利禄带着大队人马，前往他们的会堂——他们这样称呼他们的祈祷之所——夺走它们，并将犹太人逐出城外，任凭民众劫掠他们的财物。\newline{}于是，自马其顿的\Person{亚历山大}时代起就居住在该城的犹太人被逐出，被剥夺一切所有，四散流离。\newline{}其中有一位名叫\Person{阿达曼提乌斯}的医生，逃到君士坦丁堡主教\Person{阿提库斯}那里，宣信基督教，一段时间后返回\Place{亚历山大}并在那里定居。\newline{}但\Place{亚历山大}总督\Person{奥勒斯特}对这些事大为愤怒，且极其痛心如此大的一座城市突然失去如此众多的人口；于是他立即将全部情况禀报了皇帝。\newline{}济利禄也写信给他，描述犹太人的暴行；同时派人到\Person{奥勒斯特}那里调解和好之事，因为民众催促他这样做。\newline{}当\Person{奥勒斯特}拒绝听取善意表示时，济利禄向他伸出\WorkTitle{福音书}，以为对宗教的尊重会促使他放下怨恨。\newline{}然而，即使这样也未能使总督平静下来，他仍对主教怀有不可化解的敌意，随后发生了以下事件。\par

% structure: p0028 source=h2 target=subsection reason=relative-heading-rank; base=h2

\subsection{第十四章 \Place{尼特里亚}的隐修士下来，针对\Place{亚历山大}总督发动叛乱。}

有些居住在尼特里亚山中的隐修士，性情极为火烈——他们先前曾被德奥斐洛不义地武装起来反对狄奥斯库鲁斯及其弟兄们——如今又因一股热忱而激动，决心为济利禄而战。于是约有五百人离开隐修院，进入城中；他们遇见总督乘坐战车，便称他为外教偶像崇拜者，并施以许多其他辱骂之词。总督以为这是济利禄设下的圈套，便声明自己是基督徒，且在君士坦丁堡由主教阿提库斯施洗。但他们几乎不听他的声明，其中一人名叫阿摩尼乌斯，向奥瑞斯忒斯投掷石块，击中他的头部，血流如注；所有卫兵除少数外都逃散了，纷纷冲入人群，有的朝这边，有的朝那边，害怕被石头砸死。与此同时，亚历山大城的民众赶来救援总督，将其余的隐修士击退；但抓住阿摩尼乌斯后，把他交给了总督。总督立即公开施以酷刑，刑罚极其严酷，以致他死于酷刑之下；不久之后，总督向皇帝报告了所发生的事。另一方面，济利禄也向皇帝呈递了他对此事的陈述；并将阿摩尼乌斯的尸体安放在某教堂内，给他起了新名号\OriginalTerm{塔乌玛西乌斯}{Thaumasius}，命他列入殉道者行列，并在教堂中颂扬他的英勇，称他是在为虔诚而战的冲突中倒下的。但较为明智的人，虽然是基督徒，却不接受济利禄对此人的偏见评价；因为他们深知他是为自己的鲁莽行为受罚，并非因不肯否认基督而死于酷刑。济利禄自己也意识到这一点，便让这件事的记忆逐渐被沉默所抹去。但济利禄与奥瑞斯忒斯之间的敌意并未因此丝毫消减，反而因一件类似前述的事而重新燃起。\par

% structure: p0030 source=h2 target=subsection reason=relative-heading-rank; base=h2

\subsection{第十五章　论女哲学家希帕提娅}

亚历山大城有一位名叫希帕提娅的妇女，是哲学家塞昂之女。她在文学和科学上的造诣远超同时代的所有哲学家。她继承了柏拉图与普罗提诺的学派，向听众讲解哲学原理，许多人从远方来受教。由于心灵修养而获得的沉着与从容，她时常公开出现在官员面前。她也不以在男人集会上露面为羞。因她非凡的尊严与美德，所有男人更加敬佩她。然而，她仍成为当时盛行的政治嫉妒的牺牲品。由于她常与奥瑞斯忒斯会面，在基督徒民众中流传诽谤之言，说是她阻碍了奥瑞斯忒斯与主教和解。于是，其中一些人被凶残而偏执的热忱所驱使，其首领是一位名叫伯多禄的读经员，她们在她回家途中埋伏，将她从马车中拖出，带到名为\Place{Caesareum}的教堂，将她完全剥光，然后用瓦片杀害。她们将她的身体撕成碎片，把残肢带到名为\Place{Cinaron}的地方，在那里焚毁。此事不仅给济利禄，也给整个亚历山大教会带来了极大的耻辱。确实，没有什么比允许屠杀、斗殴及此类行径更远离基督信仰的精神了。此事发生在三月的大斋期内，即济利禄就任主教第四年，霍诺留第十次任执政官、德奥多西第六次任执政官之时。\par

% structure: p0032 source=h2 target=subsection reason=relative-heading-rank; base=h2

\subsection{第十六章　犹太人再次对基督徒施暴并受惩罚}

不久之后，犹太人重新对基督徒施行恶毒而不敬神的行径，并招致应得的惩罚。在叙利亚介于卡尔基斯与安提约基雅之间一个名为\Place{Inmestar}的地方，犹太人按他们惯常的方式以各种嬉戏取乐。他们这样做了许多荒唐之事，最后因醉酒而冲动，竟然嘲笑基督徒甚至基督本人；为了戏弄十字架和信赖被钉者的人，他们抓住一个基督徒男孩，把他绑在十字架上，开始嘲笑讥讽他。但没过多久，他们狂怒不已，鞭打那孩子，直到他死在他们的手下。此举引发了犹太人与基督徒之间的激烈冲突；皇帝得知此事后，命令该省总督查明并惩罚肇事者。于是，该地的犹太居民为他们亵渎神灵的嬉戏中所犯的恶行付出了代价。\par

% structure: p0034 source=h2 target=subsection reason=relative-heading-rank; base=h2

\subsection{第十七章　诺瓦提安派主教保禄为一名犹太骗子施洗时所行的奇迹}

大约在这时，\Person{克里桑斯}，诺瓦替雅奴派的主教，在主持自己教派的教会七年后，于\Person{莫纳修斯}与\Person{普林塔}担任执政官期间的八月二十六日去世。继任主教的是\Person{保禄}，他曾是一名拉丁文教师；但后来他弃绝拉丁文，投身于苦修生活；他建立了一座隐修院，过着与旷野中隐修士相差无几的生活。事实上，我本人发现他正是\Person{埃瓦格利乌斯}所说的旷野隐修士应有的人：效仿他们持续禁食、沉默、戒除肉食，并且大都也戒绝油和酒。此外，他关心穷人的需要，其程度超过任何人；他不知疲倦地探访监中之人，并为许多囚犯向法官求情，法官因其卓越的虔敬而乐于听从。但我为何要延长对他的记述呢？因为我即将提到一件非常值得记载的事迹。有一个犹太骗子，假装皈依基督教，经常受洗，并以此诡计聚敛了不少钱财。他用这种欺骗手段迷惑了许多基督教派（因为他已从亚流派和马其顿派接受了洗礼）之后，因再无他人可施其伪善，便终于来到诺瓦替雅奴派的主教\Person{保禄}面前，声称自己诚心渴望受洗，请求从他手中领受洗礼。\Person{保禄}称赞那犹太人的决心，但告诉他，除非他在信仰的根基上得到教导，并禁食祈祷多日，否则不能为他施行此礼。那犹太人不情愿地被迫禁食，却更加急切地请求受洗；\Person{保禄}见他如此迫切，不愿以更久的拖延令他沮丧，便同意了他的请求，并准备了受洗所需的一切物品。他为他买了一件白衣，吩咐注满洗礼池，然后带领那犹太人到池边准备为他施洗。但天主一种无形的力量使水突然消失。主教和在场的人自然对真正的原因毫无察觉，以为水通过下面的水道排走了——他们通常用那些水道排空洗礼池；于是这些通道被仔细堵塞，又注满洗礼池。然而，当犹太人第二次被带往洗礼池时，水又像先前一样消失了。于是\Person{保禄}对那犹太人说：\Quote{你要么是个作恶的人，可怜的家伙，要么是个已经受过洗的无知之人。}人们聚集观看这一奇迹，其中一人认出了那犹太人，并指认他不久之前已由主教\Person{阿提库斯}施洗。诺瓦替雅奴派主教\Person{保禄}的手所行的奇迹就是这样。\par

% structure: p0036 source=h2 target=subsection reason=relative-heading-rank; base=h2

\subsection{第十八章 罗马人与波斯人之间在波斯王\Person{伊斯迪格德斯}去世后再度开战。}

波斯王\Person{伊斯迪格尔德斯}去世了，他生前未曾骚扰其境内的基督徒。他的儿子，名叫\Person{瓦拉兰尼斯}，继承了王位。这位王子屈服于术士的影响，严厉迫害那里的基督徒，对他们施加各种波斯式的刑罚和折磨。因此，由于压迫，他们被迫离开本国，到罗马人中间寻求避难，恳求他们不要让他们被完全消灭。主教\Person{阿提库斯}以极大的仁慈接待了这些请愿者，并尽一切可能帮助他们；于是他将实情告知了皇帝\Person{狄奥多西}。与此同时，罗马人对波斯人的另一项不满也暴露出来。也就是说，波斯人不肯归还从罗马人中雇佣来的金矿劳工，还掠夺了罗马商人。这些事所引起的恶感因波斯基督徒逃入罗马领土而大大加剧。因为波斯王立刻派遣使节要求交出逃亡者。但罗马人绝不愿意交出他们；不仅因为他们希望保护自己的请愿者，而且因为他们愿意为基督教信仰做任何事。因此他们宁愿重新与波斯人开战，也不愿让基督徒悲惨地被消灭。盟约于是被破坏，随之而来了一场激烈的战争。关于这场战争，我认为给予一些简短的叙述并非不合时宜。罗马皇帝首先派遣了一支由将军\Person{阿尔达布里乌斯}指挥的军队；他穿过\Place{亚美尼亚}侵入波斯，蹂躏了其中一个叫做\Place{阿扎泽内}的省份。波斯将军\Person{纳尔塞乌斯}率波斯军队向他进攻，但一经交战便被打败，被迫撤退。后来他认为通过\Place{美索不达米亚}出其不意地侵入那里无人防守的罗马领土是有利的，想借此报复敌人。但\Person{纳尔塞乌斯}的这一计划并未逃过罗马将军的注意。因此，在掠夺了\Place{阿扎泽内}之后，他自己也急忙进军\Place{美索不达米亚}。因此，\Person{纳尔塞乌斯}虽然拥有一支庞大的军队，却未能入侵罗马省份；但他到达\Place{尼西比斯}——一座属于波斯人、位于两国边境的城市——后，派人去见\Person{阿尔达布里乌斯}，希望双方就战事做出相互安排，并指定交战的时间和地点。但他对使者说：\Quote{告诉\Person{纳尔塞乌斯}，罗马皇帝不会在他乐意的时候作战。}皇帝察觉到波斯人在集结全部兵力，便向军队补充兵员，并将全部希望寄托于天主以获得胜利：以下情况证明，国王并未立时从这一虔诚的信赖中获益。当\Place{君士坦丁堡}居民大为恐慌，对战争结局忧心忡忡时，来自天主的天使向\Place{比提尼亚}的一些人显现，这些人正因私事前往\Place{君士坦丁堡}，天使吩咐他们告诉民众不要惊慌，而要向天主祈祷，并确信罗马人将取得胜利。因为他们说，他们自己是由天主派遣来保护他们的。当这个消息传开时，不仅安慰了城中的居民，也使士兵们更加勇敢。正如我们所说，战事从\Place{亚美尼亚}转移到\Place{美索不达米亚}，罗马人将波斯人困在\Place{尼西比斯}城中并加以围攻；他们建造了木塔，用机械将其推进到城墙边，杀死了大量守城者以及赶来救援的人。当波斯君主\Person{瓦拉兰尼斯}得知他的\Place{阿扎泽内}省份已被摧毁，另一方面他的军队在\Place{尼西比斯}城被严密围困时，他决定亲自率领全部兵力进攻罗马人：但出于对罗马人勇气的畏惧，他恳求撒拉逊人的援助，撒拉逊人当时由一位名叫\Person{阿拉蒙达鲁斯}的好战首领统治。这位王子因此带来了大批撒拉逊援军，并劝告波斯王不必畏惧，因为他很快就会将罗马人置于他的权力之下，并将\Place{叙利亚}的\Place{安提约基雅}交到他手中。但事态并未实现这些承诺；因为天主使撒拉逊人心中充满了极大的恐慌；他们以为罗马军队正在进攻他们，又找不到其他逃生之路，便全副武装地跳入\Place{幼发拉底河}，其中近十万人溺毙。这就是恐慌的性质。\par

围攻\Place{尼西比斯}的罗马人得知波斯王正带领大量大象前来，转而惊慌起来，烧毁了他们在围攻中使用的一切器械，退回了本国。此后发生了哪些交战，另一位罗马将军\Person{阿雷奥宾杜斯}如何在单挑中杀死最勇敢的波斯人，\Person{阿尔达布里乌斯}又如何用伏击消灭了七名波斯指挥官，以及另一位罗马将军\Person{维提安}如何击败了撒拉逊残余部队，我相信我应该略过不提，以免离题太远。\par

% structure: p0039 source=h2 target=subsection reason=relative-heading-rank; base=h2

\subsection{第十九章 论信使\Person{帕拉迪乌斯}}

我将尝试说明，\Person{狄奥多西}皇帝是如何在难以置信的短时间内得知所发生之事，以及他是如何迅速获知远方发生之事的。因为在他的臣民中，有幸拥有一位身心俱具非凡精力的人，名为\Person{帕拉狄乌斯}；他骑马极其迅猛，能在三天内到达罗马和波斯王国的边界，再用同样多的时间返回\Place{君士坦丁堡}。同一人受皇帝派遣，以同样神速穿行于世间其他区域；因此一位雄辩家曾恰如其分地说：\Quote{此人以其速度证明罗马帝国的广阔疆域变得渺小。}波斯王本人也对有关这位信使的快速事迹感到惊讶；但我们只需满足于上述关于他的细节即可。\par

% structure: p0041 source=h2 target=subsection reason=relative-heading-rank; base=h2

\subsection{第二十章 波斯人再次被罗马人击败。}

当时，居住在\Place{君士坦丁堡}的罗马皇帝，因深知天主已明确赐予他胜利，心怀仁慈，尽管他的部下在战争中取得了成功，他仍渴望缔结和平；为此，他派遣了最信任的人\Person{赫利翁}，委以与波斯人订立和平条约的使命。\Person{赫利翁}到达\Place{美索不达米亚}后，在罗马人为自身安全挖有壕沟的地方，先前派遣了一位雄辩之士\Person{马克西明}作为代表，他是军队总司令\Person{阿尔达布里乌斯}的同僚，去为和平条款做初步安排。\Person{马克西明}来到波斯王面前，说他被派遣来处理此事，并非受罗马皇帝派遣，而是受他的将军们派遣；因为他说这场战争甚至不为皇帝所知，即便知道，皇帝也会认为微不足道。当波斯王欣然决定接见使团时——因为他的军队正遭受粮草匮乏——他们中有一支以\OriginalTerm{不朽者}{the Immortals}命名的队伍来到他面前。这是一支约一万名勇士的队伍——他们劝谏国王不要听取任何和平提议，直到他们向罗马人发动一次进攻，因为罗马人现已变得极其疏忽大意。国王赞同他们的建议，下令囚禁使节并派人看守，并允许\OriginalTerm{不朽者}{the Immortals}实施他们针对罗马人的计划。于是，他们到达指定地点后，分成两队，意图包围部分罗马军队。罗马人只看到一股波斯人逼近，便准备迎战，未看见另一队，因为后者突然冲入战场。但就在战斗即将开始时，天主眷顾如此安排，另一支由将军\Person{普罗科皮乌斯}率领的罗马军队从某座山后出现，发现同伴处于危险，便从后方攻击波斯人。这样，片刻前包围罗马人的那些人，自己反被包围。罗马人迅速彻底消灭了这些人后，转而对付从埋伏中冲出的那些波斯人，同样用标枪将他们全部杀死。就这样，那些被波斯人称为\OriginalTerm{不朽者}{the Immortals}的人，全部被证明是会死的，\Person{基督}对波斯人施行了这次报复，因为他们流了许多虔诚敬拜者的血。波斯王得知灾难后，假装对发生之事一无所知，下令允许使团觐见，他这样对使节说：\Quote{我同意和平，并非向罗马人屈服，而是为取悦你，我发现你是所有罗马人中最明智的。}就这样，这场因波斯受苦的基督徒而发起的战争结束了，当时是两位奥古斯都执政之年，即\Person{霍诺留}第十三次、\Person{狄奥多西}第十次执政，第三百届奥林匹克运动会第四年；随着战争结束，在波斯针对基督徒发起的迫害也终止了。\par

% structure: p0043 source=h2 target=subsection reason=relative-heading-rank; base=h2

\subsection{第二十一章 \Place{阿米达}主教\Person{阿卡西乌斯}善待波斯俘虏。}

\Place{阿米达}主教\Person{阿卡西乌}的一件高尚行为，当时极大地提升了他在所有人中的声誉。由于罗马士兵无论如何都不愿将所俘获的俘虏交还给波斯王，这些约七千人正在遭受饥荒蹂躏的\Place{阿扎泽内}濒临饿死，这令波斯王极为苦恼。于是\Person{阿卡西乌}认为这样的事绝不能等闲视之；他召集了圣职人员，对他们这样说道：\Quote{我们的天主，我的兄弟们，既不需要盘碟，也不需要杯盏；因为祂既不饮也不食，一无所缺。既然教会凭着信友的慷慨拥有许多金银器皿，我们理应将其出售，以便用所得银钱赎回俘虏，并为他们提供食物。}他说了这些话以及其他许多类似的话之后，下令将器皿熔化，用所得款项向士兵支付俘虏的赎金，并供养了他们一段时间；然后为他们备好旅途所需，将他们送回本国。这位卓越的\Person{阿卡西乌}的仁慈之举令波斯王惊叹，仿佛罗马人不仅凭战争中的勇武，也凭和平中的仁爱来战胜敌人。据说波斯王还希望\Person{阿卡西乌}能来见他，以便能一睹此人的风采；这一愿望很快遵照\Person{狄奥多西}皇帝的谕令得以实现。罗马人凭天主的恩宠取得了如此显赫的胜利后，许多以辩才著称之人都撰写了颂词献给皇帝，并公开诵读。皇后本人也创作了一首英雄体诗：她极具文学素养；她是雅典修辞学家\Person{莱昂提乌}的女儿，从小受到父亲在各类学问上的教导；\Person{阿提库}主教在皇帝与她成婚前不久为她付洗，并赋予她基督徒的名字\Person{欧多西亚}，取代了她异教的名字\Person{雅典纳伊斯}。正如我所说，许多人借此机会发表了颂词。有些人是出于被皇帝注意的欲望；另一些人则急于向大众展示才华，不愿自己辛苦取得的成就埋没于无名。\par

% structure: p0045 source=h2 target=subsection reason=relative-heading-rank; base=h2

\subsection{第二十二章 小狄奥多西皇帝的德行}

但我既不急于求皇帝的注意，也不想炫耀自己的口才，却感到有责任如实记述皇帝所拥有的非凡德行：因我深信，对如此卓越的德行保持沉默，对后世之人实为不公。首先，这位君主虽生于、长于帝位，却并未因出身和教育的环境而变得愚钝或柔弱。他展现出如此睿智，以至与他交谈的人皆认为他凭经验获得了智慧。他忍受艰辛时如此刚毅，能勇敢地承受寒暑；经常禁食，尤其每逢星期三和星期五；他这样做，是出于竭力准确遵守基督教所规定的一切形式的虔诚。他使皇宫与隐修院几乎无异：因为天未亮他便与姐妹们一同起身，以应和的方式背诵赞美天主的圣咏。通过这种训练，他熟记圣经；常与主教们讨论圣经问题，仿佛他是一名资深的受秩司铎。他收集圣书及其注释的勤勉，甚至超过昔日的\Person{托勒密·费拉德尔弗斯}。在仁慈和人性方面，他远超众人。因为\Person{尤利安}皇帝虽自称哲学家，却无法克制对讥讽他的安提约基亚人的愤怒，对\Person{特奥多罗}施加了最痛苦的折磨。相反，\Person{特奥多西}告别了亚里士多德的推理，以行动实践哲学，控制愤怒、忧伤和快乐。他从未报复过任何伤害他的人；也从未有人见他动怒。当他的一些密友问他为何从不将罪犯处死时，他回答说：\Quote{但愿我能使死者复活。}另一人问同样的问题，他答道：\Quote{处死一个人凡人并非什么伟大或困难的事；只有天主才能通过悔改使已死的人复活。}他如此习惯于施行仁慈，以至于若有人被判死刑，被带往刑场时，从未让他走到城门口之前，赦免令便已发出，命令立即返回。有一次他在君士坦丁堡的露天剧场举行野兽围猎表演，民众喊道：\Quote{让一位最勇敢的斗兽士去与激怒的野兽搏斗！}但他对他们说：\Quote{你们不知道我们惯于以人道之心观看这些表演吗？}这句话教导民众以后满足于不那么残忍的表演。他的虔诚在于他尊敬所有奉献给天主服务的人，并特别尊崇那些他确知生活圣德卓越的人。据说赫布龙的主教在君士坦丁堡去世后，皇帝希望得到他的粗毛麻布道袍；那衣服虽极其肮脏，他却穿在外袍之上，希望借此在某种程度上分享死者的圣德。某年天气非常恶劣，因民众急切要求，他不得不在赛马场举行惯常赛事；当赛马场挤满观众时，暴风雨加剧，大雪纷飞。于是皇帝非常清楚地表明了他对天主的意向：他让传令官向民众宣布：\Quote{最好、也更合宜的，是停止表演，一同向天主祈祷，使我们能免受即将来临的风暴之害。}传令官刚执行完命令，所有民众便极其欢喜地一致开始祈求并赞美天主，以致整座城市成为一大集会；皇帝本人身着朝服，走入人群之中，带头唱起赞美诗。他的期望没有落空，因为天气开始恢复往常的晴朗；上主恩赐给所有人丰厚的收成，而非预期的歉收。若遇战争，他就如同\Person{达味}一样投靠天主，深知他是战争的裁决者，并藉祈祷使战争顺利结束。因此，此刻我将叙述，在波斯战争后不久，他因信赖天主而击败了篡位者\Person{若望}（当时\Person{奥诺留}已于八月十五日，阿斯克勒皮奥多图斯与马里安任执政官之年去世）。因为我认为当时所发生的事值得提及，因皇帝派去讨伐篡位者的将领们遇到了类似于以色列人在\Person{梅瑟}带领下渡过红海时所发生的事。然而，这些事我将非常简短地叙述，把那些需要专论的大量细节留给他人。\par

% structure: p0047 source=h2 target=subsection reason=relative-heading-rank; base=h2

\subsection{第二十三章　\Person{奥诺留}皇帝死后，\Person{若望}在罗马篡位；他因\Person{特奥多西二世}的祈祷而遭毁灭}

当皇帝\Person{奥诺留}去世时，已为唯一统治者的\Person{狄奥多西}接到消息后，尽可能长时间地隐瞒真相，时而用这种说法，时而用那种说法误导百姓。但他私下派遣一支军队前往\Place{达尔马提亚}的\Place{萨洛纳}城，以便一旦西方发生任何叛乱时，能有兵力加以遏制。做好这些临时安排后，他终于公开宣布了叔父的死讯。与此同时，皇帝秘书监\Person{若望}不满足于已获得的尊荣，篡夺了最高权力，并派遣使节觐见皇帝\Person{狄奥多西}，请求承认他为帝国同僚。但那位皇帝首先逮捕了使节，随后派出曾在波斯战争中功勋卓著的军队总司令\Person{阿尔达布里乌斯}。他到达\Place{萨洛纳}后，从那里乘船前往\Place{阿奎莱亚}。起初人们认为他运气好，但后来证明运气对他不利。因为刮起了逆风，他被吹入篡位者手中。篡位者抓获他后，更加确信皇帝会因为情况紧急而选举并宣布他为皇帝，以保全其总司令的性命。皇帝得知此事后确实大为忧虑，派去讨伐篡位者的军队也同样担忧，生怕\Person{阿尔达布里乌斯}遭到篡位者的虐待。\Person{阿尔达布里乌斯}之子\Person{阿斯帕尔}得知父亲落入篡位者之手，同时明白叛军因大量蛮族的加入而增强，不知如何是好。然而在这危急时刻，虔敬皇帝的祈祷起了作用。因为天主的使者以牧羊人的形象出现，引导\Person{阿斯帕尔}及其所率军队，穿过\Place{拉文纳}附近的湖泊——当时篡位者正驻跸在该城——并在那里困住了这位军事首领。此前无人知晓能徒步渡过该湖，但天主那时使原本不可通行之处变得可以通行。于是他们仿佛行走在干地上渡过湖泊，发现城门大开，便制服了篡位者。这一事件使那位至为虔敬的皇帝有机会再次展示他对天主的虔诚。因为篡位者覆灭的消息传来时，他正在竞技场观赏赛会，他立刻对百姓说：\Quote{诸位，如今请你们离开这些娱乐，前往教堂向天主献上感恩，祂的手已推翻了篡位者。}他这样说时，赛会立即被弃置不顾，百姓全都同他一起唱着赞美诗走出赛马场，同心同声。到了教堂后，全城又成为一个会众；他们在教堂里以这些敬礼活动度过了余下的时间。\par

% structure: p0049 source=h2 target=subsection reason=relative-heading-rank; base=h2

\subsection{第二十四章 君士坦提乌斯与普拉西狄亚之子、狄奥多西之表弟瓦伦提尼安被立为皇帝。}

篡位者死后，皇帝\Person{狄奥多西}非常焦虑，不知该立谁为西方皇帝。他有一位表弟名叫\Person{瓦伦提尼安}，年纪尚幼，是其姑母\Person{普拉西狄亚}之子——\Person{普拉西狄亚}是大\Person{狄奥多西}之女，两位奥古斯都\Person{阿卡狄乌斯}与\Person{奥诺留}之妹，也是那位曾被\Person{奥诺留}立为皇帝、与其共同执政不久后便去世的\Person{君士坦提乌斯}之妻。他将这位表弟封为凯撒，派往西方，将政务托付给其母\Person{普拉西狄亚}。他自己也急忙前往意大利，以便亲自立表弟为皇帝，同时亲临其地，以劝告影响当地居民，使他们不再轻易顺从篡位者。但当他到达\Place{塞萨洛尼卡}时，因病无法继续前行，于是通过贵族\Person{赫利翁}将皇冠送往其表弟处，自己返回了\Place{君士坦丁堡}。关于这些事，我认为此处所述已足够。\par

% structure: p0051 source=h2 target=subsection reason=relative-heading-rank; base=h2

\subsection{第二十五章 \Person{阿提库斯}主教的基督徒仁爱。他将\Person{若望}之名登录于双联簿。他预知自己的死亡。}

\Emph{与此同时}，主教\Person{阿提库斯}使教会事务异常兴盛，他以智慧处理一切，并通过教导激励百姓修德。他察觉到教会即将分裂，因为若望派的人单独聚集，便下令在祈祷中提及\Person{若望}，如同对其他已故主教所行的那样。他相信这样做会使许多人回归教会。他如此慷慨，不仅供给本教区穷人，还向邻近城市的贫困者输送捐献，以供应他们的需要并增进其舒适。有一次，他派人送三百金币给\Place{尼西亚}教会司铎\Person{卡利奥皮乌斯}，同时附上以下信函。\par

\Quote{\Person{阿提库斯}致\Person{卡利奥皮乌斯}——在主内问安。}\par

\Quote{我得知在贵城有一万名急需之人，其状况理当引起虔诚者的怜悯。我说一万，是指其众多而非精确数字。既然我从那位以慷慨之手常供养忠实管家者那里收到一笔款项，又恰逢有些人因贫困所迫——这是为了考验那些拥有资财却不愿周济穷困者的人——我的朋友，请收下这三百金币，并视情况处置。我相信你会留意分给那些羞于乞讨的人，而非那些终其一生靠他人养活的人。施舍时不要因宗教立场而区别对待，要喂养饥饿者，无论他们是否与我们意见一致。}\par

阿提库斯就这样挂念着远方的穷人。他还努力消除某些人的迷信。因为他获悉，那些因犹太人的逾越节而与诺瓦提安派分离的人，已将萨巴提乌斯的遗体从罗得岛（他流放死于该岛）运回，埋葬后，常在他的墓前祈祷；他便在夜间将遗体掘出，葬于一隐秘墓穴；那些先前在此处朝拜的人发现坟墓已被打开后，便从此不再崇敬该墓。此外，他在为地方命名上展现了许多品味。对于欧克塞因海入口处一个古称法马库斯的港口，他命名为特拉佩亚；因为他不想让举行宗教集会的地方因不吉之名而蒙羞。他将君士坦丁堡郊区另一处地方称为阿基罗波利斯，原由如下：克吕索波利斯是位于博斯普鲁斯海峡顶端的一个古老港口，多位早期作家曾提及，尤其是斯特拉博、大马士革的尼各老，以及杰出的色诺芬在他的\WorkTitle{居鲁士远征记}第六卷中；又在\WorkTitle{希腊史}第一卷中论及该地时说：\Quote{亚西比德曾筑墙圍住它，并在其中设立通行税；因为所有从本都航行出来的人都习惯在那里缴纳什一税。}阿提库斯见前者正对着克吕索波利斯，位置十分宜人，便宣称它最适合称为阿基罗波利斯；此言一出，名称便牢固确立。有人对他说，不应准许诺瓦提安派在城内举行集会；他回答说：\Quote{难道你们不知道，他们在君士坦提乌斯和瓦伦斯迫害期间曾与我们一同受苦吗？况且，}他说，\Quote{他们是我们信经的见证：因为他们虽然很久以前就与教会分离，却从未在信仰上带来任何革新。}一次，他因主教祝圣之事在尼凯亚，见到诺瓦提安派主教阿斯克勒皮亚德斯年事已高，便问他：\Quote{你担任主教多少年了？}当回答是五十年时，他说：\Quote{你真是有福之人，能如此长久地从事\Quote{善工}\BibleRef{弟前 3:1}}。他又对同一位阿斯克勒皮亚德斯说：\Quote{我称赞诺瓦图斯，但绝不能赞同诺瓦提安派。}阿斯克勒皮亚德斯对这奇怪说法感到惊讶，问道：\Quote{主教，您这话是什么意思？}阿提库斯给出了区分的理由：\Quote{我称赞诺瓦图斯拒绝与那些献过祭的人共融，因为我自己也会这样做；但我不能赞扬诺瓦提安派，因为他们因极轻微的过错就将平信徒逐出共融。}阿斯克勒皮亚德斯回答说：\Quote{除了向偶像献祭，圣经中还有许多\Quote{至于死的罪}\BibleRef{若一 5:17}；为此你们只绝罚神职人员，但我们连平信徒也绝罚，将赦罪之权唯独保留于天主。}阿提库斯还预知了自己的死亡；因为他离开尼凯亚时，对那里的长老卡利奥皮乌斯说：\Quote{趁秋天之前赶到君士坦丁堡，你若还想再见我一面；因为若迟过那时，你将不见我存活。}这预言果然应验；他在自己主教任期的第二十一年，即狄奥多西第十一次执政、瓦伦提尼安凯撒第一次执政时，于十月十日逝世。狄奥多西皇帝当时正从塞萨洛尼基赶来，未能在其葬礼前抵达君士坦丁堡，因为阿提库斯在皇帝抵达前一日已安葬。不久后，同月十月二十三日，小瓦伦提尼安被宣布为奥古斯都。\newline{}\par

% structure: p0056 source=h2 target=subsection reason=relative-heading-rank; base=h2

\subsection{第二十六章 西西尼乌斯被选接替阿提库斯。}

阿提库斯去世后，关于推选继任者发生了激烈争论，有人推举此人，有人推举彼人。据说，有一方极力支持一位名叫斐理伯的长老；另一方则希望提拔同样身为长老的普洛克鲁斯；但民众普遍希望将主教职位授予西西尼乌斯。此人也是长老，但在城内未担任任何教会职务，他被任命在埃莱亚（君士坦丁堡郊区一村庄）的教堂从事神圣职务。该村庄位于城市对面的港口彼岸，按古老习俗，全体居民每年在那里聚集庆祝救主升天。所有平信徒都热切拥护此人，因为他以虔诚闻名，尤其因他照顾穷人甚至\Quote{超过自己的力量}\BibleRef{格后 8:3}。平信徒的热忱因而占了上风，西西尼乌斯于二月二十八日被祝圣，时值次年执政官任期，即狄奥多西第十二次、瓦伦提尼安第二次执政。长老斐理伯因他人被选优先于自己而深感懊恼，甚至在其\WorkTitle{基督教史}中引入此事，对受祝圣者及祝圣者都作了十分苛责的评论，对平信徒则更是严厉。但他所说之事，我无论如何不能写下。既然我不赞同他轻率地将它们写入文字，然而我认为在此略提一下他及其著作，并非不合时宜。\newline{}\par

% structure: p0058 source=h2 target=subsection reason=relative-heading-rank; base=h2

\subsection{第二十七章 西德的长老斐理伯卷帙浩繁的著作。}

\Person{斐理伯}是\Place{息德}人；\Place{息德}是\Place{潘非利亚}的一座城市。诡辩家\Person{特洛依鲁斯}也来自此地，\Person{斐理伯}自称与他有近亲关系。他是一名执事，因此得以与主教\Person{金口若望}亲密往来。他勤奋钻研文学，除了取得相当可观的文学成就外，还广泛收集了各个知识领域的书籍。他模仿亚细亚风格，撰写了许多论著，其中尝试驳斥\Person{儒利安}皇帝反对基督徒的论著，并编纂了一部\WorkTitle{基督教史}，分为三十六卷；每卷又分若干册，总数近千册，每册的内容提要本身与册子篇幅相当。这部作品他命名为\WorkTitle{基督教史}而非\WorkTitle{教会史}，并在其中汇集了大量五花八门的材料，意图显示自己并非对哲学和科学知识一无所知：因为其中包含了几何定理、天文思辨、算术计算、音乐原理，以及岛屿、山脉、森林的地理描述，还有其他各种无关紧要的内容。他将这些不相干的细节强行与主题关联，使他的著作成为一部极其松散的作品，在我看来，对无知者和有学识者都毫无用处；因为不识字的人无法欣赏其文辞的高深，真正有判断力的人则谴责其冗长的赘言。但每个人都可根据自己的品味对这些书作出判断。我只需补充一点：他混淆了所描述事件的年代顺序；因为他在叙述了\Person{德奥多西}皇帝统治时期发生的事情后，立刻又回到\Person{亚大纳削}主教的时代；他经常这样。关于\Person{斐理伯}已说得够多了：现在我们必须谈谈\Person{西息宁}主教任内所发生的事。\par

% structure: p0060 source=h2 target=subsection reason=relative-heading-rank; base=h2

\subsection{第二十八章．\Person{西息宁}祝圣\Person{普罗克鲁斯}为\Place{基齐库斯}主教，但被人民拒绝。}

\Place{基齐库斯}的主教去世后，\Person{西息宁}祝圣\Person{普罗克鲁斯}为该城的主教。但当他准备前往那里时，当地居民抢先选举了一位名叫\Person{达马修}的隐修者。他们这样做是漠视一条法律，该法律禁止他们在没有\Place{君士坦丁堡}主教批准的情况下祝圣主教；但他们声称这是特别授予\Person{阿提库}个人的特权。因此，\Person{普罗克鲁斯}仍然没有自己的主教座堂，但因在\Place{君士坦丁堡}各教堂的讲道而声名鹊起。然而，我们将在适当的地方更详细地谈论他。\Person{西息宁}在担任主教后仅两年，于\Person{希耶琉斯}与\Person{阿尔达布里乌斯}执政之年十二月二十四日去世。他以节制、生活正直和对穷人的仁慈而闻名；他性格异常和蔼、坦诚，这使得他相当不喜事务，因此被勤于活动的人视为懒惰。\par

% structure: p0062 source=h2 target=subsection reason=relative-heading-rank; base=h2

\subsection{第二十九章．\Place{安提约基雅}的\Person{聂斯脱利}升任\Place{君士坦丁堡}主教座。他对异端者的迫害。}

西西尼乌斯去世后，由于君士坦丁堡的圣职人员表现出争权好胜的精神，皇帝们决定不让该教会的任何人填补空缺的主教之位，尽管许多人热切希望祝圣斐理伯，也有不少人力挺推举普若克鲁斯。于是，他们从安提约基雅请来一位陌生人，名叫聂斯多略，原籍盖尔马尼基，以极佳的嗓音和流利的口才著称；他们认为这些资质对于教导百姓很重要。因此，过了三个月后，聂斯多略被从安提约基雅带来，有些人因其节制而大大赞扬他；但至于他在其他方面具有何种性情，凡有辨别力的人从首次讲道便能察觉。他于四月十日，在斐利克斯和托鲁斯执政官任下被祝圣，立即当着全体百姓向皇帝说出那著名的话：\Quote{我的君王，请将大地从异端中净化，我将以天堂回报你。请协助我消灭异端，我将协助你战胜波斯人。}这些言论虽然使一些对异端之名抱有无知敌意的群众极为满意，但正如我所说，那些善于从言谈预测人品的人，却毫不遗漏地察觉他轻浮的心思、暴躁和虚荣的性情，因为他竟不能抑制自己片刻就爆发出如此激烈的言辞；借用谚语\Quote{尚未尝过城中水}，便显为疯狂的迫害者。果然，在他祝圣后的第五天，他决定拆毁亚略派常私下举行敬礼的小堂，从而迫使这些人铤而走险；因为当他们看到自己小堂遭受破坏时，便向小堂放火，火势蔓延，也使许多邻近建筑化为灰烬。因此全城骚动，亚略派忿怒复仇，准备行动；但天主，城市的守护者，未让灾祸达到顶点。然而从那时起，他们给聂斯多略贴上\Quote{纵火者}的标签，这不只是异端者所为，连同信仰的人也是如此。因为他不能安静，而是千方百计骚扰那些不赞同他观点的人，不断扰乱公共安宁。他也恼怒了诺瓦替安派，因为他们的主教保禄因其虔诚而处处受到尊敬，这使他心生嫉妒；但皇帝通过劝诫制止了他的狂热。至于他在整个亚细亚、吕底亚和卡里亚给十四日派带来了怎样的灾难，以及在米利都、撒尔德因他引发的民乱中丧生了多少人，我认为应当略过不提。他因所有这些暴行以及放纵无度的口舌而遭受的惩罚，我将在稍后提及。\par

% structure: p0064 source=h2 target=subsection reason=relative-heading-rank; base=h2

\subsection{第三十章：在狄奥多西二世治下勃艮第人信奉基督。}

现在我必须叙述一件发生在此时、值得记载的史事。在莱茵河彼岸有一个称为勃艮第人的蛮族；他们过着和平的生活，因为几乎全是工匠，靠手艺为生。匈人不断侵扰这个民族，蹂躏其土地，常常杀害他们许多人。因此，在这困境中，勃艮第人决定不求助于任何人，而是将自己托付给某位天主的保护；他们认真考虑后，认为罗马人的天主大力保护敬畏他的人，于是全体一致接受了基督信仰。他们前往高卢的一座城市，请求当地主教授予基督圣洗；主教命他们禁食七日，在此期间教导他们信仰的基本原理，第八日给他们付洗后遣散。从此他们充满信心，出征攻打入侵者，没有辜负自己的希望。因为匈人的王乌普塔在夜间因饮食过量而死，勃艮第人便攻击这群无统帅的敌人；虽然他们人数少而敌人极多，却大获全胜；勃艮第人总共不过三千，却消灭了不下万敌。从那时起，该民族热诚皈依基督宗教。大约同时，亚略派主教巴尔巴斯于六月二十四日，在狄奥多西第十三次执政官任、瓦伦提尼安第三次执政官任下去世，萨巴提乌斯继任。关于这些事，已说得足够多了。\par

% structure: p0066 source=h2 target=subsection reason=relative-heading-rank; base=h2

\subsection{第三十一章：聂斯多略骚扰马其顿派。}

诚然，\Person{聂斯多略}的行为有违教会的惯例，且在其他方面也为自己招来憎恨，这从他任主教期间所发生的事便可看出。赫勒斯滂地区革尔马城的主教安多尼，受聂斯多略不容异己之例的鼓动，以贯彻宗主教意图为借口，开始迫害马其顿派。马其顿派一度忍受他的骚扰；但当安多尼变本加厉时，他们再也无法忍受他的苛待，遂陷入绝望的鲁莽，收买了两个将义置于次位、利字当先的人，刺杀了折磨他们的人。马其顿派犯下此罪后，聂斯多略借此机会变本加厉地对付他们，并说服皇帝剥夺了他们的教堂。因此，他们不仅失去了在君士坦丁堡城内旧城墙前所拥有的教堂，也失去了在基齐库斯以及赫勒斯滂乡间所拥有的许多其他教堂。当时他们中有许多人因此归顺了大公教会，宣认了\Quote{\OriginalTerm{同质}{homoousian}}信德。但正如谚语所说：\Quote{醉汉不缺酒，好斗者不乏争端}。聂斯多略的下场正是如此：他在竭力将别人逐出教会之后，自己却因以下原因被逐。\par

% structure: p0068 source=h2 target=subsection reason=relative-heading-rank; base=h2

\subsection{第三十二章　司铎阿纳斯塔修斯，聂斯多略的信仰由此被歪曲}

\Person{聂斯多略}有一位他从安提约基雅带来的同伴，名叫\Person{阿纳斯塔修斯}的司铎；聂斯多略对他极为敬重，并在处理最重要的事务时与他商议。这位阿纳斯塔修斯某日在教堂讲道时说：\Quote{谁也不可称玛利亚为\OriginalTerm{天主之母}{Theotocos}，因为玛利亚不过是一个女人；天主不可能由女人所生。}这话引起极大骚动，使许多圣职人员和教友深感不安；他们此前一直受教导承认基督是天主，并且决不因降生的安排而将祂的人性与天主性分离，他们听从宗徒的话说：\BibleQuote{纵使我们曾按血肉认识基督，但从今以后我们不再这样认识祂了。}\BibleRef{格后 5:16}又说：\BibleQuote{所以，我们应当离开基督初步的道理，努力前进，达到圆满的地步。}\BibleRef{希 6:1}正如我们所说，教会对这番言论极为反感，而聂斯多略急于维护阿纳斯塔修斯的命题——他不愿自己所敬重的人被定为亵渎——便就此发表了多次公开演讲，采取了论战的态度，完全否定了\OriginalTerm{天主之母}{Theotocos}这一称号。因此，关于这一问题的争论，有些人持一种精神，另一些人持另一种，由此引发的讨论分裂了教会，如同黑暗中的角斗者互相搏斗，各方都发出最混乱、最矛盾的断言。聂斯多略因此在大众中获得了宣扬主不过是凡人的亵渎教条的恶名，并试图将撒摩沙塔的保禄和福提诺的教条强加于教会；争论如此激烈，以至于人们认为有必要召开一次大公会议来审查这一争议。我本人读过聂斯多略的著作，发现他是个不学无术的人，并愿坦诚表达我对他的个人看法：既然我已经毫无个人敌意地提及他的过错，我也同样不会因敌对者的指控而有失偏颇地贬低他的优点。因此，我不能承认他是撒摩沙塔的保禄或福提诺的追随者，也不能承认他否认基督的天主性；但他似乎对\OriginalTerm{天主之母}{Theotocos}这个词感到害怕，如同见了某个可怕的幽灵。事实上，他在这件事上表现出的无端恐惧恰恰暴露了他极端的无知：因为他天生口才流利，被人视为有教养，但实际上他无知得可耻。事实上，他轻视对古代注释者进行精确钻研的苦差；并且，因自己口才敏捷而自高自大，不注重古人的教诲，自认为是最伟大的人。显然他不知道，在古老的抄本中，《\BibleBook{若望}{一书}》上写着：\BibleQuote{凡分裂耶稣的灵，都不是出于天主。}这段经文的残缺应归于那些企图将天主性与人性安排分开的人；或者用早期注释者的话说，有些人篡改了这封书信，旨在\Quote{将基督的人性与祂的天主性分离}。但在救主身上，人性与天主性结合，不是两个位格，而只有一个。因此，古人有鉴于此，毫不迟疑地称玛利亚为\OriginalTerm{天主之母}{Theotocos}。正如\Person{欧瑟比·庞菲洛}在他的\WorkTitle{君士坦丁传}第三卷中所写的那样：\par

\Quote{事实上，\Quote{天主与我们同在}为了我们而屈尊降生；祂诞生之地，希伯来人称之为白冷。因此，虔诚的皇后\Person{海伦纳}用最辉煌的纪念物装饰了那\OriginalTerm{天主之母}{Theotocos}的产房，以极其丰富的饰品点缀了那片圣地。}\par

奥力振在其\WorkTitle{注释}第一卷论宗徒致罗马人书时，也详尽阐明了\OriginalTerm{天主之母}{Theotocos}一词的用法。显然，聂斯多略对古代著作所知甚少，正因如此，如我所言，他只是反对这词：因为他并不像福提诺或萨莫萨特的保禄那样断言基督只是一个凡人，他自己发表的讲道集已充分证明这一点。在这些讲道中，他从未否认天主圣言的本有位格；相反，他始终维护其本质的、独立的位格与存在。他也从未像福提诺和萨莫萨特人那样，也不像摩尼教徒和孟他努的追随者那样妄图否认其自立性。事实上，我从亲自阅读聂斯多略的著作以及他的仰慕者的保证中，发现他正是如此。但他这场无谓的争论在教会中引起了不小的骚动。\par

% structure: p0072 source=h2 target=subsection reason=relative-heading-rank; base=h2

\subsection{第三十三章 逃亡奴隶亵渎大堂祭台}

正当事情处于这种状态时，教会内发生了一件暴行。原来，一位显贵的外籍仆从因受主人虐待，逃到教堂；他们竟拔剑直冲祭台。无论怎样恳求，都无法说服他们离开，以致阻碍了神圣礼仪的进行；但因他们固执地占据数日，挥舞武器，威胁任何敢于靠近的人——事实上，他们杀死了一名神职人员，又伤了另一人——最终只得自尽。一位亲历此亵渎圣所的人评论说，这种亵渎是不祥之兆，并引用了一位古代诗人的两句抑扬格诗句以支持他的看法：\par

\Quote{当圣殿被不敬之罪玷污时，\newline{}便发生这样的预兆。}\par

发出这预言的人并未因这些不祥之兆而失望：因为它们似乎预示了百姓的分裂，以及引发此事之人的被罢免。\par

% structure: p0076 source=h2 target=subsection reason=relative-heading-rank; base=h2

\subsection{第三十四章 厄弗所会议反对聂斯多略。他被罢免。}

不久之后，皇帝下令各地主教聚集在厄弗所。因此，逾越节庆典刚过，聂斯多略就在一大群支持者的护送下前往厄弗所，并发现许多主教已经到达。亚历山大里亚主教济利禄有所耽搁，直到临近五旬节才抵达。五旬节后五天，耶路撒冷主教朱文纳尔到达。当安提约基雅的若望尚未抵达时，已聚集的人开始审议此事；亚历山大里亚的济利禄开始了激烈的言词交锋，意图恐吓聂斯多略，因为他对他深怀敌意。当许多人宣告基督是天主时，聂斯多略说：\InnerQuote{我不能称那两三个月大的婴儿为天主。因此，我对你们的血是清白的，今后不再到你们中间来。}说完这些话，他便离开了集会，此后与那些持有类似见解的其他主教另行集会。于是，在场的人分裂为两派。支持济利禄的那一派组成了会议，传唤聂斯多略；但他拒绝与会，并推迟到安提约基雅的若望到达。因此，济利禄的追随者着手审查聂斯多略就争议问题所发表的公开讲道；在反复阅读后，认定它们明显亵渎了天主圣子，便罢免了他。事情发生后，聂斯多略的支持者另行组成会议，在此会议中罢免了济利禄本人，并与他一同罢免了厄弗所主教梅默农。这些事过后不久，安提约基雅主教若望到达；得知所发生的一切后，他毫不含糊地谴责济利禄是这场混乱的根源，因他如此仓促地罢免了聂斯多略。于是，济利禄联合朱文纳尔报复若望，也罢免了若望。当事情陷入如此混乱时，聂斯多略看到所引发的争端正导致共融的破坏，便痛悔地称玛利亚为\OriginalTerm{天主之母}{Theotocos}，并喊道：\InnerQuote{如果你们愿意，就让玛利亚被称为\OriginalTerm{天主之母}{Theotocos}吧，愿一切争论止息。}但尽管他作了这样的撤回，却未受到理会；因为他被罢免的决定并未撤销，而是被流放到绿洲，至今仍在那里。这次会议就此结束。这些事发生于六月二十八日，巴苏斯和安提奥库斯执政之年。若望返回自己的主教区后，召集了几位主教，罢免了也已返回教区的济利禄；但不久之后，他们放下了敌意，彼此接纳为朋友，互相恢复了对方的主教席位。但聂斯多略被罢免后，君士坦丁堡的各教会中激起了巨大的骚动。因为百姓因我们之前所称的他不幸的言论而分裂；神职人员则一致弃绝他。我们基督徒惯常对那些宣扬任何亵渎道理的人宣布这种判决，即当我们将他们的不敬公之于众，如同立在柱子上，使之受众人咒骂。\par

% structure: p0078 source=h2 target=subsection reason=relative-heading-rank; base=h2

\subsection{第三十五章 马克西米安被选为君士坦丁堡主教，尽管有人希望普罗克鲁斯接任此职。}

在这之后，关于君士坦丁堡主教的选举又有一场辩论。许多人支持我们已经提及的斐理伯，但更多的人主张普洛克鲁斯。若非一些最有影响力的人士以教会法规禁止从一教区调任至另一城市为由进行干预，普洛克鲁斯的候选本会成功。民众相信了这种说法，因此受阻。聂斯多略被废黜约四个月后，一位名叫玛克西米安的人被擢升为主教，他过苦修生活，也被列为司铎。他因自费修建墓穴以安葬虔诚者去世后的遗体而享有盛誉，但\BibleQuote{言语粗疏}\BibleRef{格后11:6}，且倾向于过安静的生活。\newline{}\par

% structure: p0080 source=h2 target=subsection reason=relative-heading-rank; base=h2

\subsection{作者对从一教区调任至另一教区之有效性的意见。\newline{}}

但由于某些人以教会法规中有一条禁令为由，阻止了普洛克鲁斯的选举，因为他先前已被任命为齐兹库斯教区的主教，我愿就此略加评论。那些当时擅自提出这种排除理由的人，在我看来并未陈述真相；他们或是出于对普洛克鲁斯的偏见，或至少自己对法规以及教会中屡次且常有裨益的先例全然无知。欧瑟比·庞斐理在《教会史》第六卷中记载：卡帕多细亚某城的主教亚历山大，为虔诚目的来到耶路撒冷，被该城居民挽留，被立为主教，作为纳尔西苏斯的继任者；此后他终生在那里管理教会。由此可见，在我们祖先看来，每当认为适宜时，将主教从一城调任至另一城是多么无足轻重的事。但若有必要证明那些阻止普洛克鲁斯受祝圣之人的说法纯属虚妄，我将在此论著中附上与此相关的法规。其内容如下：\newline{}\par

\Quote{如果有人被祝圣为主教后，因并非自身过错，而是由于人民不愿接纳，或因其他必要原因，未能前往他被任命的那座教会，那么他应享有其所获品秩的尊荣与职能，只要他不干预他可能服务的教会的事务。但他有义务服从省区会议在审查此事后所作的任何决定。}\newline{}\par

以上是法规的原文。现在，我将通过列举那些曾被调任的主教的名字，来证明许多主教曾因特殊情况的紧急需要而从一城调至另一城。佩里格内斯被祝圣为帕特雷主教；但因该城居民拒绝接纳他，罗马主教指示将他指派给科林多的教省主教区，该区因其前任主教去世而空缺；他在那里主持教会直至去世。额我略最初被立为卡帕多细亚的萨西玛主教，后来被调至纳齐安佐。梅利提乌斯首先管理塞巴斯蒂亚的教会，随后治理安提约基雅的教会。安提约基雅主教亚历山大将塞琉西亚主教多西特乌斯调至奇里乞亚的塔尔索。雷维伦提乌斯从腓尼基的阿尔卡被调走，后来至提洛。若望从吕底亚的戈尔杜姆城被调至普罗科内苏斯，并管理那里的教会。帕拉迪乌斯从赫勒诺波利斯被调至阿斯普纳；亚历山大从同城被调至阿德里亚尼。德奥斐禄从亚细亚的阿帕美亚被调至欧多克西奥波利斯（古称萨兰布里亚）。波利卡普从米西亚的塞克桑塔普里斯塔城被调至色雷斯的尼科波利斯。希耶罗斐禄从弗里吉亚的特拉佩佐波利斯被调至色雷斯的普洛提诺波利斯。奥普提穆斯从弗里吉亚的阿格达米亚被调至皮西迪亚的安提约基雅；西尔瓦努斯从色雷斯的菲利普波利斯被调至特洛亚。以上列举的主教从一教区调至另一教区的例子足以说明问题；关于西尔瓦努斯，他从色雷斯的菲利普波利斯调任至特洛亚，我认为在此处简要叙述是合适的。\newline{}\par

% structure: p0084 source=h2 target=subsection reason=relative-heading-rank; base=h2

\subsection{西尔瓦努斯（原菲利普波利斯主教，后为特洛亚主教）所行的奇迹。\newline{}}

西尔瓦努斯原是一位修辞学家，曾在诡辩派特洛伊鲁斯的学校里受教育；但为了在基督徒生活上达到完美，他开始苦修生活，并舍弃了修辞学家的披肩。君士坦丁堡主教阿提库斯留意到他，后来祝圣他为菲利普波利斯主教。他在色雷斯住了三年；但因无法忍受该地区的寒冷——他的体质虚弱多病——他恳求阿提库斯另派他人接替他的位置，声称他放弃在色雷斯的职务别无他因，只是因为寒冷。此事办妥后，西尔瓦努斯住在君士坦丁堡，在那里他实行极为严格的苦修，以至于蔑视时代的奢侈讲究，经常穿着干草编制的凉鞋出现在这座大都市拥挤的街道上。一段时间后，特洛亚的主教去世；因此该城居民前来见阿提库斯，商讨继任者的事。在他考虑应该祝圣谁给他们时，西尔瓦努斯恰巧前来拜访他，这立即让他免除了进一步的忧虑；因为他对西尔瓦努斯说：\Quote{你现在再无借口逃避教会的牧灵管理了；因特洛亚不是寒冷之地：所以天主已顾念你身体的软弱，为你提供了适宜的住所。因此，我的兄弟，请立刻前往那里。}西尔瓦努斯于是迁至该城。\newline{}\par

藉他的媒介，在此发生了一件奇迹，我即将叙述如下。在\Place{特洛阿}海岸上，最近建造了一艘巨大的货船，他们称之为\Quote{浮船}，用以运送巨大的石柱。这艘船必须下水。但尽管许多粗绳系在它上面，并动用大量人力，船却丝毫不动。这些尝试连续几天重复，结果相同，人们开始认为有魔鬼扣住了船；于是他们去找主教\Person{西尔瓦诺}，恳求他去那里祈祷。因为他们以为只有这样船才能下水。他以其特有的谦卑回答说，他只是个罪人，这工作属于某个义人，而不属于他。最后被他们不断的恳求所说服，他走近海岸，在那里祈祷后，他触摸了其中一根绳子，并鼓励其余人用力拉，船在第一次拉动下立即启动，迅速滑入海中。这个由\Person{西尔瓦诺}之手所行的奇迹，激起了全省人民对虔诚的热忱。但\Person{西尔瓦诺}非凡的德行还以其他各种方式显明。他注意到神职人员利用诉讼当事人的争执谋利，就从不任命任何神职人员为审判官；而是让诉讼人的文件交给他自己，他召来某个他所信任的、正直的虔诚信友，将案件的裁决交给他，很快便公平地解决了诉讼人之间的所有分歧；通过这种做法，\Person{西尔瓦诺}赢得了各阶层人士的极大声誉。\par

我们确实在叙述\Person{西尔瓦诺}的事迹时，相当偏离了历史的进程；但我想这不会没有益处。不过现在让我们回到我们离开的地方。\Person{马克西米安}于十月二十五日，在\Person{巴苏斯}和\Person{安提奥库斯}执政年间被祝圣后，教会的事务变得更有秩序、更安宁了。\par

% structure: p0088 source=h2 target=subsection reason=relative-heading-rank; base=h2

\subsection{第三十八章 \Place{克里特岛}许多犹太人信奉基督信仰}

大约在这一时期，居住在\Place{克里特岛}的大量犹太人因以下不幸的事件而皈依了基督信仰。某个犹太骗子假称自己是\Person{梅瑟}，奉天使命率领该岛的犹太人，并引导他们穿越海洋；因为他声称自己就是从前带领以色列人穿越红海的那一位。因此整整一年，他走遍岛上各城，说服犹太人相信这些保证。他还叫他们舍弃金钱和其他财产，承诺带领他们从干涸的海底进入福地。受这种期望的欺骗，他们忽视了一切事务，鄙弃自己所拥有的，任凭任何人随意取走。当这骗子指定的出发日期到来时，他亲自领头，所有人带着妻儿跟随其后。他带领他们来到一处俯瞰大海的海岬，从那里命令他们纵身跳入海中。那些最先到达悬崖的人照做了，立即丧生，有的撞在岩石上粉身碎骨，有的溺在水中；若非天主的眷顾使一些基督徒渔夫和商人出现在那里，更多人会丧命。这些人救起了几个几乎淹死的人，他们在危险中意识到了自己行为的疯狂。他们阻止了其余人跳下，告诉他们那些先跳者的毁灭。最后，犹太人意识到自己受了多么可怕的欺骗，首先责怪自己轻信的愚昧，然后试图抓住假\Person{梅瑟}处死。但他们未能抓住他，因为他突然消失了，这使人们普遍认为那是一个恶毒的魔鬼，为了毁灭该国的犹太人而取了人形。由于这次经历，当时克里特岛的许多犹太人放弃了犹太教，归附了基督信仰。\par

% structure: p0090 source=h2 target=subsection reason=relative-heading-rank; base=h2

\subsection{第三十九章 诺瓦提安派教会得免于火}

此后不久，诺瓦提安派的主教\Person{保禄}获得了真正爱天主之人的声誉，远超从前。因为一场可怕的火灾在\Place{君士坦丁堡}爆发，空前未有——大火吞噬了城市的大部分地区——当最大的公共谷仓、\Person{阿喀琉斯}浴场以及火路上的一切都被烧毁时，它最终逼近了位于\Place{佩拉尔古斯}附近的诺瓦提安派教堂。主教\Person{保禄}见教堂危险，便奔到祭台上，将教堂及其内一切之保全托付给天主；他不仅为教堂祈祷，也为城市祈祷，不停息。天主垂听了他，正如事件清楚证明的：尽管大火从所有门窗进入该祈祷所，却未造成任何损害。而许多邻近建筑被这吞噬性元素所吞没，教堂本身却在整场大火中完好无损，战胜了狂暴的火焰。这种情况持续了两天两夜，大火熄灭时已烧毁了城市很大一部分；但教堂保持完整，更奇妙的是，无论在木料还是墙壁上，连一丝烟痕都未看到。此事发生在八月十七日，\Person{狄奥多西}第十四次执政时，他与\Person{马克西穆斯}共同担任执政官。从此以后，诺瓦提安派每年在八月十七日以特别的感恩向天主庆祝他们教堂的保全。几乎所有人——基督徒和大部分外教人——从那时起继续视那个地方为一个特别神圣的地点，因着为其安全所行的奇迹。关于这些事，就此打住。\par

% structure: p0092 source=h2 target=subsection reason=relative-heading-rank; base=h2

\subsection{第四十章　普禄格禄斯继马西米安为君士坦丁堡主教}

马西米安平安治理教会两年又五个月，于阿勒欧宾杜斯与阿斯帕尔执政之年四月十二日去世。该日恰是逾越节前一周的第五天，即星期四。当时狄奥多西皇帝为防选举主教时常引起的教会骚乱，对这事作了明智安排：为免再因主教人选发生争执而引发动乱，他在马西米安遗体下葬前毫不延迟地指示当时在城中的众主教将普禄格禄斯置于主教席位。因为他已收到罗马主教策肋斯定予以认可此选之信函，该信函已转寄给亚历山大的济利禄、安提约基雅的若望及得撒洛尼的鲁福斯；信函中向他们保证，一位曾被提名且实为某教会之主教的人调任另一教区并无障碍。普禄格禄斯就此就任主教职，并为马西米安举行了葬礼。现在也须简要介绍他。\par

% structure: p0094 source=h2 target=subsection reason=relative-heading-rank; base=h2

\subsection{第四十一章　普禄格禄斯的卓越品德}

普禄格禄斯早年即为读经员，勤勉上学，潜心修辞学。成年后常住来于阿提库斯主教身边，被立为其秘书。他大有进步后，其恩主擢升他为执事；后升为司铎，如前所述，由西辛纽斯祝圣为基齐库斯主教。但这些事都是很早以前的了。此时他被分派了君士坦丁堡的主教席位。他是一位道德卓越堪与任何人相比的人；因受阿提库斯培养，他热心效法那位主教的一切德行。然而他比其师更践行忍耐，后者有时对异端者施以严厉；普禄格禄斯对众人温良，深信仁慈比暴力更有效地促进真理之事业。因此他决意不困扰任何异端，在自己身上恢复了教会那屡遭损害的温和与尊严。在这点上他效法了狄奥多西皇帝的榜样；因皇帝曾决意永不使用皇权对付罪犯，普禄格禄斯也决意不搅扰那些对神圣之事持有与他不同见解的人。\par

% structure: p0096 source=h2 target=subsection reason=relative-heading-rank; base=h2

\subsection{第四十二章　称赞幼年狄奥多西皇帝}

为此皇帝极为尊敬普禄格禄斯。事实上他本人就是所有真正神职人员的榜样，从不赞许那些企图迫害他人的人。我甚至敢断言，在温和方面他超越了所有曾忠实地担任司祭职的人。现今最可恰如其分地应用\BibleBook{}{户籍纪}中关于梅瑟的记载：\Quote{梅瑟为人极其谦和，胜过地上的众人}\BibleRef{户 12:3}；因为狄奥多西皇帝\Quote{谦和胜过地上的众人}。正因这温和，天主使他的敌人不战而败，正如擒获篡位者若望以及随后蛮族的溃败所清楚证明的。因为宇宙的天主在我们这个时代赐予这位最虔诚的皇帝超自然的助佑，与古时赐予义人的相似。我写这些并非出于阿谀，而是如实叙述众人皆可作证的事实。\par

% structure: p0098 source=h2 target=subsection reason=relative-heading-rank; base=h2

\subsection{第四十三章　曾为篡位者若望盟友的蛮族之灾祸}

篡位者死后，他曾召来协助攻打罗马人的蛮族准备蹂躏罗马行省。皇帝闻讯后，照常将此事交托于天主，恒切祈祷，迅速获得了所求；因为值得注意蛮族所遭的灾祸。其首领名鲁伽斯，被雷电击毙。随后瘟疫爆发，消灭了他手下大部分人；若还不够，天降烈火，烧死了许多幸存者。这使蛮族异常恐惧；与其说因他们胆敢对罗马这样勇武的民族动武，不如说因他们看到罗马人有大能的天主相助。此时，主教普禄格禄斯在教会讲道时，将厄则克耳的一段预言应用于天主在最近危难中的拯救，因而大受赞赏。预言的话如下：\par

\Quote{人子，你应讲预言攻击哥格，即罗士、默舍士和突巴耳的首领。我必用瘟疫、流血、暴雨、冰雹审判他；我必降火和硫磺在他身上，在他所有的军队身上，并在他同来的许多民族身上。我必显示为大，显示为圣，在列国眼前显示自己；他们必知道我是主。}\par

如我所说，这种预言的应用大受欢迎，并提高了普禄格禄斯的声望。此外，天主的眷顾以各种其他方式奖赏了皇帝的温和，其中之一如下。\par

% structure: p0102 source=h2 target=subsection reason=relative-heading-rank; base=h2

\subsection{第四十四章　瓦伦提尼安皇帝与狄奥多西之女欧多克西亚的婚姻}

他与皇后\Person{欧多基亚}育有一女，名为\Person{欧多克西亚}。他的表弟\Person{瓦伦提尼安}，被他任命为西方皇帝，自己请求娶她为妻。当\Person{德奥多西}皇帝同意这一提议后，他们商议应在两大帝国的边界何处举行婚礼，最后决定双方为此前往\Place{德撒洛尼基}（大约在中途）。但\Person{瓦伦提尼安}派人传话，说不愿给他添麻烦，他自己会前往\Place{君士坦丁堡}。于是，他在西方留下足够的守卫后，便为婚事前往那里，婚礼在\Person{伊西多尔}与\Person{西纳托尔}执政官之年举行；之后他携妻返回西方。这一吉祥事件发生在当时。\par

% structure: p0104 source=h2 target=subsection reason=relative-heading-rank; base=h2

\subsection{第四十五章 \Person{金口若望}的遗体被皇帝应\Person{普罗克鲁斯}之请迁至\Place{君士坦丁堡}，安奉于\Place{宗徒教堂}。}

此后不久，主教\Person{普罗克鲁斯}将因\Person{若望}主教被废而脱离教会的人重新领回教会，他通过一项明智的措施平息了他们的愤怒。此事须在此叙述。他获得皇帝许可，将\Person{若望}的遗体从埋葬地\Place{科马纳}迁至\Place{君士坦丁堡}，时值他被废后第三十五年。他庄严地抬着遗体游城后，极为尊敬地将之安放在\Place{宗徒教堂}内。由此，那位主教的爱慕者被安抚，重新与\OriginalTerm{天主教会}{catholic Church}共融。此事发生在\Person{德奥多西}皇帝第十六次执政官之年的1月27日。但令我惊讶的是，自\Person{奥力振}死后对他所发的嫉妒，竟未波及\Person{若望}。因为前者在死后约两百年被\Person{德奥斐禄}绝罚；而后者在死后第三十五年即被\Person{普罗克鲁斯}恢复共融！\Person{普罗克鲁斯}与\Person{德奥斐禄}如此不同。有洞察力和智慧的人在这些事过去和现在如何发生上不会被欺骗。\par

% structure: p0106 source=h2 target=subsection reason=relative-heading-rank; base=h2

\subsection{第四十六章 诺瓦提安派主教\Person{保禄}去世，\Person{马西安}被选为继承人。}

在\Person{若望}遗体迁葬后不久，诺瓦提安派主教\Person{保禄}于同年执政官之下的7月21日去世。他在自己的葬礼上，在某种意义上，将不同派别联合在同一教会中。因为所有派别都同来送葬，一起咏唱圣咏，因为他在世时以其正直受到众人普遍尊敬。但\Person{保禄}在临终前行了一件值得纪念的事，我认为将此记载在这部历史中是有益的，因为本书的读者了解此事或许会感兴趣。为免那重要事迹的光辉被细枝末节的叙述所掩盖，我将不赘述他在病中如何严格恪守饮食的苦修纪律，毫无偏离自己规定的路线，也未因惯常的热忱而省略任何日常敬拜功课。但这件事是什么？他意识到自己离世在即，便派人将他管辖下所有教会的长老召来，对他们说：\Quote{趁我还活着，你们要关注选举一位主教来管理你们，免得日后你们教会的安宁受到扰乱。}他们回答说这事最好不要由他们决定：\Quote{因为我们当中有人对此事有这种看法，有人有那种看法，我们绝不会推举同一个人。因此我们希望您亲自指定您愿意继任的人。}\Person{保禄}说：\Quote{那么，请你们给我一份书面声明，保证你们将选举我所指派的人。}他们写下承诺并签名盖印后，\Person{保禄}从床上坐起，写下了\Person{马西安}的名字在纸上，没有告知在场任何人他所写的内容。此人曾被他擢升为长老，并受过他的苦修训练，但当时已外出。他将此文件折叠，加上自己的封印，又让主要长老们也封印；之后他将文件交给当时住在\Place{君士坦丁堡}的\Place{斯基泰}的诺瓦提安派主教\Person{马尔库斯}，对他说：\Quote{若天主乐意让我在此世继续多活些时日，就请将托付给你妥善保管的这份寄存物还给我。但如果他以为应当召我离去，你便会在此发现我选定的主教继承人。}此后不久他便去世；死后第三天，当着许多人的面展开文件，发现里面写着\Person{马西安}的名字，众人都高呼他配得此荣誉。于是立即派遣使者去接\Person{马西安}回\Place{君士坦丁堡}。这些人以虔诚的诡计，发现他住在\Place{弗里吉亚}的\Place{提贝里奥波利斯}，便将他带回；随即祝圣并安置在主教座上，在同月21日。\par

% structure: p0108 source=h2 target=subsection reason=relative-heading-rank; base=h2

\subsection{第四十七章 皇后\Person{欧多基亚}奉\Person{德奥多西}皇帝之命前往\Place{耶路撒冷}。}

此外，\Person{德奥多西}皇帝为所获恩惠向天主献上感恩；同时以最特别的敬意尊崇基督。他也派皇后\Person{欧多基亚}前往\Place{耶路撒冷}，因为她曾许愿若活着看到女儿成婚便去那里。皇后在访问圣城时，以最贵重的礼物装饰了那里的教堂；当时以及回来后，她又用各种饰物装饰了东方其他城市的所有教堂。\par

% structure: p0110 source=h2 target=subsection reason=relative-heading-rank; base=h2

\subsection{第四十八章 \Person{塔拉西乌斯}被祝圣为\Place{卡帕多西亚的凯撒勒雅}主教。}

大约在同一时期，在\Person{狄奥多西}担任第十七任执政官时，\Person{普罗克鲁斯}主教实行了一件古时无人做过的事。\Place{卡帕多西亚}的\Place{凯撒勒雅}主教\Person{菲尔穆斯}去世后，该地居民来到\Place{君士坦丁堡}，就主教人选征求\Person{普罗克鲁斯}的意见。当\Person{普罗克鲁斯}正在考虑应选何人担任该职务时，恰逢安息日所有元老都到教堂拜访他；其中也有\Person{塔拉修斯}，此人曾治理\Place{伊利里库姆}的各族与各城。由于传闻皇帝将委任他治理东部地区，\Person{普罗克鲁斯}便按手在他头上，立他为\Place{凯撒勒雅}的主教，以取代原近卫总督。\par

此时教会事务处于如此繁荣的状态。但我们将在此结束我们的历史，祈求各地教会连同城市与民族都能安居和平；因为只要和平持续，那些想要撰写历史的人便找不到可用的材料。而我们自己，圣洁的天主之人\Person{狄奥多}，若不是喜好骚乱的人选择保持安静，我们按照您的请求所承担的工作就无法在七卷书中完成。\par

这最后一卷记载了三十二年间的事件；整部历史共七卷，涵盖了一百四十年间。它始于第271个奥林匹克纪的第一年，即\Person{君士坦丁}被宣布为皇帝之时；结束于第305个奥林匹克纪的第二年，即\Person{狄奥多西}皇帝担任第十七任执政官之时。\par

\SourceRecord{Translated by A.C. Zenos.；Nicene and Post-Nicene Fathers, Second Series；Vol. 2.；Edited by Philip Schaff and Henry Wace.；Buffalo, NY: Christian Literature Publishing Co.,；1890.；<http://www.newadvent.org/fathers/26017.htm>.}
